\documentclass[11pt, reqno]{amsart}

% The main file. Packages, environments and generic letters are loaded first, then the notation of the note.
% Each section has its own file; a long section is split into parts, each part inputting the next,
% so that this file inputs only the first part of every section; the final section, Complements, is input separately.
% The preamble: packages, page layout, draft notes, theorem environments, cross-reference names and generic letters.
% Loaded by the main file before the notation. Only the packages loaded here are available to the sections.

% GENUS workshop style: page geometry, fonts, hyperref, fancyhdr and the claim-status tags (genus-collab/meta)
\usepackage{genus}
% text width of the original note (170 mm), so that its displayed formulas still fit the line
\geometry{margin=2cm}
\usepackage[utf8]{inputenc}
\usepackage[T1]{fontenc}

%-------------------------------------------------------------------------------
% Packages
\usepackage[dvipsnames]{xcolor}
\usepackage{amssymb}
\usepackage[mathscr]{eucal}
\usepackage{amsthm}
\usepackage{amsmath}
\usepackage{mathtools}
\definecolor{highlight}{HTML}{3465a4}
% hyperref is loaded by genus.sty, with the GENUS link colours
\hypersetup{hypertexnames=false, hyperfootnotes=false, linktocpage=true, bookmarksdepth=2}
\usepackage{cleveref}
\usepackage{url}
\usepackage{array}

%-------------------------------------------------------------------------------
% Typefaces and layout
% (mathpazo and courier dropped: genus.sty sets Latin Modern)
\usepackage{microtype}
\linespread{1.06}
\setlength{\parskip}{4pt}
\setlength{\parindent}{0pt}
\setcounter{secnumdepth}{3}
% enumitem is not assumed: first-level enumerate is labelled (i), (ii), ...; for named parts use
% description or \item[(a)] with an explicit label.
\renewcommand{\labelenumi}{(\roman{enumi})}
\renewcommand{\theenumi}{(\roman{enumi})}
\numberwithin{equation}{section}

%-------------------------------------------------------------------------------
% Draft notes: visible, coloured, usable in text and in math. Set \showdraftnotesfalse to suppress them.
\newif\ifshowdraftnotes
\showdraftnotestrue
\newcommand{\draftnotetext}[1]{\textcolor{red}{\footnotesize[\textbf{Draft note:} #1]}}
\newcommand{\draftnote}[1]{\ifshowdraftnotes\ifmmode\text{\draftnotetext{#1}}\else\draftnotetext{#1}\fi\fi}

%-------------------------------------------------------------------------------
% Theorem environments
\theoremstyle{plain}
\newtheorem{thm}{Theorem}[section]
\newtheorem{lem}[thm]{Lemma}
\newtheorem{prop}[thm]{Proposition}
\newtheorem{cor}[thm]{Corollary}
\theoremstyle{definition}
\newtheorem{defn}[thm]{Definition}
\newtheorem{rmk}[thm]{Remark}
\newtheorem{notation}[thm]{Notation}
% genus.sty defines an example environment on its own counter; redefine it on the counter of thm
\let\example\undefined\let\endexample\undefined\expandafter\let\csname c@example\endcsname\undefined
\newtheorem{example}[thm]{Example}
\newtheorem{conv}[thm]{Convention}

\crefname{equation}{Eq.}{Eqs.}
\crefname{thm}{Theorem}{Theorems}
\crefname{lem}{Lemma}{Lemmas}
\crefname{prop}{Proposition}{Propositions}
\crefname{cor}{Corollary}{Corollaries}
\crefname{defn}{Definition}{Definitions}
\crefname{rmk}{Remark}{Remarks}
\crefname{notation}{Notation}{Notations}
\crefname{example}{Example}{Examples}
\crefname{conv}{Convention}{Conventions}
\crefname{section}{Section}{Sections}
\crefname{subsection}{Section}{Sections}
\crefname{appendix}{Appendix}{Appendices}
\crefname{table}{Table}{Tables}

%-------------------------------------------------------------------------------
% Generic letters
\newcommand{\bbA}{\mathbb{A}}
\newcommand{\bbC}{\mathbb{C}}
\newcommand{\bbD}{\mathbb{D}}
\newcommand{\bbE}{\mathbb{E}}
\newcommand{\bbF}{\mathbb{F}}
\newcommand{\bbK}{\mathbb{K}}
\newcommand{\bbL}{\mathbb{L}}
\newcommand{\bbP}{\mathbb{P}}
\newcommand{\bbQ}{\mathbb{Q}}
\newcommand{\bbR}{\mathbb{R}}
\newcommand{\bbT}{\mathbb{T}}
\newcommand{\bbZ}{\mathbb{Z}}
\newcommand{\cA}{\mathcal{A}}
\newcommand{\cC}{\mathcal{C}}
\newcommand{\cE}{\mathcal{E}}
\newcommand{\cF}{\mathcal{F}}
\newcommand{\cG}{\mathcal{G}}
\newcommand{\cL}{\mathcal{L}}
\newcommand{\cN}{\mathcal{N}}
\newcommand{\cO}{\mathcal{O}}
\newcommand{\cP}{\mathcal{P}}
\newcommand{\cQ}{\mathcal{Q}}
\newcommand{\cW}{\mathcal{W}}
\newcommand{\cZ}{\mathcal{Z}}
\newcommand{\mgp}[1]{\mathsf{#1}}

%-------------------------------------------------------------------------------
% General
\newcommand{\dd}{\mathrm{d}}
\newcommand{\ri}{\mathrm{i}}
\newcommand{\re}{\mathrm{e}}
\newcommand{\de}{\partial}
% \llbracket: from a font package if available, otherwise a substitute (stmaryrd is not assumed)
\providecommand{\llbracket}{[\![}
\providecommand{\rrbracket}{]\!]}
\newcommand{\fps}[1]{\llbracket #1 \rrbracket}
\newcommand{\lser}[1]{(\!(#1)\!)}
\newcommand{\Frac}{\mathrm{Frac}}
\DeclareMathOperator{\Tr}{Tr}
\DeclareMathOperator{\tr}{tr}
\DeclareMathOperator{\End}{End}
\DeclareMathOperator{\Aut}{Aut}
\DeclareMathOperator*{\Res}{Res}
\DeclareMathOperator{\val}{val}
\DeclareMathOperator{\ev}{ev}
\DeclareMathOperator{\diag}{diag}
\DeclareMathOperator{\ord}{ord}
\newcommand{\tensor}{\otimes}
\newcommand{\coef}[1]{[#1]}                    % coefficient extraction, e.g. \coef{z^n}

% The notation of the note. Loaded by the main file after the preamble.
% A section file defines a new macro only if it is unavoidable, at its top, with \providecommand and the section prefix.

%-------------------------------------------------------------------------------
% Variables, points, rings
\newcommand{\Dq}{D}                            % q d/dq = -(1-u) d/du
\newcommand{\zh}{\hat{z}}                      % z/s^4 = L^4 z
\newcommand{\tc}{t_{\mathrm{c}}}               % varpi_1/varpi_0
\newcommand{\Cstar}{\mathbb{C}_{\ast}}         % C minus {0, 1, zeta, ..., zeta^4}
\newcommand{\AL}{\mathbb{A}_{L}}               % Q[L, L^{-1}, (L^5 - 1)^{-1}]
\newcommand{\bbH}{\mathbb{H}}                  % the state space, flat basis H^a
\newcommand{\Qn}{\mathscr{Q}}                  % the degree variable of F_g
\newcommand{\logq}{\mathfrak{l}}               % the formal log q
\newcommand{\Nring}{\mathscr{N}}               % the presented differential ring
\newcommand{\Cexp}{\mathscr{C}}                % the conifold expansion \Cexp_{\alpha,\beta}
\DeclareMathOperator{\Wr}{Wr}                  % Wronskian
\DeclareMathOperator{\Gau}{Gau}                % solution module of the gauge equation
\DeclareMathOperator{\wt}{wt}                  % weight
\DeclareMathOperator{\Hess}{Hess}              % Hessian

%-------------------------------------------------------------------------------
% The formal quintic, rows, (QRR)
\newcommand{\PFfour}{\mathrm{PF}_{4}}
\newcommand{\PFfive}{\mathrm{PF}_{5}}
\newcommand{\TFT}{\boldsymbol{\omega}}         % TFT (\TFT, \hat\TFT, \TFT^{\mathrm m}_v)
\newcommand{\Qr}{\mathcal{Q}}                  % rows Q^{(i)}_a
\newcommand{\Phiact}{\Phi^{\mathrm{act}}}
\newcommand{\Phican}{\Phi^{\mathrm{can}}}
\newcommand{\PhiSP}{\Phi^{\mathrm{SP}}}
\newcommand{\bk}{\mathsf{b}}                   % \bk_k = B_{2k}/(2k(2k-1))
\newcommand{\bw}{\mathbf{w}}                   % \bw_i
\newcommand{\hw}{h_{\bw}}                      % h_{\bw}
\newcommand{\ttr}{\mathfrak{t}}                % PPZ translation \ttr_R(z) = z1 - zR^{-1}(z)1

%-------------------------------------------------------------------------------
% Cyclic systems, gauges, normalisation
\newcommand{\Bprod}{\mathcal{B}}               % prod_a b_a
\newcommand{\Pc}{P_{\mathrm{can}}}             % the canonical gauge
\newcommand{\Dn}{\mathsf{D}}                   % normalisation series
\newcommand{\cn}{\mathsf{c}}                   % freedom P -> P c(z)
\newcommand{\Ghat}{\hat{G}}
\newcommand{\Rhat}{\hat{R}}
\newcommand{\YC}{Y^{\mathrm{C}}}
\newcommand{\YCQ}{Y^{\mathrm{C}}_{\mathrm{Q}}}
\newcommand{\IC}{I^{\mathrm{C}}}
\newcommand{\Conc}{\mathcal{C}}                % concomitant
\newcommand{\Jet}{\mathcal{J}}                 % jet matrix
\newcommand{\Mcyc}{M_{\mathrm{cyc}}}
\newcommand{\sh}{\mathrm{sh}}                  % the automorphism ell -> ell + 1
\newcommand{\fc}{\mathfrak{f}}                 % free coefficients \fc_m
\newcommand{\Fg}[1]{(\mathrm{F}_{#1})}         % the window (F_g)

%-------------------------------------------------------------------------------
% The model, transport, topological recursion
\newcommand{\Wv}{W_{v}}                        % p^4/4 + v/p
\newcommand{\Mo}{M_{0}}                        % the model R-matrix in Teleman form
\newcommand{\Fs}{\mathsf{F}}                   % model row series \Fs_b(t)
\newcommand{\Fm}{F^{\mathrm{m}}}               % model free energy
\newcommand{\FTR}{F^{\mathrm{TR}}}             % TR free energy
\newcommand{\cH}{\mathcal{H}}                  % loop spaces
\newcommand{\cR}{\mathcal{R}}                  % weighted rings (appendix A)
\newcommand{\cV}{\mathcal{V}}                  % the residue (Section 8)
\newcommand{\cX}{\mathcal{X}}                  % the variation field (Section 6)
\newcommand{\sK}{\mathcal{K}}                  % the slice K
\newcommand{\sKo}{\mathcal{K}_{0}}             % its kernel part K_0
\newcommand{\fr}{\mathrm{frame}}               % \fr(k,a)
\newcommand{\se}{\mathsf{e}}                   % Euler-dilaton vector z1 + vH
\newcommand{\ML}[1]{\mathrm{ML}_{\mathcal{K}}(#1)}  % the Model Lemma on K at genus #1

%-------------------------------------------------------------------------------
% Connection problem, Gamma identity, formal stationary phase
\newcommand{\fa}{\mathfrak{a}}                 % amplitude
\newcommand{\SPn}{\mathrm{SP}^{\circ}}         % normalised stationary phase
\newcommand{\SPh}{\widehat{\mathrm{SP}}}       % divided by its constant term
\newcommand{\Gst}{\Gamma_{\mathrm{St}}}        % formal Stirling form
\newcommand{\Ginf}{G_{\infty}}                 % the orbifold value


%-------------------------------------------------------------------------------
% Title

\title[The formal quintic at the conifold point]{Conifold regularity of the formal quintic, graph by graph\protect\draftnote{Working title; to be settled by the authors.}}
\author{The GENUS collaboration}
\genusproject{solutions/conifold-gap-formal-quintic}

\emergencystretch=2em
\begin{document}

\begin{abstract}
The formal quintic of Lho and Pandharipande is the $\cO(5)$-twisted equivariant Gromov--Witten theory of $\bbP^4$ specialised at the fifth roots of unity. Its genus-$g$ potential is a Givental--Teleman sum over stable graphs of genus $g$, with the $R$-matrix normalised at $q=0$ by quantum Riemann--Roch. We expand the theory at the conifold point $u=1-5^5q=0$ in a conifold frame: the graph sum is written in a differential ring generated by the periods, and a regular period and the vanishing period at $u=0$ are substituted for the two periods at $q=0$. This is a formal change of frame; no analytic continuation is involved. We factorise the expanded $R$-matrix as $R=R'\,\Mo(\tau(u)\zh)$, where $\Mo$ is the $R$-matrix of the Landau--Ginzburg model $p^4/4+v/p$ with primitive form $\dd p$ and $\tau(u)\zh$ is the natural variable defined by the canonical coordinate. We prove that, for every $g\geq 2$ and in every conifold frame of a two-parameter rational family, in the graph sum of $R'$ on the theory obtained from the topological part by the action of $\Mo(\tau(u)\zh)$, the term of every leg-free stable graph of genus $g$ with at least one edge is regular at the conifold point, and the single-vertex graph contributes $c_g\tc^{2-2g}$ plus a regular series, where $\tc$ is the flat coordinate of the vanishing period and $c_g=(-1)^{g-1}B_{2g}/\bigl(2g(2g-2)\,5^{g-1}\bigr)$. Summing over graphs gives the conifold gap for this expansion of the genus-$g$ potential of the formal quintic, for every $g\geq2$. The proof reduces the statement, graph by graph, to a lemma about the model, proved by topological recursion on the model curve, and to the equality of two normalisations of the $R$-matrix: the canonical one at the conifold point and the quantum Riemann--Roch one at $q=0$. We prove this equality through an identity of formal stationary phase for a Mellin--Barnes integral, which gives the value at $q=0$ of its saddle series. The results hold modulo the following inputs from the literature, and every other step is proved: results on topological recursion, applied only to the model curve; results of Pandharipande, Pixton and Zvonkine on the action of $R$-matrices, including the formula for the pull-back of $\psi$ classes along the map forgetting a point; the fact that the $R$-matrix is the asymptotic solution of the quantum differential equation, and its value at $q=0$ given by quantum Riemann--Roch; the genus-$0$ mirror theorem and the genus-$0$ axioms of the twisted theory; Teleman's reconstruction, with the axioms of the twisted theory in all genera; and a push-forward formula of Arbarello and Cornalba for products of $\psi$ classes. We do not treat the quintic threefold itself.
\draftnote{Draft abstract; to be settled by the authors.}
\end{abstract}

\maketitle
\genusdisclaimer

\setcounter{tocdepth}{1}
\tableofcontents

% Section 1: Introduction (label prefix int). The statements of this section are proved in the body;
% each statement names its body version. Continued in part 2.
\section{Introduction}
\label{sec:int}

%-------------------------------------------------------------------------------
\subsection{The setting}
\label{sec:int-setting}

Let the torus $(\bbC^*)^5$ act on $\bbP^4$ with weights $\lambda_0,\dots,\lambda_4$ on $H^0(\bbP^4,\cO(1))$. The \emph{formal quintic} of Lho and Pandharipande \cite[\S0.2, (5)]{LP18}, \cite[\S0.2, (5)]{LP20} is the $\cO(5)$-twisted equivariant Gromov--Witten theory of $\bbP^4$ at the specialisation $\lambda_i=\zeta^i$, $\zeta=\re^{2\pi\ri/5}$, (see \cite[\S3.1; \S3.2, (9)]{GJR18}; the resulting hypergeometric series is studied in \cite{ZZ08}): its invariants are $N_{g,d}=\int_{[\overline{\mathcal M}_g(\bbP^4,d)]^{\mathrm{vir}}}e(R\pi_*f^*\cO(5))$, defined by virtual localisation \cite{MR1666787} before the weights are specialised, and $F_g=\sum_{d\geq0}N_{g,d}\Qn^d$ (\cref{def:set-fq}). By the genus-zero mirror theorem the $J$-function at $t=T(q)H$ is $I(q,z)/C_0(q)$, with $C_0=\sum_d\frac{(5d)!}{(d!)^5}q^d$ and the mirror map $T=I_1/C_0$ (\cref{thm:set-mirror}), and $\Qn=\re^{T(q)}$. With $u:=1-5^5q$ and $L:=u^{-1/5}=1/s$, the quantum product is cyclic (\cref{prop:set-cyclicLR}); $q=0$ is $u=1$, and $u=0$ is the \emph{conifold point}.

By Teleman's reconstruction \cite[Theorem~(2) in (1.15)]{MR2917177}\draftnote{numbers from arXiv v3; journal numbering not compared} (see \cite[\S4.1, Theorem~4.1]{GJR18}) $F_g$ is the Givental--Teleman sum over the stable graphs of genus $g$ of Givental's symplectic $R$-matrix $R(z)=1+O(z)$ acting on the topological field theory $\TFT$, in the convention of \cite[Defs.~2.2, 2.13]{PPZ15} (\cref{def:set-rows,def:set-graphsum,thm:set-GT}); at $q=0$ quantum Riemann--Roch, in the form \cite[(17)]{GJR18}, \cite[(19)--(20)]{GJR17} (which rest on \cite[\S2.2, Theorem~1]{Giv01} for the tangent weights and on \cite[Corollary~1]{MR2276766} for the twist), gives
\[
R(q=0)=\diag\bigl(h_{\bw_0},\dots,h_{\bw_4}\bigr),\qquad h_{\bw}(z)=\exp\Bigl(-\sum_{k\geq1}\bk_k\,z^{2k-1}\sum_{a\in\bw}a^{1-2k}\Bigr),\qquad \bk_k=\frac{B_{2k}}{2k(2k-1)},
\]
with $\bw_i=\{\zeta^i-\zeta^j\}_{j\neq i}\cup\{-5\zeta^i\}$, the input (QRR) (\cref{thm:set-qrr}); $h:=h_{\bw_0}$. In the conifold expansion defined below, where $L=s^{-1}$, the lowered rows $\Qr^{(i)}_a(z)=\eta(H^a,R(z)\tilde e_i)$ grow with $L$, and the natural variable is $\zh:=L^4z=z/s^4$.

\emph{Conifold frames.} Let $f_0,f_1,f_2$ be the solutions $f_k=u^k+O(u^3)$ at $u=0$ of the Picard--Fuchs operator $\PFfour$ of the quintic, and $\varpi_1=f_1+\frac7{10}f_2=u+O(u^2)$ the vanishing period (\cref{lem:set-frobbasis}). A \emph{conifold frame} is a period $\varpi_0=f_0+\alpha\varpi_1+\beta f_2$, $(\alpha,\beta)\in\bbQ^2$, and $\tc:=\varpi_1/\varpi_0=u+O(u^2)$ (\cref{def:set-frames}). The row coefficients lie in $\Nring\tensor\bbQ(\zeta)$, $\Nring$ a differential ring generated by the periods, $L^{\pm1}$ and a formal $\log q$ (\cref{def:set-confexp,cor:gau-rowsNN}); the \emph{conifold expansion} $\Cexp_{\alpha,\beta}\colon\Nring\to\bbQ\lser{s}$ expands at $u=0$ in the frame, and an expanded quantity is \emph{regular at the conifold point} if it lies in $\bbQ\fps{u}$.

\emph{The model and the natural form.} The model is the Landau--Ginzburg potential $\Wv(p)=p^4/4+v/p$ with primitive form $\dd p$, whose five critical points $p^5=v$ coalesce into the pole $p=0$ as $v\to0$ (\cref{def:set-model}). Its rows, the normalised stationary-phase series of $\int p^{\nu_b}\re^{-\Wv(p)/z}\dd p$, $\nu=(0,-1,3,2,1)$, form the model $R$-matrix $\Mo(t)$, $t=z/\varsigma^4$, $\varsigma=v^{1/5}$ (\cref{prop:set-modrows,def:set-M0}). Let $\tau(u):=1/{}_2F_1(1,\frac45;\frac95;u)=\bigl(\sum_{n\geq0}\frac{4u^n}{5n+4}\bigr)^{-1}$ and $v(u):=u\,\tau(u)^{-5/4}$ (\cref{def:set-tau}). At $v=v(u)$ the canonical coordinates of the model are those of the formal quintic measured from the conifold point, and $z/\varsigma^4=\tau(u)\zh=-\frac54\zeta^iz/(u_i-u_i^{\mathrm c})$ (\cref{lem:set-canon,prop:gau-equalprod}); so $M:=\Mo(\tau(u)\zh)$ is the $R$-matrix of the model at $v(u)$ in the variable of the formal quintic. Writing the expanded rows as $\Qr=R'M$, $R':=\Qr M^{-1}$, gives the \emph{natural form} of the graph sum: the outer graph sum $R'.(M.\TFT)$ has the vertices $M.\TFT$ and the edges and leaves of $R'$ (\cref{def:set-outer}).

%-------------------------------------------------------------------------------
\subsection{The conifold gap}
\label{sec:int-gap}

The conifold gap is a prediction of the topological B-model: near a conifold point the genus-$g$ potential, expanded in the flat coordinate of the vanishing period, has the universal leading pole of the $c=1$ string at the self-dual radius and no other pole \cite{Bershadsky:1993cx} (see \cite[\S1]{HKQ09}), \cite{Ghoshal:1995wm}, \cite[\S2.3]{HK07}, \cite[(1.2)]{HKQ09}. For the quintic threefold it was left open after \cite{GJR18,CGL21,CGL26} (see \cite[Remark~1.6]{GJR18}), was an input of high-genus computations \cite{HKQ09,LR22}, and is proved in \cite[Theorem~2]{CGYZ26} by master-space localisation. It holds for the resolved conifold \cite[\S1.1, footnote~1]{Br25}, for local $\bbP^2$ in all genera \cite{Br25} (after a check to genus $7$ in \cite{CI21}, see \cite[\S1.2]{Br25}), and for the topological recursion of toric mirror curves at a generic one-node degeneration \cite{FCCZ26}.

The potentials of the formal quintic satisfy holomorphic anomaly equations \cite[Theorems~1 and~2, \S0.4]{LP20}. By \cite[Remark~1.9]{CGYZ26}, the approach of \cite{GJR18} reduces the gap of the quintic to two key statements, the conifold gap for the formal quintic and the conifold regularity of its correlators with special insertions, and the authors of \cite{CGYZ26} expect their method to prove the first as well. We prove, conditionally on the inputs of \cref{sec:int-inputs}, a graph-local statement for the formal quintic (Theorem~A); summing over the graphs gives its conifold gap in every genus $g\geq2$ (Corollary~B), with the constant of \cite[Conjecture~1.5]{CGYZ26} (\cref{rmk:int-constant}).

The proof of \cite{CGYZ26} also factorises an $R$-matrix, that of a master-space theory \cite[Proposition~3.1]{CGYZ26}; here the $R$-matrix of the formal quintic itself is factorised, through a universal local model. \draftnote{The authors decide whether to cite their companion notes and how to relate the note to the sentence of \cite[\S1.1]{Br25} on $R$-matrix methods.}

The normalisation at $q=0$ is carried to the conifold point by an identity of formal stationary phase (Theorem~Gamma), through a change of chart as in Givental's treatment of projective spaces \cite[\S10]{MR1901075}. The analogous identity for the Landau--Ginzburg integral of Guo, Janda and Ruan is stated without proof in \cite[Remark~7.2]{GJR18}, and the value at the orbifold point is \cite[Lemma~7.7, (37)]{GJR18}; for Gamma factors in quantum cohomology see also \cite[Lemma~4.20]{MR2553377} and \cite{GGI16}. \draftnote{Re-check the latest versions of arXiv:2609.13921 and arXiv:2509.19298 before submission.}

% Section 1, part 2: the statements. Continued in part 3.

%-------------------------------------------------------------------------------
\subsection{Results}
\label{sec:int-results}

For a leg-free stable graph $\Gamma$, a matrix series $X$ and a cohomological field theory $\Omega$, $\mathrm{Cont}_\Gamma[E_X,q_X;\Omega]$ is the contraction of the vertex tensors of $\Omega$ along the edges $E_X$ of $X$, with the leaf point $q_X$ of $X$, as in \cite[Defs.~2.2, 2.13]{PPZ15} when $X$ is symplectic with $X(0)=1$, and with the edge of \cref{def:set-outer} in general, including the factor $|\Aut\Gamma|^{-1}$ (\cref{def:set-graphsum,def:set-outer,def:flh-outer}). The \emph{vertex graph} has one vertex, of genus $g$, and no edge; its term contains the leaf decorations.

\begin{thm}[Theorem~A]\genusstatus{AI-proposed, unchecked}
\label{thm:int-main}
Let $g\geq2$ and let $(\alpha,\beta)\in\bbQ^2$ define a conifold frame (\cref{def:set-frames}), with vanishing period $\varpi_1=u+O(u^2)$ and $\tc=\varpi_1/\varpi_0$. Let $M:=\Mo(\tau(u)\zh)$, exactly in $u$, and $R':=\Qr M^{-1}$, where $\Qr$ is the matrix of the lowered rows of \cref{def:set-rows} expanded in the frame (\cref{def:set-confexp}). For a leg-free stable graph $\Gamma$ of genus $g$ let
\[
T_\Gamma:=\mathrm{Cont}_\Gamma\bigl[E_{R'},q_{R'};M.\TFT\bigr]
\]
be its outer term (\cref{def:set-outer}), where $\TFT$ is the topological field theory of the formal quintic (\cref{def:set-frob}), expanded in the frame. Then:
\begin{enumerate}
\item $T_\Gamma\in\bbQ\fps{u}$ for every $\Gamma$ with at least one edge;
\item for the vertex graph, $T_\Gamma\in c_g\tc^{2-2g}+\bbQ\fps{u}$, where
\[
c_g=\frac{(-1)^{g-1}B_{2g}}{2g(2g-2)\,5^{g-1}}=\frac{|B_{2g}|}{2g(2g-2)\,5^{g-1}} .
\]
\end{enumerate}
\end{thm}

\Cref{thm:int-main} is proved as \cref{thm:mai-main}; for one $g$ it is Property N($g$) of \cref{def:set-N}. Here $c_2=1/1200$, $c_3=1/25200$ (\cref{rmk:set-cg}), and the hatted terms are $G_\Gamma=u^{2g-2}T_\Gamma$ (\cref{lem:flh-hat}). The statement is about each graph separately; we do not know how to obtain it from a statement about the singular part of the sum (\cref{rmk:mai-graphlocal}).

\begin{cor}[Corollary~B: the conifold gap of the formal quintic]\genusstatus{AI-proposed, unchecked}
\label{cor:int-gap}
Let $g\geq2$ and $(\alpha,\beta)\in\bbQ^2$. The conifold expansion $F_g^{\mathrm{cf}}\in\bbQ\lser{s}$ of the genus-$g$ potential of the formal quintic in the frame $(\alpha,\beta)$ (\cref{def:mai-Fcf}) satisfies
\[
F_g^{\mathrm{cf}}\in c_g\tc^{2-2g}+\bbQ\fps{u} .
\]
Since $\bbQ\fps{u}=\bbQ\fps{\tc}$, the series $F_g^{\mathrm{cf}}$ has no fractional power of $u$, no term $\tc^k$ with $3-2g\leq k\leq-1$, and leading term $c_g\tc^{2-2g}$.
\end{cor}

Here $F_g^{\mathrm{cf}}$ is the image under $\Cexp_{\alpha,\beta}$ of the genus-$g$ graph sum written in $\Nring$, whose evaluation is $F_g$ by (T) and (QRR); by the group property of the action \cite[Def.~2.13, Prop.~2.14]{PPZ15} it is the sum of the $T_\Gamma$ (\cref{cor:mai-gap}).

The proof compares two normalisations of the $R$-matrix. The cyclic systems of the formal quintic and of the model at $v(u)$ are related by a gauge regular at $u=0$, unique up to a scalar factor $\cn(z)$ (Theorem~U, \cref{thm:gau-U}); the \emph{canonical gauge} $\Pc$ is fixed at the conifold point by the gauge identity $\Pc\YC(v(u))=\YCQ$ between the vanishing solutions (\cref{thm:gau-G,def:gau-can}), the rows $\Qr$ at $q=0$ by (QRR). Both $\Qr$ and $\Pc\Qr^{\mathrm m}(v(u))$, $\Qr^{\mathrm m}$ the model rows, are formal WKB solutions of the same system with the same exponentials, so
\[
\Qr\,\Dn=\Pc\,\Qr^{\mathrm m}(v(u)),\qquad \Dn=\diag(\Dn_0,\dots,\Dn_4),\qquad \Dn_i(z)=\Dn_0(\zeta^{-i}z),\qquad \Dn_0\in\bbQ\fps{z}
\]
(\cref{thm:gau-R0,cor:gau-colours}; that $\Dn_0$ has rational coefficients is part of \cref{thm:con-glue}).

\begin{thm}[Theorem~C: the normalisation]\genusstatus{AI-proposed, unchecked}
\label{thm:int-D}
In every conifold frame, $\Dn_0=1$ in $\bbQ\fps{z}$. Hence $\Dn=1$, and the window $\Fg{g}$, that is $\Dn_0\equiv1\bmod z^{3g-2}$ (\cref{def:gau-Fg}), holds for every $g\geq2$. Equivalently, $\Qr=\Pc\,\Qr^{\mathrm m}(v(u))$: the actual gauge is the canonical gauge, and the $R$-matrix of the formal quintic, normalised at $q=0$ by (QRR), is in the canonical normalisation at the conifold point.
\end{thm}

\Cref{thm:int-D} is \cref{thm:gam-D} with \cref{cor:mai-actual}; its consequences for the natural form are \cref{cor:mai-natfact}. It is reduced to one value at $q=0$ of a saddle series:

\begin{thm}[Theorem~Gamma]\genusstatus{AI-proposed, unchecked}
\label{thm:int-gamma}
Let $c(z;L)=\sum_{n\geq0}c_n(L)z^n$, $c_0(L)=1$, be the full-Stirling saddle series of \cref{def:gam-c}: the normalised stationary-phase series, divided by its constant term ($\SPh$ of \cref{def:fsp-wick}), at $\sigma=L$ of
\[
\Gamma\Bigl(1+\frac{5\sigma}{z}\Bigr)\prod_{j=0}^{4}\Gamma\Bigl(\frac{\zeta^j-\sigma}{z}\Bigr)(-q)^{\sigma/z}\,\dd\sigma,\qquad q=q(L)=\frac{1-L^{-5}}{5^5},
\]
with every Gamma factor replaced by its formal Stirling form (\cref{def:fsp-stirling}). Then every $c_n$ lies in $\AL=\bbQ[L,L^{-1},(L^5-1)^{-1}]$, and:
\begin{enumerate}
\item[(a)] no $c_n$ has a pole at $L=1$, that is, at $q=0$, and $c(z;1)=h(z)$;
\item[(b)] no $c_n$ has a pole at $L=0$, that is, at $q=\infty$, and
\[
c(z;0)=\Ginf(z):=\exp\Bigl(\sum_{l\geq1}\frac{(-1)^{l+1}B_{5l+1}(\frac15)}{l(5l+1)}\,z^{5l}\Bigr),
\]
the series given by \cite[Lemma~7.7, (37)]{GJR18} at $m=5$, $k=1$, $\lambda=1$ in the notation there.
\end{enumerate}
\end{thm}

\Cref{thm:int-gamma} is formed by parts (i), (ii) and the first assertion of (iv) of \cref{thm:gam-main}, with \cref{thm:gam-ginf}. Part (a) is proved by a local Mellin inversion (\cref{thm:gam-lmi}) from the formal stationary-phase lemmas of \cref{sec:fsp} alone; by $c_n(L)=L^{4n}\mathrm{ph}^{\mathrm{SP}}_n(L^{-5})$ (\cref{prop:gam-hatflat}), Theorem~P (\cref{thm:con-Pb}) and the glue (\cref{thm:con-glue}) it is equivalent to $\Dn=1$ (\cref{cor:gam-equiv,cor:con-A2}). Part (b) says that $c(z;L)$ takes at $L=0$ the series displayed on the right of \cite[Lemma~7.7, (37)]{GJR18}, read as a series in $z$, which is obtained there for the Landau--Ginzburg series of that paper (we do not claim that the two series coincide); it identifies the free coefficients of the canonical colour-$0$ function (\cref{cor:gam-free}) and is not used for \cref{thm:int-main}.

\begin{rmk}
\label{rmk:int-constant}
The constant $c_g$ is that of \cite[Conjecture~1.5]{CGYZ26} for the quintic threefold, in the same normalisation $\varpi_1=u+O(u^2)$; the two theories have different potentials, so nothing about the quintic threefold follows (\cref{sec:sgn-lit}, \cref{rmk:mai-constant}). \draftnote{Numbers from arXiv:2609.13921v2, read in the HTML rendering; to be re-read in the PDF.}
\end{rmk}

% Section 1, part 3: inputs, what is not claimed, strategy, plan of the note.

%-------------------------------------------------------------------------------
\subsection{Inputs}
\label{sec:int-inputs}

The results are conditional on the following inputs. Statement numbers refer to the arXiv versions named in the bibliography.
\begin{description}
\item[(Q)] \emph{Quoted theorems.} Topological recursion, applied only to the model curve $x=\Wv(p)$, $y=p$ and its rescalings $x=\kappa\Wv$ (\cref{sec:mod}): \cite[Def.~4.2, (4-9), with the kernel (4-3); Theorems~4.2 and~4.6; Theorem~5.1, (5-14), with \S5.3, (5-27)--(5-30)]{Eynard:2007kz}; the theorem \cite[Theorem~4.1, with (4-1)--(4-3)]{DuninBarkowski:2012bw}, with its graph weights \cite[\S2.3, (2-22)--(2-26)]{DuninBarkowski:2012bw} (the dilaton leaves (2-23) and the factor $1/|\Aut|$ of (2-26)), as recalled in \cite[\S2.6, (2-30)--(2-32), (2-34), (2-35)]{Dunin-Barkowski:2015caa}, with matrices in the left-action convention of \cite[Remark~2.1]{Dunin-Barkowski:2015caa} (upper index the row), and with the action of \cite[\S2.3]{Dunin-Barkowski:2015caa} read with the factor $1/|\Aut\Gamma|$ and with the dilaton leaves of \cite[Def.~2.13]{PPZ15}, that is, as the action of \cite[Defs.~2.2, 2.13]{PPZ15}; it is applied to $x=4\Wv$, $y=p$, whose critical points $p^5=v$ have the pairwise distinct critical values $5p_c^4$. Its hypotheses, namely symplecticity, \cite[(2-32)]{Dunin-Barkowski:2015caa} and the edge condition \cite[(4-1)]{DuninBarkowski:2012bw}, are proved for this curve in \cref{prop:mod-Rtr}; the form \cite[(2-33)]{Dunin-Barkowski:2015caa} of the edge condition is misprinted (it has $+\sum_k$ in place of $\delta_{ij}-\sum_k$) and is not used, nor are \cite[Prop.~4.3]{Dunin-Barkowski:2015caa} and the sentence preceding \cite[Theorem~8.1]{Dunin-Barkowski:2015caa}; and \cite[Theorem~7.1: (7.6), (7.7) with its factor $1/\hbar^2$, (7.9), (7.10), with the kernel (1.1) and the assumptions of \S1.1]{ABDKS25}. The $R$-matrix action: \cite[\S0.2, (1); Def.~0.3, in particular (iii); (5); Remark~2.1 and (6); Def.~2.2, with the convention stated after it (p.~20); Def.~2.5, with the convention stated after it; Prop.~2.12, with the formula (8) in its proof; Def.~2.13; Prop.~2.14]{PPZ15}; Props.~2.12 and~2.14 are used as statements over $\bbQ$ and transferred to the other coefficient rings in \cref{def:set-graphsum}; (8) enters only through \cref{lem:set-pushins}.
\item[(QRR)] \emph{The $R$-matrix along the line: flatness and the value at the large-radius point} (\cref{thm:set-qrr}). The $R$-matrix $R$ of (T) has coefficients in $\End(\bbH)\otimes\bbQ(\zeta)\fps q$ and satisfies: (\emph{flatness}) for every colour $i$ the column $R\,\tilde e_i\,\re^{u_i/z}$ solves the quantum differential equation $z\partial_Tv=H*v$ along the line (\cref{def:set-rows}~(b)); (\emph{value}) $R(q=0)=\diag(h_{\bw_0},\dots,h_{\bw_4})$ in the normalised canonical basis at $q=0$. Flatness is a theorem about the $R$-matrix of the cohomological field theory, not a convention; the convention (Givental's) is only that the legs of the graph sum carry $R^{-1}(\psi)$. Flatness is used in \cref{prop:set-cyclicLR}~(iii) and in the last sentence of \cref{cor:set-cyclicPF5}~(b), on which \cref{lem:gau-colrows}, \cref{thm:gau-R0}~(iii), \cref{cor:gau-rowsNN}, \cref{thm:con-glue} and \cref{lem:mai-Psi} rest. The value is used for every colour, in the form $\Phi_i(z,1)=h_{\bw_i}(z)$ (\cref{eq:set-qrr}): in \cref{lem:gau-colrows} and, for colour $0$, in the glue \cref{thm:con-glue}. Sources for flatness: \cite[Lemma~5.2, (16)]{GJR18}, which is \cref{prop:set-cyclicLR}~(iii) at $\lambda=1$; \cite[\S4.2, (12)--(13)]{GJR17}; in general \cite[Section~1, Proposition (a)--(c)]{Giv01} and \cite[Props.~7.7, 8.2~(ii), 8.4, 8.7~(ii)]{MR2917177}. Sources for the value: \cite[(17), in the proof of Lemma~5.2]{GJR18}; \cite[(19)--(20), in the proof of Lemma~4.6, at $t=0$]{GJR17}, which state it for this theory; \cite[\S2.2, Theorem~1 and the Remark after it]{Giv01}; \cite[Theorem~1, Corollary~1]{MR2276766}, which give the tangent factor and the Euler-class twist separately.
\item[(M)] \emph{The genus-zero mirror theorem} \cite[Theorem~11.6]{MR1408320}\draftnote{numbers from arXiv v2; IMRN numbering not compared}, in the Calabi--Yau case $n=4$, $r=1$, $l_1=5$, $\lambda'=0$, with the set-up of \cite[\S9]{MR1408320} and \cite[Prop.~11.4]{MR1408320}: at $\lambda_i=\zeta^i$, $J(T(q)H,z)=I(q,z)/C_0(q)$ with $T=I_1/C_0$ (\cref{thm:set-mirror}).
\item[(G0)] \emph{The genus-zero axioms of the twisted theory}: the genus-zero twisted classes form a cohomological field theory with unit for $\eta$ \cite[\S4.1]{GJR18}, \cite[Def.~0.3]{PPZ15}, polynomial and $S_5$-symmetric in $\lambda$, with the string, divisor and splitting axioms, the genus-zero topological recursion relation and the quantum differential equation \cite[\S5.2, the display before (14), of which (14) is the form in the $q$-direction]{GJR18} (\cref{thm:set-G0}).
\item[(T)] \emph{Teleman's reconstruction, with the axioms of the twisted theory in all genera} (\cref{thm:set-GT}). (i) At $\lambda_i=\zeta^i$ the classes $\Omega_{g,n}$ of \cref{def:set-fq}, in all genera, form a cohomological field theory with flat unit $1$ for $\eta$ \cite[\S4.1]{GJR18}, \cite[Def.~0.3]{PPZ15}, and satisfy the divisor equation: for $2g-2+n>0$ the part of degree $d$ of $\pi_*\Omega_{g,n+1}(\gamma_1,\dots,\gamma_n,H)$ is $d$ times that of $\Omega_{g,n}(\gamma_1,\dots,\gamma_n)$, where $\pi$ forgets the last point; so the theory at $\Qn=\re^{T(q)}$ is the theory at the point $t=T(q)H$ with $\Qn=1$. (ii) This theory, which is semisimple (\cref{cor:set-cyclicPF5}~(a)), is the graph sum of \cref{def:set-graphsum}, $\Omega=R.\TFT$, for an $R$-matrix $R$ with coefficients in $\End(\bbH)\otimes\bbQ(\zeta)\fps q$, symplectic and with $R|_{z=0}=1$ \cite[Theorem~(2) in (1.15)]{MR2917177}, \cite[\S4.1, Theorem~4.1]{GJR18}; this is the $R$-matrix of \cref{def:set-rows}, to which (QRR) refers. In \cref{thm:int-main,thm:int-D} (T) identifies that $R$-matrix, whose rows define $T_\Gamma$ and $\Dn$; in \cref{cor:int-gap} it identifies the total with $F_g$.
\end{description}
In addition, the formula \cite[(1.7)]{AC96}, with $\kappa_0=2g-2+n$ \cite[p.~5]{AC96}, is used in part (i) of the dilaton lemma \cref{lem:set-dilaton}. The definitions use \cite[\S0.2, (5)]{LP18}, \cite[\S0.2, (5), (10)]{LP20} and \cite{MR1666787}; \cite{ZZ08} is cited for the series of \cref{lem:set-periods} only. The inputs are read with the following definitions and conventions of their sources: the normalisation of the $J$-function \cite[Section~6, (10)]{MR2276766} in (M) (\cref{thm:set-mirror}); the free energies \cite[Def.~4.3, (4-10)]{Eynard:2007kz} (\cref{conv:mod-kernel,cor:mod-trF}) and the set-up \cite[(5-1), the assumption before (5-4), (5-13)]{Eynard:2007kz} of \cite[Theorem~5.1]{Eynard:2007kz} (\cref{thm:mod-var}); the recursion \cite[(2-26)--(2-29)]{Dunin-Barkowski:2015caa}, to which the recollection of \cite[Theorem~4.1]{DuninBarkowski:2012bw} in \cite[\S2.6]{Dunin-Barkowski:2015caa} refers (\cref{conv:mod-kernel}); and the reading of the integrals \cite[(2-31), (2-32)]{Dunin-Barkowski:2015caa} as formal Laplace series (\cref{lem:mod-laplace,prop:mod-Rtr}): termwise Gaussian integration of the holomorphic part, as in \cite[(3-57)]{DuninBarkowski:2012bw}, and the value $\mathsf G_{-2}=-\vartheta$ for the polar term $w^{-2}\dd w$ (cf.\ \cite[(3-64)]{DuninBarkowski:2012bw}); and the edge condition \cite[(4-1)]{DuninBarkowski:2012bw} read in the coordinates $w_i$ of \cite[\S2.6]{Dunin-Barkowski:2015caa}, that is, as \eqref{eq:mod-edge} (\cref{prop:mod-Rtr}~(iii)).

\emph{Where the inputs enter.} (M) and (G0) enter only in \cref{sec:set} and its appendix \cref{sec:set-lrproof} (\cref{thm:set-mirror,prop:set-cyclicLR,cor:set-cyclicPF5}); (QRR) at the places named above: flatness in \cref{prop:set-cyclicLR}~(iii) and \cref{cor:set-cyclicPF5}~(b), the value in \cref{lem:gau-colrows} and \cref{thm:con-glue}; (T) in \cref{thm:set-GT} and \cref{cor:mai-gap}; topological recursion only in \cref{sec:mod}; the details are in \cref{sec:mai-thm}. \Cref{sec:gau} uses no quoted input beyond (M), (G0) and (QRR), which enter through the rows of \cref{sec:set} (flatness, \cref{def:set-rows}~(b), in \cref{thm:gau-R0}~(iii)--(iv)) and in \cref{lem:gau-colrows} (the value; hence in \cref{cor:gau-rowsNN}~(iii) and \cref{cor:gau-colours}); the transport lemmas of \cref{sec:tra} use \cite[Defs.~2.2, 2.13]{PPZ15} and the dilaton lemma, and Theorem~R adds no input; Theorem~P(a) (\cref{thm:con-Pa}) and part (a) of \cref{thm:int-gamma} use no quoted input.

\emph{Conditionality.} Theorem~A holds modulo the quoted theorems (Q) of \cite{Eynard:2007kz}, \cite{DuninBarkowski:2012bw} (as recalled in \cite{Dunin-Barkowski:2015caa}), \cite{ABDKS25} and \cite{PPZ15} with the statement numbers above, the flatness of the $R$-matrix and its value at the large-radius point (QRR), the genus-$0$ mirror theorem (M) and the genus-$0$ axioms (G0) of the twisted theory, Teleman's reconstruction with the axioms of the twisted theory in all genera (T), and the formula \cite[(1.7)]{AC96}; every other step is proved in the note.

In particular Theorem~U (over every field of characteristic $0$, \cref{thm:gau-U}) and every lemma of the chain are proved here; Stirling's series enters only through its formal proof as the stationary phase of Euler's integral (\cref{lem:fsp-stirling}), and the implicit function theorem only in its formal version (\cref{lem:fsp-ift}). \Cref{cor:int-gap} uses, beyond \cref{thm:int-main}, only (T), part of \cite{PPZ15} (listed after \cref{cor:mai-gap}) and proved statements. \Cref{thm:int-D} uses the $R$-matrix that (T) identifies only through $R|_{z=0}=1$, its coefficient ring $\bbQ(\zeta)\fps q$ and (QRR), that is, flatness (\cref{def:set-rows}~(b)) and the value at $q=0$, and the genus-zero data of \cref{sec:set}, which rest on (M) and (G0). \Cref{thm:int-gamma} uses no quoted input; the identity \cite[24.4.1]{DLMF} of Bernoulli polynomials used in part (b) is derived from their generating function in the proof of \cref{lem:fsp-dlmf5118}. No finite computation enters a proof; the computations that accompanied the work are summarised in \cref{sec:evi}.

All other citations are further sources, comparisons, conventions (\cref{sec:sgn}) or context, and enter no proof; among them \cite[Corollary~1.2]{Hock25} (\cref{rmk:mod-hock}), \cite[Theorem~4.4, (4-23); Corollary~4.1, (4-26)--(4-29); (A-1), (A-5), (A-25); Theorem~4.7, (4-30); (5-15)]{Eynard:2007kz}, \cite[(26)]{LP20}, \cite[Corollary~11.7]{MR1408320} (a misprinted form of (M)), \cite[Props.~6.1--6.2]{GR19}, \cite[\S2.4.3, (17); \S2.4.4, (19)--(20)]{LP20} and \cite{CGYZ26}.

%-------------------------------------------------------------------------------
\subsection{What is not claimed}
\label{sec:int-notclaimed}

\begin{enumerate}
\item Nothing is claimed about the quintic threefold or its graph sums; the formal quintic is a different theory. The conifold gap of the quintic threefold is \cite[Theorem~2]{CGYZ26}.
\item Nothing is claimed for frames other than those of \cref{def:set-frames} and the rescalings of their two periods (\cref{sec:sgn-lit}), for $g=1$, or away from $u=0$.
\item Nothing is claimed about the analytic continuation of the series $F_g(q)$ to a neighbourhood of $u=0$. The conifold expansion $\Cexp_{\alpha,\beta}$ is a formal change of frame on the ring $\Nring$ (\cref{def:set-confexp}), applied to the graph-sum representative $F_g^{\Nring}$ of $F_g$ (\cref{def:mai-Fcf}); this is the sense of ``conifold expansion'' and of ``conifold gap'' in \cref{cor:int-gap}.
\item The second key statement of \cite[Remark~1.9]{CGYZ26} is not addressed (the Model Lemma, \cref{def:red-ML}, concerns the model), and \cref{cor:int-gap} is not claimed to be literally the first, which is named there but not stated.
\item No novelty is claimed for the formal stationary-phase lemmas of \cref{sec:fsp}, proved for completeness, for the formula of $\Ginf$ \cite[Lemma~7.7, (37)]{GJR18}, or for the change of chart \cite[\S10]{MR1901075}.
\item Topological recursion is not applied to a mirror curve of the quintic or of the formal quintic, and \cite{FCCZ26} is not used.
\end{enumerate}

%-------------------------------------------------------------------------------
\subsection{Strategy}
\label{sec:int-strategy}

\emph{Reduction to the model (\cref{sec:gau,sec:tra,sec:red}).} In a conifold frame the rows satisfy a cyclic system with one simple pole at $u=0$ (\cref{prop:gau-quintic}), from \cref{prop:set-cyclicLR} and the cyclic form of $\PFfour$ (\cref{thm:gau-cyc}); the model pulled back by $v(u)$ has the same product of coefficients, which fixes $\tau$ (\cref{prop:gau-equalprod}). Theorem~U gives the regular gauge between them, and the gauge identity singles out the symplectic gauge $\Pc$ (\cref{thm:gau-G,prop:gau-symp}). Its edges lie in $\sKo\tensor\sKo$ for a slice $\sK$ defined by the vanishing solution of the model (\cref{prop:tra-can}), and Theorem~R (\cref{thm:red-main,cor:red-can}) reduces the regularity of the edge terms of the flat sum of $\Pc.\Omega^{\mathrm{mod}}_{v(u)}$, and the vertex law, to the Model Lemma on $\sK$ (\cref{def:red-ML}).

\emph{The model (\cref{sec:mod}).} The model cohomological field theory is topological recursion on its curve (\cref{thm:mod-tr}); insertions from $\sK$ become legs with pairings polynomial in $v$ (\cref{prop:mod-LD,thm:mod-V}), and the free energies are $c_hv^{2-2h}$ (\cref{thm:mod-gap}); this is the Model Lemma (\cref{thm:mod-MLK}).

\emph{The normalisation (\cref{sec:flh,sec:con,sec:gam}).} The comparison theorem (\cref{thm:flh-main}) identifies $T_\Gamma$ with the flat term, graph by graph, if $\Dn_0\equiv1\bmod z^{3g-2}$; Theorem~P (\cref{thm:con-Pa,thm:con-Pb}) and the glue (\cref{thm:con-glue}) turn $\Dn=1$ into part (a) of \cref{thm:int-gamma} (\cref{thm:gam-D}). \Cref{thm:mai-main} assembles the parts.

%-------------------------------------------------------------------------------
\subsection{Plan of the note}
\label{sec:int-plan}

\Cref{sec:set} defines the objects and states (M), (G0), (QRR) and (T); \cref{sec:gau} proves Theorem~U, the gauge identity and symplecticity; \cref{sec:tra} proves the transport lemmas and \cref{sec:red} Theorem~R; \cref{sec:mod} proves the Model Lemma; \cref{sec:flh} compares the flat and the outer sums; \cref{sec:con} proves Theorem~P and the glue, and \cref{sec:gam} \cref{thm:int-gamma,thm:int-D}. \Cref{sec:mai} proves \cref{thm:int-main,cor:int-gap}, and \cref{sec:cmp} records the natural factorisation and universality. The appendices prove the large-radius cyclic structure (\cref{sec:set-lrproof}), the cyclic form of $\PFfour$ (\cref{sec:gau-cycproof}), the $u$-adic convergence lemma (\cref{sec:tra-convproof}) and the lemmas of formal stationary phase (\cref{sec:fsp}), and record the sign conventions (\cref{sec:sgn}) and the finite computations (\cref{sec:evi}).


\subsection*{Acknowledgements}

This work was supported by the UKRI grant UKRI2394 and by the Leverhulme Trust grant LIP-2023-007.



% Section 2, part 1 of 8; label prefix: set. The main file inputs only this part; each part inputs the next.
% Part 1: introduction; the formal quintic, genus-zero data, PF5 = z^5 D PF4 - 1, the quoted inputs (M) and (G0).
% Part 2: the large-radius cyclic structure (statement; the proof of (i)-(ii) is in the appendix on the large-radius cyclic structure).
% Part 3: its corollary; points; Frobenius data.
% Part 4: the R-matrix and its rows, (QRR), the twist sign, the rotation of h_w; the graph sum, (T).
% Part 5: the Frobenius basis at the conifold point, the conifold frames, the conifold expansion, tau, the canonical coordinates.
% Part 6: the model, its rows and the T-system, M_0 and its symplecticity.
% Part 7: the model cohomological field theory, its homogeneity, the dilaton lemma.
% Part 8: the hatted theory, the outer graph sum; Property N(g) and the constants c_g.
% No new macros. Local letters are declared where they occur.
\section{The formal quintic, the model and the outer graph sum}
\label{sec:set}

This section fixes the objects of the note and states the quoted inputs (M), (G0) (\cref{sec:set-fq}), (QRR) (\cref{sec:set-rows}) and (T) (\cref{sec:set-graph}). Four facts used in the definitions are proved later, without using the definitions that cite them: \cref{lem:gau-confexp}, \cref{cor:gau-rowsNN}, \cref{lem:gau-colrows}, and the symplecticity of $R'$ (from \cref{thm:gam-D}, whose proof does not use $R'$). Printed numbers are values of the definitions; no finite computation enters a proof.

%-------------------------------------------------------------------------------
\subsection{The formal quintic and its genus-zero data}
\label{sec:set-fq}

Throughout, $\zeta\coloneqq\re^{2\pi\ri/5}$, colours $i,j$ and flat indices $a,b$ run over $\bbZ/5$, and $\overline{\mathcal M}_{g,n}$ is the moduli space of stable curves of genus $g$ with $n$ marked points.

\begin{defn}[The formal quintic]
\label{def:set-fq}
Let $\bbT=(\bbC^*)^5$ act on $\bbP^4$ with weights $\lambda_0,\dots,\lambda_4$ on $H^0(\bbP^4,\cO(1))$, with fixed points $p_0,\dots,p_4$, and let $H\coloneqq c_1^{\bbT}(\cO(1))$; so $H|_{p_i}=\lambda_i$, the tangent weights at $p_i$ are $\lambda_i-\lambda_j$ ($j\neq i$), and $\cO(5)|_{p_i}$ has the weight $5\lambda_i$ \cite[(10)]{LP20}. For $g,d\geq0$ the \emph{formal quintic invariant} is $N_{g,d}\coloneqq\int_{[\overline{\mathcal M}_g(\bbP^4,d)]^{\mathrm{vir}}}e\bigl(R\pi_*f^*\cO(5)\bigr)$ \cite[Section~0.2, (5)]{LP18}, \cite[Section~0.2, (5)]{LP20} (whenever the moduli space is non-empty), defined by virtual localisation \cite{MR1666787}, the Euler class being taken on each fixed locus as the quotient of the top Chern classes of the two terms of $R\pi_*f^*\cO(5)$. It is a rational function of the weights, homogeneous of degree $0$ (the rank $5d+1-g$ equals the virtual dimension), specialised at $\lambda_i=\zeta^i\lambda$, where it does not depend on $\lambda$; we put $\lambda=1$. (The second term, $\bbE^\vee\otimes\cO(5)|_{p_i}$ on the contracted components, has an invertible top Chern class there, as $5\zeta^i\neq0$. The other denominators of the localisation formula are products of the weights $(a\lambda_i+b\lambda_j)/d-\lambda_k$ ($a+b=d$, non-zero before the specialisation), $(\lambda_i-\lambda_j)/d_1+(\lambda_i-\lambda_k)/d_2$ and of classes $(\lambda_i-\lambda_j)/d-\psi$ ($j,k\neq i$); none vanishes at $\lambda_i=\zeta^i$, as a proper convex combination of two distinct fifth roots of unity has modulus $<1$.) This is the specialisation of \cite[Section~0.2]{LP18} and \cite[Section~0.2]{LP20}; see \cite[Section~3.1; Section~3.2, (9)]{GJR18}. The resulting hypergeometric series (\cref{lem:set-periods}) is studied in \cite{ZZ08}. For $g\geq2$ the \emph{genus-$g$ potential} is $F_g\coloneqq\sum_{d\geq0}N_{g,d}\,\Qn^d\in\bbQ\fps{\Qn}$, the term $d=0$ included; along the mirror map $\Qn=\re^{T(q)}\in q\bbQ\fps q$ (\cref{def:set-genus0}) it becomes a power series in $q$, again written $F_g$.

More generally, for $2g-2+n>0$ and $\gamma_1,\dots,\gamma_n\in H^*_{\bbT}(\bbP^4)$, with the stabilisation map $\rho$,
\[
\Omega_{g,n}(\gamma_1,\dots,\gamma_n)\coloneqq\sum_{d\geq0}\Qn^d\,\rho_*\Bigl(e\bigl(R\pi_*f^*\cO(5)\bigr)\prod_{j}\mathrm{ev}_j^*\gamma_j\cap[\overline{\mathcal M}_{g,n}(\bbP^4,d)]^{\mathrm{vir}}\Bigr)\in H^*(\overline{\mathcal M}_{g,n})\fps{\Qn},
\]
defined and specialised in the same way; with the \emph{twisted pairing} $\eta(a,b)\coloneqq\int_{\bbP^4}ab\,e_{\bbT}(\cO(5))$ these are the classes of the cohomological field theory of the formal quintic \cite[Section~4.1]{GJR18}, and $F_g=\int_{\overline{\mathcal M}_g}\Omega_{g,0}$. At $\lambda_i=\zeta^i$, $\prod_j(H-\lambda_j)=0$ becomes $H^5=1$, and the \emph{state space} is $\bbH\coloneqq\bbQ[H]/(H^5-1)$, with the flat basis $H^a$ ($a\in\bbZ/5$), $H^0=1$. Restriction to $p_i$ is $H\mapsto\zeta^i$, and an element of $\bbH\otimes\bbQ(\zeta)$ is determined by its five restrictions. As $e_{\bbT}(\cO(5))=5H$ and $e_{\bbT}(T_{p_i}\bbP^4)=\prod_{j\neq i}(\lambda_i-\lambda_j)=5\zeta^{4i}$, localisation gives
\[
\eta(H^a,H^b)=\sum_{i\in\bbZ/5}\frac{5\,\zeta^{i(a+b+1)}}{5\,\zeta^{4i}}=5\,\delta_{a+b\equiv3\bmod5},\qquad\eta^{cc'}=\tfrac15\,\delta_{c+c'\equiv3\bmod5},
\]
$(\eta^{cc'})$ the inverse matrix. The $N_{g,d}$ are rational: permuting the coordinates of $\bbC^5$ permutes the weights and preserves $H$ and $\cO(5)$ with their linearisations, so the rational function is symmetric, and its value at $\lambda_i=\zeta^i$ is fixed by the Galois group of $\bbQ(\zeta)$.
\end{defn}

\begin{defn}[Genus-zero data]
\label{def:set-genus0}
Let $\Dq\coloneqq q\,\dd/\dd q$, and let $\logq$ be an indeterminate with $\Dq\logq=1$, a formal $\log q$; $\Dq$ acts on $\bbQ\fps q[\logq]$, and $z$ is a constant for $\Dq$. Put
\[
C_0\coloneqq\sum_{d\geq0}\frac{(5d)!}{(d!)^5}q^d,\qquad\tilde I_1\coloneqq5\sum_{d\geq1}\frac{(5d)!}{(d!)^5}\Bigl(\sum_{m=d+1}^{5d}\frac1m\Bigr)q^d,\qquad I_1\coloneqq C_0\logq+\tilde I_1,
\]
the \emph{mirror map} $T\coloneqq I_1/C_0=\logq+\tilde I_1/C_0$, $C_1\coloneqq\Dq T$, and $\re^T\coloneqq q\exp(\tilde I_1/C_0)\in q\bbQ\fps q$; thus $C_0,C_1\in1+q\bbQ\fps q$. Put $L\coloneqq(1-5^5q)^{-1/5}$, so that $\Dq L=5^4qL^6=(L^6-L)/5$. The \emph{I-function} $I(q,z)$ is the $\bbH\otimes\bbQ(\zeta)$-valued series, $z$ times a series in $z^{-1}$ with coefficients in $\bbQ(\zeta)\fps q[\logq]$, with the restrictions
\[
I|_{p_i}=I_i(q,z)\coloneqq z\,q^{\zeta^i/z}\sum_{d\geq0}q^d\,\frac{\prod_{k=1}^{5d}(5\zeta^i+kz)}{\prod_{j\in\bbZ/5}\prod_{k=1}^{d}(\zeta^i-\zeta^j+kz)},\qquad q^{\zeta^i/z}\coloneqq\exp(\zeta^i\logq/z).
\]
The fifth-order and the fourth-order operators are
\begin{equation}
\label{eq:set-PF5}
\PFfive\coloneqq(z\Dq)^5-1-q\prod_{k=1}^5(5z\Dq+kz),
\end{equation}
\begin{equation}
\label{eq:set-PF4}
\PFfour\coloneqq\Dq^4-5q\prod_{k=1}^4(5\Dq+k)=(1-5^5q)\Dq^4-6250q\Dq^3-4375q\Dq^2-1250q\Dq-120q .
\end{equation}
With $u\coloneqq1-5^5q$ (\cref{def:set-points}), $\PFfour=u\Dq^4-(1-u)\bigl(2\Dq^3+\frac75\Dq^2+\frac25\Dq+\frac{24}{625}\bigr)$, the coefficients being the elementary symmetric functions of $\frac15,\dots,\frac45$. The eigenvalues of $H*$, the canonical coordinates and the normalisations $\Delta_i$ are consequences of (M) and (G0); they are stated in \cref{cor:set-cyclicPF5}~(a).
\end{defn}

\begin{lem}[$\PFfive=z^5\Dq\PFfour-1$]\genusstatus{AI-proposed, unchecked}
\label{lem:set-pf5}
As operators on any differential $\bbQ$-algebra that contains $q$ with $\Dq q=q$ and in which $z$ is a constant (for instance any differential $\bbQ\fps q$-algebra), $\PFfive=z^5\,\Dq\circ\PFfour-1$.
\end{lem}

\begin{proof}
As $\Dq\circ q=q\circ(\Dq+1)$, $\Dq\circ\PFfour=\Dq^5-q\prod_{k=1}^5(5\Dq+k)$; and as $z$ is a constant, $\PFfive=z^5\bigl(\Dq^5-q\prod_{k=1}^5(5\Dq+k)\bigr)-1$.
\end{proof}

\begin{thm}[The genus-zero mirror theorem (M); quoted]\genusstatus{AI-proposed, unchecked}
\label{thm:set-mirror}
For $d\geq0$ and $n\geq1$, with $n\geq3$ if $d=0$, let $E_{0,n,d}\coloneqq R^0\pi_*f^*\cO(5)$ on $\overline{\mathcal M}_{0,n}(\bbP^4,d)$, a vector bundle of rank $5d+1$ because $R^1\pi_*f^*\cO(5)=0$ in genus $0$, and let
\[
\langle\psi^{k_1}\gamma_1,\dots,\psi^{k_n}\gamma_n\rangle_{0,n,d}\coloneqq\int_{[\overline{\mathcal M}_{0,n}(\bbP^4,d)]}e_{\bbT}(E_{0,n,d})\prod_{j=1}^n\mathrm{ev}_j^*\gamma_j\,\psi_j^{k_j}
\]
be the genus-zero twisted invariants. The \emph{J-function}, in the normalisation of \cite[Section~6, (10)]{MR2276766} with the twisted pairing, is
\[
J(t,z)\coloneqq z+t+\sum_{n,d}\frac{\Qn^d}{n!}\sum_{a\in\bbZ/5}\Bigl\langle t,\dots,t,\frac{H^a}{z-\psi}\Bigr\rangle_{0,n+1,d}H^\vee_a,\qquad t\in H^*_{\bbT}(\bbP^4),
\]
where $H^\vee_a$ is the $\eta$-dual basis and the sum runs over the $(n,d)$ for which $\overline{\mathcal M}_{0,n+1}(\bbP^4,d)$ is defined. For $t=\tau H$ the divisor equation (\cref{thm:set-G0}~(b)) gives $J(\tau H,z)=\re^{\tau H/z}\bigl(z+\sum_{d\geq1}(\Qn\re^{\tau})^d\sum_a\langle H^a/(z-\psi)\rangle_{0,1,d}H^\vee_a\bigr)$. Since $T=\logq+\tilde I_1/C_0$ contains the formal logarithm, $J(T(q)H,z)$ at $\Qn=1$ is \emph{defined} by this formula, with $\re^{dT}\coloneqq(\re^T)^d$, $\re^T=q\exp(\tilde I_1/C_0)$ (\cref{def:set-genus0}), and $\re^{TH/z}\coloneqq\sum_{k\geq0}(TH)^k/(k!\,z^k)$:
\[
J(TH,z)\coloneqq\re^{TH/z}\Bigl(z+\sum_{d\geq1}\re^{dT}\sum_{a\in\bbZ/5}\Bigl\langle\frac{H^a}{z-\psi}\Bigr\rangle_{0,1,d}H^\vee_a\Bigr).
\]
Then at $\lambda_i=\zeta^i$ (see also (d) of the translation below),
\[
\text{(M)}\qquad J\bigl(T(q)H,z\bigr)=\frac{I(q,z)}{C_0(q)},
\]
an identity of $z$ times series in $z^{-1}$ with coefficients in $\bbH\otimes\bbQ(\zeta)\fps q[\logq]$.
\end{thm}

\emph{The pin.} (M) is \cite[Theorem~11.6]{MR1408320} (set-up in Section~9, pp.~23--27; the proof rests on Proposition~11.4) in the Calabi--Yau case $n=4$, $r=1$, $l_1=5$, $\lambda'=0$, at $\lambda_i=\zeta^i$. In this paragraph the letters are Givental's: $\bbC\bbP^n$ with the torus of $\bbC^{n+1}$ (weights $\lambda_\alpha$), the bundle $\oplus_a\cO(l_a)$ with fibre weights $\lambda'_a$, the space $Y_{2,d}$ of genus-zero stable maps of degree $d$ with markings $x_0,x_1$, the Euler class $E'_d$ of the bundle of sections vanishing at $x_0$, $c=-\psi_0$, and for $r=1$ (p.~27) $Z_i\coloneqq\sum_{d\geq0}q^d\int_{Y_{2,d}}e_0^*\bigl(\prod_{j\neq i}(p-\lambda_j)\bigr)(\hbar+c)^{-1}E'_d$, the term $d=0$ being $1$. With $\mathfrak f\coloneqq\sum_d\prod_a(l_ad)!/(d!)^{n+1}q^d$ and $\mathfrak g_l\coloneqq\sum_{d\geq1}q^d\prod_a(l_ad)!/(d!)^{n+1}\sum_{m=1}^{ld}\frac1m$, Theorem~11.6 states, for $\sum_al_a=n+1$, that
\[
Z_i^*\coloneqq\sum_{d\geq0}q^d\,\frac{\prod_a\prod_{m=1}^{l_ad}(l_a\lambda_i-\lambda'_a+m\hbar)}{\prod_\alpha\prod_{m=1}^d(\lambda_i-\lambda_\alpha+m\hbar)}
\]
coincides with $Z_i$ after (1) the substitution $Q=q\exp\{\sum_al_a(\mathfrak g_{l_a}-\mathfrak g_1)/\mathfrak f\}$, (2) the multiplication of $Z_i(Q(q),\hbar)$ by $\mathcal E_i\coloneqq\exp\{(\mathfrak f\hbar)^{-1}[\sum_a(l_a\lambda_i-\lambda'_a)\mathfrak g_{l_a}-(\sum_\alpha(\lambda_i-\lambda_\alpha))\mathfrak g_1]\}$, and (3) the multiplication by $\mathfrak f$. With $\mathfrak s_i^*(t)\coloneqq\re^{\lambda_it/\hbar}Z_i^*(\re^t,\hbar)$, $\mathfrak s_i(T)\coloneqq\re^{\lambda_iT/\hbar}Z_i(\re^T,\hbar)$ and $\re^T=Q(\re^t)$, this reads $\mathfrak s_i^*(t)=\mathfrak f\,\mathcal E_i\,\re^{-\lambda_i(T-t)/\hbar}\,\mathfrak s_i(T)$; as $T-t=\sum_al_a(\mathfrak g_{l_a}-\mathfrak g_1)/\mathfrak f$ and $\sum_al_a=n+1$, the terms in $\lambda_i$ cancel in the exponent, and $\mathfrak s_i^*(t)=\mathfrak f\exp\{(\mathfrak g_1\sum_\alpha\lambda_\alpha-\sum_a\lambda'_a\mathfrak g_{l_a})/(\hbar\mathfrak f)\}\,\mathfrak s_i(T)$. (\cite[Corollary~11.7]{MR1408320} prints $l_\alpha$ for $\lambda_\alpha$ here.)

\emph{Translation.} Put $n=4$, $r=1$, $l_1=5$, $\lambda'_1=0$, $p=H$, $\hbar=z$, $\re^t=q$, $t=\logq$. (a) $\mathfrak f=C_0$ and $5(\mathfrak g_5-\mathfrak g_1)=\tilde I_1$, so (1) is the mirror map of \cref{def:set-genus0}. (b) The factor above is $C_0\exp\{e_1(\lambda)\mathfrak g_1/(zC_0)\}$, $e_1(\lambda)=\sum_\alpha\lambda_\alpha$, which is $C_0$ at $\lambda_i=\zeta^i$. (c) $\mathfrak s_i^*(t)=I_i(q,z)/z$. (d) $\mathfrak s_i(T)=J(TH,z)|_{p_i}/z$: the $\eta$-dual of $\gamma\mapsto\int_{\overline{\mathcal M}_{0,1}(\bbP^4,d)}\beta\,\mathrm{ev}^*\gamma$ is $\mathrm{ev}_*(\beta)/e_{\bbT}(\cO(5))$, and $E_{0,1,d}$ is an extension of $\mathrm{ev}^*\cO(5)$ by the sections vanishing at the marking ($f^*\cO(5)$ is globally generated), so the definition of $J(TH,z)$ reads $J(TH,z)=\re^{TH/z}\bigl(z+\sum_{d\geq1}\re^{dT}\mathrm{ev}_*\bigl(E'_d/(z-\psi)\bigr)\bigr)$; on the other side $E'_d$ is pulled back along the map forgetting $x_1$, the string equation gives $\pi_*(\hbar-\psi_0)^{-1}=\hbar^{-1}(\hbar-\psi_0)^{-1}$ ($d\geq1$), and $\int_{\bbP^4}\prod_{j\neq i}(H-\lambda_j)\,\gamma=\gamma|_{p_i}$, so $Z_i(q,z)=1+z^{-1}\sum_{d\geq1}q^d\,\mathrm{ev}_*\bigl(E'_d/(z-\psi)\bigr)|_{p_i}$. By (a)--(d) the relation reads $C_0\,J(TH,z)|_{p_i}=I_i(q,z)$ for every $i$, which is (M).

\emph{Hypotheses.} The case is Calabi--Yau, $l_1=5=n+1$; the classes are polynomial in $\lambda$, $\lambda'$ and $\hbar$, so $\lambda'=0$ is a specialisation; and in genus zero $H^1(C,f^*\cO(5))=0$, so Givental's class $E_d$ is the Euler class of $E_{0,n,d}$ (convexity is all that is needed). After removing $\re^{\lambda_iT/z}$ both sides are, coefficientwise in $q$ and $z^{-1}$, polynomials in the weights, while $T$ and $C_0$ do not depend on them; so the specialisation $\lambda_i=\zeta^i$ is a ring homomorphism, and $e_1(\zeta^0,\dots,\zeta^4)=0$ gives (M).

\begin{thm}[The genus-zero axioms (G0); quoted]\genusstatus{AI-proposed, unchecked}
\label{thm:set-G0}
At $\lambda_i=\zeta^i$ the classes $\Omega_{0,n}$ of \cref{def:set-fq} form the genus-zero part of a cohomological field theory with unit $1$ for the pairing $\eta$ \cite[Section~4.1]{GJR18}, \cite[Def.~0.3]{PPZ15}, and the following hold.
\begin{itemize}
\item[(a)] For $d\geq1$ the invariant $\langle H^{a_1},\dots,H^{a_n}\rangle_{0,n,d}$, before the specialisation, lies in $\bbQ[\lambda_0,\dots,\lambda_4]$, is symmetric in the $\lambda_j$ and homogeneous of degree $\sum_ja_j-n$: the moduli space is proper of dimension $5d+n+1$, $E_{0,n,d}$ has rank $5d+1$, and a permutation of the coordinates of $\bbC^5$ permutes the weights and preserves $H$ and $\cO(5)$ with their linearisations.
\item[(b)] The string and divisor equations hold, also with descendant insertions at the other points (with the usual correction terms, from $\psi_j=\pi^*\psi_j+D_j$): $E_{0,n+1,d}$ is the pull-back of $E_{0,n,d}$ along the map $\pi$ forgetting a point, and $\int_dH=d$.
\item[(c)] The splitting axiom for $\eta$ holds (by the normalisation sequence along a boundary divisor, $e(E)$ restricted to the divisor times $e(\mathrm{ev}_{\mathrm{node}}^*\cO(5))$ is $e(E_1)e(E_2)$), hence the WDVV equations, and so does the genus-zero topological recursion relation.
\item[(d)] The quantum differential equation \cite[Section~5.2, the display before (14), of which (14) is the form in the $q$-direction]{GJR18}: let
\[
\langle\!\langle x,y,w\rangle\!\rangle(t)\coloneqq\sum_{n,d}\frac{\Qn^d}{n!}\langle x,y,w,t,\dots,t\rangle_{0,n+3,d},\qquad x*y\coloneqq\sum_{c,c'}\langle\!\langle x,y,H^c\rangle\!\rangle\,\eta^{cc'}H^{c'},
\]
the quantum product at $t$, $(\eta^{cc'})$ being the inverse matrix of $(\eta(H^c,H^{c'}))$, and let $J_a\coloneqq\partial J/\partial t^a$ be the derivative in the direction $H^a$. Then $z\,\partial_{t^b}J_a=\sum_c(H^b*H^a)^cJ_c$, $J_a=H^a+O(z^{-1})$, and $J_0=J/z$ (string equation). At $t=T(q)H$, $\Qn=1$ these series are understood as in \cref{thm:set-mirror}: by (b), $\langle\!\langle x,y,w\rangle\!\rangle(TH)\coloneqq\eta(xy,w)+\sum_{d\geq1}\re^{dT}\langle x,y,w\rangle_{0,3,d}$, and $J_a(TH,z)$ is $\partial J/\partial t^a$ at $t=\tau H$ with $(\Qn\re^{\tau})^d$ replaced by $\re^{dT}$ and $\re^{\tau H/z}$ by $\re^{TH/z}$. (d) follows from (b), (c) and the genus-zero topological recursion relation.
\end{itemize}
\end{thm}

% Section 2, part 2 of 8: the large-radius cyclic structure (statement).

Along the line $t=T(q)H$, with $\Qn=1$, functions of $t$ are functions of $q$: by the divisor equation (G0)(b) the insertions of $TH$ are resummed, so that the functions lie in $\bbQ[T]\fps{\re^T}$, and $T\mapsto T(q)$, $\re^T\mapsto q\exp(\tilde I_1/C_0)$ (\cref{def:set-genus0}) maps them to $\bbQ\fps q[\logq]$, intertwining $\partial_T$ with $C_1^{-1}\Dq$, and $\partial_T\coloneqq C_1^{-1}\Dq$ is the derivative in the coordinate $T$ of $H$, i.e.\ $\partial_{t^1}$ along the line (not the derivation $\de_t$ of \cref{def:gau-conc}). By (G0), the quantum product $*$ at $t=TH$ is commutative and associative, has the unit $1$, and is Frobenius for $\eta$: $\eta(x*y,w)=\langle\!\langle x,y,w\rangle\!\rangle=\eta(x,y*w)$.

\begin{prop}[The large-radius cyclic structure]\genusstatus{AI-proposed, unchecked}
\label{prop:set-cyclicLR}
Assume (M) and (G0), and for (iii) the flatness part of (QRR), \cref{def:set-rows}~(b). Along the line $t=T(q)H$, at $\lambda_i=\zeta^i$ and $\Qn=1$, the following hold.
\begin{enumerate}
\item In the flat basis, $H*H^a=\kappa^{\mathrm{LR}}_aH^{a+1}$ for $a\in\bbZ/5$, where
\[
\kappa^{\mathrm{LR}}=(\kappa^{\mathrm{LR}}_0,\dots,\kappa^{\mathrm{LR}}_4)=\Bigl(1,\ \frac{L^5}{C_0^2C_1^3},\ 1,\ \frac{C_0}{C_1},\ \frac{C_0}{C_1}\Bigr).
\]
\item Consequently, for every solution $v$ of the quantum differential equation $z\partial_Tv=H*v$, in any differential $\bbQ\fps q$-algebra in which $z$ is a constant, the lowered components $v_a\coloneqq\eta(H^a,v)$ satisfy $z\Dq v_a=C_1\kappa^{\mathrm{LR}}_a\,v_{a+1}$ for all $a\in\bbZ/5$.
\item For every colour $i$ the rows of \cref{def:set-rows} (defined in \cref{sec:set-rows}) satisfy $(z\Dq+\zeta^iL)\Qr^{(i)}_a=C_1\kappa^{\mathrm{LR}}_a\,\Qr^{(i)}_{a+1}$ for $a\in\bbZ/5$. Equivalently, for $X_a\coloneqq\zeta^{-i}\Qr^{(i)}_a$,
\[
X_0=L\Phi_i/C_0,\qquad X_{a+1}=(C_1\kappa^{\mathrm{LR}}_a)^{-1}(z\Dq+\zeta^iL)X_a\quad(a=0,\dots,3),\qquad(z\Dq+\zeta^iL)X_4=C_0X_0 .
\]
\end{enumerate}
\end{prop}

The superscript LR distinguishes $\kappa^{\mathrm{LR}}$ from the coefficients $\kappa_a$ of \cref{prop:gau-quintic}. Parts (i) and (ii) are proved in \cref{sec:set-lrproof}, (iii) after \cref{def:set-rows}. Below, $\Mcyc^{\mathrm{LR}}\coloneqq\partial_T^2\circ(\kappa^{\mathrm{LR}}_1)^{-1}\circ\partial_T^2\circ C_0^{-1}$, and $\Mcyc^{\mathrm{LR}}=\kappa^{\mathrm{LR}}_3\,\PFfour$ by \eqref{eq:set-McycLR}.

% Section 2, part 3 of 8: the corollary of prop:set-cyclicLR; points; Frobenius data.

\begin{cor}[Eigenvalues, idempotents and the fifth-order equation]\genusstatus{AI-proposed, unchecked}
\label{cor:set-cyclicPF5}
Assume (M) and (G0), and for the last sentence of (b) the flatness part of (QRR), \cref{def:set-rows}~(b).
\begin{itemize}
\item[(a)] $(H*)^5=(L/C_1)^5$. The eigenvalues of $H*$ are $\mu_i\coloneqq\zeta^iL/C_1$ ($i\in\bbZ/5$), labelled by $\mu_i(0)=\zeta^i$, the classical eigenvalue at $p_i$. The idempotents are
\[
e_i=\frac15\sum_{a=0}^4\Bigl(\prod_{b<a}\kappa^{\mathrm{LR}}_b\Bigr)\mu_i^{-a}H^a,\qquad\Delta_i\coloneqq\frac1{\eta(e_i,e_i)}=\frac{C_0^2\,\zeta^{3i}}{L^2}.
\]
The canonical coordinates along the line are functions $u_i$ with $\partial_Tu_i=\mu_i$, that is $\Dq u_i=\zeta^iL$; they are determined up to additive constants, the constants of $\Dq$ on $\bbQ(\zeta)\fps q[\logq]$ being $\bbQ(\zeta)$ (as in the proof of \cref{lem:set-periods}; for the conifold chart see \cref{lem:set-canon}).
\item[(b)] For every solution $v$ of $z\partial_Tv=H*v$ in a differential $\bbQ\fps q$-algebra in which $z$ is a constant, $C_0v_0$ solves $\PFfive$. In particular, for the column $v=R(z)\tilde e_i\re^{u_i/z}$ of \cref{def:set-rows}, the function $\re^{u_i/z}\,z\,L\,\Phi_i(\zh,u)$ solves $\PFfive$: this is the WKB form of colour $i$.
\end{itemize}
\end{cor}

\begin{proof}
(a) By \cref{prop:set-cyclicLR}~(i), $(H*)^5H^a=\prod_b\kappa^{\mathrm{LR}}_b\,H^a=(L/C_1)^5H^a$, so the characteristic polynomial of $H*$ is $x^5-(L/C_1)^5$, with the pairwise distinct roots $\zeta^iL/C_1$, as $L/C_1\in1+q\bbQ\fps q$. The equation $H*x=\mu_ix$ for $x=\sum_ax_aH^a$ reads $x_a=x_{a-1}\kappa^{\mathrm{LR}}_{a-1}/\mu_i$, consistently around the cycle as $\prod_b\kappa^{\mathrm{LR}}_b=\mu_i^5$. As $*$ is commutative and associative with unit $1$, the eigenlines are ideals spanned by idempotents $e_i$ with $\sum_ie_i=1$; for the coefficient $x_0^{(i)}$ of $H^0$ in $e_i$, the coefficient of $H^a$ in $\sum_ie_i=1$ is $\bigl(\prod_{b<a}\kappa^{\mathrm{LR}}_b\bigr)(C_1/L)^a\sum_ix_0^{(i)}\zeta^{-ia}=\delta_{a0}$, whose unique solution is $x_0^{(i)}=1/5$. By the Frobenius property and $\eta(1,H^a)=5\,\delta_{a3}$, $\eta(e_i,e_i)=\eta(1,e_i)=\kappa^{\mathrm{LR}}_0\kappa^{\mathrm{LR}}_1\kappa^{\mathrm{LR}}_2\,\mu_i^{-3}=L^2/(C_0^2\zeta^{3i})$. Finally $\partial_Tu_i=\mu_i$ and $\Dq=C_1\partial_T$ give $\Dq u_i=\zeta^iL$.

(b) By \cref{prop:set-cyclicLR}~(ii), $v_1=z\partial_Tv_0$, $v_2=z^2(\kappa^{\mathrm{LR}}_1)^{-1}\partial_T^2v_0$, $v_3=z^3\partial_T\bigl((\kappa^{\mathrm{LR}}_1)^{-1}\partial_T^2v_0\bigr)$ and, by \eqref{eq:set-McycLR},
\[
v_4=z^4(\kappa^{\mathrm{LR}}_3)^{-1}\Mcyc^{\mathrm{LR}}(C_0v_0)=z^4\,\PFfour(C_0v_0).
\]
The closing relation $z\partial_Tv_4=\kappa^{\mathrm{LR}}_4v_0$ becomes $z^5C_1^{-1}\Dq\PFfour(C_0v_0)=(C_0/C_1)v_0$, i.e.\ $\PFfive(C_0v_0)=0$ by \cref{lem:set-pf5}. The last statement concerns the rows and is proved after \cref{def:set-rows}.
\end{proof}

\begin{defn}[Points and variables]
\label{def:set-points}
Put $u\coloneqq1-5^5q$ and $s\coloneqq u^{1/5}$, so that $L=u^{-1/5}=s^{-1}$ and $\Dq=q\,\dd/\dd q=-(1-u)\,\dd/\dd u$, and put $\zh\coloneqq z/s^4=L^4z$, the hatted variable. The \emph{large-radius point} is $q=0$ ($u=1$, $L=s=1$, $\zh=z$), the \emph{conifold point} is $u=0$ ($q=5^{-5}$), and the \emph{orbifold point} is $q=\infty$ ($L=0$). Near $q=0$ the functions of the note lie in $\bbQ\fps q[\logq]$; near the conifold point the note uses their conifold expansions (\cref{def:set-confexp}), which are Laurent series in $s$, with $\Dq=-\frac{1-s^5}5s^{-4}\frac{\dd}{\dd s}$. The conifold expansion substitutes the periods of a conifold frame for $C_0$ and $I_1$ in a representative in the ring $\Nring$; it is a formal change of frame, and no analytic continuation is involved.
\end{defn}

\begin{defn}[Frobenius data and the topological field theory]
\label{def:set-frob}
Along the line $t=T(q)H$ let $e_i$ be the idempotents and $\Delta_i=1/\eta(e_i,e_i)$ (\cref{cor:set-cyclicPF5}~(a)), $\eta$ the pairing of \cref{def:set-fq}. The \emph{normalised canonical basis} is $\tilde e_i\coloneqq\Delta_i^{1/2}e_i$, so $\eta(\tilde e_i,\tilde e_j)=\delta_{ij}$, with the square roots $\zeta^{1/2}\coloneqq\zeta^3$ and $\Delta_i^{1/2}\coloneqq\zeta^{4i}\,C_0/L\in\bbQ(\zeta)\fps q$, so that $(\Delta_i^{1/2})^2=\zeta^{8i}C_0^2/L^2=\Delta_i$. At $q=0$, $\Delta_i(0)=W^0_i\coloneqq\prod_{j\neq i}(\lambda_i-\lambda_j)/(5\lambda_i)=\zeta^{3i}$. The \emph{topological field theory of the formal quintic} is the TFT $\TFT$ on $\bbH\otimes\bbQ(\zeta)\fps q$ with $\TFT_{g,n}(\tilde e_i,\dots,\tilde e_i)\coloneqq\Delta_i^{(2g-2+n)/2}=(\zeta^{4i}C_0/L)^{2g-2+n}$ and $\TFT_{g,n}$ vanishing on mixed colours; its unit is $1_{\TFT}\coloneqq\sum_i\Delta_i^{-1/2}\tilde e_i=\sum_ie_i=1$. We also regard $\TFT$ as a TFT on the colour space, with the orthonormal basis $\tilde e_i$ ($i\in\bbZ/5$) and the unit $1_{\TFT}$. By \cref{thm:set-GT}, since $R|_{z=0}=1$, $\TFT$ is the degree-zero part of the cohomological field theory along the line.
\end{defn}

\input{sections/s02_setup_4}



% Section 3, part 1 of 8; label prefix: gau. The main file inputs only this part; each part inputs the next.
% Part 1: introduction; the cyclic form of PF4 (objects, statements, the conifold expansion).
% Part 2: the quintic system and the pulled-back model; Theorem U (notation, statement).
% Part 3: the commutant, lifting, the scalar equation.
% Part 4: proof of Theorem U, the leading term, the Z/5-degree, graded gauges.
% Part 5: (R1), the diagonal, the vanishing solutions.
% Part 6: the gauge identity (G), the canonical gauge, symplecticity, HPF.
% Part 7: polynomiality, formal WKB uniqueness, colour covariance of the model columns and of the actual rows.
% Part 8: the normalisation series D, the rows in the ring, the colours of D, (F_g).
% The proofs of the cyclic form of PF4 are in the appendix file a2_cyclicPF4. No new macros.
\section{Cyclic systems, Theorem U and the canonical gauge}
\label{sec:gau}

We compare the cyclic systems of the formal quintic at the conifold point and of the model: both have one simple pole and equal products (\cref{sec:gau-cyc,sec:gau-systems}), their gauges form a free module of rank one (Theorem~U, \cref{sec:gau-U}), (G) singles out the symplectic canonical gauge $\Pc$ (\cref{sec:gau-van}), and $\Pc$ applied to the model columns is the matrix of rows times the normalisation series $\Dn$ (\cref{sec:gau-R0}).

\emph{Conventions.} $\Dq=q\,\dd/\dd q=-(1-u)\,\dd/\dd u$ (\cref{def:set-genus0,def:set-points}); for a series in $u$ a prime denotes $\dd/\dd u$; $z$ is a constant, and $z\de_u$ is $z\,\dd/\dd u$ at fixed $z$, so $z\Dq=-(1-u)z\de_u$. $E_{ab}$ ($a,b\in\bbZ/5$) are the matrix units and $\mathbf e_a$ the standard basis vectors; solutions, gauges and rows are in lowered components (\cref{def:set-rows}).

%-------------------------------------------------------------------------------
\subsection{The cyclic form of the fourth-order operator}
\label{sec:gau-cyc}

In this subsection $\PFfour=\sum_{k=0}^4A_k\Dq^k$ is written in the variable $u$ (\cref{def:set-genus0}):
\[
A_4=u,\qquad A_3=-2(1-u),\qquad A_2=-\tfrac75(1-u),\qquad A_1=-\tfrac25(1-u),\qquad A_0=-\tfrac{24}{625}(1-u),
\]
and we use $\Dq(1-u)=1-u$ and $\Dq u=-(1-u)$ throughout. The proofs of \cref{lem:gau-confexp,thm:gau-cyc,cor:gau-cycsys} are in \cref{sec:gau-cycproof}.

\begin{defn}[Objects of the cyclic-form theorem]
\label{def:gau-conc}
\begin{enumerate}
\item $\bbK:=\bbQ\lser{u}$ with the derivation $\Dq$, and $\bbK[\ell]$ is the polynomial ring in an indeterminate $\ell$ with the unique extension of $\Dq$ for which $\Dq\ell=1-u^{-1}$ (so $\ell$ plays the role of $\log u$). $\bbK[\Dq]$ is the ring of differential operators $\sum_ka_k\Dq^k$, $a_k\in\bbK$, with $\Dq\circ a-a\circ\Dq=\Dq(a)$; it acts on every differential ring that extends $(\bbK,\Dq)$.
\item $f_0,f_1,f_2\in\bbQ\fps u$ are the solutions of $\PFfour$ with $f_k=u^k+O(u^3)$, and $\varpi_1=f_1+\gamma f_2$, $\gamma=\frac7{10}$, is the vanishing period (\cref{lem:set-frobbasis}). $y_3:=\varpi_1\ell+\rho_3\in\bbK[\ell]$, $\rho_3\in u^3\bbQ\fps u$, is the logarithmic solution of \cref{lem:set-frobbasis}~(iii), whose indeterminate $\ell$ is the one above; thus $\PFfour y_3=0$ in $\bbK[\ell]$. Put $\mathcal S:=\operatorname{span}_\bbQ(f_0,f_1,f_2,y_3)\subset\bbK[\ell]$. An element of $\mathcal S$ lies in $\bbK$ if and only if its coefficient of $y_3$ vanishes, so $\mathcal S\cap\bbK=\operatorname{span}(f_0,f_1,f_2)$, the space of power-series solutions of $\PFfour$.
\item For $\cP=\sum_{k\leq4}a_k\Dq^k$ with coefficients $a_k$ in a differential ring $(\cR_0,\Dq)$, for instance $\cP\in\bbK[\Dq]$, the \emph{formal adjoint} is $\cP^*:=\sum_k(-\Dq)^k\circ a_k$, and the \emph{concomitant} of $\cP$ is the bilinear differential expression
\[
\Conc_\cP(f,g):=\sum_{k=1}^4\sum_{j=0}^{k-1}(-1)^j\,\Dq^j(a_kg)\,\Dq^{k-1-j}f .
\]
We write $\Conc:=\Conc_{\PFfour}$ and expand it as $\Conc(f,g)=\sum_{a,b=0}^3\Dq^af\,\Jet_{ab}\,\Dq^bg$ with $\Jet_{ab}\in\bbQ[u]\subset\bbK$; $\Jet=(\Jet_{ab})$ is the \emph{jet matrix}. As the $A_k$ lie in $\bbQ[u]$, $\PFfour$, $\PFfour^*$ and $\Conc$ act on every differential ring containing $(\bbQ[u],\Dq)$: on $\bbK[\ell]$, $\bbQ\fps q[\logq]$ ($u=1-5^5q$), $\Nring$ ($u=L^{-5}$) and $\bbQ\lser s$ (\cref{def:set-confexp,lem:gau-confexp,cor:gau-rowsNN}).
\item For a basis $\mathbf x=(x_0,\dots,x_3)$ of $\mathcal S$, the \emph{Gram matrix} is $\Conc_{\mathbf x}:=(\Conc(x_i,x_j))_{i,j=0}^3$ and the \emph{Wronskian matrix} is $\Wr_{\mathbf x}:=(\Dq^ax_i)_{a,i=0}^3$ (row $a$, column $i$), whose determinant is the Wronskian $\Wr(x_0,\dots,x_3)$ in the sense of \cref{lem:set-wronskian}. The basis $\mathbf f:=(f_0,f_1,f_2,y_3)$ is the standard one.
\item $\sh$ is the $\bbK$-algebra automorphism of $\bbK[\ell]$ with $\sh(\ell)=\ell+1$.
\item For $t\in\bbK$ with $\Dq t\neq0$, $\de_t:=(\Dq t)^{-1}\Dq$, a derivation of $\bbK[\ell]$. For a basis $\mathbf x$ of $\mathcal S$ (with $\Conc_{\mathbf x}$ invertible, \cref{lem:gau-conc}) the \emph{transversality sums} are $\Upsilon_k:=\sum_{i,j}(\Conc_{\mathbf x}^{-1})_{ij}\,x_i\,\Dq^kx_j$ and, for $y_0\in\bbK^\times$, $\widetilde\Upsilon_k:=\sum_{i,j}(\Conc_{\mathbf x}^{-1})_{ij}\,(x_i/y_0)\,\de_t^k(x_j/y_0)$ ($k\geq0$), elements of the fraction field of $\bbK[\ell]$.
\item For $y_0,y_1\in\bbK$ with $y_0\neq0$ and $t:=y_1/y_0$ not constant, the \emph{cyclic operator} is
\begin{equation}
\label{eq:gau-Mcyc}
\Mcyc(y_0,y_1):=\frac1{y_0}\circ\Dq\circ\frac1{\Dq t}\circ\Dq\circ\bigl(u\,y_0^2(\Dq t)^2\bigr)\circ\Dq\circ\frac1{\Dq t}\circ\Dq\circ\frac1{y_0}\ \in\bbK[\Dq].
\end{equation}
\end{enumerate}
\end{defn}

\begin{lem}[The conifold expansion is a ring map]\genusstatus{AI-proposed, unchecked}
\label{lem:gau-confexp}
Let $(\alpha,\beta)\in\bbQ^2$, $\varpi_0=f_0+\alpha\varpi_1+\beta f_2$ and $\tc=\varpi_1/\varpi_0$ (\cref{def:set-frames}), and extend $\Dq$ from $\bbK$ to $\bbQ\lser s$, $s^5=u$, by $\Dq s=-(1-s^5)/(5s^4)$. The assignment of \cref{def:set-confexp} on the generators of $\Nring$ extends to a unique ring homomorphism $\Cexp_{\alpha,\beta}\colon\Nring\to\bbQ\lser s$, and $\Cexp_{\alpha,\beta}\circ\Dq=\Dq\circ\Cexp_{\alpha,\beta}$. For the elements $C_0=c_0$, $I_1=c_0\logq+i_0$, $T=I_1/C_0$ and $C_1=\mathsf W/c_0^2$ of $\Nring$ (\cref{def:set-confexp}), which evaluate to the functions of \cref{def:set-genus0},
\[
\Cexp_{\alpha,\beta}(C_0)=\varpi_0,\qquad \Cexp_{\alpha,\beta}(I_1)=\varpi_1,\qquad \Cexp_{\alpha,\beta}(\mathsf W)=\varpi_0\Dq\varpi_1-\varpi_1\Dq\varpi_0=\varpi_0^2\,\Dq\tc ,
\]
so the conifold expansion sends $C_0$, $I_1$, $T$ and $C_1$ to $\varpi_0$, $\varpi_1$, $\tc$ and $\Dq\tc$.
\end{lem}

\begin{thm}[The cyclic form of $\PFfour$, (CYC)]\genusstatus{AI-proposed, unchecked}
\label{thm:gau-cyc}
Let $y_0\in\mathcal S\cap\bbK$ be non-zero and not proportional to $\varpi_1$, and let $\varpi_1$ be taken at any non-zero scale. Then, as operators in $\bbK[\Dq]$,
\begin{equation}
\label{eq:gau-cyc}
\PFfour=\Mcyc(y_0,\varpi_1)=\frac1{y_0}\circ\Dq\circ\frac1{\Dq t}\circ\Dq\circ\bigl(u\,y_0^2(\Dq t)^2\bigr)\circ\Dq\circ\frac1{\Dq t}\circ\Dq\circ\frac1{y_0},\qquad t=\frac{\varpi_1}{y_0}.
\end{equation}
More generally $\PFfour=\Mcyc(y_0,y_1)$ for every basis $(y_0,y_1)$ of an isotropic plane in $\mathcal S\cap\bbK$; these are the planes that contain $\varpi_1$ (\cref{lem:gau-iso}~(iii)). $\Mcyc(y_0,y_1)$ is unchanged under $y_0\mapsto\mu y_0$, $y_1\mapsto\mu'y_1$ ($\mu,\mu'\in\bbQ^\times$) and $y_1\mapsto y_1+ay_0$ ($a\in\bbQ$). In particular $\PFfour=\Mcyc(\varpi_0,\varpi_1)$ in every conifold frame of \cref{def:set-frames}.
\end{thm}

\begin{cor}[The cyclic system is $\PFfive$]\genusstatus{AI-proposed, unchecked}
\label{cor:gau-cycsys}
Let $(y_0,y_1)$ be as in \cref{thm:gau-cyc}, $t=y_1/y_0$, and
\[
\kappa=(\kappa_0,\dots,\kappa_4):=\Bigl(\Dq t,\ \frac{L^5}{y_0^2(\Dq t)^2},\ \Dq t,\ y_0,\ y_0\Bigr),\qquad L^5:=u^{-1},
\]
units of $\bbK$. Let $(\cR,\Dq)$ be a differential ring extending $(\bbK,\Dq)$ and $z\in\cR$ with $\Dq z=0$.
\begin{enumerate}
\item For $Y_0,\dots,Y_4\in\cR$ the cyclic system
\begin{equation}
\label{eq:gau-system}
z\Dq Y_a=\kappa_aY_{a+1}\qquad(a\in\bbZ/5)
\end{equation}
holds if and only if $y_0Y_0$ solves $\PFfive$ and $Y_{a+1}=z\Dq Y_a/\kappa_a$ for $a=0,\dots,3$.
\item If $z$ is invertible in $\cR$, then for every $w\in\cR$ the same holds with $z\Dq$ replaced by $z\Dq+w$ in \eqref{eq:gau-system}, in the relations $Y_{a+1}=z\Dq Y_a/\kappa_a$ and in $\PFfive$.
\end{enumerate}
\end{cor}

% Section 3, part 2 of 8: the quintic system and the pulled-back model; Theorem U (notation, statement).

%-------------------------------------------------------------------------------
\subsection{The quintic system and the pulled-back model}
\label{sec:gau-systems}

\begin{defn}[Cyclic systems with one simple pole; regular gauges]
\label{def:gau-onepole}
Let $\Bbbk$ be a field of characteristic $0$. A \emph{cyclic system with one simple pole} over $\Bbbk\fps u$ is a system $z\de_uY=AY$ with $A=\sum_{a\in\bbZ/5}b_aE_{a,a+1}$ whose coefficients have the \emph{one-pole shape}: $b_a\in\Bbbk\fps u^\times$ for $a\neq1$ and $ub_1\in\Bbbk\fps u^\times$. Its \emph{residue} is $r:=(ub_1)(0)\in\Bbbk^\times$ and its \emph{product} is $\Bprod:=\prod_ab_a\in\Bbbk\lser u$. It is \emph{self-dual} if $b_0=b_2$ and $b_3=b_4$, i.e.\ $\eta A=A^T\eta$ for $\eta=(5\,\delta_{a+b\equiv3})_{a,b}$ (the entries $(c,d)$ are $5b_{d-1}\delta_{c+d\equiv4}$ and $5b_{c-1}\delta_{c+d\equiv4}$). For a second such system $\tilde A=\sum_a\tilde b_aE_{a,a+1}$, with residue $\tilde r$, the condition (EP) of \emph{equal products} is $\prod_ab_a=\prod_a\tilde b_a$. The \emph{gauge equation} from $\tilde A$ to $A$ is $zP'=AP-P\tilde A$ for $P=\sum_{k\geq0}z^kP_k$, $P_k\in\mathfrak{gl}_5(\Bbbk\fps u)$; then $z\de_u(P\tilde Y)=AP\tilde Y$ whenever $z\de_u\tilde Y=\tilde A\tilde Y$. $\Gau(A,\tilde A)$ denotes the set of all solutions $P\in\mathfrak{gl}_5(\Bbbk\fps u)\fps z$; it is a $\Bbbk\fps z$-module, since $z(P\cn)'=(AP-P\tilde A)\cn$ for $\cn\in\Bbbk\fps z$. A \emph{regular gauge} is an element of $\Gau(A,\tilde A)$ with $P_0\in\mathrm{GL}_5(\Bbbk\fps u)$. $P$ is \emph{graded} if $P_k=\sum_bp_{k,b}E_{b+k,b}$ for every $k$.
\end{defn}

\begin{prop}[The quintic cyclic system]\genusstatus{AI-proposed, unchecked}
\label{prop:gau-quintic}
Let $(\alpha,\beta)\in\bbQ^2$ give a conifold frame (\cref{def:set-frames}), $\tc=\varpi_1/\varpi_0$, and
\[
\kappa:=\Bigl(\Dq\tc,\ \frac{L^5}{\varpi_0^2(\Dq\tc)^2},\ \Dq\tc,\ \varpi_0,\ \varpi_0\Bigr),\qquad L^5=u^{-1},
\]
the vector of \cref{cor:gau-cycsys} for $(y_0,y_1)=(\varpi_0,\varpi_1)$.
\begin{enumerate}
\item $\kappa_a=\Cexp_{\alpha,\beta}(C_1\kappa^{\mathrm{LR}}_a)$ for $a\in\bbZ/5$, where $C_1\kappa^{\mathrm{LR}}=(C_1,L^5/(C_0^2C_1^2),C_1,C_0,C_0)$ (\cref{prop:set-cyclicLR}) is represented in $\Nring$ by $(\mathsf W/c_0^2,\ L^5c_0^2/\mathsf W^2,\ \mathsf W/c_0^2,\ c_0,\ c_0)$. Hence the conifold expansion carries the forward relations $(z\Dq+\zeta^iL)X_a=C_1\kappa^{\mathrm{LR}}_aX_{a+1}$ ($a=0,\dots,3$) of elements of $\Nring\lser z\otimes\bbQ(\zeta)$ to the relations $(z\Dq+\zeta^iL)\Cexp(X_a)=\kappa_a\Cexp(X_{a+1})$, $L=s^{-1}$; by \cref{cor:gau-cycsys}~(ii) with $w=\zeta^is^{-1}$ the closing relation $(z\Dq+\zeta^iL)\Cexp(X_4)=\kappa_4\Cexp(X_0)$ is then equivalent to the equation $\PFfive$, with $z\Dq$ replaced by $z\Dq+\zeta^iL$, for $\varpi_0\Cexp(X_0)$; for the rows it is obtained in \cref{thm:gau-R0}~(iv) and \cref{cor:gau-colours}.
\item The cyclic system \eqref{eq:gau-system} is $z\de_uY=AY$ with $A=\sum_ab_aE_{a,a+1}$, $b_a:=-\kappa_a/(1-u)$,
\[
b=\Bigl(\tc',\ -\frac1{u\varpi_0^2(1-u)^3\tc'^2},\ \tc',\ -\frac{\varpi_0}{1-u},\ -\frac{\varpi_0}{1-u}\Bigr),
\]
a self-dual cyclic system with one simple pole over $\bbQ\fps u$, with residue $r=-1$, product $\prod_ab_a=-1/(u(1-u)^5)$, and $b_a(0)=1,1,-1,-1$ for $a=0,2,3,4$.
\item More generally, for $y_0\in\mathcal S\cap\bbK$ not proportional to $\varpi_1$ and $t=\varpi_1/y_0$, the coefficients $-\kappa_a/(1-u)$ of \cref{cor:gau-cycsys} have the one-pole shape if and only if $y_0(0)\neq0$.
\end{enumerate}
\end{prop}

\begin{proof}
(i) By \cref{lem:gau-confexp}, $\Cexp(c_0)=\varpi_0$, $\Cexp(\mathsf W/c_0^2)=\Dq\tc$ and $\Cexp(L^5)=u^{-1}$, which gives the five entries (for $\kappa_1$, $u^{-1}\varpi_0^2/(\varpi_0^4(\Dq\tc)^2)$). As $\Cexp$ is a ring homomorphism commuting with $\Dq$, with $\Cexp(L)=s^{-1}$, it carries each forward relation to the relation with $\kappa_a$; the statement on the closing relation is \cref{cor:gau-cycsys}~(ii) with $w=\zeta^is^{-1}$ and $(y_0,y_1)=(\varpi_0,\varpi_1)$. (ii) As $\Dq=-(1-u)\,\dd/\dd u$, $z\Dq Y_a=\kappa_aY_{a+1}$ is $z\de_uY_a=b_aY_{a+1}$, and $\Dq\tc=-(1-u)\tc'$, $L^5=u^{-1}$ give $b$. As $\varpi_0(0)=\tc'(0)=1$ (\cref{def:set-frames}), $b_0,b_2,b_3,b_4$ are units with the stated values and $(ub_1)(0)=-1$; self-duality is visible, and $\prod_ab_a=\tc'^2\cdot\bigl(-u^{-1}\varpi_0^{-2}(1-u)^{-3}\tc'^{-2}\bigr)\cdot\varpi_0^2(1-u)^{-2}$. (iii) If $y_0(0)\neq0$, then $t=u/y_0(0)+O(u^2)$, so $t'(0)\neq0$ and the computation of (ii) applies to $(y_0,t)$; if $y_0(0)=0$, $-y_0/(1-u)$ is not a unit.
\end{proof}

\begin{prop}[Equal products and $v(u)$]\genusstatus{AI-proposed, unchecked}
\label{prop:gau-equalprod}
Let $v\in u\bbQ\fps u$ with $v'(0)\neq0$. The model system $z\,\dd Y_b/\dd v=b^{\mathrm m}_bY_{b+1}$, $b^{\mathrm m}=(1,-v^{-1},1,-1,-1)$ (\cref{prop:set-modrows}), pulled back along $v=v(u)$ is $z\de_uY=\tilde AY$ with $\tilde A=v'(u)\sum_bb^{\mathrm m}_b(v(u))E_{b,b+1}$, i.e.
\[
\tilde b=\bigl(v',\ -v'/v,\ v',\ -v',\ -v'\bigr),\qquad \prod_a\tilde b_a=-v'^5/v .
\]
\begin{enumerate}
\item If $v=u+O(u^2)$, $\tilde A$ is a self-dual cyclic system with one simple pole over $\bbQ\fps u$, with residue $\tilde r=-1$ and $\tilde b_a(0)=1,1,-1,-1$ for $a=0,2,3,4$.
\item For $v=u+O(u^2)$, the products of $\tilde A$ and of the quintic system of \cref{prop:gau-quintic} are equal if and only if $v'^5/v=1/(u(1-u)^5)$, if and only if $v=u\,\tau(u)^{-5/4}$ (\cref{def:set-tau}).
\item (Necessity of (EP).) If $A,\tilde A$ are cyclic systems with one simple pole over $\Bbbk\fps u$ and $\Gau(A,\tilde A)\neq0$, then $\prod_ab_a=\prod_a\tilde b_a$.
\end{enumerate}
In particular, for $v=v(u)$ and every conifold frame, the quintic system and the pulled-back model are two self-dual cyclic systems with one simple pole satisfying (EP), and $b_a(0)=\tilde b_a(0)$ for $a\neq1$, $r=\tilde r=-1$.
\end{prop}

\begin{proof}
The chain rule gives $\tilde A$, and $\prod_a\tilde b_a=v'^5\cdot(-v^{-1})\cdot(-1)^2$. (i) $v'$ is a unit with $v'(0)=1$, and $-uv'/v$ a unit with value $-1$ at $u=0$. (ii) The first equivalence is \cref{prop:gau-quintic}~(ii). Write $v=uw$, $w\in1+u\bbQ\fps u$; elements of $1+u\bbQ\fps u$ have unique $n$-th roots there (binomial series; $c^n\neq1$ for $c\in1+u\bbQ\fps u\smallsetminus\{1\}$). As $v'(0)=1$, $v'^5=w(1-u)^{-5}$ is equivalent to $v'=w^{1/5}(1-u)^{-1}$. In $\bbQ\lser s$, $s^5=u$, with $v^{4/5}:=s^4w^{4/5}$ and $\frac{\dd}{\dd u}v^{4/5}=\frac45s^{-1}w^{-1/5}v'$, this reads $\frac{\dd}{\dd u}v^{4/5}=\frac45s^{-1}(1-u)^{-1}$, which is also the derivative of $s^4/\tau=\sum_{n\geq0}\frac4{5n+4}u^{n+4/5}$ (\cref{def:set-tau}). As $\dd/\dd u$ has kernel $\bbQ$ on $\bbQ\lser s$ and $v^{4/5}-s^4/\tau\in s^4\bbQ\fps u$, the equation is equivalent to $w^{4/5}=\tau^{-1}$, i.e.\ $w=\tau^{-5/4}$. (iii) If $P\neq0$ and $P_n$ is its first non-zero coefficient, the coefficient of $z^n$ of the gauge equation is $AP_n=P_n\tilde A$, so $A^5P_n=P_n\tilde A^5$; as $A^5=(\prod_ab_a)1$ and $\tilde A^5=(\prod_a\tilde b_a)1$, (EP) follows. The last sentence combines (i), (ii) and \cref{prop:gau-quintic}~(ii).
\end{proof}

%-------------------------------------------------------------------------------
\subsection{Theorem U}
\label{sec:gau-U}

Throughout this subsection $\Bbbk$ is a field of characteristic $0$, $A=\sum_ab_aE_{a,a+1}$ and $\tilde A=\sum_a\tilde b_aE_{a,a+1}$ are cyclic systems with one simple pole over $\Bbbk\fps u$ satisfying (EP), and a prime is $\dd/\dd u$. The following notation is local to this subsection.
\begin{itemize}
\item $\mathfrak{gl}^{(m)}:=\bigoplus_a\Bbbk\lser uE_{a,a+m}$ for $m\in\bbZ/5$; an element of $\mathfrak{gl}^{(m)}$ has \emph{shift} $m$, and $\mathfrak{gl}_5(\Bbbk\lser u)=\bigoplus_m\mathfrak{gl}^{(m)}$.
\item $\Bprod^{(m)}_a:=\tilde b_a\tilde b_{a+1}\cdots\tilde b_{a+m-1}$ for $0\leq m\leq5$ (so $\Bprod^{(0)}_a=1$ and $\Bprod^{(5)}_a=\Bprod:=\prod_a\tilde b_a$), and $\mathcal Y:=\tilde A^{-1}$.
\item $\Theta_U:=\diag(d^U_0,\dots,d^U_4)$ with $d^U_0:=1$ and $d^U_{a+1}:=d^U_a\tilde b_a/b_a$; $\operatorname{dlog}\Theta_U:=\Theta_U^{-1}\Theta_U'$; $\mathfrak d:=(\det\Theta_U/\det\Theta_U(0))^{-1/5}$.
\item If $\Bbbk$ contains a primitive fifth root of unity $\zeta$, put $\mathrm{gr}:=\diag(0,1,2,3,4)$, $\zeta^{i(\mathrm{gr}+1)}:=\diag(\zeta^{i(a+1)})_a$, and define the \emph{colour rotation} on matrices and on vectors by
\[
\operatorname{Rot}_iP(u,z):=\zeta^{i(\mathrm{gr}+1)}P(u,\zeta^{-i}z)\zeta^{-i(\mathrm{gr}+1)},\qquad \operatorname{Rot}_iY(u,z):=\zeta^{i(\mathrm{gr}+1)}Y(u,\zeta^{-i}z).
\]
\end{itemize}
By (EP), $d^U_4\tilde b_4/b_4=\prod_a\tilde b_a/b_a=1=d^U_0$, so the recursion of $d^U$ holds for every $a\in\bbZ/5$. The $d^U_a$ are units of $\Bbbk\fps u$ ($\tilde b_1/b_1=(u\tilde b_1)/(ub_1)$), so $\mathfrak d\in1+u\Bbbk\fps u$ is a binomial series and $\mathfrak d'=-\frac15\tr(\operatorname{dlog}\Theta_U)\,\mathfrak d$; and $(A\Theta_U)_{a,a+1}=b_ad^U_{a+1}=d^U_a\tilde b_a=(\Theta_U\tilde A)_{a,a+1}$, so $A\Theta_U=\Theta_U\tilde A$.

\begin{thm}[Theorem U]\genusstatus{AI-proposed, unchecked}
\label{thm:gau-U}
Let $A$, $\tilde A$ be cyclic systems with one simple pole over $\Bbbk\fps u$ with equal products.
\begin{enumerate}
\item[(a)] (Existence.) There is $P^\natural\in\Gau(A,\tilde A)$ with $P^\natural_0=\mathfrak d\,\Theta_U$, a diagonal matrix with unit entries.
\item[(b)] (Uniqueness.) For every $P^\natural\in\Gau(A,\tilde A)$ with $P^\natural_0=\mathfrak d\,\Theta_U$, the map $\Bbbk\fps z\to\Gau(A,\tilde A)$, $\cn\mapsto P^\natural\cn(z)$, is bijective. Consequently every solution $P$ of the gauge equation with coefficients in $\mathfrak{gl}_5(\Bbbk\fps u)$ is $P^\natural\cn(z)$ for a unique $\cn\in\Bbbk\fps z$; its $z^0$-term is $P_0=\cn(0)\,\mathfrak d\,\Theta_U$, a diagonal matrix, and $P$ is a regular gauge if and only if $\cn(0)\neq0$. If $P$ is one regular gauge, the regular gauges are exactly the $P\cn(z)$ with $\cn\in\Bbbk\fps z^\times$. For the pair $(\tilde A,\tilde A)$: $\Gau(\tilde A,\tilde A)=\Bbbk\fps z\cdot1$.
\item[(c)] (Graded gauges.) Graded regular gauges exist; if $P$ is one, the graded regular gauges are exactly the $P\cn(z)$ with $\cn\in\Bbbk\fps{z^5}^\times$.
\item[(d)] (The field.) If all $b_a$, $\tilde b_a$ lie in $\bbQ\lser u$, then $P^\natural$ can be chosen in $\mathfrak{gl}_5(\bbQ\fps u)\fps z$; for every field $\Bbbk\supseteq\bbQ$ the solutions with coefficients in $\mathfrak{gl}_5(\Bbbk\fps u)$ are then $P^\natural\Bbbk\fps z$, and $P^\natural\cn$ has rational coefficients if and only if $\cn\in\bbQ\fps z$.
\end{enumerate}
\end{thm}

% Section 3, part 3 of 8: the commutant, lifting, the scalar equation.

\begin{lem}[The commutant]\genusstatus{AI-proposed, unchecked}
\label{lem:gau-comm}
Let $F$ be a field of characteristic $0$, $\tilde b_0,\dots,\tilde b_4\in F^\times$, $\tilde A=\sum_a\tilde b_aE_{a,a+1}$, and define $\mathfrak{gl}^{(m)}$, $\Bprod^{(m)}_a$, $\Bprod$ and $\mathcal Y$ over $F$ as above.
\begin{enumerate}
\item[(a)] $\tilde A^m=\sum_a\Bprod^{(m)}_aE_{a,a+m}$ for $0\leq m\leq5$, so $\tilde A^5=\Bprod\,1$; $\mathcal Y=\sum_a\tilde b_a^{-1}E_{a+1,a}$, $\mathcal Y^j=\sum_a(\Bprod^{(j)}_a)^{-1}E_{a+j,a}$ for $0\leq j\leq5$, and $\tilde A^m=\Bprod\,\mathcal Y^{5-m}$ for $1\leq m\leq5$; $\det(x1-\tilde A)=x^5-\Bprod$.
\item[(b)] For $X=\sum_a\xi_aE_{a,a+m}\in\mathfrak{gl}^{(m)}$, $[\tilde A,X]=\sum_a(\tilde b_a\xi_{a+1}-\xi_a\tilde b_{a+m})E_{a,a+m+1}\in\mathfrak{gl}^{(m+1)}$.
\item[(c)] For every $m$, $\ker(\operatorname{ad}\tilde A)\cap\mathfrak{gl}^{(m)}=F\tilde A^m$; hence $\ker(\operatorname{ad}\tilde A)=F[\tilde A]=F[\mathcal Y]=\bigoplus_{j=0}^4F\mathcal Y^j$, of dimension $5$.
\item[(d)] For $m\in\{1,\dots,5\}$ and $X=\sum_a\xi_aE_{a,a+m}\in\mathfrak{gl}^{(m)}$ put $\lambda_m(X):=\sum_a\xi_a/\Bprod^{(m)}_a$ (a local letter). Then $\lambda_m(\tilde A^m)=5$, $\operatorname{ad}\tilde A(\mathfrak{gl}^{(m-1)})=\ker\lambda_m$, $\mathfrak{gl}^{(m)}=F\tilde A^m\oplus\ker\lambda_m$, and $\mathfrak{gl}_5(F)=F[\mathcal Y]\oplus\operatorname{im}(\operatorname{ad}\tilde A)$.
\item[(e)] The projection $\operatorname{pr}$ onto $F[\mathcal Y]$ along $\operatorname{im}(\operatorname{ad}\tilde A)$ is $F$-linear and preserves the shift; for $X\in\mathfrak{gl}^{(m)}$, $1\leq m\leq5$,
\[
\operatorname{pr}(X)=\tfrac15\lambda_m(X)\,\tilde A^m=\tfrac15\Bigl[\sum_a\xi_a\Bprod^{(5-m)}_{a+m}\Bigr]\mathcal Y^{5-m}.
\]
Equivalently, $\operatorname{pr}(X)$ is the unique element of $F[\mathcal Y]$ with $\tr(\operatorname{pr}(X)Z)=\tr(XZ)$ for all $Z\in F[\mathcal Y]$.
\item[(f)] Let $F=\Bbbk\lser u$ and let $\tilde A$ have the one-pole shape. Then $\mathcal Y\in\mathfrak{gl}_5(\Bbbk\fps u)$; an element $\sum_{j=0}^4w_j\mathcal Y^j$ ($w_j\in\Bbbk\lser u$) lies in $\mathfrak{gl}_5(\Bbbk\fps u)$ if and only if every $w_j$ lies in $\Bbbk\fps u$; the commutant of $\tilde A$ in $\mathfrak{gl}_5(\Bbbk\fps u)$ is the ring $\Bbbk\fps u[\mathcal Y]:=\bigoplus_{j=0}^4\Bbbk\fps u\,\mathcal Y^j$; and for $X\in\mathfrak{gl}_5(\Bbbk\fps u)$ the coefficient of $\mathcal Y^j$ in $\operatorname{pr}(X)$ lies in $u^{-1}\Bbbk\fps u$ for $j=1,\dots,4$ and in $\Bbbk\fps u$ for $j=0$.
\item[(g)] (Not used.) $x^5-\Bprod$ is separable, so $\tilde A$ is regular semisimple; in the situation of (f), $x^5-\Bprod$ is irreducible over $\Bbbk\lser u$, so $F[\mathcal Y]$ is a field.
\end{enumerate}
\end{lem}

\begin{proof}
(a) Induction on $m$, with $\Bprod^{(m)}_a\tilde b_{a+m}=\Bprod^{(m+1)}_a$ and $E_{a,a+5}=E_{aa}$; $\tilde A\sum_c\tilde b_c^{-1}E_{c+1,c}=1$ gives $\mathcal Y$, and $\mathcal Y^j$ follows likewise. Reindexing, $\mathcal Y^{5-m}=\sum_a(\Bprod^{(5-m)}_{a+m})^{-1}E_{a,a+m}$, and $\Bprod^{(m)}_a\Bprod^{(5-m)}_{a+m}=\Bprod$. In the Leibniz expansion of $\det(x1-\tilde A)$ only the identity and the even $5$-cycle $a\mapsto a+1$ contribute ($\pi(a)=a+1$ forces $\pi(a+1)=a+2$, and so on), the latter with $\prod_a(-\tilde b_a)=-\Bprod$.
(b) $\tilde AX=\sum_a\tilde b_a\xi_{a+1}E_{a,a+1+m}$ and $X\tilde A=\sum_a\xi_a\tilde b_{a+m}E_{a,a+m+1}$.
(c) By (b) $\operatorname{ad}\tilde A$ shifts the (direct) grading by $1$, so it suffices to treat $X\in\mathfrak{gl}^{(m)}$: $[\tilde A,X]=0$ means $\xi_{a+1}=\xi_a\tilde b_{a+m}/\tilde b_a$, whose solutions are determined by $\xi_0$, and $\xi_a=\Bprod^{(m)}_a$ is one. With $\tilde A^m=\Bprod\mathcal Y^{5-m}$ and $\mathcal Y=\Bprod^{-1}\tilde A^4$, $F[\tilde A]=F[\mathcal Y]$.
(d) For $Z\in\ker(\operatorname{ad}\tilde A)$, $\tr(Z[\tilde A,Y])=0$. For $X\in\mathfrak{gl}^{(m)}$, (a) gives $\tr(X\tilde A^{5-m})=\sum_a\xi_a\Bprod^{(5-m)}_{a+m}=\Bprod\lambda_m(X)$, so $\lambda_m$ vanishes on $\operatorname{ad}\tilde A(\mathfrak{gl}^{(m-1)})$, which has dimension $4$ by (c); as $\lambda_m(\tilde A^m)=5\neq0$, this image is $\ker\lambda_m$. Summing over $m$ gives the last assertion.
(e) $X=\frac15\lambda_m(X)\tilde A^m+\bigl(X-\frac15\lambda_m(X)\tilde A^m\bigr)$ with the second term in $\ker\lambda_m$; the second form uses (a). $X-\operatorname{pr}(X)\in\operatorname{im}(\operatorname{ad}\tilde A)$ is trace-orthogonal to $F[\mathcal Y]$ by (d), and the trace form is non-degenerate on $F[\tilde A]$: in the basis $1,\tilde A,\dots,\tilde A^4$ its Gram matrix has $5$ at $(0,0)$, $5\Bprod$ at $(i,5-i)$ and $0$ elsewhere ($\tilde A^n$ has shift $n$).
(f) The entries of $\mathcal Y$ are the $\tilde b_a^{-1}$, integral since $\tilde b_1^{-1}=u(u\tilde b_1)^{-1}$. The supports $\{(a+j,a)\}$ of $\mathcal Y^0,\dots,\mathcal Y^4$ partition the $25$ positions, and the entry $(a+j,a)$ of $\sum_jw_j\mathcal Y^j$ is $w_j(\Bprod^{(j)}_a)^{-1}$. For $j\leq4$ the window $\{a,\dots,a+j-1\}$ contains the pole index $1$ at most once, so $(\Bprod^{(j)}_a)^{-1}$ is integral, and a unit for $a=2$; this gives the criterion. $\Bbbk\fps u[\mathcal Y]$ is a ring since $\mathcal Y^5=\Bprod^{-1}1$, $\Bprod^{-1}\in u\Bbbk\fps u$, and by (c) it is the integral commutant. For integral $X\in\mathfrak{gl}^{(m)}$, $1\leq m\leq4$, the coefficient $\frac15\sum_a\xi_a\Bprod^{(5-m)}_{a+m}$ involves at most four consecutive weights, at most one of them $\tilde b_1$, so it lies in $u^{-1}\Bbbk\fps u$; for $m=5$ it is $\frac15\tr X$.
(g) $\gcd(x^5-\Bprod,5x^4)=1$. In the situation of (f), with $\Bprod^\circ:=u\Bprod$ (a unit) and $x=1/X$, $x^5-\Bprod=-\Bprod^\circ u^{-1}X^{-5}\bigl(X^5-u(\Bprod^\circ)^{-1}\bigr)$, an Eisenstein polynomial at $u$.
\end{proof}

\begin{lem}[Lifting]\genusstatus{AI-proposed, unchecked}
\label{lem:gau-lift}
Let $\tilde A$ have the one-pole shape, let $X\in\mathfrak{gl}_5(\Bbbk\fps u)$ have shift $m+1$ with $0\leq m\leq4$, and suppose $\operatorname{pr}(X)=0$. Then there is $\mathcal Z\in\mathfrak{gl}_5(\Bbbk\fps u)$ of shift $m$ with $[\tilde A,\mathcal Z]=X$, and every solution $\mathcal Z'\in\mathfrak{gl}_5(\Bbbk\fps u)$ of $[\tilde A,\mathcal Z']=X$ is $\mathcal Z+w$ with $w\in\Bbbk\fps u[\mathcal Y]$. If $\tilde A$ and $X$ have rational coefficients, $\mathcal Z$ can be chosen rational.
\end{lem}

\begin{proof}
Write $X=\sum_a\xi_aE_{a,a+m+1}$ and try $\mathcal Z=\sum_a\rho_a\Bprod^{(m)}_aE_{a,a+m}$ (local letters). By \cref{lem:gau-comm}~(b), $[\tilde A,\mathcal Z]=\sum_a(\rho_{a+1}-\rho_a)\Bprod^{(m+1)}_aE_{a,a+m+1}$, so the equation is $\rho_{a+1}-\rho_a=\sigma_a:=\xi_a/\Bprod^{(m+1)}_a$, solvable if and only if $\sum_a\sigma_a=\lambda_{m+1}(X)=0$, i.e.\ $\operatorname{pr}(X)=0$ (\cref{lem:gau-comm}~(e)). Each $\sigma_a$ is integral, with $\sigma_a(0)=0$ when the window $\{a,\dots,a+m\}$ contains $1$. Put $\rho_1:=0$, $\rho_{a+1}:=\rho_a+\sigma_a$ ($a=1,\dots,4$); the relation for $a=0$ follows from $\sum_a\sigma_a=0$. The entry $\rho_a\Bprod^{(m)}_a$ is integral when the window $\{a,\dots,a+m-1\}$ avoids $1$, and otherwise ($a=2-m,\dots,1$) when $\rho_a(0)=0$: for $a=1$ by choice, and for $2-m\leq a\leq0$ because $\rho_a=-\sum_{c=a}^0\sigma_c$ with every window $\{c,\dots,c+m\}$ containing $1$. Two integral solutions differ by an element of $\Bbbk\fps u[\mathcal Y]$ (\cref{lem:gau-comm}~(c), (f)); the construction uses only field operations.
\end{proof}

\begin{lem}[The scalar equation]\genusstatus{AI-proposed, unchecked}
\label{lem:gau-scal}
Let $A$, $\tilde A$ have the one-pole shape and satisfy (EP).
\begin{enumerate}
\item[(a)] For $w=\sum_{j=0}^4w_j\mathcal Y^j$ with $w_j\in\Bbbk\lser u$,
\[
\operatorname{pr}\bigl(w'+\operatorname{dlog}\Theta_U\,w\bigr)=\sum_{j=0}^4\bigl(w_j'+\chi^U_jw_j\bigr)\mathcal Y^j,\qquad \chi^U_j:=-\frac j5\,\frac{\Bprod'}{\Bprod}+\frac15\tr(\operatorname{dlog}\Theta_U).
\]
In particular $\operatorname{pr}((\mathcal Y^j)')=-\frac j5(\Bprod'/\Bprod)\mathcal Y^j$ and $\operatorname{pr}(\operatorname{dlog}\Theta_U\,\mathcal Y^j)=\frac15\tr(\operatorname{dlog}\Theta_U)\,\mathcal Y^j$. Moreover $\chi^U_j=\frac j{5u}+\chi^\circ_j$ with $\chi^\circ_j:=-\frac j5(\Bprod^\circ)'/\Bprod^\circ+\frac15\tr(\operatorname{dlog}\Theta_U)\in\Bbbk\fps u$, where $\Bprod^\circ:=u\Bprod\in\Bbbk\fps u^\times$.
\item[(b)] For $Z\in\mathfrak{gl}_5(\Bbbk\fps u)$ the equation $\operatorname{pr}(w'+\operatorname{dlog}\Theta_U\,w)=-\operatorname{pr}(Z)$ has solutions $w\in\Bbbk\fps u[\mathcal Y]$; their components $w_1,\dots,w_4$ are unique, and $w_0$ is unique up to $w_0\mapsto w_0+\cn\,\mathfrak d$, $\cn\in\Bbbk$.
\item[(c)] For $Z=0$ the solutions in $\Bbbk\fps u[\mathcal Y]$ are exactly the $\cn\,\mathfrak d\,1$, $\cn\in\Bbbk$.
\item[(d)] If the data are rational, so are the solutions for $\cn\in\bbQ$; $\mathfrak d\in1+u\bbQ\fps u$.
\end{enumerate}
\end{lem}

\begin{proof}
(a) $(\mathcal Y^j)'$ has shift $5-j$, with entries $((\Bprod^{(j)}_a)^{-1})'$ at $(a+j,a)$; by \cref{lem:gau-comm}~(e),
\[
\operatorname{pr}((\mathcal Y^j)')=\frac15\Bigl[\sum_a\bigl((\Bprod^{(j)}_a)^{-1}\bigr)'\Bprod^{(j)}_a\Bigr]\mathcal Y^j=-\frac15\Bigl[\sum_a\frac{(\Bprod^{(j)}_a)'}{\Bprod^{(j)}_a}\Bigr]\mathcal Y^j ,
\]
and $\sum_a(\Bprod^{(j)}_a)'/\Bprod^{(j)}_a=\sum_a\sum_{i<j}\tilde b_{a+i}'/\tilde b_{a+i}=j\Bprod'/\Bprod$, each index occurring $j$ times. Likewise $\operatorname{dlog}\Theta_U\mathcal Y^j$ has entries $\frac{(d^U_{a+j})'}{d^U_{a+j}}(\Bprod^{(j)}_a)^{-1}$, and \cref{lem:gau-comm}~(e) gives $\frac15\tr(\operatorname{dlog}\Theta_U)\mathcal Y^j$. As $\operatorname{pr}$ is $\Bbbk\lser u$-linear and fixes $\Bbbk\lser u[\mathcal Y]$, the formula follows. Finally $\Bprod'/\Bprod=(\Bprod^\circ)'/\Bprod^\circ-u^{-1}$, and $\tr(\operatorname{dlog}\Theta_U)$ is integral since the $d^U_c$ are units.

(b) By \cref{lem:gau-comm}~(f), $\operatorname{pr}(Z)=\sum_j\phi^Z_j\mathcal Y^j$ with $u\phi^Z_j\in\Bbbk\fps u$ for $j\geq1$ and $\phi^Z_0\in\Bbbk\fps u$ (local letters), and by (a) the equation splits into $w_j'+\chi^U_jw_j=-\phi^Z_j$. For $j=1,\dots,4$, multiplied by $u$, the coefficient of $u^n$ reads
\[
\Bigl(n+\frac j5\Bigr)w_{j,n}=-\coef{u^n}(u\phi^Z_j)-\sum_{l=0}^{n-1}\coef{u^{n-1-l}}\chi^\circ_j\cdot w_{j,l},
\]
with $n+j/5\neq0$: exactly one solution $w_j\in\Bbbk\fps u$. For $j=0$, $\chi^U_0=-\mathfrak d'/\mathfrak d$, so the equation is $\mathfrak d\,(w_0/\mathfrak d)'=-\phi^Z_0$, with solutions $w_0=\mathfrak d\bigl(\cn-\int_0^u\phi^Z_0/\mathfrak d\,\dd u\bigr)$, $\cn\in\Bbbk$ ($\int_0^u$ the formal antiderivative without constant term), and no others. (c) is (b) with $\phi^Z=0$. (d) The pivots are rational and $\mathfrak d$ is a binomial series of an element of $1+u\bbQ\fps u$.
\end{proof}

\input{sections/s03_gauge_4}



% Section 4, part 1 of 4; label prefix: tra. Part 2: frames, the affine slice, bivectors.
% Part 3: pointer to the u-adic convergence lemma. Part 4: dressings, edges, Lemmas E, T, F, the canonical gauge.
% No new macros. Local letters (declared where used): phi (a map of finite type), chi (sequences), pi_r, N, Theta.

\section{The vanishing tower and transport}
\label{sec:tra}

This section fixes the two coefficient settings in which the model is used, defines the vanishing tower $\Pi_i$, the slice $\sK$ and its frames, and proves the transport lemmas: for every symplectic dressing $P$ of the model that annihilates the vanishing period, in particular for the canonical gauge $\Pc$ of \cref{def:gau-can}, the edges lie in $\sKo\otimes\sKo$, the leaf point lies in $\sK$, its dilaton rescaling lies in the affine slice $-z+\sK$, and every graph term is a finite sum. Every infinite sum is justified $u$-adically. In this section and the next, $q$ denotes an element of a loop space, as in Givental's formalism; the parameter $q$ of the formal quintic does not occur.

\subsection{The vanishing tower and the slice}
\label{sec:tra-tower}

\begin{conv}[loop spaces, the two settings, forms, insertions]
\label{conv:tra-forms}
The state space $\bbH$ has the flat basis $H^a$, $a\in\bbZ/5$, with unit $1 = H^0$ and $H = H^1$, and the pairing $\eta(H^a,H^b) = 5\,\delta_{a+b\equiv 3}$ of \cref{def:set-fq}; as a matrix, $\eta$ is $5$ times the permutation matrix of $a\mapsto 3-a$, so $\eta^{-1} = \eta/25$. The charges and signs are those of \cref{def:set-model}, $\nu = (\nu_0,\dots,\nu_4) = (0,-1,3,2,1)$ and $\varepsilon = (1,-1,-1,1,1)$; put $w_b := \nu_b+1$. The \emph{weight} of $v^mz^kH^b$ is $5m+4k+w_b$ and the weight of a scalar $v^mz^k$ ($k\in\bbZ$) is $5m+4k$; $V_w$ denotes the $\bbQ$-span of the $v^mz^kH^b$ ($m,k\geq 0$) of weight $w$, which is the weight space of \cref{def:mod-weights}. In terms of the exponent weights of \cref{lem:set-modhom}, $4k+w_b = 5\wt(k,b)+5$.
\begin{description}
\item[Graded setting] Coefficients in $\bbQ[v]$. Each $V_w$ is finite-dimensional and consists of polynomials in $z$; put $\cH^{\mathrm{gr}}_+ := \bigoplus_wV_w = \bbH\otimes\bbQ[v][z]$. Givental's loop space $\bbH\otimes\bbQ[v]\lser{z^{-1}}$, with positive half $\cH^{\mathrm{gr}}_+$, is consistent in this setting, and every sum below is finite.
\item[$u$-adic setting] $v$ is a non-zero element of $u\bbQ\fps{u}$, written $v(u)$ when its dependence on $u$ is stressed; from \cref{def:tra-dressing} on, $v'(0)\neq 0$; for the formal quintic it is the series $v(u) = u\tau(u)^{-5/4}$ of \cref{def:set-tau}. The loop space and its halves are
\[
\cH^u := \bbH\otimes\bbQ\lser{z}\fps{u},\qquad \cH^u_+ := \bbH\otimes\bbQ\fps{u,z},\qquad \cH^u_- := \bbH\otimes z^{-1}\bbQ[z^{-1}]\fps{u}.
\]
Here $\bbQ\lser{z}\fps{u}$ is the power-series ring in $u$ over the field $\bbQ\lser{z}$: its elements are the series $\sum_{k\in\bbZ}c_kz^k$ with $c_k\in\bbQ\fps{u}$ and $c_k\to 0$ $u$-adically as $k\to-\infty$, and products are finite modulo every power of $u$. It contains the one-variable objects of this section: the leaf points and the entries of the dressings are power series in $z$, and the vanishing solution has infinitely many negative powers of $z$; the edges, in two variables, are treated in \cref{lem:tra-twoslot,lem:tra-E}. Every $f\in\cH^u$ splits uniquely, coefficientwise in $u$, as $f = f_{\geq 0}+f_{<0}$ with $f_{\geq 0}\in\cH^u_+$ and $f_{<0}\in\cH^u_-$; $f_{\geq k}$ is the part with $z$-powers $\geq k$. The form
\[
\Omega_{\mathrm s}(f,g) := \Res_z\,\eta(f(-z),g(z))\,\dd z,
\]
with the residue taken coefficientwise in $u$, is $\bbQ\fps{u}$-bilinear and antisymmetric on $\cH^u$ (the substitution $z\mapsto -z$ is an automorphism of $\bbQ\lser{z}$ which changes the sign of the residue, and $\eta$ is symmetric), and $\cH^u_+$, $\cH^u_-$ are isotropic and complementary. The substitution $v\mapsto v(u)$ maps every $V_w$ into $\cH^u_+$, and every scalar series $\sum_{n\geq 1}c_nv^{(w+4n)/5}z^{-n}$ of scalar weight $w$ ($c_n\in\bbQ$) into $z^{-1}\bbQ[z^{-1}]\fps{u}$, since its $z^{-n}$-coefficient has $u$-order at least $(w+4n)/5$.
\end{description}
\emph{Insertions and components.} An insertion $\psi^kH^b$ is identified with the vector $z^kH^b$. For $q$ in $\cH^{\mathrm{gr}}_+$ or $\cH^u_+$ we write $q_{(k,b)} := \coef{z^k}q^b$ for its flat components, so $q = \sum_{k,b}q_{(k,b)}z^kH^b$. Solutions are written in lowered components: a vector $J$ has the components $Y_a := \eta(H^a,J)$, i.e.\ $Y = \eta J$, and for $n = \sum_an^aH^a$ we put $\langle Y,n\rangle := \sum_aY_an^a = \eta(J,n)$. Matrices $P$ (gauges and dressings) act on lowered components; the \emph{vector form} of $P$ is $P^{\mathrm{vec}} := \eta^{-1}P\eta$, the conjugate of $P$ by the permutation $a\mapsto 3-a$, which acts on vectors. Since $\eta^{-1} = \eta/25$, the condition $P(-z)^T\eta P(z) = \eta$ holds for $P$ if and only if it holds for $P^{\mathrm{vec}}$. The vanishing solution $\YC$ of \cref{def:gau-van} has the vector form $J^{\mathrm C} := \eta^{-1}\YC$; thus $\eta(J^{\mathrm C}(-z),q(z)) = \langle\YC(-z),q(z)\rangle$.

\emph{Summable families; maps of finite type.} A family $(x_l)_{l\in L}$ in $\cH^u_+$ is \emph{summable} if for every flat index $(k,b)$ and every $N$ only finitely many $x_l$ have $\coef{z^k}x_l^b\notin u^N\bbQ\fps{u}$; its sum is taken coefficientwise and is finite modulo every power of $u$. The same definition is used for bivectors in $\bbH\otimes\bbH\otimes\bbQ\fps{\psi,\psi'}\fps{u}$, with the indices $(k,l,a,b)$ of the coefficients of $\psi^k\psi'^lH^a\otimes H^b$. A family $(y_l)$ in $\bbQ\lser{u}$ is summable if all $y_l$ lie in one $u^{N_0}\bbQ\fps{u}$ and for every $N$ only finitely many lie outside $u^N\bbQ\fps{u}$. Every regrouping of a summable family (a partition into subfamilies) consists of summable families; the family of their sums is summable (for a fixed coefficient index and $N$ the finitely many members outside $u^N\bbQ\fps{u}$ lie in finitely many subfamilies, and every other sub-sum lies in $u^N\bbQ\fps{u}$); and the sum of their sums is the total sum, because coefficientwise and modulo $u^N$ this is a statement about finite sums; the sum of a summable family of elements of $\bbQ\fps{u}$ lies in $\bbQ\fps{u}$. A $\bbQ\fps{u}$-linear map $\phi\colon\cH^u_+\to\bbQ\lser{u}$ (a local letter) is \emph{of finite type} if $\phi(\cH^u_+)$ lies in some $u^{N_0}\bbQ\fps{u}$ and for every $N$ there are a finite set $F_N$ of flat indices and an integer $M_N$ such that $\phi(x)\in u^N\bbQ\fps{u}$ whenever $\coef{z^k}x^b\in u^{M_N}\bbQ\fps{u}$ for all $(k,b)\in F_N$, the other coefficients of $x$ being arbitrary; the same definition is used for maps on bivectors.
\end{conv}

\begin{lem}[the summation lemma]\genusstatus{AI-proposed, unchecked}
\label{lem:tra-sum}
If $\phi$ is of finite type and $(x_l)_{l\in L}$ is a summable family in $\cH^u_+$, then $(\phi(x_l))_{l\in L}$ is summable and $\phi(\sum_lx_l) = \sum_l\phi(x_l)$. The same holds for maps of finite type on bivectors and summable families of bivectors.
\end{lem}

\begin{proof}
Fix $N$, with $F_N$, $M_N$ as in \cref{conv:tra-forms}. By summability only a finite set $L_0\subset L$ of indices has a coefficient $\coef{z^k}x_l^b$, $(k,b)\in F_N$, outside $u^{M_N}\bbQ\fps{u}$; so $\phi(x_l)\in u^N\bbQ\fps{u}$ for $l\notin L_0$, and, as $\sum_{l\notin L_0}x_l$ has all its coefficients with index in $F_N$ in $u^{M_N}\bbQ\fps{u}$, linearity gives $\phi(\sum_lx_l)-\sum_{l\in L_0}\phi(x_l)\in u^N\bbQ\fps{u}$. Since all $\phi(x_l)$ lie in $u^{N_0}\bbQ\fps{u}$, $(\phi(x_l))$ is summable, and $\phi(\sum_lx_l)\equiv\sum_l\phi(x_l)$ modulo every $u^N$. Bivectors are treated in the same way.
\end{proof}

The maps of finite type used below are $x\mapsto\coef{z^k}x^b$, the $\Pi_i$, a model correlator in one insertion (\cref{lem:tra-conv}), $\Omega_{\mathrm s}(f,\cdot)$ for $f\in\cH^u_-$, and, coefficientwise in the other slot, $\Pi_i\otimes 1$ (\cref{lem:tra-twoslot}) and the contraction of \cref{lem:tra-E}.

\begin{defn}[the vanishing tower]
\label{def:tra-tower}
Let $\YC$ be the vanishing solution of the model (\cref{def:gau-van}), in lowered components. For $q$ in $\cH^{\mathrm{gr}}_+$ or in $\cH^u_+$ and $i\geq 0$ put
\begin{equation}
\label{eq:tra-tower}
\Pi_i(q) := -\coef{z^i}\,\langle\YC(-z),q(z)\rangle .
\end{equation}
$\Pi_0$ is the \emph{vanishing covector}. In the $u$-adic setting $\YC$ is taken at $v = v(u)$, where it lies in $\bbQ[z^{-1}]\fps{u}^5$ (\cref{lem:gau-vanshape}), and the product is taken in $\bbQ\lser{z}\fps{u}$ (\cref{conv:tra-forms}).
\end{defn}

\emph{The shift rule and the values.} The series $\YC$ has no positive power of $z$. Since $\coef{z^i}(z^k\YC_b(-z)) = (-1)^{k-i}\coef{z^{-(k-i)}}\YC_b$,
\begin{equation}
\label{eq:tra-shift}
\Pi_i(z^kH^b) = \Pi_0(z^{k-i}H^b)\ \ (k\geq i),\qquad \Pi_i(z^kH^b) = 0\ \ (k<i),\qquad \Pi_0(z^mH^b) = (-1)^{m+1}\coef{z^{-m}}\YC_b .
\end{equation}
By the closed form of \cref{lem:gau-vanshape}, $\coef{z^{-n}}\YC_b = \varepsilon_b(-1)^nv^{n-r}/(4^r\,r!\,(n-r)!)$ with $r = \lfloor n/5\rfloor$ if $n\equiv\nu_b+1\pmod 5$, and $\coef{z^{-n}}\YC_b = 0$ otherwise. As $\nu_b+1\equiv 1-b\pmod 5$, this gives
\begin{equation}
\label{eq:tra-values}
\Pi_0(z^mH^b) = -\varepsilon_b\,\frac{v^{m-r}}{4^r\,r!\,(m-r)!}\quad\bigl(r = \lfloor m/5\rfloor\bigr)\ \ \text{if } b+m\equiv 1\pmod 5,\qquad \Pi_0(z^mH^b) = 0\ \ \text{otherwise}.
\end{equation}
In particular:
\begin{enumerate}
\item $\Pi_0(z^mH) = \pi_rv^{4r}$ if $m = 5r$, and $\Pi_0(z^mH) = 0$ if $5\nmid m$, where $\pi_r := 1/(4^r\,r!\,(4r)!)$, $\pi_0 = 1$;
\item $\Pi_0(H^a) = 0$ for $a\neq 1$;
\item the $v$-exponent of a non-zero $\Pi_0(z^mH^b)$ is $m-\lfloor m/5\rfloor\geq 4m/5$; so in the $u$-adic setting $\ord_u\Pi_0(z^mH^b)\geq 4m/5$.
\end{enumerate}
For example $\Pi_0(H) = 1$, $\Pi_0(z1) = -v$ and $\Pi_0(z^6\,1) = -v^5/480$.

\emph{Homogeneity and finite type.} The coefficient of $z^{-n}$ in $\YC_b$ is a multiple of $v^{n-\lfloor n/5\rfloor}$, of scalar weight $5(n-\lfloor n/5\rfloor)-4n = w_b$ when $n\equiv\nu_b+1\pmod 5$ (as $0\leq\nu_b+1\leq 4$); so $\YC_b$ is homogeneous of scalar weight $w_b$, and in the graded setting $\Pi_i$ maps $V_w$ into $\bbQ v^{(w-4i)/5}$ (zero unless $(w-4i)/5$ is a non-negative integer). In the $u$-adic setting $\Pi_i(q) = \sum_{k\geq i,\,b}q_{(k,b)}\,\Pi_0(z^{k-i}H^b)$ converges by (iii), and modulo $u^N$ it depends only on the $q_{(k,b)}$ with $4(k-i)/5<N$, taken modulo $u^N$: $\Pi_i$ maps $\cH^u_+$ to $\bbQ\fps{u}$ and is of finite type.

\emph{The tower as symplectic pairings.} For $k\geq 1$, $z^{-k}J^{\mathrm C}\in\cH^u_-$ by \cref{lem:gau-vanshape}, and for $q\in\cH^u_+$
\begin{equation}
\label{eq:tra-pairing}
\Omega_{\mathrm s}(z^{-k}J^{\mathrm C},q) = (-1)^{k+1}\,\Pi_{k-1}(q),
\end{equation}
because $(z^{-k}J^{\mathrm C})(-z) = (-1)^kz^{-k}J^{\mathrm C}(-z)$ and $\eta(J^{\mathrm C}(-z),q(z)) = \langle\YC(-z),q(z)\rangle$, so the left side is $(-1)^k\coef{z^{k-1}}\langle\YC(-z),q(z)\rangle$.

\begin{lem}[the operator $\Theta$]\genusstatus{AI-proposed, unchecked}
\label{lem:tra-theta}
Let $A$ be $\bbQ[v]$ (graded setting) or $\bbQ\fps{u}$ with $v = v(u)$ ($u$-adic setting). For a sequence $\chi = (\chi_j)_{j\geq 0}$ in $A$, finitely supported in the graded setting, put $\chi(z) := \sum_j\chi_jz^j$ and $\Theta\chi := (\Pi_i(\chi(z)H))_{i\geq 0}$. Then $\Theta = 1+N$, where
\[
(N\chi)_i := \sum_{r\geq 1}\pi_r\,v^{4r}\,\chi_{i+5r},
\]
and $\Theta$ is bijective on the set of these sequences, with inverse $\sum_{r\geq 0}(-N)^r$. The operator $N$, and hence $\Theta^{-1}$, maps sequences vanishing outside $[0,L)$ to sequences vanishing outside $[0,L)$; in the graded setting both map sequences with $\chi_j\in\bbQ v^{(w-4j)/5}$ for all $j$ to sequences of the same kind.
\end{lem}

\begin{proof}
By the shift rule \eqref{eq:tra-shift} and the value (i) above, $\Pi_i(\chi(z)H) =\sum_{j\geq i}\chi_j\Pi_0(z^{j-i}H) = \sum_{r\geq 0}\pi_rv^{4r}\chi_{i+5r}$, which is $\chi_i+(N\chi)_i$ (a convergent sum in the $u$-adic setting, as $\ord_uv^{4r}\geq 4r$). In both settings $(1+N)\sum_{r\leq R}(-N)^r\chi = \chi-(-N)^{R+1}\chi = \sum_{r\leq R}(-N)^r(1+N)\chi$. Graded setting: if $\chi$ vanishes outside $[0,L)$, then $N\chi$ vanishes outside $[0,L-5)$, so $N^r\chi = 0$ for $5r\geq L$ and the sums are finite. $u$-adic setting: $\ord_u(N^r\chi)_i\geq 4r+\min_j\ord_u\chi_j$ uniformly in $i$, so $((-N)^r\chi)_{r\geq 0}$ is summable coefficientwise, $N$ commutes with its sum (the double family $(\pi_kv^{4k}((-N)^r\chi)_{i+5k})_{k\geq 1,\,r\geq 0}$ is summable at each $i$), and $(-N)^{R+1}\chi\to 0$. So $\sum_r(-N)^r$ is a two-sided inverse of $\Theta$. The statement on supports holds for $N$ by its definition, and in the graded setting $v^{4r}\chi_{i+5r}$ has scalar weight $20r+(w-4i-20r) = w-4i$ if $\chi_j$ has weight $w-4j$ for all $j$; both properties pass to $\Theta^{-1}$.
\end{proof}

\begin{defn}[the slice $\sK$ and $\sKo$]
\label{def:tra-slice}
In either setting, with $A$ as in \cref{lem:tra-theta} and $\cH^A_+ := \cH^{\mathrm{gr}}_+$, resp.\ $\cH^A_+ := \cH^u_+$,
\begin{equation}
\label{eq:tra-slice}
\sK := \{q\in\cH^A_+ : \Pi_i(q) = 0\ \text{for all } i\geq 1\},\qquad \sKo := \sK\cap\ker\Pi_0 .
\end{equation}
Equivalently, $\sKo$ is the set of $n\in\cH^A_+$ with $\langle\YC(-z),n(z)\rangle$ in $z^{-1}\bbQ[z^{-1}]\fps{u}$ in the $u$-adic setting, and in $z^{-1}\bbQ[v]\fps{z^{-1}}$ in the graded setting. The \emph{affine slice} is $-z+\sK := -z1+\sK$; its description is \cref{cor:tra-slice}. When the setting must be named we write $\sK^A$, $\sKo^A$; the \emph{homogeneous} elements of $\sK$ are the elements of the $\sK\cap V_w$.
\end{defn}

% Section 4, part 2 of 4: frames, the affine slice, coefficients of frames, frame expansion of bivectors.
% Local letters: chi^q (the H-part of q), pr_{K_0}, f_j^{(k,a)} (the coefficients of frame(k,a)), n(c).

\begin{lem}[frames]\genusstatus{AI-proposed, unchecked}
\label{lem:tra-frames}
Let $A$ be as in \cref{lem:tra-theta}, and $A\{z\} := \bbQ[v][z]$ in the graded, $A\{z\} := \bbQ\fps{u,z}$ in the $u$-adic setting.
\begin{description}
\item[(a)] $\cH^A_+ = \sKo^A\oplus A\{z\}H$. In the graded setting $V_w = (\sKo\cap V_w)\oplus(A\{z\}H\cap V_w)$ for every $w$.
\item[(b)] For every $(k,a)$ with $a\neq 1$ there is exactly one element of $\sKo^A$ of the form $z^kH^a+\chi(z)H$ with $\chi(z)\in A\{z\}$; it is
\[
\fr(k,a) = z^kH^a+\sum_{j<k}f_jz^jH,\qquad f_j = c_{k,a,j}\,v^{(4(k-j)+\nu_a+1)/5},\quad c_{k,a,j}\in\bbQ,
\]
where $c_{k,a,j} = 0$ unless the exponent is an integer, which is then $\geq 1$. No element of $\sKo^A$ has the form $z^kH+\sum_{j<k}\chi_jz^jH$.
\item[(c)] (Frame coordinates.) The map $\sKo^A\to\prod_{k\geq 0,\,a\neq 1}A$, $n\mapsto(n_{(k,a)})$, is injective. In the $u$-adic setting it is bijective; in the graded setting its image is the set of finitely supported families in $\bbQ[v]$, and on $\sKo\cap V_w$ the $(k,a)$-member is a rational multiple of $v^{(w-4k-\nu_a-1)/5}$. The inverse sends $(c_{(k,a)})$ to the sum of the summable family $(c_{(k,a)}\fr(k,a))$. In particular the projection $\mathrm{pr}_{\sKo}$ of $\cH^A_+$ onto $\sKo^A$ along $A\{z\}H$ is $\mathrm{pr}_{\sKo}(q) = \sum_{k,\,a\neq 1}q_{(k,a)}\fr(k,a)$.
\item[(d)] The frames of the $u$-adic setting are the frames of the graded setting at $v = v(u)$.
\end{description}
Moreover $\fr(k,a) = z^kH^a$ if and only if $k<\nu_a+1$; these ten \emph{trivial frames} are those with $(k,a) = (0,0)$, $(0,2)$, $(1,2)$, $(2,2)$, $(3,2)$, $(0,3)$, $(1,3)$, $(2,3)$, $(0,4)$, $(1,4)$. For example
\begin{gather*}
\fr(1,0) = z1+vH =: \se,\qquad \fr(2,4) = z^2H^4+\tfrac{v^2}{2}H,\qquad \fr(3,3) = z^3H^3+\tfrac{v^3}{6}H,\\
\fr(4,2) = z^4H^2-\tfrac{v^4}{24}H,\qquad \fr(6,0) = z^6\,1+vz^5H-\tfrac{v^5}{120}H ;
\end{gather*}
$\se$ is the Euler--dilaton vector.
\end{lem}

\begin{proof}
\emph{The decomposition (a).} For $q\in\cH^A_+$ put $\chi^q := \Theta^{-1}((\Pi_i(q))_{i\geq 0})$ (\cref{lem:tra-theta}) and $\mathrm{pr}_{\sKo}(q) := q-\chi^q(z)H$. By the definition of $\Theta$, $\Pi_i(\chi^q(z)H) = (\Theta\chi^q)_i = \Pi_i(q)$, so $\mathrm{pr}_{\sKo}(q)\in\sKo^A$ and $q = \mathrm{pr}_{\sKo}(q)+\chi^q(z)H$; in the graded setting $(\Pi_i(q))_i$ is finitely supported and $\chi^q(z)$ is a polynomial. If $\chi(z)H\in\sKo^A$, then $\Theta\chi = 0$, hence $\chi = 0$ by the injectivity of $\Theta$: the sum is direct. For $q\in V_w$, $(\Pi_i(q))_i$ vanishes beyond the $z$-degree of $q$ and $\Pi_i(q)\in\bbQ v^{(w-4i)/5}$ (the paragraph on homogeneity after \cref{def:tra-tower}); by \cref{lem:tra-theta} $\chi^q_j\in\bbQ v^{(w-4j)/5}$, so $\chi^q(z)H$ and $\mathrm{pr}_{\sKo}(q)$ lie in $V_w$. This proves (a).

\emph{Frames: (b) and (d).} For $a\neq 1$ put $\fr(k,a) := \mathrm{pr}_{\sKo}(z^kH^a)$. The sequence $(\Pi_i(z^kH^a))_i$, which is $(\Pi_0(z^{k-i}H^a))_i$, vanishes for $i>k$ by the shift rule \eqref{eq:tra-shift} and at $i = k$ by the value (ii) after \cref{def:tra-tower}; so it vanishes outside $[0,k)$, hence so does $\chi^{z^kH^a}$ (\cref{lem:tra-theta}), and $\fr(k,a) = z^kH^a+\sum_{j<k}f_jz^jH$ with $f_j := -\chi^{z^kH^a}_j$. Uniqueness: the difference of two such elements lies in $\sKo^A\cap A\{z\}H = 0$. An element $z^kH+\sum_{j<k}\chi_jz^jH\neq 0$ lies in $A\{z\}H$, hence not in $\sKo^A$, by (a). Exponents: in the graded setting $z^kH^a\in V_w$ with $w = 4k+\nu_a+1$, so by the proof of (a) $f_j\in\bbQ v^{(w-4j)/5} = \bbQ v^{(4(k-j)+\nu_a+1)/5}$, and $f_j = 0$ unless this exponent is a non-negative integer, i.e.\ unless $j\equiv k-\nu_a-1\pmod 5$; for $j\leq k-1$ and $a\neq 1$ (so $\nu_a\geq 0$) the exponent is at least $(4+1)/5 = 1$. For (d): the substitution $v\mapsto v(u)$ commutes with every $\Pi_i$ on polynomials in $z$ over $\bbQ[v]$, so a graded frame at $v(u)$ lies in the $u$-adic $\sKo$, has the form of (b), and is the $u$-adic frame by uniqueness.

By \cref{lem:tra-theta} the frame condition $\Pi_i(\fr(k,a)) = 0$ ($i\geq 0$) reads $f_i+(Nf)_i = -\Pi_0(z^{k-i}H^a)$, i.e.
\[
f_i = -\Pi_0(z^{k-i}H^a)-\sum_{r\geq 1}\pi_rv^{4r}f_{i+5r}\qquad(0\leq i<k),
\]
from which the examples follow with \eqref{eq:tra-values}; e.g.\ for $\fr(6,0)$, $f_5 = -\Pi_0(z1) = v$ and $f_0 = -\Pi_0(z^6\,1)-\pi_1v^4f_5 = -v^5/120$. Trivial frames: if $k\geq\nu_a+1$, the index $j_0 := k-\nu_a-1$ satisfies $0\leq j_0<k$ and $j_0+5r>k$ for $r\geq 1$ (as $\nu_a\leq 3$), so the recursion gives $f_{j_0} = -\Pi_0(z^{\nu_a+1}H^a) = \varepsilon_av^{\nu_a+1}/(\nu_a+1)!\neq 0$; if $k<\nu_a+1$, no $j<k$ satisfies $j\equiv k-\nu_a-1\pmod 5$, since $1\leq k-j\leq\nu_a$, so all $f_j$ vanish. Counting the $k<\nu_a+1$ for $a = 0,2,3,4$ ($\nu_a = 0,3,2,1$) gives the ten trivial frames.

\emph{Frame coordinates: (c).} Let $c = (c_{(k,a)})_{k\geq 0,\,a\neq 1}$ be a family in $A$, finitely supported in the graded setting, and write $f^{(k,a)}_j$ for the coefficient $f_j$ of $\fr(k,a)$. The coefficient $\coef{z^j}(\cdot)^b$ of $c_{(k,a)}\fr(k,a)$ is, for $b\neq 1$, $c_{(j,b)}$ if $(k,a) = (j,b)$ and $0$ otherwise; for $b = 1$ it is $c_{(k,a)}f^{(k,a)}_j$, of $u$-order at least $(4(k-j)+\nu_a+1)/5\geq(4(k-j)+1)/5$ by (b) and $\ord_uv\geq 1$. So the family $(c_{(k,a)}\fr(k,a))$ is summable; let $n(c)$ be its sum. As each $\Pi_i$ is of finite type, \cref{lem:tra-sum} gives $\Pi_i(n(c)) = 0$: $n(c)\in\sKo^A$, with $n(c)_{(k,a)} = c_{(k,a)}$ for $a\neq 1$. Conversely, for $n\in\sKo^A$ and $c_{(k,a)} := n_{(k,a)}$, $n-n(c)$ has no component with $a\neq 1$, so it lies in $\sKo^A\cap A\{z\}H = 0$ by (a). Hence $n\mapsto(n_{(k,a)})$ is bijective with inverse $c\mapsto n(c)$. For any $q$, $\mathrm{pr}_{\sKo}(q)$ and $\sum_{k,\,a\neq 1}q_{(k,a)}\fr(k,a)$ are elements of $\sKo^A$ with the components of $q$ for $a\neq 1$ ($\mathrm{pr}_{\sKo}$ changes only $H$-components); so they are equal, and
\begin{equation}
\label{eq:tra-decomp}
q = \sum_{k,\,a\neq 1}q_{(k,a)}\fr(k,a)+\chi^q(z)H,\qquad \chi^q_j = q_{(j,1)}-\sum_{k>j,\,a\neq 1}q_{(k,a)}f^{(k,a)}_j,
\end{equation}
the last sum being $u$-adically convergent (finite in the graded setting). In the graded setting an element of $\sKo$ is a polynomial in $z$, so its frame coordinates are finitely supported, and on $\sKo\cap V_w$ the coordinate $n_{(k,a)} = \coef{z^k}n^a$ is the coefficient of a monomial $v^mz^kH^a$ of weight $w$, a rational multiple of $v^{(w-4k-\nu_a-1)/5}$; as $\Pi_i$ maps $V_w$ to $\bbQ v^{(w-4i)/5}$, $\sKo\cap\cH^{\mathrm{gr}}_+ = \bigoplus_w(\sKo\cap V_w)$.
\end{proof}

\begin{cor}[the affine slice]\genusstatus{AI-proposed, unchecked}
\label{cor:tra-slice}
$\sK^A = \sKo^A\oplus AH$, the decomposition being $q = (q-\Pi_0(q)H)+\Pi_0(q)H$. Moreover $-z1\in\sK^A$, so the affine slice $-z+\sK$ is the set $\{q\in\cH^A_+ : \Pi_i(q) = 0,\ i\geq 1\}$, i.e.\ $\sK^A$ itself; its points are written uniquely as $q = -z1+n+\mu H$ with $n\in\sKo^A$ and $\mu\in A$ (a local letter), and
\[
\Pi_0(-z1+n+\mu H) = v+\mu .
\]
\end{cor}

\begin{proof}
By the shift rule $\Pi_i(H) = 0$ for $i\geq 1$, and $\Pi_0(H) = 1$; so $H\in\sK^A$, $AH\cap\sKo^A = 0$, and $q-\Pi_0(q)H\in\sKo^A$ for $q\in\sK^A$. By \eqref{eq:tra-values} and the shift rule, $\Pi_0(-z1) = v$, $\Pi_1(-z1) = -\Pi_0(1) = 0$ and $\Pi_i(-z1) = 0$ for $i\geq 2$. Hence $-z1\in\sK^A$, and for $q = -z1+y$, $y = n+\mu H\in\sK^A$, linearity gives $\Pi_0(q) = v+\mu$.
\end{proof}

\begin{lem}[coefficients of frames]\genusstatus{AI-proposed, unchecked}
\label{lem:tra-coeff}
Over $\bbQ\fps{u}$, every $n\in\sKo$ is the sum of the summable family $(n_{(k,a)}\fr(k,a))$, where $\fr(k,a)\in\sKo\cap V_{4k+\nu_a+1}$ is homogeneous over $\bbQ[v]$, evaluated at $v(u)$, and every non-zero coefficient $f_j$ of $\fr(k,a)$ has the $v$-exponent $\wt(k,a)-\wt(j,1) = (4(k-j)+\nu_a+1)/5\geq 1$; every $q\in\sK$ is such a sum plus $\Pi_0(q)H$. So every insertion used in \cref{thm:red-main} is a $\bbQ\fps{u}$-combination, in the sense of summable families, of homogeneous elements of $\sKo$ and of $H$.
\end{lem}

\begin{proof}
\Cref{lem:tra-frames}~(b)-(d) and \cref{cor:tra-slice}, with $\wt(k,a)-\wt(j,1) = ((4k+\nu_a-4)-(4j-1-4))/5$.
\end{proof}

\begin{lem}[frame expansion of bivectors]\genusstatus{AI-proposed, unchecked}
\label{lem:tra-twoslot}
For $E\in\bbH\otimes\bbH\otimes\bbQ\fps{\psi,\psi'}\fps{u}$ write $E^{ab}_{kl} := \coef{\psi^k\psi'^l}E^{ab}$ (first index: first slot), put $E_{\cdot,(l,b)} := \sum_{k,a}E^{ab}_{kl}\psi^kH^a$, an element of $\cH^u_+$ in the variable $\psi$, and
\[
(\Pi_i\otimes 1)E := \sum_{l,b}\Pi_i(E_{\cdot,(l,b)})\,\psi'^lH^b ;
\]
$(1\otimes\Pi_i)E$ is defined in the same way with the two slots exchanged. The following are equivalent: (i) $(\Pi_i\otimes 1)E = 0$ and $(1\otimes\Pi_i)E = 0$ for all $i\geq 0$; (ii) $E$ is the sum of the summable family $(E^{\alpha\beta}\,\fr_\alpha(\psi)\otimes\fr_\beta(\psi'))$ over $\alpha = (k,a)$, $\beta = (l,b)$ with $a,b\neq 1$, where $E^{\alpha\beta} := E^{ab}_{kl}$. The coefficients in (ii) are unique: they are the flat components of $E$ with $a,b\neq 1$.
\end{lem}

\begin{proof}
The family in (ii) is summable for any coefficients in $\bbQ\fps{u}$: the coefficient of $\psi^{k'}\psi'^{l'}H^{a'}\otimes H^{b'}$ of the $(\alpha,\beta)$-term is a product of a coefficient of $\fr_\alpha$ and one of $\fr_\beta$; for $a'\neq 1$ only $\alpha = (k',a')$ contributes, for $a' = 1$ the $u$-order is at least $(4(k-k')+1)/5$ (the exponent of \cref{lem:tra-frames}~(b), with $\nu_a\geq0$ and $\ord_uv\geq1$), and likewise in the second slot. The same computation shows that the coefficient of $\psi^k\psi'^lH^a\otimes H^b$, $a,b\neq 1$, of the sum is $E^{\alpha\beta}$: this is the uniqueness.
(ii)$\Rightarrow$(i): coefficientwise in the other slot, $\Pi_i\otimes 1$ and $1\otimes\Pi_i$ are of finite type and kill every term; apply \cref{lem:tra-sum}.
(i)$\Rightarrow$(ii): for each $(l,b)$, $E_{\cdot,(l,b)}\in\sKo$ by the first half of (i), so $E_{\cdot,(l,b)} = \sum_\alpha E^{ab}_{kl}\fr_\alpha(\psi)$ over $\alpha = (k,a)$, $a\neq 1$ (\cref{lem:tra-frames}~(c)); regrouping the double family $\bigl(E^{ab}_{kl}\,\fr_\alpha(\psi)\otimes\psi'^lH^b\bigr)$ over $(\alpha,(l,b))$, summable by the computation of the first paragraph in the first slot (the second slot is a flat monomial), gives the summable family $(\fr_\alpha\otimes G_\alpha)_\alpha$ (\cref{conv:tra-forms}) and $E = \sum_\alpha\fr_\alpha(\psi)\otimes G_\alpha(\psi')$ with $G_\alpha := \sum_{l,b}E^{ab}_{kl}\psi'^lH^b$. By \cref{lem:tra-sum}, $0 = (1\otimes\Pi_i)E = \sum_\alpha\Pi_i(G_\alpha)\fr_\alpha(\psi)$, an element of $\sKo$ with frame coordinates $(\Pi_i(G_\alpha))_\alpha$, which therefore vanish. So $G_\alpha\in\sKo$ and $G_\alpha = \sum_\beta E^{\alpha\beta}\fr_\beta(\psi')$; the resulting double family is summable, with sum $E$.
\end{proof}

``Edges in $\sKo\otimes\sKo$'' means (ii), a completed tensor product made precise by summable families.

% Section 4, part 3 of 4: the u-adic convergence lemma (stated and proved in an appendix).

All infinite sums of this section and of \cref{sec:red,sec:flh} are justified by \cref{lem:tra-conv}, stated and proved in \cref{sec:tra-convproof}: frames and their expansions inside correlators, edges and leaves are summable (parts (a), (b)); the vanishing series and its pairings are ring operations in $\bbQ\lser{z}\fps{u}$ (part (c)); leaves of $\psi$-degree $\geq 2$ are bounded by the dimension (part (d)); the vertex graph resums under the Model Lemma (part (e)); and the dilaton re-expansion of a vertex tensor converges when the dilaton factor is $1$ at $u = 0$, and in general not otherwise (part (f)).

\input{sections/s04_transport_4}



% Section 5 (one part); label prefix: red.
% No new macros. Local letters: y^circ (the K_0-part of the leaf, as in sec:tra), t_alpha (frames at the half-edges).

\section{Reduction to the model}
\label{sec:red}

This section proves Theorem R: for a dressing $P$ of the model with (Sp) and (A) (\cref{def:tra-dressing}), every term of the flat graph sum with an edge is regular, and the vertex graph is $c_gT_P^{2-2g}$ up to a regular series, provided that the Model Lemma on $\sK$ holds in the genera that occur; the Model Lemma is proved in \cref{sec:mod} (\cref{thm:mod-MLK}). The proof uses only the results of \cref{sec:tra} and \cref{lem:tra-conv}, no gauge equation and no datum of the quintic beyond $T_P$.

\subsection{The flat sum and the Model Lemma}
\label{sec:red-sum}

\begin{defn}[the flat sum]
\label{def:red-flatsum}
Let $P$ be a dressing with (Sp), with edge $E_P$, leaf point $q_T$ and $\tilde q_T = d\,q_T$ (\cref{def:tra-edge}). For a leg-free stable graph $\Gamma$ of genus $g\geq 2$, $\mathrm{Cont}_\Gamma$ is $|\Aut\Gamma|^{-1}$ times the contraction along $\Gamma$, as in \cref{def:set-graphsum}, of $E_P$ on every edge with the vertex tensor
\[
d^{\,2g_v-2+n_v}\,\dd^{n_v}\overline F_{g_v}(\tilde q_T)
\]
at every vertex, of genus $g_v$ and valence $n_v$, one slot per half-edge; here $\overline F_h$ is the dilaton-shifted ancestor potential of the model (\cref{def:set-modcohft,def:set-graphsum}) and $\dd^n\overline F_h$ its $n$-th derivative. The \emph{flat sum} is $F^P_g := \sum_\Gamma\mathrm{Cont}_\Gamma$, the sum over the leg-free stable graphs of genus $g$, with no further factor. The \emph{vertex graph} has one vertex and no edge, and
\begin{equation}
\label{eq:red-vertex}
\mathrm{Cont}_{\mathrm{vertex}} = d^{\,2g-2}\,\overline F_g(\tilde q_T) = d^{\,2g-2}\sum_{m\geq 0}\frac1{m!}\langle(dT_{\geq 2})^m\rangle_g ,
\end{equation}
with $dT_{\geq 2} = \tilde q_T+z1\in z^2\cH^u_+$ as in \eqref{eq:tra-leaf};
an \emph{edge graph} is a graph with at least one edge; $\Gamma$ is \emph{regular} if $\mathrm{Cont}_\Gamma\in\bbQ\fps{u}$. For $P = \Pc$ we write $F_\Gamma := \mathrm{Cont}_\Gamma$, the \emph{flat term} of $\Gamma$; it is compared with the outer term $T_\Gamma$ in \cref{thm:flh-main}.
\end{defn}

The vertex tensor is defined by this homogeneous form, which makes sense for every dressing with (Sp), since its leaf $dT_{\geq 2}$ has $\psi$-degree at least $2$; every $\mathrm{Cont}_\Gamma$ is a finite sum (\cref{lem:tra-F}). When $d(0) = 1$, as for $\Pc$ (\cref{prop:tra-can}), it equals the Taylor expansion $\dd^{n_v}\overline F_{g_v}(q_T)$ at the base point $-z1$ (\cref{lem:tra-F}); when $d(0)\neq 1$ that expansion has in general no $u$-adic limit (\cref{lem:tra-conv}~(f2)).

\begin{defn}[the Model Lemma on $\sK$]
\label{def:red-ML}
For $h\geq 2$ let $\Fm_h(v) := \langle\ \rangle_h(v)$ be the genus-$h$ free energy of the model, a rational multiple of $v^{2-2h}$ by \cref{lem:set-modhom}; put $\Fm_1(v) := -\frac{1}{12}\log v$ and $\Fm_0(v) := -\frac52v^2\log v$, of which only the derivatives of order $\geq 1$, resp.\ $\geq 3$, are used. For $h\geq 0$, $\ML{h}$ is the following statement.
\begin{enumerate}
\item For $n\geq 1$ with $2h-2+n>0$ and homogeneous $X_1,\dots,X_n\in\sK$ (graded setting, coefficients in $\bbQ[v]$),
\[
\langle X_1\cdots X_n\rangle_h(v)-(\Fm_h)^{(n)}(v)\prod_{i=1}^n\Pi_0(X_i)\ \in\ \bbQ[v] .
\]
\item For $h\geq 2$, $\Fm_h = c_hv^{2-2h}$, with $c_h$ as in \eqref{eq:set-cg}.
\end{enumerate}
By multilinearity over $\bbQ[v]$ and polarisation, (i) is equivalent to the function form: for every $y\in\sK$ (graded setting) and every $n$ with $2h-2+n>0$, the Taylor coefficient $\frac1{n!}\bigl(\langle y^n\rangle_h-(\Fm_h)^{(n)}(v)\Pi_0(y)^n\bigr)$ of $\overline F_h-\Fm_h\circ\Pi_0$ at the base point lies in $\bbQ[v]$; in short, $\overline F_h|_{-z+\sK} = \Fm_h\circ\Pi_0+\mathrm{Reg}_h$ with a regular remainder $\mathrm{Reg}_h$, an identity of the Taylor coefficients at the base point of order $n>2-2h$ (all orders for $h\geq2$, $n\geq1$ for $h=1$, $n\geq3$ for $h=0$).
\end{defn}

\begin{lem}[genus bound for vertices of edge graphs]\genusstatus{AI-proposed, unchecked}
\label{lem:red-genus}
Let $\Gamma$ be a leg-free stable graph of genus $g$ with at least one edge. Then every vertex of $\Gamma$ carries at least one half-edge and has genus at most $g-1$.
\end{lem}

\begin{proof}
$\Gamma$ is connected and has an edge, so no vertex is isolated. The genus of $\Gamma$ is the sum of the vertex genera plus the first Betti number of $\Gamma$. If some vertex had genus $g$, all other vertices would have genus $0$ and $\Gamma$ would be a tree; a tree with an edge has at least two vertices of valence $1$, so a vertex of genus $0$ would carry a single half-edge, against stability. Hence every vertex has genus at most $g-1$.
\end{proof}

\subsection{Theorem R}
\label{sec:red-thm}

\begin{thm}[Theorem R]\genusstatus{AI-proposed, unchecked}
\label{thm:red-main}
Fix $g\geq 2$ and a dressing $P$ of the model at $v(u)$, $v(0) = 0$, $v'(0)\neq 0$, with (Sp) and (A) (\cref{def:tra-dressing}).
\begin{description}
\item[(a)] If $\ML{h}$ holds for $0\leq h\leq g-1$, then $\mathrm{Cont}_\Gamma\in\bbQ\fps{u}$ for every leg-free stable graph $\Gamma$ of genus $g$ with at least one edge.
\item[(b)] If $\ML{g}$ holds, then $\mathrm{Cont}_{\mathrm{vertex}}\in c_g\,T_P^{2-2g}+\bbQ\fps{u}$.
\end{description}
The gauge equation is not used.
\end{thm}

\begin{proof}
\emph{Step 1 (edges and leaves).} By \cref{cor:tra-edges} and \cref{lem:tra-twoslot}, $E_P = \sum_{\alpha,\beta}E_P^{\alpha\beta}\fr_\alpha(\psi)\otimes\fr_\beta(\psi')$, a summable family with $E_P^{\alpha\beta}\in\bbQ\fps{u}$; so every half-edge insertion is a frame, a homogeneous element of $\sKo$ evaluated at $v(u)$ (\cref{lem:tra-coeff}). By \cref{lem:tra-T}~(c), (d), the leaf is $dT_{\geq 2} = \tilde q_T+z1 = y^\circ+\delta_PH$ with $y^\circ\in\sKo$, whose frame coordinates lie in $\bbQ\fps{u}$ and vanish for $k\leq 1$, and $\delta_P\in u^2\bbQ\fps{u}$.

\emph{Step 2 (vertex factors).} By \cref{lem:red-genus}, every vertex of an edge graph carries $n\geq 1$ half-edges and has genus $h\leq g-1$. By \cref{def:red-flatsum}, the factor of such a vertex whose half-edges carry the frames $t_{\alpha_1},\dots,t_{\alpha_n}$ is
\[
d^{\,2h-2+n}\sum_{m\geq 0}\frac1{m!}\langle t_{\alpha_1}\cdots t_{\alpha_n}(dT_{\geq 2})^m\rangle_h ,
\]
a finite sum ($m\leq 3h-3+n$, \cref{lem:tra-conv}~(d)). Expanding multilinearly, with $y^\circ$ expanded in frames (\cref{lem:tra-coeff}), every term is a product of elements of $\bbQ\fps{u}$ (powers of $d$ and $\delta_P$, frame coordinates of $y^\circ$) and of a model correlator of frames and $H$ containing the frame $t_{\alpha_1}$. As $\Pi_0(t_{\alpha_1}) = 0$, $\ML{h}$~(i) says that this correlator lies in $\bbQ[v]$, hence in $\bbQ\fps{u}$ at $v = v(u)$.

\emph{Step 3 (edge graphs).} By \cref{lem:tra-conv}~(b4), applied to the edges $E_P$, which have property (i) of \cref{lem:tra-twoslot} by \cref{cor:tra-edges}, and to the leaves $dT_{\geq 2}\in z^2\cH^u_+\cap\sK$, the vertex tensors of \cref{def:red-flatsum} being the vertex factors of \cref{lem:tra-conv}~(b3) with $y=dT_{\geq2}$ times the units $d^{\,2g_w-2+n_w}$ of $\bbQ\fps{u}$, of product $d^{\,2g-2}$, which multiplies the contraction and every member of the family, the contraction along $\Gamma$ is the sum of the summable family obtained by expanding every edge in frames and every leaf as in Step 1; by \cref{lem:tra-conv}~(b2) no frame is split into its flat components. By Steps 1 and 2 every member of the family lies in $\bbQ\fps{u}$, hence so does its sum: $\mathrm{Cont}_\Gamma\in\bbQ\fps{u}$. This proves (a).

\emph{Step 4 (the vertex graph).} By \eqref{eq:red-vertex}, $\mathrm{Cont}_{\mathrm{vertex}} = d^{\,2g-2}\,\overline F_g(\tilde q_T)$. By \cref{lem:tra-T}~(c), (d), $\tilde q_T\in -z1+z^2\cH^u_+$, $\Pi_i(\tilde q_T) = 0$ for $i\geq 1$, and $\Pi_0(\tilde q_T)-v = \delta_P\in u^2\bbQ\fps{u}$. So \cref{lem:tra-conv}~(e) ($g\geq 2$, $v'(0)\neq 0$) gives under $\ML{g}$ that $\overline F_g(\tilde q_T)\in c_g\,(v+\delta_P)^{2-2g}+\bbQ\fps{u}$. Finally $v+\delta_P = dT_P$, so $d^{\,2g-2}c_g(dT_P)^{2-2g} = c_gT_P^{2-2g}$, and $d^{\,2g-2}\bbQ\fps{u} = \bbQ\fps{u}$ as $d$ is a unit. This proves (b); since $v'(0)\neq 0$, $T_P = (v+\delta_P)/d$ is $u$ times a unit, so $T_P^{2-2g}\in\bbQ\lser{u}$.
\end{proof}

\begin{cor}[Theorem R for the canonical gauge]\genusstatus{AI-proposed, unchecked}
\label{cor:red-can}
Let $(\alpha,\beta)\in\bbQ^2$ give a conifold frame (\cref{def:set-frames}) and let $g\geq 2$. If $\ML{h}$ holds for $0\leq h\leq g-1$, every edge graph of genus $g$ of the flat sum of $\Pc.\Omega^{\mathrm{mod}}_{v(u)}$ is regular: $F_\Gamma\in\bbQ\fps{u}$. If moreover $\ML{g}$ holds, the vertex graph satisfies $F_{\mathrm{vertex}}\in c_g\,\tc^{2-2g}+\bbQ\fps{u}$.
\end{cor}

\begin{proof}
By \cref{prop:tra-can}, $\Pc$ is a dressing of the model at $v(u) = u\tau(u)^{-5/4}$, with $v(0) = 0$ and $v'(0) = 1$, satisfying (Sp) and (A), and $T_{\Pc} = \tc$. Apply \cref{thm:red-main}.
\end{proof}

By \cref{thm:mod-MLK} the hypotheses $\ML{h}$ hold for every $h\geq 0$; \cref{thm:flh-main} transfers the conclusion of \cref{cor:red-can} from the flat terms $F_\Gamma$ to the outer terms $T_\Gamma$, and \cref{thm:mai-main} assembles the result.

% Section 6, part 1 (parts 2-10 follow by \input). Label prefix: mod. No new macros.
\section{The model: topological recursion on the model curve}
\label{sec:mod}

Topological recursion is applied here, and only here, to the curve $x=\Wv(p)$, $y=p$ on $\bbP^1$, $v\neq 0$, and to its rescalings $(\kappa\Wv,\mu p)$ (\cref{lem:mod-scale}); nothing about the quintic enters. The Model Lemma on $\sK$ (\cref{thm:mod-MLK}) follows from the identification with topological recursion (\cref{thm:mod-tr}), a leg dictionary (\cref{prop:mod-LD,prop:mod-dii}), a lattice theorem (\cref{thm:mod-lattice}), L-regularity (\cref{thm:mod-Lreg}), the Gap Theorem (\cref{thm:mod-gap}) and Theorem~V (\cref{thm:mod-V}). The inputs (Q) are those of \cref{sec:int-inputs}, each cited where used.

\subsection{Conventions; the model as topological recursion}
\label{sec:mod-tr}

\begin{conv}[Topological recursion]
\label{conv:mod-kernel}
Let $x,y$ be meromorphic functions on $\bbP^1$ such that $\dd x$ has only simple zeros, the \emph{branch points} $\varsigma_b$, and $\dd y\neq 0$ there; let $\iota_b$ be the local involution of $x$ near $\varsigma_b$ ($x\circ\iota_b=x$, $\iota_b\neq\mathrm{id}$). Put $\omega_{0,1}:=y\,\dd x$ and $\omega_{0,2}(p_1,p_2):=\dd p_1\,\dd p_2/(p_1-p_2)^2$. For $2h-2+n+1>0$ the correlators $\omega_{h,n+1}$ are defined by the recursion, with $J=(p_1,\dots,p_n)$,
\begin{equation}
\label{eq:mod-kernel}
\begin{gathered}
\omega_{h,n+1}(p_0,J)=\sum_b\Res_{q\to\varsigma_b}K_b(p_0,q)\Big[\omega_{h-1,n+2}(q,\iota_bq,J)+{\sum}'\,\omega_{h_1,|I_1|+1}(q,I_1)\,\omega_{h_2,|I_2|+1}(\iota_bq,I_2)\Big],\\
K_b(p_0,q):=\frac{1}{2}\,\frac{\int_{\iota_bq}^{q}\omega_{0,2}(p_0,\cdot)}{\big(y(q)-y(\iota_bq)\big)\,\dd x(q)},
\end{gathered}
\end{equation}
where ${\sum}'$ runs over $h_1+h_2=h$ and $I_1\sqcup I_2=J$ with the two terms containing $\omega_{0,1}$ omitted. Here $q$ is a point of the curve. This is the recursion of \cite[(2-26)--(2-29)]{Dunin-Barkowski:2015caa}, with the kernel (2-28), to which the quoted theorem of \cref{thm:mod-tr} refers. For $h\geq 2$, on a curve where $\omega_{h,1}$ has zero residues at the branch points, so that the choice of $\Phi_{y\dd x}$ does not matter (for the model curve, \cref{lem:mod-xfun}), the free energy is defined by $(2-2h)\FTR_h:=\sum_b\Res_{\varsigma_b}\Phi_{y\dd x}\,\omega_{h,1}$, where $\dd\Phi_{y\dd x}=y\,\dd x$; only this definition is used (it is that of \cite[(2.1)]{Hock25}; \cite{Dunin-Barkowski:2015caa} defines no free energy). The kernel printed in \cite[(4-3), (4-9)]{Eynard:2007kz} equals the kernel of \cite[(1.1)]{ABDKS25} with $\omega^{(0)}_{1,0}=-y\,\dd x$, and both are $-K_b$. Since $\omega_{h,n}$ is built with $2h-2+n$ kernels, the correlators $W^{(h)}_n$ of \cite{Eynard:2007kz} and $\omega^{(h)}_{n,0}$ of \cite{ABDKS25} are
$\omega_{h,n}=(-1)^n\omega^{(h)}_{n,0}$ for all $(h,n)$, $\omega_{h,n}=(-1)^nW^{(h)}_n$ for $2h-2+n\geq0$ (at $(h,n)=(0,1)$, $W^{(0)}_1=0$ \cite[(4-7)]{Eynard:2007kz}), and $\FTR_h=-F^{(h)}$, with $F^{(h)}$ of \cite[Def.~4.3 (4-10)]{Eynard:2007kz}; in particular our $\omega_{0,3}$ is the negative of \cite[Thm.~4.1 (4-20)]{Eynard:2007kz}. The same dictionary is recorded in \cref{sec:sgn-tr} (\cref{tab:sgn-tr}).
\end{conv}

By \cite[Thm.~4.2, Thm.~4.6]{Eynard:2007kz}, translated by the dictionary above, every $\omega_{h,n}$ with $2h-2+n>0$ is a symmetric meromorphic multidifferential whose poles, in each variable, lie at the branch points only; in particular it has no poles on the diagonals.

\begin{defn}[The model curve]
\label{def:mod-curve}
Let $v\in\bbC$, $v\neq 0$, fix a fifth root $\varsigma$ of $v$, and let $x:=\Wv(p)=p^4/4+v/p$ and $y:=p$ on $\bbP^1$ with the coordinate $p$. Then $\dd x=(p^5-v)\,\dd p/p^2$; the branch points are $\varsigma_b:=\zeta^{-b}\varsigma$, $b\in\bbZ/5$; $x''(\varsigma_b)=5\varsigma_b^2\neq 0$; the critical values $x(\varsigma_b)=\frac54\zeta^b\varsigma^4$ are pairwise distinct; $\dd y=\dd p\neq 0$; and $p=0$, $p=\infty$ are the poles of $x$ and of $y\,\dd x$. The five points of a fibre $x^{-1}(X)$ are the roots of $T^5-4XT+4v$. The weights are $p:1$, $v:5$, $x:4$, $\bar p:=v/p:4$, $\log p:0$. The correlators $\omega_{h,n}$ and $\FTR_h$ of this curve are those of \cref{conv:mod-kernel}; their homogeneity in $v$ is \cref{lem:mod-scale}~(ii).
\end{defn}

\begin{defn}[Pairings with legs]
\label{def:mod-pair}
The \emph{leg class} is $\cL:=\bbQ[v][p]+\bbQ(v)[x]$. For $2h-2+n>0$ and legs $\phi_1,\dots,\phi_n\in\cL$ put
\[
\langle\phi_1\cdots\phi_n\rangle^{\mathrm{TR}}_h:=\sum_{b_1,\dots,b_n}\Res_{p_1=\varsigma_{b_1}}\cdots\Res_{p_n=\varsigma_{b_n}}\prod_{l=1}^n\phi_l(p_l)\,\omega_{h,n}(p_1,\dots,p_n),
\]
the \emph{residue pairing}; for one slot we write $\langle\phi,w\rangle:=\sum_b\Res_{\varsigma_b}\phi\,w$. For legs in $\bbQ[v][p]$ the residue pairing equals the pairing $(2\pi\ri)^{-1}\oint_{|p|=R}$ on a circle with $R>|\varsigma|$, since such a leg is holomorphic at $p=0$, where $\omega_{h,n}$ is holomorphic too. Let $\mathfrak r$ be a ray from $0$ to $\infty$ avoiding the branch points. The \emph{log leg} $-\log p$ is paired by
\[
\langle\cdots,-\log p,\cdots\rangle^{\mathrm{TR}}_h:=\int_{p'\in\mathfrak r}\omega_{h,n}(\dots,p',\dots),
\]
and several log legs by iterated integrals along rays; in a mixed pairing the slots with legs in $\cL$ are paired by residues and the other slots along rays. These integrals converge, may be taken in any order, and do not depend on the rays. Indeed $\omega_{h,n}$ is a meromorphic multidifferential on $(\bbP^1)^n$ whose polar divisor lies in $\bigcup_{i,b}\{p_i=\varsigma_b\}$: another component would give, for generic values of the other variables, a pole in one variable off the branch points, against the paragraph after \cref{conv:mod-kernel}. So $\omega_{h,n}$ is jointly holomorphic on products of rays and near their end points, hence bounded there by a constant times $\prod_i|\dd p_i|/(1+|p_i|)^2$; in each variable its residues at the branch points vanish (\cref{lem:mod-xfun} with $f=1$, whose proof does not use pairings), and moving a ray across $\varsigma_b$ changes the integral by $2\pi\ri$ times such a residue. Moreover, if $\eta$ is a meromorphic $1$-form on $\bbP^1$ with poles only at the branch points and zero residues there, then $\eta=\dd G$ with $G$ rational and holomorphic at $0$ and $\infty$, and
\[
\int_{\mathfrak r}\eta=G(\infty)-G(0)=\sum_b\Res_{\varsigma_b}(-\log p)\,\eta,
\]
for any branch of $\log p$ near each $\varsigma_b$: indeed $\Res_{\varsigma_b}(\log p\,\dd G)=-\Res_{\varsigma_b}G\,\dd p/p$, and the residue theorem for the rational form $G\,\dd p/p$, whose other poles are $0$ and $\infty$, gives $-\sum_b\Res_{\varsigma_b}G\,\dd p/p=G(0)-G(\infty)$. So a log leg may equally be paired by residues at the branch points.
\end{defn}

\begin{defn}[Dual descendant functions]
\label{def:mod-xi}
Let $D_x:=\dd/\dd x=\frac{p^2}{p^5-v}\frac{\dd}{\dd p}$ acting on functions, and for $a\in\bbZ/5$
\[
\xi_a:=\varepsilon'_a\,v^{\alpha_a}\frac{p^{m_a}}{p^5-v},\quad (m_a,\varepsilon'_a,\alpha_a)=(4,-1,0),\ (0,1,1),\ (1,1,0),\ (2,-1,0),\ (3,-1,0)\quad(a=0,\dots,4),
\]
and $\xi_{a,d}:=D_x^d\xi_a$ for $d\geq 0$. (The signs $\varepsilon'_a$ are local, not the $\varepsilon_b$ of \cref{def:set-model}.) Each $\xi_{a,d}$ is rational with poles only at the branch points, holomorphic at $p=0$, and $O(p^{-1-4d})$ at $p=\infty$.
\end{defn}

\begin{defn}[Ancestor vector of a leg]
\label{def:mod-ancvec}
For a leg $\phi\in\cL$, or $\phi=-\log p$, or a $\bbQ[v]$-linear combination of these, put $q_\phi(z):=\sum_{a,d}\langle\phi,\dd\xi_{a,d}\rangle H^az^d$; a leg $\phi$ with $q_\phi=X$ is called a \emph{TR leg} of $X$. Here $\langle-\log p,\dd\xi_{a,d}\rangle=\xi_{a,d}(\infty)-\xi_{a,d}(0)$ by \cref{def:mod-pair}.
\end{defn}

The ancestor correlators $\langle\prod_i\psi^{d_i}H^{a_i}\rangle_h(v)$ of $\Omega^{\mathrm{mod}}_v$ are those of \cref{def:set-modcohft}, extended multilinearly to vectors $q_1,\dots,q_n\in\bbH[z]$ (\cref{conv:tra-forms}).

\begin{thm}[The model as topological recursion]\genusstatus{AI-proposed, unchecked}
\label{thm:mod-tr}
Let $v\neq 0$, $n\geq1$ and $2h-2+n>0$. The correlators of the model curve are
\begin{equation}
\label{eq:mod-tr}
\omega_{h,n}(p_1,\dots,p_n)=(-5)^{h-1}\sum_{a_1,\dots,a_n}\sum_{d_1,\dots,d_n\geq 0}\Big\langle\prod_{i=1}^n\psi^{d_i}H^{a_i}\Big\rangle_h(v)\,\prod_{i=1}^n\dd\xi_{a_i,d_i}(p_i),
\end{equation}
a finite sum, the brackets being the ancestor correlators of $\Omega^{\mathrm{mod}}_v$. Equivalently, for legs $\phi_1,\dots,\phi_n$ as in \cref{def:mod-ancvec}, $\langle q_{\phi_1}\cdots q_{\phi_n}\rangle_h=(-5)^{1-h}\langle\phi_1\cdots\phi_n\rangle^{\mathrm{TR}}_h$.
\end{thm}

% Section 6, part 2.
Throughout \cref{lem:mod-laplace,lem:mod-scale,prop:mod-Rtr}, $\kappa\neq 0$ is a complex constant and $x_\kappa:=\kappa\Wv$; the curve $(\bbP^1,x_\kappa,y=p)$ has the branch points $\varsigma_j$ of \cref{def:mod-curve}. The cases used are $\kappa=4$ and $\kappa=1$. The letter $\vartheta$ denotes the Laplace variable of \cite[(2-31), (2-32)]{Dunin-Barkowski:2015caa} (called $\zeta$ there).

\begin{lem}[Formal Laplace series]\genusstatus{AI-proposed, unchecked}
\label{lem:mod-laplace}
Let $c$ be a branch point of $x_\kappa$ and $w$ a local coordinate at $c$ with $x_\kappa=x_\kappa(c)-w^2/2$. For a $1$-form $\omega=\sum_{m\geq -M}a_mw^m\,\dd w$ meromorphic at $c$ with $\Res_c\omega=0$ put
\begin{equation}
\label{eq:mod-laplace}
\begin{gathered}
\mathsf L_c[\omega](\vartheta):=\sum_ma_m\,\mathsf G_m(\vartheta),\\
\mathsf G_{2n}:=(2n-1)!!\,\vartheta^{-n}\ (n\geq 0),\qquad \mathsf G_{-2n}:=\frac{(-1)^n\vartheta^n}{(2n-1)!!}\ (n\geq 1),\qquad \mathsf G_m:=0\ (m\text{ odd}),
\end{gathered}
\end{equation}
a formal Laurent series in $\vartheta^{-1}$ (the coefficient $a_{-1}$ vanishes). Then:
\begin{enumerate}
\item if $F$ is meromorphic at $c$ and $F\,\dd x_\kappa$ has zero residue at $c$, then $\mathsf L_c[\dd F]=-\vartheta\,\mathsf L_c[F\,\dd x_\kappa]$;
\item $w$ is determined by $x_\kappa$ up to $w\mapsto -w$, which multiplies $\mathsf L_c$ by $-1$;
\item let $k\in\bbC^\times$ and $v=k^5v'$, so that $x_{\kappa,v}(kp')=k^4x_{\kappa,v'}(p')$ and $\varsigma_j(v)=k\varsigma_j(v')$, and normalise the coordinates $w_j$ at $\varsigma_j$ by $\frac{\dd w_j}{\dd p}(\varsigma_j)=\sqrt{-5\kappa}\,\varsigma_j$ with one square root for all $v$; then $w_j(v)(kp')=k^2w_j(v')(p')$ and $\mathsf L_{\varsigma_j(v)}[p^m\dd p](\vartheta)=k^{m-1}\,\mathsf L_{\varsigma_j(v')}[p'^m\dd p'](k^4\vartheta)$;
\item if $\omega$ is holomorphic at $c$, then $\vartheta\frac{\dd}{\dd\vartheta}\mathsf L_c[\omega]=\vartheta\,\mathsf L_c\big[(x_\kappa-x_\kappa(c))\,\omega\big]+\frac12\mathsf L_c[\omega]$.
\end{enumerate}
\end{lem}

The integrals $\int\omega\,e^{(x_\kappa-x_\kappa(c))\vartheta}$ in \cite[(2-31), (2-32)]{Dunin-Barkowski:2015caa} are read as $\sqrt{2\pi/\vartheta}\,\mathsf L_c[\omega]$. For the holomorphic part this is the termwise Gaussian integral. The only polar term that occurs is $w^{-2}\dd w$, in $\dd\xi^i$ at its own branch point; as $w^{-2}\re^{-\vartheta w^2/2}\dd w$ has no residue, its integral along any path avoiding $0$ is $-\sqrt{2\pi\vartheta}=\sqrt{2\pi/\vartheta}\,\mathsf G_{-2}$, the only value for which $R_{\mathrm{TR}}^{-1}=1+O(\vartheta^{-1})$.

\begin{proof}
All four statements are linear, so it suffices to check monomials. (i) For $F=w^m$, $m\neq -2$ (for $F=w^{-2}$ the form $F\,\dd x_\kappa$ has a residue), $\dd F=mw^{m-1}\dd w$ and $F\,\dd x_\kappa=-w^{m+1}\dd w$, so the claim is $m\mathsf G_{m-1}=\vartheta\mathsf G_{m+1}$: both sides are $(2n+1)!!\,\vartheta^{-n}$ for $m=2n+1\geq 1$, $(-1)^n\vartheta^{n+1}/(2n-1)!!$ for $m=-(2n+1)$, and $0$ for even $m$. (ii) $w^2=2(x_\kappa(c)-x_\kappa)$, and $w\mapsto -w$ multiplies $w^m\dd w$ by $(-1)^{m+1}$, while only even $m$ contribute. (iii) $w_j(v)(kp')$ and $k^2w_j(v')(p')$ are coordinates at $\varsigma_j(v')$ with the same square $2k^4(x_{\kappa,v'}(\varsigma_j(v'))-x_{\kappa,v'}(p'))$ and the same derivative $k^2\sqrt{-5\kappa}\,\varsigma_j(v')$ there, so they agree. If $p'^m\dd p'=\sum_nb_nw'^n\dd w'$ with $w'=w_j(v')$, then $p^m\dd p=\sum_nk^{m-1-2n}b_nw^n\dd w$ with $w=k^2w'$, and $\mathsf G_{2s}(\vartheta)=k^{4s}\mathsf G_{2s}(k^4\vartheta)$, $\mathsf G_{-2s}(\vartheta)=k^{-4s}\mathsf G_{-2s}(k^4\vartheta)$ give the claim. (iv) On $w^{2n}\dd w$ the left side is $-n(2n-1)!!\,\vartheta^{-n}$; as $(x_\kappa-x_\kappa(c))w^{2n}\dd w=-\frac12w^{2n+2}\dd w$, the right side is $-\frac12(2n+1)!!\,\vartheta^{-n}+\frac12(2n-1)!!\,\vartheta^{-n}$, the same; odd powers give $0$ on both sides.
\end{proof}

\begin{lem}[Scaling of the curve]\genusstatus{AI-proposed, unchecked}
\label{lem:mod-scale}
\begin{enumerate}
\item For constants $\kappa,\mu\neq 0$ let $\omega^{(\kappa,\mu)}_{h,n}$ and $F^{(\kappa,\mu)}_h$ be the correlators and free energies of $(\bbP^1,\kappa\Wv,\mu p,\omega_{0,2})$. Then $\omega^{(\kappa,\mu)}_{h,n}=(\kappa\mu)^{2-2h-n}\omega_{h,n}$ for $2h-2+n>0$, and $F^{(\kappa,\mu)}_h=(\kappa\mu)^{2-2h}\FTR_h$ for $h\geq 2$.
\item For $k\in\bbC^\times$ let $m_k(p):=kp$. Then $m_k^*\,\omega^{(k^5v)}_{h,n}=k^{5(2-2h-n)}\,\omega^{(v)}_{h,n}$ and $\FTR_h(k^5v)=k^{5(2-2h)}\FTR_h(v)$, where the superscript is the parameter of the curve. Consequently a pairing of $\omega_{h,n}$ with the legs $p^{m_1},\dots,p^{m_j}$ in $j$ slots and with $-\log p$ in the other $n-j$ slots is a function $C(v)$ with $C(k^5v)=k^{5e}C(v)$, $e=2-2h-n+\sum_lm_l/5$.
\end{enumerate}
\end{lem}

\begin{proof}
(i) The branch points, the involutions $\iota_b$ and $\omega_{0,2}$ are unchanged, and the kernel becomes $K_b/(\kappa\mu)$. By induction on $2h-2+n$, starting from the exponent $0=2-0-2$ of $\omega_{0,2}$: every term of \eqref{eq:mod-kernel} for $\omega_{h,n+1}$ is one kernel times $\omega_{h-1,n+2}$, of exponent $2-2h-n$, or times $\omega_{h_1,|I_1|+1}\omega_{h_2,|I_2|+1}$, of exponent $(1-2h_1-|I_1|)+(1-2h_2-|I_2|)=2-2h-n$, so the total exponent is $2-2h-(n+1)$. For the free energy, $\Phi_{y\dd x}$ becomes $\kappa\mu\,\Phi_{y\dd x}$. (ii) The automorphism $m_k$ of $\bbP^1$ satisfies $m_k^*W_{k^5v}=k^4\Wv$, $m_k^*p=kp$ and $m_k^*\omega_{0,2}=\omega_{0,2}$, and maps the branch points of $\Wv$ to those of $W_{k^5v}$. The recursion is natural under isomorphisms of curves, so $m_k^*\omega^{(k^5v)}_{h,n}$ is the correlator of $(\bbP^1,k^4\Wv,kp,\omega_{0,2})$, which is $k^{5(2-2h-n)}\omega^{(v)}_{h,n}$ by (i); likewise for $\FTR_h$. For the pairings, pull the residues back along $m_k$: a leg $p_l^{m_l}$ contributes $k^{m_l}$, and a ray goes to a ray, along which the pairing does not depend on the choice (\cref{def:mod-pair}).
\end{proof}

\begin{prop}[The $R$-matrix of the curve]\genusstatus{AI-proposed, unchecked}
\label{prop:mod-Rtr}
Let $v\neq 0$ and take at the branch points of $x_\kappa$ the coordinates $w_j$ of \cite[\S2.6]{Dunin-Barkowski:2015caa}, $x_\kappa=x_\kappa(\varsigma_j)-w_j^2/2$, normalised by $\frac{\dd w_j}{\dd p}(\varsigma_j)=\sqrt{-5\kappa}\,\varsigma_j$ with one square root for all $j$. Let $e^{\mathrm{TR}}_j$, $\Delta_j$ and $\xi^j$ be the corresponding data (2-30), (2-34) of \cite[\S2.6]{Dunin-Barkowski:2015caa}: $\Delta_j^{-1/2}=\frac{\dd y}{\dd w_j}(\varsigma_j)$ and $\xi^j(q)=\int_\infty^q\omega_{0,2}(w_j,\cdot)/\dd w_j|_{w_j=0}$; and let $R_{\mathrm{TR}}$ be the matrix series defined by (2-31) read through \cref{lem:mod-laplace}, $R_{\mathrm{TR}}^{-1}(\vartheta^{-1})^j{}_i:=-\vartheta^{-1}\mathsf L_{\varsigma_j}[\dd\xi^i]$ (upper index the row: the left-action convention of \cite[Remark~2.1]{Dunin-Barkowski:2015caa}, in which (2-32) is $R^{-1}$ applied to the unit \cite[Section~3]{Dunin-Barkowski:2015caa}). Then, as formal series in $\vartheta^{-1}$:
\begin{enumerate}
\item $R_{\mathrm{TR}}^{-1}(\vartheta^{-1})^j{}_i=\Mo(t)_{ij}$ with $t=-1/(\kappa\varsigma^4\vartheta)$, for all $i,j$; consequently
\begin{equation}
\label{eq:mod-Rtr}
R_{\mathrm{TR}}(\psi)=\Mo\big(\psi/(\kappa\varsigma^4)\big)
\end{equation}
in the basis $e^{\mathrm{TR}}_j$, and $R_{\mathrm{TR}}$ is symplectic in this orthonormal basis;
\item (2-32) holds: $\sum_kR_{\mathrm{TR}}^{-1}(\vartheta^{-1})^i{}_k\Delta_k^{-1/2}=\mathsf L_{\varsigma_i}[\dd y]$, both sides being $\varsigma^{-1}\Rhat^{\mathrm m,(i)}_0(t)/\sqrt{-5\kappa}$. In the frozen basis $\bar e_j:=e^{\mathrm{TR}}_j/\sqrt{-5\kappa}$ the unit $\sum_k\Delta_k^{-1/2}e^{\mathrm{TR}}_k$ is $\hat 1:=\sum_k\bar e_k/\varsigma_k$, and the dilaton leaves of the graph formula carry the translation $\psi\big(\hat 1-\Mo^{-1}(\psi/(\kappa\varsigma^4))\hat 1\big)=\psi\,\hat 1-\psi R_{\mathrm{TR}}^{-1}(\psi)\hat 1$ of \cite[Def.~2.13]{PPZ15}.
\item (the edge condition) For a double series $\Omega=\sum_{m,n\geq-M}a_{mn}\,w_i^m\,w_j'^{\,n}\,\dd w_i\,\dd w_j'$ in the coordinates $w_i$ at $\varsigma_i$ in the first point and $w_j$ at $\varsigma_j$ in the second, with $a_{-1,n}=a_{m,-1}=0$, put $(\mathsf L_i\otimes\mathsf L_j)[\Omega](\vartheta_1,\vartheta_2):=\sum_{m,n}a_{mn}\mathsf G_m(\vartheta_1)\mathsf G_n(\vartheta_2)$, and let $B^{\mathrm{reg}}_{ij}:=\omega_{0,2}-\delta_{ij}\,\dd w_i\,\dd w_i'/(w_i-w_i')^2$, which is holomorphic at $(\varsigma_i,\varsigma_j)$. Then, with $\psi_a:=\vartheta_a^{-1}$,
\begin{equation}
\label{eq:mod-edge}
(\psi_1+\psi_2)\,(\mathsf L_i\otimes\mathsf L_j)\bigl[B^{\mathrm{reg}}_{ij}\bigr]=\delta_{ij}-\sum_kR_{\mathrm{TR}}^{-1}(\psi_1)^i{}_k\,R_{\mathrm{TR}}^{-1}(\psi_2)^j{}_k\qquad\text{for all }i,j .
\end{equation}
This is the edge condition \cite[(4-1), with (3-2) and (3-57)]{DuninBarkowski:2012bw}, written in the coordinates $w_i$ of \cite[\S2.6]{Dunin-Barkowski:2015caa} (the coordinate of \cite[(3-1)]{DuninBarkowski:2012bw} is normalised by $x=z^2+a_i$ instead): the coefficient of $\psi_1^{d_1}\psi_2^{d_2}$ in $(\mathsf L_i\otimes\mathsf L_j)[B^{\mathrm{reg}}_{ij}]$ is $(2d_1-1)!!\,(2d_2-1)!!\,a_{2d_1,2d_2}$, the edge weight (3-57), and the right side divided by $\psi_1+\psi_2$ is the edge of \cite[Def.~2.2]{PPZ15} for $R_{\mathrm{TR}}$ in the orthonormal basis $e^{\mathrm{TR}}_j$ (\cref{sec:sgn-R}). So the expansion of $\omega_{0,2}$ at the branch points is the edge of $R_{\mathrm{TR}}$. The form printed in \cite[(2-33)]{Dunin-Barkowski:2015caa} has $+\sum_k$ in place of $\delta_{ij}-\sum_k$ (so for $i\neq j$ it differs from \eqref{eq:mod-edge} by a sign) and is not used.
\end{enumerate}
\end{prop}

The sign of the dilaton leaves is that of \cite[\S2.3]{DuninBarkowski:2012bw} and of \cite[Prop.~2.12, Def.~2.13]{PPZ15}; \cite[\S2.3]{Dunin-Barkowski:2015caa} prints the opposite sign. With legs $R^{-1}(\psi)$ the sign is fixed by the requirement that the action preserve the unit: it is that of \cite[(2-23)]{DuninBarkowski:2012bw} and of \cite[Prop.~2.12]{PPZ15} (\cref{sec:sgn-R}).

% Section 6, part 3.
\begin{proof}[Proof of \cref{prop:mod-Rtr}]
Write $\sqrt{-5\kappa}$ for the root of the statement and $\mathsf L_j:=\mathsf L_{\varsigma_j}$; indices run over $\bbZ/5$.

\emph{Step 1 (integration by parts).} By (2-34), $\dd\xi^i(q)=\frac{\dd p}{\dd w_i}(\varsigma_i)\,\dd q/(\varsigma_i-q)^2$, so $\xi^i(q)=-1/\big(\sqrt{-5\kappa}\,\varsigma_i(q-\varsigma_i)\big)$. It has at most a simple pole at $\varsigma_j$, so $\xi^i\dd x_\kappa$ is holomorphic there, and \cref{lem:mod-laplace}~(i) gives $R_{\mathrm{TR}}^{-1}(\vartheta^{-1})^j{}_i=\mathsf L_j[\xi^i\dd x_\kappa]$.

\emph{Step 2 (partial fractions).} Since $\dd x_\kappa=\kappa(p^5-v)\,\dd p/p^2$ and $(p^5-v)/(p-\varsigma_i)=\sum_{k=0}^4\varsigma_i^kp^{4-k}$, with the local letters $J^{(j)}_m:=\mathsf L_j[p^m\dd p]$, $\tilde J^{(j)}_m:=\sqrt{-5\kappa}\,J^{(j)}_m$:
\begin{equation}
\label{eq:mod-Rtr-pf}
R_{\mathrm{TR}}^{-1}(\vartheta^{-1})^j{}_i=\frac15\sum_{k=0}^4\varsigma_i^{k-1}\,\tilde J^{(j)}_{2-k}.
\end{equation}

\emph{Step 3 (the five series are model rows).} Fix $j$, write $J_m:=J^{(j)}_m$, $z:=-1/(\kappa\vartheta)$ (so $\vartheta x_\kappa=-\Wv/z$) and $t:=z/\varsigma^4$. \emph{(Mod2).} \Cref{lem:mod-laplace}~(i) with $F=p^m$ gives $mzJ_{m-1}-J_{m+3}+vJ_{m-2}=0$; in particular $J_3=vJ_{-2}$. \emph{(Mod1).} By \cref{lem:mod-laplace}~(iii) with $k=\varsigma$, $v'=1$, $J_m(v;\vartheta)=\varsigma^{m-1}J_m(1;\varsigma^4\vartheta)$, so $5v\,\partial_vJ_m=(m-1)J_m+4\vartheta\partial_\vartheta J_m$. By \cref{lem:mod-laplace}~(iv), $\vartheta\partial_\vartheta J_m=\vartheta\kappa\big(\tfrac14J_{m+4}+vJ_{m-1}\big)-\tfrac54\vartheta\kappa\varsigma_j^4J_m+\tfrac12J_m$; eliminating $J_{m+4}=(m+1)zJ_m+vJ_{m-1}$ by (Mod2), using $\kappa\vartheta=-1/z$ and inserting the result into (Mod1) gives $z\partial_vJ_m=-J_{m-1}+J_m/\varsigma_j$ (as $\varsigma_j^5=v$). Put $\tilde Y_m:=e^{-\Wv(\varsigma_j)/z}J_m$ (a formal symbol, with the Leibniz rule); as $\partial_v[\Wv(\varsigma_j(v))]=1/\varsigma_j$, $z\partial_v\tilde Y_m=-\tilde Y_{m-1}$. Hence $Y_b:=\varepsilon_b\tilde Y_{\nu_b}$ satisfy $z\partial_vY_b=b^{\mathrm m}_bY_{b+1}$ (for $b=1$ use $Y_2=-\tilde Y_3=-v\tilde Y_{-2}$ by (Mod2)), the derivation of \cref{prop:set-modrows}~(a) for these series. Write $Y_b=e^{U_j/z}\varsigma^{\nu_b-1}F_b$ with $U_j:=-\Wv(\varsigma_j)$ and $F_b:=\varepsilon_b\varsigma^{1-\nu_b}J_{\nu_b}$. By the homogeneity above, $F_b=\varepsilon_bJ_{\nu_b}(1;\varsigma^4\vartheta)$ is a power series in $t$, and the computation of \cref{prop:set-modrows}~(b) shows that $(F_b)_b$ solves the $T$-system of colour $j$. (Here $F_b\in\bbC\fps t$. The computation and the uniqueness statement of \cref{prop:set-modrows}~(b) are formal and hold verbatim over $\bbC$ in place of $\bbQ(\zeta)$: the recursion has the pivot $4n\zeta^{4j}\neq0$.) \emph{Initial values.} $J_m=\varsigma_j^{m-1}/\sqrt{-5\kappa}+O(\vartheta^{-1})$ and $(\varsigma_j/\varsigma)^{\nu_b-1}=\zeta^{j(b+1)}$, so $\sqrt{-5\kappa}F_b(0)=\varepsilon_b\zeta^{j(b+1)}=\Rhat^{\mathrm m,(j)}_b(0)$ (\cref{def:set-M0}). \emph{Uniqueness.} By \cref{lem:gau-colmod}~(i), (ii), $\Rhat^{\mathrm m,(j)}=\zeta^jF^{(j)}$ also solves the colour-$j$ system, and by \cref{prop:set-modrows}~(b) (pivot $4n\zeta^{4j}$) a power-series solution is determined by the constant term of its component $b=0$. Hence $\tilde J^{(j)}_{m_b}=\varepsilon_b\varsigma^{m_b-1}\Rhat^{\mathrm m,(j)}_b(t)$ for $(m_0,\dots,m_4)=(0,-1,-2,2,1)$, $m_b\equiv -b\pmod 5$ (for $b=2$, $\tilde J_{-2}=\tilde J_3/v$).

\emph{Step 4 (assembly).} As $b$ runs over $\bbZ/5$, $k:=2-m_b$ runs over $0,\dots,4$, and $\varsigma_i^{k-1}\varsigma^{m_b-1}=\zeta^{-i(b+1)}$. By Step 3 and \eqref{eq:mod-Rtr-pf}, $R_{\mathrm{TR}}^{-1}(\vartheta^{-1})^j{}_i=\frac15\sum_b\varepsilon_b\zeta^{-i(b+1)}\,\Rhat^{\mathrm m,(j)}_b(t)=\sum_b(\Psi_0^{-1})_{ib}\Rhat^{\mathrm m,(j)}_b(t)=\Mo(t)_{ij}$, since $(\Psi_0)_{bi}=\varepsilon_b\zeta^{i(b+1)}$ has the inverse $\frac15\varepsilon_b\zeta^{-i(b+1)}$.

\emph{Step 5 (inverse).} With the upper index as row, Step 4 reads $R_{\mathrm{TR}}^{-1}(\psi)=\Mo\big(-\psi/(\kappa\varsigma^4)\big)^T$, $\psi=\vartheta^{-1}$. By \cref{lem:set-M0sympl}, $\Mo(-t)^T=\Mo(t)^{-1}$, which gives \eqref{eq:mod-Rtr} and symplecticity. This proves (i).

(ii) By (2-30), $\Delta_k^{-1/2}=1/(\sqrt{-5\kappa}\varsigma_k)$. By (i), and as $1/\varsigma_k=\zeta^k/\varsigma$ (\cref{def:mod-curve}) and $\sum_k\zeta^{-kb}=5\delta_{b,0}$, the left side of (2-32) is $\sum_k\Mo(t)_{ki}/(\sqrt{-5\kappa}\,\varsigma_k)=\Rhat^{\mathrm m,(i)}_0(t)/(\sqrt{-5\kappa}\,\varsigma)$, and the right side is $J^{(i)}_0=\Rhat^{\mathrm m,(i)}_0(t)/(\sqrt{-5\kappa}\varsigma)$ by Step 3. The unit is $\sum_k\Delta_k^{-1/2}e^{\mathrm{TR}}_k=\sum_k\bar e_k/\varsigma_k$. Given (2-32), in the graph formula of \cite[Thm.~4.1]{DuninBarkowski:2012bw} a dilaton leaf is a marked point of its vertex carrying $\psi^k$, $k\geq 2$, with the weight $\coef{z^{k-1}}\bigl(-R(-z)\bigr)$ applied to the unit \cite[\S2.3]{DuninBarkowski:2012bw}, in total $\psi\bigl(1-R(-\psi)\bigr)$ applied to the unit; by \cref{lem:set-pushins}~(ii) at every vertex these points give the vertex integrals of the translation by this series in the push-forward form of \cite[Def.~2.5]{PPZ15} and \cite[\S2.3]{Dunin-Barkowski:2015caa}; and $R(-z)$ there, which also decorates the ordinary leaves, is $R^{-1}(z)$ of \cite[\S2.3]{Dunin-Barkowski:2015caa}. Finally $R_{\mathrm{TR}}^{-1}(\psi)=\Mo(\psi/(\kappa\varsigma^4))^{-1}$ by (i).

(iii) Put $f:=1/x_\kappa'(p)=p^2/(\kappa(p^5-v))$ and, for a bidifferential $\Omega$, $\mathsf D\Omega:=\dd_1(\Omega/\dd x_{\kappa,1})+\dd_2(\Omega/\dd x_{\kappa,2})$, where $\Omega/\dd x_{\kappa,a}$ is a function of the $a$-th point (local letters). \emph{Step A:} $\mathsf D\omega_{0,2}=\sum_k\dd\xi^k(p_1)\,\dd\xi^k(p_2)$. Indeed
\[
\mathsf D\omega_{0,2}=\Bigl[\frac{f'(p_1)+f'(p_2)}{(p_1-p_2)^2}+\frac{2\bigl(f(p_2)-f(p_1)\bigr)}{(p_1-p_2)^3}\Bigr]\dd p_1\,\dd p_2 .
\]
Fix $p_2\neq0,\infty$ off the branch points. As a $1$-form in $p_1$ this is regular at $p_1=p_2$ (the bracket tends to $f'''(p_2)/6$) and holomorphic at $p_1=0$ and $p_1=\infty$ ($f$ has a double zero at $0$ and is $O(p^{-3})$ at $\infty$). At $p_1=\varsigma_k$, where $f=r_k/(p-\varsigma_k)+O(1)$ with $r_k:=1/x_\kappa''(\varsigma_k)$, the simple poles of the two terms cancel, and the principal part is $-r_k\,\dd p_1\,\dd p_2/\bigl((p_1-\varsigma_k)^2(\varsigma_k-p_2)^2\bigr)$. By Step~1, $\dd\xi^k(p_1)\,\dd\xi^k(p_2)=\bigl(\tfrac{\dd p}{\dd w_k}(\varsigma_k)\bigr)^2\dd p_1\,\dd p_2/\bigl((\varsigma_k-p_1)^2(\varsigma_k-p_2)^2\bigr)$, which is holomorphic in $p_1$ off $\varsigma_k$ and has the same principal part there, because $x_\kappa=x_\kappa(\varsigma_k)-w_k^2/2$ gives $x_\kappa''(\varsigma_k)\bigl(\tfrac{\dd p}{\dd w_k}(\varsigma_k)\bigr)^2=-1$. The difference of the two sides is a holomorphic $1$-form in $p_1$ on $\bbP^1$, hence zero. \emph{Step B.} If $\Omega$ is holomorphic at $(\varsigma_i,\varsigma_j)$, then $(\mathsf L_i\otimes\mathsf L_j)[\mathsf D\Omega]=-(\vartheta_1+\vartheta_2)(\mathsf L_i\otimes\mathsf L_j)[\Omega]$: apply \cref{lem:mod-laplace}~(i) in the first point to $F=\Omega/\dd x_{\kappa,1}$, which has at most a simple pole at $\varsigma_i$ while $F\,\dd x_{\kappa,1}=\Omega$ has no residue, coefficientwise in the second point, and likewise in the second point. \emph{Step C.} By the definition of $R_{\mathrm{TR}}$, $\mathsf L_j[\dd\xi^k]=-\vartheta\,R_{\mathrm{TR}}^{-1}(\vartheta^{-1})^j{}_k$. For $i\neq j$, $B^{\mathrm{reg}}_{ij}=\omega_{0,2}$, and Steps~A and~B give $-(\vartheta_1+\vartheta_2)(\mathsf L_i\otimes\mathsf L_j)[\omega_{0,2}]=\vartheta_1\vartheta_2\sum_kR_{\mathrm{TR}}^{-1}(\psi_1)^i{}_kR_{\mathrm{TR}}^{-1}(\psi_2)^j{}_k$; divide by $-\vartheta_1\vartheta_2$. For $i=j$ use the coordinate $w=w_i$, in which $1/(\dd x_\kappa/\dd w)=-1/w$: the formula of Step~A with $-1/w$ in place of $f$ gives $\mathsf D\bigl(\dd w\,\dd w'/(w-w')^2\bigr)=\dd w\,\dd w'/(w^2w'^2)$, as $(w^{-2}+w'^{-2})/(w-w')^2-2/(ww'(w-w')^2)=1/(w^2w'^2)$. So $\mathsf DB^{\mathrm{reg}}_{ii}=\sum_k\dd\xi^k\,\dd\xi^k-\dd w\,\dd w'/(w^2w'^2)$, an identity of double Laurent series in $w,w'$ (both sides have poles only on $w=0$ and $w'=0$, and $(\mathsf L_i\otimes\mathsf L_i)$ of a product $\alpha(w)\beta(w')$ is $\mathsf L_i[\alpha]\,\mathsf L_i[\beta]$), and by Step~B and $\mathsf G_{-2}=-\vartheta$, $-(\vartheta_1+\vartheta_2)(\mathsf L_i\otimes\mathsf L_i)[B^{\mathrm{reg}}_{ii}]=\vartheta_1\vartheta_2\bigl(\sum_kR_{\mathrm{TR}}^{-1}(\psi_1)^i{}_kR_{\mathrm{TR}}^{-1}(\psi_2)^i{}_k-1\bigr)$; divide by $-\vartheta_1\vartheta_2$. (At $\psi_2=-\psi_1$, \eqref{eq:mod-edge} gives the symplecticity of (i) again.)
\end{proof}

\begin{proof}[Proof of \cref{thm:mod-tr}]
\emph{Step 1 (the quoted theorem).} \cite[Thm.~4.1]{DuninBarkowski:2012bw}, as recalled in \cite[\S2.6, (2-30)--(2-32), (2-34), (2-35)]{Dunin-Barkowski:2015caa}, with the action of \cite[\S2.3]{Dunin-Barkowski:2015caa} (legs $R^{-1}(\psi)$, edges as in \cite[Def.~2.2]{PPZ15}, dilaton leaves with the translation of \cite[Def.~2.13]{PPZ15}, in the push-forward form of \cite[Def.~2.5]{PPZ15} by the proof of \cref{prop:mod-Rtr}~(ii)), gives for a compact curve with simple branch points, $\dd y\neq 0$ there, and the Bergman kernel: if the matrix $R$ of (2-31) is symplectic in the orthonormal basis, (2-32) holds, and the expansion of the Bergman kernel at the branch points is the edge of $R$, in the form \cite[(4-1)]{DuninBarkowski:2012bw}, that is \eqref{eq:mod-edge} (the printed \cite[(2-33)]{Dunin-Barkowski:2015caa} has $+\sum_k$ in place of $\delta_{ij}-\sum_k$), then for $2h-2+n>0$
\begin{equation}
\label{eq:mod-235}
\omega_{h,n}=\sum_{i_1,\dots,i_n}\sum_{d_1,\dots,d_n\geq 0}\int_{\overline{\mathcal M}_{h,n}}\alpha_{h,n}(e_{i_1},\dots,e_{i_n})\prod_{l=1}^n\psi_l^{d_l}\;\prod_{l=1}^n\dd\big((\dd/\dd x)^{d_l}\xi^{i_l}\big)(p_l),
\end{equation}
where $\alpha=R.\TFT^{\mathrm{TR}}$ and $\TFT^{\mathrm{TR}}_{h,n}(e_j,\dots,e_j)=\Delta_j^{(2h-2+n)/2}$, zero on mixed arguments. We apply it to $(\bbP^1,4\Wv,p,\omega_{0,2})$ (for orientation, $4\Wv$ is the superpotential \cite[(8-1)]{Dunin-Barkowski:2015caa} with $n=4$, $a_1=\dots=a_4=0$ and $a_5=4v$): the critical points $p^5=v$ are simple with pairwise distinct critical values, and $\dd y\neq0$ there (\cref{def:mod-curve}). The three hypotheses are \cref{prop:mod-Rtr} at $\kappa=4$: symplecticity is (i), (2-32) is (ii), and the edge condition is (iii); none of them is taken from the literature. (The compatibility condition \cite[Prop.~4.3]{Dunin-Barkowski:2015caa} is equivalent to (2-32); it holds by item (2) there, as $\dd y/\dd(4\Wv)=p^2/(4(p^5-v))$ has poles only at the zeros of $\dd\Wv$, and is not used.) \Cref{prop:mod-Rtr} gives $R^{4}(\psi)=\Mo(\psi/(4\varsigma^4))$, $\Delta^{4}_j=-20\varsigma_j^2$, $\xi^{(4,j)}=-1/(\sqrt{-20}\,\varsigma_j(p-\varsigma_j))$; so \eqref{eq:mod-235} holds for $(\bbP^1,4\Wv,p)$ with $\alpha^{4}=R^{4}.\TFT^{4}$.

\emph{Step 2 (transfer to $\kappa=1$).} Take $\sqrt{-20}=2\sqrt{-5}$ and identify $e^4_j\leftrightarrow e^1_j$. Then $\xi^{(4,j)}=\xi^{(1,j)}/2$, $\dd/\dd x_4=\frac14\dd/\dd x_1$, $\TFT^4_{h,n}=2^{2h-2+n}\TFT^1_{h,n}$, the unit of $\TFT^4$ is half that of $\TFT^1$, and $R^4(\psi)=R^1(\psi/4)$ with $R^1(\psi):=\Mo(\psi/\varsigma^4)$. In the graph sum a vertex with $n_v$ special points gives $2^{2h_v-2+n_v}$, each dilaton leaf $2^{-1}$, and $\sum_v(2h_v-2+n_v)=2h-2+n+\#\{\text{dilaton leaves}\}$, so the TFT contributes $2^{2h-2+n}$. The substitution $R\mapsto R(\cdot/c)$ multiplies every leg, dilaton and edge term by $c^{-D}$, $D$ its cohomological degree on $\overline{\mathcal M}_{h,n}$ (for edges, $E_{R(\cdot/c)}(\psi',\psi'')=c^{-1}E_R(\psi'/c,\psi''/c)$), so $[R(\cdot/c).\TFT]_{h,n}=\sum_Dc^{-D}[R.\TFT]^{(D)}_{h,n}$. In \eqref{eq:mod-235} only $D=3h-3+n-\sum_ld_l$ contributes, and leg $l$ gives $4^{-d_l}2^{-1}$ times the leg for $\kappa=1$. Hence $\omega^{(4,1)}_{h,n}=2^{2h-2+n}4^{-(3h-3+n)}2^{-n}\,\mathrm{RHS}^1=4^{2-2h-n}\,\mathrm{RHS}^1$, and $\omega^{(4,1)}_{h,n}=4^{2-2h-n}\omega_{h,n}$ by \cref{lem:mod-scale}~(i). So \eqref{eq:mod-235} holds on $x=\Wv$ with $\alpha^1=\Mo(\cdot/\varsigma^4).\TFT^1$, $\Delta^1_j=-5\varsigma_j^2$, $\xi^{(1,j)}=-1/(\sqrt{-5}\,\varsigma_j(p-\varsigma_j))$.

\emph{Step 3 (the frozen basis and the factor $(-5)^{h-1}$).} With $\bar e_j:=e^1_j/\sqrt{-5}$ and $\tilde\xi^j:=\sqrt{-5}\,\xi^{(1,j)}=-1/(\varsigma_j(p-\varsigma_j))$ the summand of \eqref{eq:mod-235} is unchanged, $\TFT^1_{h,n}(\bar e_j,\dots,\bar e_j)=(-5)^{h-1}\varsigma_j^{2h-2+n}$, $\eta(\bar e_j,\bar e_k)=-\delta_{jk}/5$, and the unit is $\hat 1$ (\cref{prop:mod-Rtr}~(ii)). Identifying $\bar e_j$ with the normalised canonical vector $\tilde e_j$ of $\TFT^{\mathrm m}_v$ ($\hat\Delta_j^{1/2}=\varsigma_j$, \cref{def:set-modcohft}, the formal root $\varsigma$ there being specialised to the fifth root of $v$ fixed in \cref{def:mod-curve}), $\TFT^1_{h,n}=(-5)^{h-1}\TFT^{\mathrm m}_{v,h,n}$, with metric $-\frac15$ times that of $\TFT^{\mathrm m}_v$ and the same unit. For a CohFT $\Omega$ with metric $\eta$ and $\lambda\neq 0$, $\lambda^{h-1}\Omega_{h,n}$ is a CohFT for $\eta/\lambda$ with the same unit, and $R.(\lambda^{h-1}\Omega)=\lambda^{h-1}(R.\Omega)$: each vertex gives $\lambda^{h_v-1}$, each edge $\lambda$, and $\sum_v(h_v-1)+\#\{\text{edges}\}=h-1$. With $\lambda=-5$: $\alpha^1=(-5)^{h-1}\big(\Mo(\cdot/\varsigma^4).\TFT^{\mathrm m}_v\big)$ on the frozen basis.

\emph{Step 4 (the flat basis).} By \cref{def:set-modcohft}, $\Omega^{\mathrm{mod}}_v=(\Psi^{\mathrm m}_v)_*(\Mo(z/\varsigma^4).\TFT^{\mathrm m}_v)$, with $\eta(H^b,\tilde e_i)=\varepsilon_b\varsigma^{\nu_b-1}\zeta^{i(b+1)}$ (\cref{def:set-M0}, \cref{prop:set-modrows}~(b)). As $H^{3-a}/5$ is $\eta$-dual to $H^a$, $\sum_i\tilde e_i\otimes\tilde\xi^i=\sum_aH^a\otimes\xi'_a$ with $\xi'_a:=\frac15\sum_i\eta(H^{3-a},\tilde e_i)\tilde\xi^i$. Using $\nu_{3-a}=2-\nu_a$ and $\zeta^{-ia}=(\varsigma_i/\varsigma)^a$,
\[
\xi'_a=-\frac{\varepsilon_{3-a}}{5}\,\varsigma^{-\nu_a-a}\sum_i\frac{\varsigma_i^a}{p-\varsigma_i},\qquad\sum_i\frac{\varsigma_i^a}{p-\varsigma_i}=\begin{cases}5p^4/(p^5-v),&a=0,\\ 5vp^{a-1}/(p^5-v),&1\leq a\leq 4,\end{cases}
\]
by partial fractions. With $\nu_a+a=0,0,5,5,5$ and $\varepsilon_{3-a}=1,-1,-1,1,1$ this gives $\xi'_a=\xi_a$ of \cref{def:mod-xi}, and \eqref{eq:mod-235} becomes \eqref{eq:mod-tr}; the sum is finite since a correlator vanishes unless $\sum_ld_l\leq 3h-3+n$. Pairing each slot with a leg (\cref{def:mod-pair}) gives the second form.
\end{proof}

\input{sections/s06_model_4}



\section{Flat and hatted graph sums}
\label{sec:flh}

This section proves the comparison theorem, \cref{thm:flh-main}: for every $g\geq 2$ and every leg-free stable graph $\Gamma$ of genus $g$, the outer term $T_\Gamma$ of the natural form (\cref{def:set-outer}) equals the term $F_\Gamma$ of $\Gamma$ in the flat sum of $\Pc.\Omega^{\mathrm{mod}}_{v(u)}$ (\cref{def:red-flatsum}), provided the normalisation series $\Dn$ of \cref{thm:gau-R0} satisfies $\Fg{g}$ (\cref{def:gau-Fg}), which is proved for every $g$ in \cref{thm:gam-D}. The argument is graph by graph; it uses no total potential and no topological recursion.

Throughout, a conifold frame (\cref{def:set-frames}) is fixed, the rows are expanded in it (\cref{def:set-confexp}, \cref{cor:gau-rowsNN}), and
\[
M\coloneqq\Mo(\tau(u)\zh),\qquad \zh=z/s^4,
\]
is read, exactly in $u$, as a power series in $z$ with coefficients in $\bbQ(\zeta)\lser{s}$. By \cref{def:set-M0} and \cref{lem:set-M0sympl}, $M(0)=1$ and $M(-z)^TM(z)=1$. As in \cref{def:set-outer}, $R'\coloneqq\Qr M^{-1}$, where $\Qr$ is the matrix of the lowered rows of \cref{def:set-rows} (flat index $a$, colour $i$). Two local letters denote diagonal matrices acting on lowered components: $\Lambda_s\coloneqq\diag(s^{1-\nu_a})$ and $\Lambda_\varsigma\coloneqq\diag(\varsigma^{1-\nu_a})$; they are not the leg map of \cref{def:mod-weights}.

\subsection{Outer sums and four lemmas}
\label{sec:flh-lemmas}

\begin{defn}[outer sums with a non-unital constant part]
\label{def:flh-outer}
Let $\cA$ be $\bbQ\fps{u}$ or $\bbQ(\zeta)\lser{s}$; the second ring contains the first through $u=s^5$, and both carry the $u$-adic topology.
\begin{enumerate}
\item Let $\Omega$ be a CohFT with flat unit $1$ \cite[Def.~0.3]{PPZ15} on a free $\cA$-module $\mathbb V$ of finite rank with a perfect symmetric pairing $\eta$ (its Gram matrix is invertible). Its correlators and its dilaton-shifted potentials (\cref{lem:set-dilaton}) are
\[
\langle\psi^{k_1}a_1\cdots\psi^{k_n}a_n\rangle^{\Omega}_h=\int_{\overline{\mathcal M}_{h,n}}\Omega_{h,n}(a_1\otimes\cdots\otimes a_n)\,\psi_1^{k_1}\cdots\psi_n^{k_n},\qquad\overline F^{\Omega}_h(q)=\sum_{n}\frac{1}{n!}\langle(q+z1)^n\rangle^{\Omega}_h ,
\]
a point $q\in\mathbb V\fps{z}$ being inserted through $z^ka\mapsto\psi^ka$. For a leg-free stable graph $\Gamma$, with vertices $v$ of genus $g_v$ and valence $n_v$, a bivector series $E\in(\mathbb V\otimes\mathbb V)\fps{\psi,\psi'}$ that is symmetric under the exchange of its two slots together with $\psi$ and $\psi'$, $E(\psi,\psi')^T=E(\psi',\psi)$, and a point $q$, let $\mathrm{Cont}_\Gamma[E,q;\Omega]$ be $|\Aut\Gamma|^{-1}$ times the contraction of the vertex tensors $\dd^{n_v}\overline F^{\Omega}_{g_v}(q)$, one slot for each half-edge, with one copy of $E$ for each edge, its two slots inserted at the two half-edges of that edge. As in \cref{def:set-graphsum}, the factor $|\Aut\Gamma|^{-1}$ is part of $\mathrm{Cont}_\Gamma$. The potentials of a CohFT are always shifted by its own flat unit.
\item If $q+z1\in z^2\mathbb V\fps{z}$, every contraction is a finite sum: at a vertex each insertion of $q+z1$ has $\psi$-degree at least $2$ and raises the dimension by one, so only finitely many such insertions, and finitely many monomials of $E$, survive the dimension count. If $q+z1\in z\epsilon\,1+z^2\mathbb V\fps{z}$ with $\epsilon\in u\bbQ\fps{u}$, the vertex tensors are the Taylor expansions at the base point $-z1$; by the dilaton equation the terms with $m$ insertions of $z\epsilon 1$ are rational multiples of $\epsilon^m$ times terms without them, so these expansions converge $u$-adically (\cref{lem:set-dilaton}~(iii); for the flat model correlators this is \cref{lem:tra-conv}~(f)). Only these two cases occur below.
\item Let $(\mathbb V',\eta',1')$ be a second such module, of the same rank as $\mathbb V$, and $O(z)=\sum_{k\geq0}O_kz^k$, with $O_k\colon\mathbb V\to\mathbb V'$ linear, a symplectic series: $O(-z)^*O(z)=1$, where $O^*\colon\mathbb V'\to\mathbb V$ is the adjoint for $\eta$ and $\eta'$; thus $O^{-1}(z)=O(-z)^*$. Assume
\[
\text{(unit)}\qquad O_0^{-1}1'=\varkappa^{-1}1\quad\text{for a unit }\varkappa\text{ of }\cA .
\]
The edge and the leaf point of $O$ are
\[
E_O(\psi,\psi')\coloneqq\frac{\eta^{-1}-\bigl(O^{-1}(\psi)\otimes O^{-1}(\psi')\bigr)\eta'^{-1}}{\psi+\psi'},\qquad q_O(z)\coloneqq O^{-1}(z)(-z1'),
\]
where $\eta^{-1}\in\mathbb V\otimes\mathbb V$ and $\eta'^{-1}\in\mathbb V'\otimes\mathbb V'$ are the inverse bivectors. By symplecticity the numerator vanishes at $\psi'=-\psi$, so $E_O$ is a power series in $\psi,\psi'$, symmetric in the sense of (i) \cite[Rem.~2.1 and (6)]{PPZ15}; by (unit), $q_O\in-z\varkappa^{-1}1+z^2\mathbb V\fps{z}$.
\item For a unit $\varkappa$ of $\cA$ the CohFT $\varkappa.\Omega$ on $(\mathbb V,\eta)$ is $(\varkappa.\Omega)_{g,n}\coloneqq\varkappa^{2g-2+n}\,\Omega_{g,n}$. Its flat unit is $\varkappa^{-1}1$; the gluing axioms hold because the exponent $2g-2+n$ is additive under the gluing maps.
\end{enumerate}
\end{defn}

If $\mathbb V'=\mathbb V$ and $O(0)=1$, then $\varkappa=1$, $q_O+z1=\ttr_O(z)$ is the translation of \cite[Def.~2.13]{PPZ15}, and $\mathrm{Cont}_\Gamma[E_O,q_O;\Omega]$ is, for $\Gamma$ of genus $g\geq2$, the term of $\Gamma$ in $\int_{\overline{\mathcal M}_g}(O.\Omega)_{g,0}$ for the unit-preserving action $O.\Omega$ (\cref{def:set-graphsum}, \cref{lem:set-pushins}~(iii)). Two series with a non-unital constant part occur below. The canonical gauge $\Pc$ satisfies (unit) with $\varkappa=d$ (\cref{lem:flh-inputs}); as its constant term is an isometry, its edge and leaf point are those of \cref{def:tra-edge}. The outer matrix $R'$, from the colour space of $\TFT$ to the state space of the formal quintic, satisfies (unit) with $\varkappa=1$ (the condition involves only the constant term, the isometry $\Psi_{\mathrm Q}$, and makes sense without symplecticity), and $T_\Gamma=\mathrm{Cont}_\Gamma[E_{R'},q_{R'};M.\TFT]$, with the edge of \cref{def:set-outer}. The definitions do not decide whether $R'$ is symplectic; under $\Fg{g}$ it agrees modulo $z^{3g-2}$ with the symplectic series $\Pc\Psi^{\mathrm m}$ (\cref{lem:flh-inputs}~(ii), (iii)), which is all that the proof of \cref{thm:flh-main} uses.

\begin{lem}[frame change]\genusstatus{AI-proposed, unchecked}
\label{lem:flh-frame}
Let $A\colon(\widetilde{\mathbb V},\tilde\eta)\to(\mathbb V,\eta)$ be an isometry independent of $z$, $\widetilde\Omega$ a CohFT with flat unit $\tilde 1$ on $\widetilde{\mathbb V}$, and $A_*\widetilde\Omega$ the CohFT on $(\mathbb V,\eta)$ with $(A_*\widetilde\Omega)_{g,n}\coloneqq\widetilde\Omega_{g,n}\circ(A^{-1})^{\otimes n}$, whose flat unit is $A\tilde1$. For every leg-free stable graph $\Gamma$ and all $E$ and $q$ as in \cref{def:flh-outer}~(i),
\[
\mathrm{Cont}_\Gamma[E,q;A_*\widetilde\Omega]=\mathrm{Cont}_\Gamma\bigl[(A^{-1}\otimes A^{-1})E,\;A^{-1}q;\;\widetilde\Omega\bigr].
\]
\end{lem}

\begin{proof}
Since $A^{-1}(q+zA\tilde1)=A^{-1}q+z\tilde1$, insertion by insertion $\overline F^{A_*\widetilde\Omega}_h(q)=\overline F^{\widetilde\Omega}_h(A^{-1}q)$, and hence $\dd^n\overline F^{A_*\widetilde\Omega}_h(q)[w_1,\dots,w_n]=\dd^n\overline F^{\widetilde\Omega}_h(A^{-1}q)[A^{-1}w_1,\dots,A^{-1}w_n]$. Every term of $\mathrm{Cont}_\Gamma$ is a full contraction of the vertex tensors with the edge bivectors; writing the identity as $AA^{-1}$ on every half-edge turns the contraction with $E$ into the contraction with $(A^{-1}\otimes A^{-1})E$. The factor $|\Aut\Gamma|^{-1}$ is the same on both sides.
\end{proof}

\begin{lem}[TFT dilaton invariance]\genusstatus{AI-proposed, unchecked}
\label{lem:flh-dil}
Let $\Omega$ be a CohFT with flat unit $1$ on $(\mathbb V,\eta)$ over $\cA$, let $\varkappa\in 1+u\bbQ\fps{u}$, let $E$ be as in \cref{def:flh-outer}~(i), and let $q\in\mathbb V\fps{z}$ with $q+z\varkappa^{-1}1\in z^2\mathbb V\fps{z}$. Then for every leg-free stable graph $\Gamma$
\[
\mathrm{Cont}_\Gamma[E,q;\varkappa.\Omega]=\mathrm{Cont}_\Gamma[E,q;\Omega].
\]
On the left the potentials are shifted by the unit $\varkappa^{-1}1$ of $\varkappa.\Omega$ and the contraction is a finite sum; on the right they are shifted by $1$, and $q+z1\in z(1-\varkappa^{-1})1+z^2\mathbb V\fps{z}$ with $1-\varkappa^{-1}\in u\bbQ\fps{u}$, so the vertex tensors are the $u$-adically convergent Taylor expansions of \cref{def:flh-outer}~(ii).
\end{lem}

\begin{proof}
By multilinearity, $\langle(q+z\varkappa^{-1}1)^n\rangle^{\varkappa.\Omega}_h=\varkappa^{2h-2+n}\langle(q+z\varkappa^{-1}1)^n\rangle^{\Omega}_h=\varkappa^{2h-2}\langle(\varkappa q+z1)^n\rangle^{\Omega}_h$. Hence $\overline F^{\varkappa.\Omega}_h(q)=\varkappa^{2h-2}\,\overline F^{\Omega}_h(\varkappa q)$ and
\[
\dd^n\overline F^{\varkappa.\Omega}_h(q)=\varkappa^{2h-2+n}\,\bigl(\dd^n\overline F^{\Omega}_h\bigr)(\varkappa q),
\]
both sides being finite sums because $\varkappa q+z1\in z^2\mathbb V\fps{z}$. The dilaton identity of \cref{lem:set-dilaton}~(iii) at the point $\varkappa q$ reads $\dd^n\overline F^{\Omega}_h(q)=\varkappa^{2h-2+n}\,\dd^n\overline F^{\Omega}_h(\varkappa q)$ for $2h-2+n>0$, which holds at every vertex of a stable graph; it applies over both rings $\cA$ because $\varkappa(0)=1$, so that its two sides are, coefficient by coefficient, $u$-adically convergent series in $1-\varkappa^{-1}\in u\bbQ\fps{u}$ times correlators of $\Omega$. So the vertex tensors of the two sides agree, and the edge factors are the same.
\end{proof}

\begin{lem}[the $R$-action commutes with the TFT rescaling]\genusstatus{AI-proposed, unchecked}
\label{lem:flh-comm}
Let $\Omega$ be a CohFT with flat unit $1$ on $(\mathbb V,\eta)$ over $\cA$, $\varkappa$ a unit of $\cA$, and $R(z)=1+O(z)$ an $\End(\mathbb V)$-valued series with $R(-z)^*R(z)=1$. For the unit-preserving action of \cite[Def.~2.13]{PPZ15},
\[
R.(\varkappa.\Omega)=\varkappa.(R.\Omega).
\]
\end{lem}

\begin{proof}
By \cite[Def.~2.13]{PPZ15}, $R.\Omega$ is the action of \cite[Def.~2.2]{PPZ15} on the translation \cite[Def.~2.5]{PPZ15} of $\Omega$ by $\ttr_R(z)=z1-zR^{-1}(z)1$, where $1$ is the unit of $\Omega$. The unit of $\varkappa.\Omega$ is $\varkappa^{-1}1$, so its translation is $\varkappa^{-1}\ttr_R$. In the graph sum for the genus-$g$, $n$-point class, a vertex of genus $g_v$ with $n_v$ half-edges and legs and $m$ translation points contributes $(\varkappa.\Omega)_{g_v,n_v+m}=\varkappa^{2g_v-2+n_v+m}\Omega_{g_v,n_v+m}$, its $m$ translation insertions contribute $\varkappa^{-m}$, and the legs and edges do not involve $\varkappa$. So every graph term is $\varkappa^{\sum_v(2g_v-2+n_v)}=\varkappa^{2g-2+n}$ times the corresponding term of $R.\Omega$, since $\sum_v(2g_v-2)+2E+n=2g-2+n$ for a connected stable graph of genus $g$ with $E$ edges and $n$ legs.
\end{proof}

\begin{lem}[hatting]\genusstatus{AI-proposed, unchecked}
\label{lem:flh-hat}
Let $g\geq 2$. For every leg-free stable graph $\Gamma$ of genus $g$, the hatted term and the outer term of \cref{def:set-outer} satisfy
\[
G_\Gamma=s^{10(g-1)}\,T_\Gamma=u^{2g-2}\,T_\Gamma .
\]
\end{lem}

We recall the hatted data of \cref{def:set-hat} in the form in which they are used. The element $\zeta^{-4i}\,\Cexp_{\alpha,\beta}(\Delta_i^{1/2})=s\varpi_0$ does not depend on $i$ ($\Delta_i^{1/2}=\zeta^{4i}C_0/L$, $\Cexp(C_0)=\varpi_0$, $\Cexp(L)=s^{-1}$); put $\hat x\coloneqq\varpi_0$. The hatted TFT is $\hat\TFT=s^{-1}.\TFT$, whose genus-$0$ three-point value at $\tilde e_i$ is $\zeta^{4i}\hat x$ and whose unit is $s$ times the unit $1_{\TFT}=\sum_i\Delta_i^{-1/2}\tilde e_i$ of $\TFT$; the hatted unit of the state space is $\hat1=s1$; the hatted rows are $\hat\Qr(u,\zh)=\Lambda_s\Qr(u,s^4\zh)$; and $\Lambda_s$ preserves $\eta$, because $\eta$ pairs $H^a$ with $H^{3-a}$ and $\nu_a+\nu_{3-a}=2$. With $\hat M(\zh)\coloneqq\Mo(\tau\zh)=M(s^4\zh)$ and $\hat R'(\zh)\coloneqq\hat\Qr(\zh)\hat M(\zh)^{-1}=\Lambda_sR'(s^4\zh)$, which is the map $R'(s^4\zh)$ written in the lowered components with respect to the $\hat H^a$ (in these, $\hat1=s1$ has the components $5\,\delta_{a3}$; $\Lambda_s$ is not applied to vectors as an isometry of $\bbH$), the hatted term is $G_\Gamma=\mathrm{Cont}_\Gamma[E_{\hat R'},q_{\hat R'};\hat M.\hat\TFT]$, with the edge and the leaf point formed in the variable $\zh$ and with the unit $\hat1$.

\begin{proof}
\emph{Edges and leaves.} Since $\Lambda_s$ preserves $\eta$, the edge of $\hat R'$ is $E_{\hat R'}(\psi,\psi')=s^4E_{R'}(s^4\psi,s^4\psi')$: its coefficient of $\psi^a\psi'^b$ is that of $E_{R'}$ times $s^{4(a+b+1)}$. With $\hat1=s1$, $q_{\hat R'}(\zh)=R'^{-1}(s^4\zh)(-\zh s1)=s^{-3}q_{R'}(s^4\zh)$, so the leaf insertion $q_{\hat R'}+\zh\,s1_{\TFT}$ equals $s^{-3}(q_{R'}+z1_{\TFT})$ at $z=s^4\zh$: its coefficient of $\zh^k$ is that of $q_{R'}+z1_{\TFT}$ at $z^k$ times $s^{4k-3}$.

\emph{Vertices.} The classes of $\TFT$ have degree $0$, and in the graph sum of \cite[Def.~2.2]{PPZ15} with the translation of \cite[Def.~2.13]{PPZ15} a term of cohomological degree $D$ uses coefficients of the $R$-matrix of total $z$-degree $D$ (the $\psi$-exponent at a leg, the $\psi$-exponents plus one at an edge, the $z$-exponent minus one at a translation point). Since $\hat M(\zh)=M(s^4\zh)$, the degree-$D$ part of $(\hat M.\TFT)_{h,m}$ is $s^{4D}$ times that of $(M.\TFT)_{h,m}$. By \cref{lem:flh-comm}, $\hat M.\hat\TFT=\hat M.(s^{-1}.\TFT)=s^{-1}.(\hat M.\TFT)$. So at a vertex of genus $g_v$ with $n_v$ half-edges and $l_v$ leaf insertions, the degree-$D_v$ part of the vertex class is multiplied by $s^{4D_v-(2g_v-2+n_v+l_v)}$.

\emph{The count.} Take a term of $\Gamma$ with edge monomials $\psi^{a_e}\psi'^{b_e}$, leaf insertions of degrees $k_l$ and vertex classes of degrees $D_v$, and let $E$ and $n_{\mathrm l}$ be the numbers of edges and of leaf insertions. At each vertex the integrand has the degree $\dim\overline{\mathcal M}_{g_v,n_v+l_v}$, so
\[
\sum_vD_v+\sum_e(a_e+b_e)+\sum_lk_l=\sum_v(3g_v-3+n_v+l_v)=3g-3-E+n_{\mathrm l} .
\]
By the two previous paragraphs the term is multiplied by $s^N$, where, using $\sum_v(2g_v-2+n_v)=2g-2$,
\begin{align*}
N&=\sum_e4(a_e+b_e+1)+\sum_l(4k_l-3)+\sum_v\bigl(4D_v-(2g_v-2+n_v+l_v)\bigr)\\
&=4(3g-3-E+n_{\mathrm l})+4E-3n_{\mathrm l}-(2g-2)-n_{\mathrm l}=10(g-1).
\end{align*}
Hence $G_\Gamma=s^{10(g-1)}T_\Gamma=u^{2g-2}T_\Gamma$.
\end{proof}

The lemma uses no property of the rows beyond \cref{def:set-hat,def:set-outer}. Once the hatted data are known to be power series in $u$ (\cref{cor:mai-natfact}), it shows that $T_\Gamma$ has a pole of order at most $2g-2$ at $u=0$.

\begin{lem}[the inputs of the comparison]\genusstatus{AI-proposed, unchecked}
\label{lem:flh-inputs}
Let $\Qr^{\mathrm m}$ be the matrix of the model columns at $v(u)$ (\cref{thm:gau-R0}), $\Psi^{\mathrm m}\coloneqq\Qr^{\mathrm m}|_{z=0}$, $\Psi_{\mathrm Q}\coloneqq\Qr|_{z=0}$, $P_0\coloneqq\Pc|_{z=0}$ and $d\coloneqq(P_0)_{33}$, and identify the colour bases of the formal quintic and of the model through the index $i$, as in \cref{cor:gau-colours}.
\begin{enumerate}
\item $\Qr^{\mathrm m}=\Psi^{\mathrm m}M$ and $\Omega^{\mathrm{mod}}_{v(u)}=(\Psi^{\mathrm m})_*\bigl(M.\TFT^{\mathrm m}_{v(u)}\bigr)$; $\Psi^{\mathrm m}$ is an isometry, independent of $z$, from the colour space with the pairing $\delta$ to $(\bbH,\eta)$, and it maps the unit of $\TFT^{\mathrm m}_{v(u)}$ to $H^0$.
\item $\Qr\Dn=\Pc\Qr^{\mathrm m}$ with $\Dn=\diag(\Dn_i)$ and $\Dn_i(z)=\Dn_0(\zeta^{-i}z)$. Under $\Fg{g}$, $R'\equiv\Pc\Psi^{\mathrm m}\bmod z^{3g-2}$.
\item $\Pc$ is symplectic; so $E_{\Pc}$ is a power series and $\Pc\Psi^{\mathrm m}$ is symplectic. Moreover
\[
R'(-z)^T\eta^{-1}R'(z)=M(z)\bigl(\Dn(z)\Dn(-z)\bigr)^{-1}M(z)^{-1};
\]
hence $R'$ is symplectic, $R'(-z)^T\eta^{-1}R'(z)=1$, if and only if $\Dn_0(z)\Dn_0(-z)=1$, and under $\Fg{g}$ it is symplectic modulo $z^{3g-2}$.
\item $\Pc$ satisfies (unit) with $\varkappa=d$, where $d\in\bbQ\fps{u}$ and $d(0)=1$; $R'$ satisfies (unit) with $\varkappa=1$; $\Psi_{\mathrm Q}=P_0\Psi^{\mathrm m}$; and $\TFT=d.\TFT^{\mathrm m}_{v(u)}$ on the identified colour spaces.
\end{enumerate}
\end{lem}

\begin{proof}
(i) By \cref{prop:set-modrows}~(b), \cref{def:set-M0} and \cref{lem:gau-colmod}, the colour-$i$ model column at $v(u)$, with its exponential stripped, has the entries $\varsigma^{\nu_b-1}\Rhat^{\mathrm m,(i)}_b(z/\varsigma^4)$, and $z/\varsigma^4=\tau\zh$ because $\varsigma=s\tau^{-1/4}$ (\cref{def:set-tau}). As $\Rhat^{\mathrm m}=\Psi_0\Mo$ (\cref{def:set-M0}), $\Qr^{\mathrm m}=\Lambda_\varsigma^{-1}\Psi_0\Mo(\tau\zh)=\Psi^{\mathrm m}M$ with $\Psi^{\mathrm m}=\Lambda_\varsigma^{-1}\Psi_0$. From $\Psi_0^T\eta^{-1}\Psi_0=1$ (\cref{def:set-M0}) and $\Lambda_\varsigma^{-1}\eta^{-1}\Lambda_\varsigma^{-1}=\eta^{-1}$ (again $\nu_a+\nu_{3-a}=2$) we get $\Psi^{\mathrm m,T}\eta^{-1}\Psi^{\mathrm m}=1$: on lowered components the pairing is $\eta^{-1}$, so $\Psi^{\mathrm m}$ is an isometry. By \cref{def:set-modcohft}, $\Omega^{\mathrm{mod}}_{v(u)}=(\Psi^{\mathrm m})_*(\Mo(z/\varsigma^4).\TFT^{\mathrm m}_{v(u)})$, and $\Mo(z/\varsigma^4)=M$. The action of $M$ preserves the unit of $\TFT^{\mathrm m}$ \cite[Prop.~2.12]{PPZ15}, and $\Omega^{\mathrm{mod}}_{v(u)}$ has the flat unit $H^0$ (\cref{lem:set-dilaton}~(ii)); so $\Psi^{\mathrm m}$ maps the unit of $\TFT^{\mathrm m}$ to $H^0$.

(ii) The identity is \cref{thm:gau-R0}~(iv) with \cref{cor:gau-colours}. Under $\Fg{g}$, $\Dn_0\equiv1\bmod z^{3g-2}$, hence $\Dn\equiv1\bmod z^{3g-2}$; from $\Qr=\Pc\Psi^{\mathrm m}M\Dn^{-1}$ we get $R'=\Pc\Psi^{\mathrm m}\,(M\Dn^{-1}M^{-1})\equiv\Pc\Psi^{\mathrm m}\bmod z^{3g-2}$.

(iii) By \cref{prop:gau-symp}, $\Pc(-z)^T\eta\,\Pc(z)=\eta$, equivalently $\Pc(-z)^T\eta^{-1}\Pc(z)=\eta^{-1}$, since $\eta^{-1}=\eta/25$ (\cref{conv:tra-forms}). So $E_{\Pc}$ is a power series (\cref{def:flh-outer}~(iii), \cref{def:tra-edge}), and $S\coloneqq\Pc\Psi^{\mathrm m}$ (a local letter) satisfies $S(-z)^T\eta^{-1}S(z)=\Psi^{\mathrm m,T}\eta^{-1}\Psi^{\mathrm m}=1$, because $\Psi^{\mathrm m}$ is an isometry independent of $z$ (by (i)): $\Pc\Psi^{\mathrm m}$ is symplectic. By the proof of (ii), $R'=S\,K$ with $K(z)\coloneqq M(z)\Dn(z)^{-1}M(z)^{-1}$. From $M(-z)^TM(z)=1$ we get $M(-z)^T=M(z)^{-1}$ and $M(-z)^{-T}=M(z)$; as $\Dn$ is diagonal, $K(-z)^T=M(z)\Dn(-z)^{-1}M(z)^{-1}$, and
\[
R'(-z)^T\eta^{-1}R'(z)=K(-z)^T\,S(-z)^T\eta^{-1}S(z)\,K(z)=K(-z)^TK(z)=M(z)\bigl(\Dn(-z)\Dn(z)\bigr)^{-1}M(z)^{-1}.
\]
This is $1$ if and only if $\Dn_i(z)\Dn_i(-z)=1$ for every $i$, that is, as $\Dn_i(z)=\Dn_0(\zeta^{-i}z)$, if and only if $\Dn_0(z)\Dn_0(-z)=1$. Under $\Fg{g}$, $\Dn\equiv1\bmod z^{3g-2}$, so the right side is $\equiv1\bmod z^{3g-2}$.

(iv) The canonical gauge is graded (\cref{def:gau-can}), so $P_0$ is diagonal and, in vector form, maps $H^0$ to $dH^0$ (\cref{def:tra-dressing}: $d=(P_0)_{33}=(P_0^{\mathrm{vec}})_{00}$, the vector form being that of \cref{conv:tra-forms}): this is (unit) with $\varkappa=d$. By \cref{lem:gau-diag}, $p_{0,a}(0)=1$ for every $a$ under (N0), which holds in the conifold frames; so $d(0)=1$, and $d\in\bbQ\fps{u}$ by \cref{thm:gau-U}. Since $M(0)=1$, the constant term of $R'$ is $\Psi_{\mathrm Q}$. Its entries are the conifold expansions of the representatives of \cref{cor:gau-rowsNN}~(i) at $n=0$: there row $0$ of colour $0$ is $A^{\mathrm{act}}_0c_0^{-1}=L/C_0$ (\cref{thm:gau-R0}~(iii)), the rows $a\geq1$ follow from $\Qr_{a+1,0}=(C_1\kappa^{\mathrm{LR}}_a)^{-1}L\,\Qr_{a,0}$ with $C_1\kappa^{\mathrm{LR}}=(C_1,L^5/(C_0^2C_1^2),C_1,C_0,C_0)$, and colour $j$ is $\zeta^{j(a+1)}$ times colour $0$. So the representative in $\Nring\otimes\bbQ(\zeta)$ of $(\Psi_{\mathrm Q})_{aj}$ (denoted $\Psi_{\mathrm Q}$ in \cref{lem:mai-Psi}, its expansion being $\Psi_{\mathrm{cf}}$ there) is $\zeta^{j(a+1)}\mathfrak m_a$ with $(\mathfrak m_0,\dots,\mathfrak m_4)=(L/C_0,\,L^2/(C_0C_1),\,C_0C_1/L^2,\,C_0/L,\,1)$ (local letters). With $\Delta_j^{-1/2}=\zeta^jL/C_0$ (\cref{cor:gau-rowsNN}~(ii)),
\[
\textstyle\sum_j\Delta_j^{-1/2}(\Psi_{\mathrm Q})_{aj}=\frac L{C_0}\,\mathfrak m_a\sum_j\zeta^{j(a+2)}=5\,\delta_{a3}=\eta(H^a,H^0)
\]
in $\Nring\otimes\bbQ(\zeta)$, since $\mathfrak m_3=C_0/L$. Its conifold expansion says that $\Psi_{\mathrm Q}$ maps the unit $1_{\TFT}=\sum_j\Delta_j^{-1/2}\tilde e_j$ of $\TFT$ (\cref{def:set-frob}) to $H^0=1$: this is (unit) with $\varkappa=1$. At $z^0$ the identity of (ii) reads $\Psi_{\mathrm Q}=P_0\Psi^{\mathrm m}$, because $\Dn(0)=1$ (\cref{thm:gau-R0}); in particular $\Psi_{\mathrm Q}$ is an isometry, $P_0$ and $\Psi^{\mathrm m}$ being isometries ($P_0^T\eta P_0=\eta$ is the part of order $z^0$ of \cref{prop:gau-symp}). Therefore $1_{\TFT}=\Psi_{\mathrm Q}^{-1}H^0=(\Psi^{\mathrm m})^{-1}P_0^{-1}H^0=d^{-1}(\Psi^{\mathrm m})^{-1}H^0$ (in lowered components $H^0$ is $5e_3$ and $(P_0)_{33}=d$), which is $d^{-1}$ times the unit of $\TFT^{\mathrm m}$ by (i). Both TFTs are diagonal in their orthonormal colour bases and vanish on mixed colours, and such a TFT is determined by its unit $\sum_iy_i\tilde e_i$: its genus-$g$ value at $\tilde e_i^{\otimes n}$ is $y_i^{-(2g-2+n)}$ (with $y_i=\Delta_i^{-1/2}$ for $\TFT$, \cref{def:set-frob}, and $y_i=\hat\Delta_i^{-1/2}$ for $\TFT^{\mathrm m}$, \cref{def:set-modcohft}). Hence $\TFT_{g,n}=d^{2g-2+n}\,\TFT^{\mathrm m}_{g,n}$, that is $\TFT=d.\TFT^{\mathrm m}$.
\end{proof}

The symplecticity of $R$ in \cref{def:set-rows}~(a) is an identity between power series in $q$ and is not used in (iii): an identity between evaluations under $\ev$ of elements of $\Nring\otimes\bbQ(\zeta)$ survives the conifold expansion only where $\ev$ is known to be injective, which is neither claimed nor needed. This is why (iii) carries the factor $\Dn$ and why the unit identity $\sum_j\Delta_j^{-1/2}(\Psi_{\mathrm Q})_{aj}=5\,\delta_{a3}$ in (iv) is proved in $\Nring\otimes\bbQ(\zeta)$ itself; the other identities at $z^0$ in (iv), $\Psi_{\mathrm Q}=P_0\Psi^{\mathrm m}$ and the isometry, are identities in $\bbQ(\zeta)\lser s$ obtained from (ii).

\subsection{The comparison theorem}
\label{sec:flh-thm}

\begin{thm}[the comparison theorem]\genusstatus{AI-proposed, unchecked}
\label{thm:flh-main}
Let $g\geq 2$ and let $\Gamma$ be a leg-free stable graph of genus $g$. Let $F_\Gamma$ be its term in the flat sum of $\Pc.\Omega^{\mathrm{mod}}_{v(u)}$ (\cref{def:red-flatsum}), $T_\Gamma$ its outer term and $G_\Gamma$ its hatted term (\cref{def:set-outer}, with $M=\Mo(\tau(u)\zh)$ exactly in $u$). Assume $\Fg{g}$. Then
\begin{equation}
\label{eq:flh-main}
T_\Gamma=F_\Gamma,\qquad\text{hence}\qquad G_\Gamma=u^{2g-2}F_\Gamma .
\end{equation}
\end{thm}

\begin{proof}
\emph{Step 1: only $R'$ through $z^{3g-3}$ enters.} At a vertex of genus $g_v$ with $n_v$ half-edges and $l_v$ leaf insertions, the classes of $M.\TFT$ have degree $\geq0$, so only insertions of total $\psi$-degree at most $3g_v-3+n_v+l_v$ contribute; summed over the vertices this bound is $3g-3-E+n_{\mathrm l}$, where $E$ and $n_{\mathrm l}$ are the numbers of edges and of leaf insertions. Since $R'$ satisfies (unit) with $\varkappa=1$ (\cref{lem:flh-inputs}~(iv)), every leaf insertion $q_{R'}+z1_{\TFT}$ has degree at least $2$. Hence a monomial $\psi^a\psi'^b$ of an edge occurs only with $a+b\leq 3g-3-E-n_{\mathrm l}\leq 3g-4$, and a leaf insertion of degree $k$ only with $k\leq 3g-1-E-n_{\mathrm l}\leq 3g-2$. With $N(\psi,\psi')\coloneqq1-\bigl(R'^{-1}(\psi)\otimes R'^{-1}(\psi')\bigr)\eta^{-1}$, $\eta^{-1}$ being the inverse bivector of $(\bbH,\eta)$ as in \cref{def:set-outer}, the edge of \cref{def:set-outer} has the numerator $\tilde N\coloneqq N(\psi,\psi')-\frac12N(-\psi',\psi')-\frac12N(\psi,-\psi)$ (a local letter), which vanishes at $\psi'=-\psi$. Writing $\tilde N=\sum_{m,n}\tilde N_{mn}\bigl(\psi^m-(-\psi')^m\bigr)\psi'^n$ and dividing by $\psi+\psi'$, the coefficient of $\psi^a\psi'^b$ in $E_{R'}$ is a combination of the coefficients of $N$ of total degree $a+b+1$, so it involves $R'^{-1}$ only through $z^{a+b+1}$; the coefficient of $z^k$ in $q_{R'}$ involves $R'^{-1}$ only through $z^{k-1}$. So $T_\Gamma$ depends on $R'$ only modulo $z^{3g-2}$, through the part of $N$ of total degree at most $3g-3$ and the part of $q_{R'}$ of degree at most $3g-2$.

\emph{Step 2: the outer data.} Put $R''\coloneqq\Pc\Psi^{\mathrm m}$; it is symplectic (\cref{lem:flh-inputs}~(iii)), so the two correction terms of its numerator vanish, and $E_{R''}$ is the edge $E_O$ of \cref{def:flh-outer}~(iii) for $O=R''$, a power series by the argument of \cite[Rem.~2.1 and (6)]{PPZ15} recalled there. By \cref{lem:flh-inputs}~(ii), $R'\equiv R''\bmod z^{3g-2}$, so the parts of total degree at most $3g-3$ of the numerators of $E_{R'}$ and $E_{R''}$ agree, and $q_{R'}\equiv q_{R''}\bmod z^{3g-1}$. Hence the coefficients of $E_{R'}$ of total degree at most $3g-4$, the only ones that enter by Step 1, equal those of $E_{R''}$; by Step 1, $T_\Gamma=\mathrm{Cont}_\Gamma[E_{R''},q_{R''};M.\TFT]$. Since $\Psi^{\mathrm m}$ is an isometry independent of $z$, $R''^{-1}=(\Psi^{\mathrm m})^{-1}\Pc^{-1}$ gives $E_{R''}=((\Psi^{\mathrm m})^{-1}\otimes(\Psi^{\mathrm m})^{-1})E_{\Pc}$ and $q_{R''}=(\Psi^{\mathrm m})^{-1}q_{\Pc}$; and $q_{R''}+zd^{-1}1_{\TFT^{\mathrm m}}=O(z^2)$, where $1_{\TFT^{\mathrm m}}$ is the unit of $\TFT^{\mathrm m}$, because $R''_0=P_0\Psi^{\mathrm m}$ maps $d^{-1}1_{\TFT^{\mathrm m}}$ to $H^0$ (\cref{lem:flh-inputs}~(i), (iv)).

\emph{Step 3: the chain.}
\begin{align*}
T_\Gamma&=\mathrm{Cont}_\Gamma[E_{R''},q_{R''};M.\TFT]&&\\
&=\mathrm{Cont}_\Gamma[E_{R''},q_{R''};M.(d.\TFT^{\mathrm m})]&&\text{(\cref{lem:flh-inputs}~(iv))}\\
&=\mathrm{Cont}_\Gamma[E_{R''},q_{R''};d.(M.\TFT^{\mathrm m})]&&\text{(\cref{lem:flh-comm}, $M(0)=1$)}\\
&=\mathrm{Cont}_\Gamma[E_{R''},q_{R''};M.\TFT^{\mathrm m}]&&\text{(\cref{lem:flh-dil}, $\varkappa=d$, $d(0)=1$)}\\
&=\mathrm{Cont}_\Gamma[E_{\Pc},q_{\Pc};(\Psi^{\mathrm m})_*(M.\TFT^{\mathrm m})]&&\text{(\cref{lem:flh-frame}, $A=\Psi^{\mathrm m}$, Step 2)}\\
&=\mathrm{Cont}_\Gamma[E_{\Pc},q_{\Pc};\Omega^{\mathrm{mod}}_{v(u)}]&&\text{(\cref{lem:flh-inputs}~(i))}.
\end{align*}
The hypothesis of \cref{lem:flh-dil} on the point holds by Step 2; the unit of $d.(M.\TFT^{\mathrm m})$ is $d^{-1}1_{\TFT^{\mathrm m}}=1_{\TFT}$, so the first three lines are finite sums.

\emph{Step 4: the flat term.} The last contraction is formed with the leaf point $q_{\Pc}=q_T$ of \cref{def:tra-edge} and with the Taylor expansions of the potentials of $\Omega^{\mathrm{mod}}_{v(u)}$ at $-z1$ in the direction $q_T+z1\in z(1-d^{-1})1+z^2\bbH\fps{z}$. Since $d(0)=1$, these expansions converge $u$-adically and equal the homogeneous vertex tensors $d^{2g_v-2+n_v}\dd^{n_v}\overline F_{g_v}(dq_T)$ of \cref{def:red-flatsum} (\cref{lem:tra-conv}~(f), \cref{lem:set-dilaton}). So the last line is $F_\Gamma$. Finally $G_\Gamma=u^{2g-2}T_\Gamma$ by \cref{lem:flh-hat}.
\end{proof}

By Step 1, at genus $g$ the outer terms involve the normalisation series only through $z^{3g-3}$: $\Fg{g}$, the vanishing of $\coef{z^k}\Dn_0$ for $1\leq k\leq 3g-3$, is all of $\Dn_0=1$ that genus $g$ sees (\cref{rmk:con-genus}).

\begin{rmk}[where the comparison can fail]
\label{rmk:flh-fail}
(i) \emph{The normalisation.} In general the proof of \cref{lem:flh-inputs}~(ii) gives $R'=\Pc\Psi^{\mathrm m}\,M\Dn^{-1}M^{-1}$. If $\Dn=\cn(z)1$ is scalar (then $\cn=\Dn_0$ lies in $\bbQ\fps{z^5}$, by $\Dn_i(z)=\Dn_0(\zeta^{-i}z)$ and \cref{thm:con-glue}) with $\cn(z)\cn(-z)=1$, then $R'=(\Pc\cn(z)^{-1})\Psi^{\mathrm m}$ with the symplectic graded regular gauge $\Pc\cn(z)^{-1}$ (\cref{prop:gau-symp}), and the proof identifies $T_\Gamma$ with the term of the flat sum of that gauge, a different graph decomposition (\cref{rmk:mai-stored}, \cref{rmk:mai-graphlocal}). (ii) \emph{The split.} For another symplectic $M'$ with $M'(0)=1$ the outer sum $(\Qr M'^{-1}).(M'.\TFT)$ has the same total (by the regrouping in the proof of \cref{cor:mai-gap}) but in general different graph terms (no example is recorded here). (iii) \emph{The rows formula.} The colour-basis form of the edge, $E_{R'}(\psi,\psi')=\bigl(1-R'(-\psi)^T\eta^{-1}R'(-\psi')\bigr)/(\psi+\psi')$, uses the symplecticity of $R'$, which holds modulo $z^{3g-2}$ under $\Fg{g}$ and exactly once $\Dn=1$; it is not inferred from the symplecticity of $R$ at large radius.
\end{rmk}


% Section 8 (part 1): the connection problem. Label prefix: con.
% No new macros. Letters used in one proof only are declared there and are local to it.

\section{The connection problem}
\label{sec:con}

The canonical gauge $\Pc$ of \cref{sec:gau} is normalised at the conifold point $u=0$, the rows of the formal quintic at $q=0$; the normalisation series $\Dn_0$ of \cref{thm:gau-R0} measures the difference, and $\Fg{g}$ for every $g$ (\cref{def:gau-Fg}) is the statement $\Dn_0=1$. This section reduces $\Dn_0=1$ to one identity at $q=0$. We define the full-Stirling saddle series $c(z;L)$ of a Mellin--Barnes summand and the hatted series $\PhiSP$ (\cref{sec:con-form}, after the conventions of \cref{sec:con-problem}). Theorem~P (\cref{sec:con-P}) shows that $\PhiSP$ solves HPF with polynomial coefficients (part~(a)) and, by the residue of a $1$-form at its pole, that it is $\Phican$ (part~(b)). The glue (\cref{sec:con-glue}) gives $\Dn_0(z)=\PhiSP(z,1)/h(z)$; the value $\PhiSP(z,1)=c(z;1)$ is determined in \cref{sec:gam}. We work in colour~$0$; the other colours follow by \cref{cor:gau-colours}.

\subsection{Conventions}
\label{sec:con-problem}

A colour-$0$ WKB function is a series $\Phi(\zh,u)=\sum_{n\geq 0}\mathrm{ph}_n(u)\zh^n$ with $\mathrm{ph}_0=1$ (\cref{def:set-points,def:set-rows}). We write $\PFfive^{L}$ for $\PFfive$ with $z\Dq$ replaced by $z\Dq+L$, that is, $\PFfive^{L}=\re^{-u_0/z}\,\PFfive\,\re^{u_0/z}$ with $\Dq u_0=L$; it is the operator $\PFfive^{(0)}$ of \cref{lem:gau-HPF,lem:gau-wkb}. At fixed $z$ one has $u\,\dd/\dd u=\theta_u-\tfrac45\theta_{\zh}$ and $z/u=\zh/s$, so that for $e\in\bbZ$ and $F$ in $\bbQ\lser{s}\fps{\zh}$
\begin{equation}
\label{eq:con-chain}
(z\Dq+L)(s^eF)=s^{e-1}\mathscr O_eF,\qquad \mathscr O_e=1+\zh(1-u)\bigl(\tfrac45\theta_{\zh}-\tfrac e5-\theta_u\bigr),
\end{equation}
the chain rule of \cref{lem:gau-HPF}. Write $\PFfive=u(z\Dq)^5-1-(1-u)\sum_{m=1}^5e_mz^m(z\Dq)^{5-m}$, with $e_m$ the elementary symmetric functions of $\tfrac15,\dots,\tfrac55$. Applying \eqref{eq:con-chain} five times to $s^{-1}\Phi$ gives $\PFfive^{L}(s^{-1}\Phi)=s^{-1}\,\mathrm{HPF}(\Phi)$, where $\mathrm{HPF}(\Phi)$ is the left side of \eqref{eq:gau-HPF}; so $L\re^{u_0/z}\Phi$ solves $\PFfive$ if and only if $\Phi$ solves HPF, in $\bbQ\lser{s}\fps{\zh}$. The free coefficients of $\Phi$ are $\fc_m=[u^{4m}]\mathrm{ph}_{5m}$ (\cref{lem:gau-HPF}).

\begin{rmk}[Genus bookkeeping]
\label{rmk:con-genus}
By \cref{def:gau-Fg}, $\Fg{g}$ says that $\Dn_0-1$ has no terms $z^1,\dots,z^{3g-3}$; the coefficient of $z^n$ enters $\Fg{g}$ first at $g=\lceil(n+3)/3\rceil$. The identity $\Dn_0(z)\Dn_0(-z)=1$, the symplecticity of $R'$ (\cref{lem:flh-inputs}~(iii)), would halve the conditions. Deriving it from the symplecticity of $R$ (\cref{def:set-rows}~(a)) would need that symplecticity for the representatives of the rows in $\Nring\otimes\bbQ(\zeta)$ (\cref{cor:gau-rowsNN}), which is not claimed. So it is not used; it follows from $\Dn_0=1$ (\cref{thm:gam-D}).
\end{rmk}

% Section 8 (part 3): the full-Stirling series, its rationality, the 1-form and PhiSP.

\subsection{The full-Stirling series, its rationality, the \texorpdfstring{$1$}{1}-form and \texorpdfstring{$\PhiSP$}{the full-Stirling WKB function}}
\label{sec:con-form}

\begin{defn}[The full-Stirling saddle series]
\label{def:gam-c}
Let $\Cstar:=\bbC\setminus\{0,1,\zeta,\zeta^2,\zeta^3,\zeta^4\}$ and $q(L):=(1-L^{-5})/5^5$; no fifth root is taken. For $L\in\Cstar$ fix an open disc $\bbD_L$ around $L$ containing neither $0$ nor any $\zeta^j$, holomorphic branches of $\log(5\sigma)$ and of $\log(\zeta^j-\sigma)$, $j=0,\dots,4$, on $\bbD_L$, and put
\[
\lambda:=-5\log(5L)+\sum_{j=0}^4\log(\zeta^j-L),\qquad\text{so that}\qquad \re^{\lambda}=\frac{\prod_j(\zeta^j-L)}{(5L)^5}=-q(L)
\]
(the letter $\lambda$ is local to this section and \cref{sec:gam}). On $\bbD_L$ define
\begin{align*}
\mathsf S_0(\sigma)&:=\lambda\sigma+5\sigma\log(5\sigma)-5\sigma+\sum_{j=0}^4\bigl[(\zeta^j-\sigma)\log(\zeta^j-\sigma)-(\zeta^j-\sigma)\bigr],\\
\mathsf S_1(\sigma)&:=\tfrac12\log(5\sigma)-\tfrac12\sum_{j=0}^4\log(\zeta^j-\sigma),\qquad
\mathsf S_{2k}(\sigma):=\bk_k\Bigl[(5\sigma)^{1-2k}+\sum_{j=0}^4(\zeta^j-\sigma)^{1-2k}\Bigr]\quad(k\geq1),\\
A_{\mathrm{St}}(\sigma,z)&:=\exp\Bigl(\mathsf S_1(\sigma)+\sum_{k\geq1}z^{2k-1}\mathsf S_{2k}(\sigma)\Bigr).
\end{align*}
Then $\mathsf S_0'(L)=0$ and $\mathsf S_0''=5/\sigma+\sum_j(\zeta^j-\sigma)^{-1}=5/(\sigma(1-\sigma^5))$, which does not vanish at $L$. The \emph{full-Stirling saddle series} is the normalised stationary-phase series (\cref{def:fsp-wick}), with constant term $1$, in the variable $\sigma$ with the measure $\dd\sigma$:
\[
c(z;L):=\SPh_L[\mathsf S_0;A_{\mathrm{St}}](z)=\sum_{n\geq0}c_n(L)z^n,\qquad c_0(L)=1 .
\]
For colour $i$ the series is defined in the same way at the critical point $\sigma=\zeta^iL$, with a disc around $\zeta^iL$, and with the branches and the logarithm of $-q$ chosen there as for $i=0$, so that $\zeta^iL$ is critical.
\end{defn}

With $\log(a/z):=\log a-\log z$ and $\Gst(1+w):=w\,\Gst(w)$ (\cref{def:fsp-stirling}) one has, identically in $\sigma\in\bbD_L$,
\[
\re^{\lambda\sigma/z}\,\Gst(1+5\sigma/z)\prod_{j=0}^4\Gst\bigl((\zeta^j-\sigma)/z\bigr)=(2\pi)^3z^2\exp\Bigl(\mathsf S_0/z+\mathsf S_1+\sum_kz^{2k-1}\mathsf S_{2k}\Bigr),
\]
because the coefficient of $(\log z)/z$ is $-[5\sigma+\sum_j(\zeta^j-\sigma)]=-\sum_j\zeta^j=0$ and that of $\log z$ is $-\frac12+\frac52=2$. So $c(z;L)$ is the normalised saddle series at $\sigma=L$ of the Mellin--Barnes summand $\Gamma(1+5\sigma/z)\prod_j\Gamma((\zeta^j-\sigma)/z)(-q)^{\sigma/z}\,\dd\sigma$, with all six Gamma factors replaced by $\Gst$ and $(-q)^{\sigma/z}:=\re^{\lambda\sigma/z}$; in the Euler integrals the twist factor $\Gamma(1+5\sigma/z)$ carries the measure $\dd y$ and the five factors $\Gamma((\zeta^j-\sigma)/z)$ carry $\dd x/x$.

\emph{Taylor data.} Consider the local power sums $p_k(L):=\sum_j(\zeta^j-L)^{-k}$. Their generating function is
\[
\sum_{k\geq1}p_k(L)\,t^{k-1}=\sum_j(\zeta^j-L-t)^{-1}=\frac{5(L+t)^4}{1-(L+t)^5},
\]
so that $p_k\in(1-L^5)^{-k}\bbQ[L]$, and one has
\begin{gather*}
\mathsf S_0(L+t)-\mathsf S_0(L)=\sum_{k\geq1}\frac{5(-1)^{k+1}L^{-k}+p_k(L)}{k(k+1)}\,t^{k+1},\qquad
-1/\mathsf S_0''(L)=-L(1-L^5)/5,\\
\mathsf S_1(L+t)-\mathsf S_1(L)=\sum_{k\geq1}\frac{(-1)^{k+1}L^{-k}+p_k(L)}{2k}\,t^{k},\\
[t^r]\,\mathsf S_{2k}(L+t)=\bk_k\binom{1-2k}{r}\bigl[5^{1-2k}L^{1-2k-r}+(-1)^rp_{2k-1+r}(L)\bigr].
\end{gather*}

\begin{prop}[Rationality]\genusstatus{AI-proposed, unchecked}
\label{prop:gam-rational}
The series $c(z;L)$ does not depend on the disc $\bbD_L$ nor on the branches of the logarithms, and every $c_n$ is the value at $L$ of an element of $\AL=\bbQ[L,L^{-1},(L^5-1)^{-1}]$. More precisely, $c(z;L)$ depends only on the Taylor coefficients of order at least $2$ of the phase at $L$ and on the normalised amplitude $A_{\mathrm{St}}/A_{\mathrm{St}}(L,0)$; $\lambda$ and the sign convention for $q$ never enter it, they only fix the dictionary $L\leftrightarrow q$.
\end{prop}

\begin{proof}
(1) $\mathsf S_0$ and $A_{\mathrm{St}}$ are holomorphic on $\bbD_L$, $\mathsf S_0'(L)=0$ and $\mathsf S_0''(L)=5/(L(1-L^5))\neq0$. (2) By \cref{lem:fsp-univ}, $c_n(L)$ is a polynomial over $\bbQ$ in the Taylor coefficients at $L$ of $\mathsf S_0$ of orders $3,\dots,2n+2$, in those of $\mathsf S_1-\mathsf S_1(L)$ and of the $\mathsf S_{2k}$ (which determine the normalised amplitude), and in $-1/\mathsf S_0''(L)$. (3) These are the expressions displayed above: polynomials in $L^{-1}$, in the $p_k\in(1-L^5)^{-k}\bbQ[L]$ and in $-L(1-L^5)/5$. (4) Hence $c_n\in\AL$. (5) The Taylor coefficients of order at least $2$ of $\mathsf S_0$, and of order at least $1$ of $\mathsf S_1$ and $\mathsf S_{2k}$, are rational in $\sigma$ and carry no branch, and $\exp(\mathsf S_1(\sigma)-\mathsf S_1(L))$ is the branch equal to $1$ at $L$. Replacing $\log(5\sigma)$ by $\log(5\sigma)+2\pi\ri n$ and $\log(\zeta^j-\sigma)$ by $\log(\zeta^j-\sigma)+2\pi\ri n_j$ changes $\mathsf S_0'$ by $2\pi\ri(5n-\sum_jn_j)$, which the formula for $\lambda$ compensates; then $\mathsf S_0$ changes by the constant $2\pi\ri\sum_j\zeta^jn_j$ and $\mathsf S_1$ by a constant, which the normalised series does not see (\cref{lem:fsp-const}). A logarithm $\lambda+2\pi\ri k$ of $-q$ with $k\neq0$ gives $\mathsf S_0'(L)=2\pi\ri k\neq0$, so $L$ is critical only for $\lambda$; a smaller disc changes nothing. So $c(z;L)$ depends only on $L$, and only through the data named in the statement.
\end{proof}

\begin{lem}[One $1$-form]\genusstatus{AI-proposed, unchecked}
\label{lem:con-form}
Let $\Xi:=-\log(1-u)\in u\bbQ\fps{u}$ and $\cG(x):=\sum_{m\geq1}x^{5m-1}/(m(5m-1))$.
\begin{enumerate}
\item[(a)] In the colour-$0$ summand
\[
q^{\sigma/z}\,\Gamma(1+5\sigma/z)\Big/\prod_{j=0}^4\Gamma\bigl(1+(\sigma-\zeta^j)/z\bigr)
\]
expand $\Gamma(1+5\sigma/z)$ and the five factors $1/\Gamma(1+(\sigma-\zeta^j)/z)$ by Stirling's series, and put $x=1/\sigma$. The exponent is $S(\sigma)/z$ with $S(\sigma):=\sigma\log q+5\sigma\log(5\sigma)-\sum_j(\sigma-\zeta^j)\log(\sigma-\zeta^j)$: the terms in $(\log z)/z$ and the terms linear in $\sigma/z$ cancel. Moreover $S(1/x)=-\cW(x)$ with $\cW(x):=\Xi/x+\cG(x)$, and $\dd\cW/\dd x=-(\Xi+\log(1-x^5))/x^2$; formally, the critical points are $x^5=u$ (an identity of symbols, with $\log(1-x^5)$ and $\log(1-u)$ as series; it is used below only through (c)).
\item[(b)] The half-densities and the measure combine to $(5\sigma/(\sigma^5-1))^{1/2}\,\dd\sigma=-\sqrt5\,(1-x^5)^{-1/2}\,\dd x$. Up to a factor independent of $x$, the summand times $\dd\sigma$ is $\re^{-\cW/z}\fa\,\dd x$ with
\[
\log\fa(x,z)=-\tfrac12\log(1-x^5)+\sum_{k\geq1}\bk_k(xz)^{2k-1}\Bigl[5^{1-2k}-5\sum_{P\geq0}\binom{2k-2+5P}{5P}x^{5P}\Bigr];
\]
$\fa$ does not depend on $u$, $\fa(0,z)=1$, and the coefficient of $x^a$ in $\fa$ is a polynomial in $z$ of degree at most $a$.
\item[(c)] In the hatted variables $x=sX$, $z=s^4\zh$ one has $\cW/z=\hat\cW/\zh$ with
\[
\hat\cW(X;u)=\sum_{m\geq1}\frac{u^{m-1}}m\Bigl[\frac1X+\frac{X^{5m-1}}{5m-1}\Bigr],\qquad \hat\cW'(X)=-\frac1{uX^2}\log\frac{1-uX^5}{1-u};
\]
$X=1$ is critical for every $u$, $\hat\cW''(1)=5/(1-u)$, $\hat\cW(1)=\frac54\tau(u)^{-1}$, and $[t^k]\hat\cW(1+t)=\sum_{m\geq1}u^{m-1}[(-1)^k+\binom{5m-1}{k}/(5m-1)]/m$ lies in $\bbQ\fps{u}$. The hatted amplitude $\hat\fa(X;\zh):=\fa(sX,s^4\zh)$ has
\[
\log\hat\fa=-\tfrac12\log(1-uX^5)+\sum_{k\geq1}\bk_k(uX\zh)^{2k-1}\Bigl[5^{1-2k}-5\sum_{P\geq0}\binom{2k-2+5P}{5P}u^PX^{5P}\Bigr];
\]
it lies in $\bbQ[X,\zh]\fps{u}$, and $\hat\fa(1;0)=(1-u)^{-1/2}$.
\item[(d)] The series $c(z;L)$ of \cref{def:gam-c} (the first Mellin--Barnes summand) and the normalised saddle series at $\sigma=L$ of the second summand of (a), with all Gamma factors replaced by their Stirling forms and the logarithm of $q$ chosen so that $L$ is critical, coincide.
\end{enumerate}
\end{lem}

% Section 8 (part 4): proof of the 1-form lemma; PhiSP and the residue; Theorem P(a), statement and part (i).

\begin{proof}
(a) Replace each $\Gamma(1+y)$ by its Stirling form (\cref{def:fsp-stirling})
\[
\Gst(1+y)=y\,\Gst(y)=\sqrt{2\pi}\,y^{y+1/2}\re^{-y}\exp\Bigl(\sum_{k\geq1}\bk_ky^{1-2k}\Bigr),
\]
with $y=5\sigma/z$ for the numerator and $y=(\sigma-\zeta^j)/z$ for the factors of the denominator, and $\log(a/z)=\log a-\log z$. The terms in $(\log z)/z$ are $-(\log z/z)\,[5\sigma-\sum_j(\sigma-\zeta^j)]=0$, those linear in $\sigma/z$ are $[-5\sigma+\sum_j(\sigma-\zeta^j)]/z=0$, and $q^{\sigma/z}=\re^{\sigma\log q/z}$; what remains of order $1/z$ is $S(\sigma)/z$. At $\sigma=1/x$ write $\log(1/x-\zeta^j)=\log(1-\zeta^jx)-\log x$; with $\sum_j(1/x-\zeta^j)=5/x$ and $\log q+5\log5=\log(1-u)$ this gives $S(1/x)=-\Xi/x-\sum_j(1/x-\zeta^j)\log(1-\zeta^jx)$. Expanding $\log(1-\zeta^jx)=-\sum_n\zeta^{jn}x^n/n$ and summing over $j$, only the terms with $5\mid n$ survive in $\sum_j\log(1-\zeta^jx)$ and only those with $5\mid n+1$ in $\sum_j\zeta^j\log(1-\zeta^jx)$, and
\[
\sum_j(1/x-\zeta^j)\log(1-\zeta^jx)=-\sum_{m\geq1}\frac{x^{5m-1}}m+5\sum_{m\geq1}\frac{x^{5m-1}}{5m-1}=\cG(x).
\]
Then $\cG'(x)=\sum_mx^{5m-2}/m=-x^{-2}\log(1-x^5)$ gives $\dd\cW/\dd x$, which vanishes if and only if $\log(1-x^5)=\log(1-u)$.

(b) The half-densities are $(5\sigma/z)^{1/2}$ from $\Gamma(1+5\sigma/z)$ and $((\sigma-\zeta^j)/z)^{-1/2}$ from the denominator; their product is $z^2(5\sigma/\prod_j(\sigma-\zeta^j))^{1/2}=z^2(5\sigma/(\sigma^5-1))^{1/2}$, and at $\sigma=1/x$ one has $5\sigma/(\sigma^5-1)=5x^4/(1-x^5)$ and $\dd\sigma=-\dd x/x^2$. The Stirling series contribute $\sum_k\bk_k[(5\sigma/z)^{1-2k}-\sum_j((\sigma-\zeta^j)/z)^{1-2k}]=\sum_k\bk_k(xz)^{2k-1}[5^{1-2k}-\sum_j(1-\zeta^jx)^{1-2k}]$, and $\sum_j(1-\zeta^jx)^{-r}=5\sum_{5\mid i}\binom{r-1+i}{i}x^i$ (expand and sum over $j$). The constant $-(2\pi)^{-2}\sqrt5$ and the power $z^2$ do not depend on $x$; what is left is $\re^{-\cW/z}\fa\,\dd x$. Every monomial of $\log\fa$ has $z$-degree at most its $x$-degree, hence so has every monomial of $\fa$.

(c) With $x=sX$ and $z=s^4\zh$ one has $\Xi/(xz)=(\Xi/u)/(X\zh)$ with $\Xi/u=\sum_mu^{m-1}/m$, and $\cG(sX)/(s^4\zh)=\sum_mu^{m-1}X^{5m-1}/(m(5m-1)\zh)$; this is $\hat\cW/\zh$. Termwise in $u$,
\[
\hat\cW'(X)=\sum_m\frac{u^{m-1}}m(X^{5m-2}-X^{-2})=(uX^2)^{-1}\sum_m\frac{u^m}m(X^{5m}-1),
\]
which is the formula stated and vanishes at $X=1$; differentiating the logarithm at $X=1$ gives $\hat\cW''(1)=5/(1-u)$; and $\hat\cW(1)=5\sum_mu^{m-1}/(5m-1)=\frac54\sum_{n\geq0}4u^n/(5n+4)=\frac54\tau^{-1}$ (\cref{def:set-tau}). The Taylor coefficients at $X=1$ follow from those of $(1+t)^{-1}$ and $(1+t)^{5m-1}$. Finally $xz=uX\zh$ and $x^5=uX^5$ in (b).

(d) By \cref{lem:fsp-univ,lem:fsp-const}, a normalised saddle series depends only on the Taylor coefficients of order at least $2$ of the phase at the critical point and on the normalised amplitude. For $y_j:=(\zeta^j-\sigma)/z$ the formal reflection identity \cref{lem:fsp-refl} gives $\log\Gst(y_j)+\log\Gst(1-y_j)=\log2\pi+(y_j-\frac12)(\log y_j-\log(-y_j))$, and $\log y_j-\log(-y_j)=\pi\ri\,\epsilon_j$ with an odd integer $\epsilon_j$ fixed by the branches. Since $1-y_j=1+(\sigma-\zeta^j)/z$, the product $\prod_j\Gst(1+(\sigma-\zeta^j)/z)^{-1}$ is a constant times $\re^{c/z}\re^{\pi\ri(\sum_j\epsilon_j)\sigma/z}\prod_j\Gst((\zeta^j-\sigma)/z)$, with $c$ independent of $\sigma$. Moreover $q^{\sigma/z}=\re^{(\lambda+\pi\ri\epsilon)\sigma/z}$ with an odd integer $\epsilon$, because the logarithms of $q$ and of $-q$ differ by an odd multiple of $\pi\ri$. So the second summand, in Stirling form, is a constant times $\re^{c/z}$ times the first with $\lambda$ replaced by $\lambda+\pi\ri e$, where $e:=\epsilon+\sum_j\epsilon_j$ is even. If $L$ is critical for both, then $e=0$ (step (5) of the proof of \cref{prop:gam-rational}). The two integrands then have the same phase up to an additive constant and the same amplitude up to a constant factor, so their normalised series coincide.
\end{proof}

\begin{defn}[The full-Stirling WKB function and the residue]
\label{def:con-PhiSP}
The \emph{full-Stirling WKB function} is the normalised stationary-phase series, with constant term $1$, of $\re^{-\hat\cW/\zh}\hat\fa\,\dd X$ at $X=1$:
\[
\PhiSP(\zh,u):=\SPh_{X=1}[-\hat\cW;\hat\fa](\zh)=\sum_{n\geq0}\mathrm{ph}^{\mathrm{SP}}_n(u)\zh^n,\qquad\mathrm{ph}^{\mathrm{SP}}_0=1 .
\]
Equivalently,
\[
\PhiSP=\mathrm{SP}_{X=1}\bigl[\re^{-\hat\cW/\zh}\hat\fa\,\dd X\bigr]\Big/\bigl((2\pi\zh/5)^{1/2}\re^{-\hat\cW(1)/\zh}\bigr),
\]
the full stationary phase of \cref{def:fsp-wick} divided by $(2\pi\zh/5)^{1/2}\re^{-\hat\cW(1)/\zh}$, which is the Gaussian factor $(2\pi\zh)^{1/2}\hat\cW''(1)^{-1/2}$ times the critical exponential $\re^{-\hat\cW(1)/\zh}$ times the amplitude value $\hat\fa(1;0)=(1-u)^{-1/2}$; in every comparison below this $\zh$-dependent factor is common to both sides. By \cref{lem:con-form}~(c) and \cref{lem:fsp-univ}, and since the covariance $(1-u)/5$ and the unit $\hat\fa(1;0)=(1-u)^{-1/2}$ lie in $\bbQ\fps u$, every $\mathrm{ph}^{\mathrm{SP}}_n$ lies in $\bbQ\fps{u}$. The \emph{residue} of the flat $1$-form is
\[
\cV:=\Res_{x=0}\re^{-\cW/z}\fa\,\dd x:=\sum_{k\geq1}\frac{(-\Xi/z)^k}{k!}\,[x^{k-1}]\bigl(\re^{-\cG/z}\fa\bigr).
\]
Since $\Xi\in u\bbQ\fps{u}$, $\cG\in x^4\bbQ\fps{x}$ and the coefficient of $x^a$ in $\fa$ has $z$-degree at most $a$, each $u$-order of $\cV$ is a finite sum, and $\cV\in z^{-1}\bbQ[z^{-1}]\fps{u}$.
\end{defn}

\subsection{Theorem P}
\label{sec:con-P}

The ring $\AL$ carries the derivation $\Dq$ with $\Dq L=(L^6-L)/5$; it contains $q=q(L)$, and $\Dq q=q$, so that $\Dq=q\,\dd/\dd q$ on $\AL$. Let $\jmath\colon\AL\to\bbQ\lser{s}$ be the ring homomorphism with $\jmath(L)=s^{-1}$; it is injective, because $\bbQ(L)\to\bbQ\lser{s}$, $L\mapsto s^{-1}$, is a field embedding.

\begin{thm}[Theorem P(a)]\genusstatus{AI-proposed, unchecked}
\label{thm:con-Pa}
Let $c(z;L)=\sum_nc_n(L)z^n$ be the full-Stirling saddle series, with $c_n\in\AL$ (\cref{def:gam-c,prop:gam-rational}).
\begin{enumerate}
\item $\PFfive^{L}\bigl(L\,c(z;L)\bigr)=0$ in $\AL\fps{z}$.
\item In $\bbQ\lser{s}\fps{z}$ one has $\jmath(c)(z)=\PhiSP(s^{-4}z,s^5)$; that is, $\jmath(c_n)=s^{-4n}\mathrm{ph}^{\mathrm{SP}}_n$ for every $n$.
\item $\PhiSP$ solves HPF in $\bbQ\fps{u}\fps{\zh}$.
\item Every $\mathrm{ph}^{\mathrm{SP}}_n$ is a polynomial in $u$; hence $c_n(L)=L^{4n}\mathrm{ph}^{\mathrm{SP}}_n(L^{-5})$ in $\AL$, and $c_n\in\bbQ[L,L^{-1}]$.
\end{enumerate}
\end{thm}

The proof of Theorem~P(a) uses no input from outside the note; part (ii) does not use (iv).

\begin{proof}[Proof of (i)]
In this proof $\sigma=L+t$, $\mathsf S_0$ stands for the formal phase $\mathsf S_0(L+t)-\mathsf S_0(L)\in t^2\AL\fps t$ at $\sigma=L$ (\cref{def:gam-c}; the constant does not matter, \cref{lem:fsp-const}), and $a_0:=A_{\mathrm{St}}/A_{\mathrm{St}}(L,0)=\exp\bigl(\mathsf S_1(L+t)-\mathsf S_1(L)+\sum_kz^{2k-1}\mathsf S_{2k}(L+t)\bigr)\in\AL\fps t\fps z$ for the normalised amplitude, both given by the Taylor data of \cref{def:gam-c}. For $\chi\in\bbQ[\sigma,z]$ put $c\langle\chi\rangle(z;L):=\SPn_L[\mathsf S_0;\chi a_0]$, so that $c=c\langle1\rangle$; the map $\chi\mapsto c\langle\chi\rangle$ extends $\AL[z]$-linearly to $\AL[\sigma,z]$ (\cref{def:fsp-wick}). The letters $\chi$, $a$, $a_0$, $b$, $b_z$, $\mathsf E$, $\mathsf T_z$, $\rho$, $\mathsf l_\rho$, $\mathsf U_\rho$, $\mathsf m_\rho$, $w$, $\Bbbk$ are local to this proof.

\emph{Step 1 (the shift identity).} We apply \cref{lem:fsp-generic}. First, $\mathsf S_0'(L)=0$ and $\mathsf S_0''(\sigma)=5/(\sigma(1-\sigma^5))=-1/\delta(\sigma)$, $\delta(\sigma)=(\sigma^6-\sigma)/5$. Next, $A_{\mathrm{St}}(\sigma,0)^2=\re^{2\mathsf S_1}=5\sigma/\prod_j(\zeta^j-\sigma)=5\sigma/(1-\sigma^5)=-\mathsf r(\sigma)$, so the quotient of $A_{\mathrm{St}}(\sigma,0)$ by $A_{\mathrm{St}}(L,0)$ is $(\mathsf r(\sigma)/\mathsf r(L))^{1/2}$, the root with constant term $1$ at $L$. Finally $b:=\chi\,A_{\mathrm{St}}(\sigma,z)/A_{\mathrm{St}}(\sigma,0)=\chi\exp\bigl(\sum_kz^{2k-1}\mathsf S_{2k}(\sigma)\bigr)$ lies in $\bbA_\sigma\fps{z}$, because $\sigma^{-1}$ and the power sums $\sum_j(\zeta^j-\sigma)^{-r}\in(1-\sigma^5)^{-r}\bbQ[\sigma]$ lie in $\bbA_\sigma$. So $\chi a_0=(\mathsf r(\sigma)/\mathsf r(L))^{1/2}b$, and \eqref{eq:fsp-generic} gives
\begin{equation}
\label{eq:con-shiftD}
(z\Dq+L)\bigl(L\,c\langle\chi\rangle\bigr)=L\,c\langle\sigma\chi\rangle\quad\text{in }\AL\lser{z},\qquad\chi\in\bbQ[\sigma,z].
\end{equation}

\emph{Step 2 ($\PFfive$).} $\PFfive=(z\Dq)^5-1-q\prod_{k=1}^5(5z\Dq+kz)$, with $q$ acting by multiplication on the left. Iterating \eqref{eq:con-shiftD} and using the $\AL[z]$-linearity of $\chi\mapsto c\langle\chi\rangle$,
\[
\PFfive^{L}(L\,c)=L\Bigl(c\langle\sigma^5-1\rangle-c\bigl\langle q\textstyle\prod_{k=1}^5(5\sigma+kz)\bigr\rangle\Bigr).
\]

\emph{Step 3 (the shift).} Put $a:=(\sigma^5-1)a_0$, and let $\mathsf T_z$ be the continuous ring endomorphism $f(t,z)\mapsto f(t+z,z)$ of $\AL\fps{t,z}$. The translate of $a$ by $z$ (\cref{lem:fsp-transl}~(b)) is $b_z=\re^{\mathsf E}\,\mathsf T_z(a)$ with
\[
\mathsf E:=\bigl(\mathsf S_0(L+t+z)-\mathsf S_0(L+t)\bigr)/z=\sum_{n\geq1}\mathsf S_0^{(n)}(L+t)\,z^{n-1}/n!\in\AL\fps{t,z},
\]
which has no constant term, because $\mathsf S_0'(L)=0$. We claim
\begin{equation}
\label{eq:con-bz}
b_z=q\prod_{k=1}^5(5\sigma+kz)\,a_0\qquad\text{in }\AL\fps{t,z},\quad\sigma=L+t .
\end{equation}
The ring $\AL\fps{t,z}$ embeds in $\Bbbk\fps{t,z}$, where $\Bbbk:=\AL\otimes_{\bbQ}\bbQ(\zeta)=\bbQ(\zeta)[L,L^{-1},(L^5-1)^{-1}]$ (in this proof only; $\bbQ(\zeta)$ is free over $\bbQ$, so $\AL\to\Bbbk$ is injective), in which $5L$ and the $\zeta^j-L$ are units; we prove \eqref{eq:con-bz} there. We write $\mathsf T_z$ also for the endomorphism $f(t,z)\mapsto f(t+z,z)$ of $\Bbbk\fps{t,z}$; it is continuous and preserves the ideal $(t,z)$, so it commutes with $\exp$ and $\log$. Let $\rho(\sigma)=\alpha\sigma+\beta$ with $\alpha\in\bbZ$, $\beta\in\bbQ(\zeta)$ and $\rho_0:=\rho(L)\in\Bbbk^\times$, so that $\rho(L+t)=\rho_0+\alpha t$ is a unit of $\Bbbk\fps t$, again written $\rho$, and put
\begin{gather*}
\mathsf l_\rho:=\log(1+\alpha t/\rho_0)\in t\Bbbk\fps t,\qquad\mathsf U_\rho:=\rho\,\mathsf l_\rho-\alpha t\in t^2\Bbbk\fps t,\\
\mathsf m_\rho:=\sum_{k\geq1}\bk_k\,z^{2k-1}\rho^{1-2k},\qquad w:=z/\rho\in z\Bbbk\fps{t,z};
\end{gather*}
these are the Taylor parts at $\sigma=L$ of $\log\rho$ (orders $\geq1$) and of $\rho\log\rho-\rho$ (orders $\geq2$), and the Stirling tail $\mu(\rho/z)$ of \cref{def:fsp-stirling}. Then
\begin{equation}
\label{eq:con-factor}
\frac{\mathsf T_z\mathsf U_\rho-\mathsf U_\rho}z+(\mathsf T_z-1)\bigl(-\tfrac12\mathsf l_\rho+\mathsf m_\rho\bigr)=\alpha\,\mathsf l_\rho+\log\mathrm{Poch}_\alpha(z/\rho)\qquad\text{in }\Bbbk\fps{t,z}.
\end{equation}
Indeed $\mathsf T_z\rho=\rho\,(1+\alpha w)$, so $\mathsf T_z\mathsf l_\rho=\mathsf l_\rho+\log(1+\alpha w)$ and $(\mathsf T_z\mathsf U_\rho-\mathsf U_\rho)/z=\alpha\,\mathsf l_\rho+(w^{-1}+\alpha)\log(1+\alpha w)-\alpha$ (here $w^{-1}\log(1+\alpha w)$ is the power series $\sum_{n\geq1}(-1)^{n+1}\alpha^nw^{n-1}/n$, as in \cref{def:fsp-stirling}), while $\mathsf m_\rho=\sum_k\bk_kw^{2k-1}$ and $\mathsf T_z\mathsf m_\rho=\sum_k\bk_kw^{2k-1}(1+\alpha w)^{1-2k}$. So the left side of \eqref{eq:con-factor} is $\alpha\,\mathsf l_\rho+\log\mathrm{St}_\alpha(w)-\log\mathrm{St}_0(w)$, and $\mathrm{St}_\alpha=\mathrm{Poch}_\alpha\,\mathrm{St}_0$ in $\bbQ\fps w$ by \cref{lem:fsp-stirfe}; the substitution $w=z/\rho$ is a continuous ring homomorphism $\bbQ\fps w\to\Bbbk\fps{t,z}$.

Apply \eqref{eq:con-factor} to $\rho_{\mathrm N}:=5\sigma$ ($\alpha=5$) and to $\rho_j:=\zeta^j-\sigma$ ($\alpha=-1$), $j=0,\dots,4$. Since $[t^{k+1}]\,\mathsf U_\rho=(-1)^{k+1}\alpha^{k+1}\rho_0^{-k}/(k(k+1))$ and $[t^k]\,\mathsf l_\rho=(-1)^{k+1}\alpha^k\rho_0^{-k}/k$ for $k\geq1$, the Taylor data of \cref{def:gam-c} give, coefficient by coefficient, $\mathsf S_0(L+t)-\mathsf S_0(L)=\mathsf U_{\rho_{\mathrm N}}+\sum_j\mathsf U_{\rho_j}$, $\mathsf S_1(L+t)-\mathsf S_1(L)=\frac12\mathsf l_{\rho_{\mathrm N}}-\frac12\sum_j\mathsf l_{\rho_j}$ and $\sum_kz^{2k-1}\mathsf S_{2k}(L+t)=\mathsf m_{\rho_{\mathrm N}}+\sum_j\mathsf m_{\rho_j}$. Summing \eqref{eq:con-factor} over the six factors and adding $(\mathsf T_z-1)\mathsf l_{\rho_{\mathrm N}}=\log(1+5z/\rho_{\mathrm N})$, with $\mathrm{Poch}_5(w)(1+5w)=\prod_{k=1}^5(1+kw)$ and $\mathrm{Poch}_{-1}(w)=(1-w)^{-1}$,
\[
\mathsf E+(\mathsf T_z-1)\log a_0=5\,\mathsf l_{\rho_{\mathrm N}}-\sum_j\mathsf l_{\rho_j}+\sum_{k=1}^5\log\Bigl(1+\frac{kz}{\rho_{\mathrm N}}\Bigr)-\sum_j\log\Bigl(1-\frac z{\rho_j}\Bigr).
\]
Exponentiate, using $\re^{\mathsf l_\rho}=\rho/\rho_0$, $\prod_j(\zeta^j-x)=1-x^5$ and $\prod_j(\zeta^j-L)/(5L)^5=-q$ (\cref{def:gam-c}):
\[
\re^{\mathsf E}\,\frac{\mathsf T_z(a_0)}{a_0}=\frac{\prod_{k=1}^5(5\sigma+kz)}{(5L)^5}\cdot\frac{\prod_j(\zeta^j-L)}{\prod_j(\zeta^j-\sigma-z)}=\frac{q\prod_{k=1}^5(5\sigma+kz)}{(\sigma+z)^5-1}.
\]
Multiplying by $\mathsf T_z(\sigma^5-1)\,a_0=((\sigma+z)^5-1)\,a_0$ gives \eqref{eq:con-bz}. By \cref{lem:fsp-transl}~(b), with $p=z$ and $a\in\AL\fps t\fps z$, $\SPn_L[\mathsf S_0;b_z]=\SPn_L[\mathsf S_0;a]$, that is, $c\langle q\prod_k(5\sigma+kz)\rangle=c\langle\sigma^5-1\rangle$; by Step~2, $\PFfive^{L}(L\,c)=0$.
\end{proof}

\input{sections/s08_connection_5}



% Section 9, part 1 of 4 (continued in parts 2-4); label prefix: gam. No new macros.
% Local letters are declared where they are used. Discs are written \bbD: the letter D is the operator q d/dq.

\section{The Gamma identity and the normalisation}
\label{sec:gam}

This section proves Theorem Gamma (\cref{thm:gam-main}): the full-Stirling saddle series $c(z;L)$ of \cref{def:gam-c} has no pole at $L=1$, that is at $q=0$, and $c(z;1)=h(z)$, the series $h=h_{\bw_0}$ of (QRR). By Theorem~P and the glue of \cref{sec:con} this is equivalent to $\Dn=1$, and it gives $\Fg{g}$ for every $g\geq 2$ (\cref{thm:gam-D}). The proof is local: near a point at which exactly one Gamma factor of the Mellin--Barnes summand degenerates, only that factor is written as an Euler integral; after a change of chart the two-variable stationary phase depends holomorphically on a parameter that vanishes at the point, where a Gaussian computation evaluates it exactly (Theorem~LMI, \cref{thm:gam-lmi}). At $L=0$, that is at $q=\infty$, the same theorem gives the free coefficients of \cref{lem:gau-HPF} in closed form (\cref{thm:gam-ginf}, \cref{cor:gam-free}).

\subsection{The saddle series}
\label{sec:gam-series}

We use $c(z;L)=\SPh_L[\mathsf S_0;A_{\mathrm{St}}]$ and its data as in \cref{def:gam-c}; $c_n\in\AL$ does not depend on the choices of disc and branches (\cref{prop:gam-rational}). At each of the three points $L=1$, $L=\zeta^i$ and $L=0$ exactly one of the six Gamma factors of the summand degenerates, namely $\Gamma((1-\sigma)/z)$, $\Gamma((\zeta^i-\sigma)/z)$ and the twist factor $\Gamma(1+5\sigma/z)$ respectively, and \cref{thm:gam-lmi} applies at each of them; it is used at $L=1$ and $L=0$, and the points $\zeta^i$ are reached by rotation (\cref{thm:gam-main}~(iii)).

\subsection{Local Mellin inversion}
\label{sec:gam-lmi}

Data that are holomorphic at the base point enter only through their Taylor series.

\begin{lem}[The delta lemma with a source]\genusstatus{AI-proposed, unchecked}
\label{lem:gam-deltasrc}
Let $a_*\in\bbC^\times$ and $c\in\bbC$, let $P\in s^2\bbC\fps{s}$, and let $A\in\bbC\fps{s}\fps{z}$. In two variables $(s,m)$ consider the phase $F(s,m):=P(s)+a_*sm$, of exactly this form (no term in $m$ alone and no term $s^km$ with $k\geq 2$), and the amplitude $\re^{cm}A(s,z)$. Then $(0,0)$ is a non-degenerate critical point of $F$, $\det\Hess F(0,0)=-a_*^2$, and, exactly,
\begin{equation}
\label{eq:gam-deltasrc}
\SPn_{(0,0)}\bigl[F;\,\re^{cm}A\bigr]=\exp\bigl(P(\delta)/z\bigr)\,A(\delta,z),\qquad \delta:=-cz/a_*,
\end{equation}
where both factors are re-expanded in $z$: $P(\delta)/z=\sum_{k\geq 2}\coef{s^k}P\,(-c/a_*)^kz^{k-1}$ lies in $z\bbC\fps{z}$, and $A(\delta,z)=\sum_{n\geq 0}z^n\sum_{j+l=n}\coef{s^jz^l}A\,(-c/a_*)^j$. For $c=0$ the series is $A(0,z)$.
\end{lem}

\begin{proof}
The gradient $(P'(s)+a_*m,\,a_*s)$ vanishes at $(0,0)$, and the Hessian there is $H=\bigl(\begin{smallmatrix}P''(0)&a_*\\ a_*&0\end{smallmatrix}\bigr)$, $\det H=-a_*^2\neq0$; the covariance $C=-H^{-1}$ has $C_{ss}=0$, $C_{sm}=-1/a_*$, $C_{mm}=P''(0)/a_*^2$, and the part of order at least three of $F$ is $F_3(s)=P(s)-P''(0)s^2/2$. Write $\re^{F_3(s)/z}A(s,z)=\sum_{a\geq 0}g_a(z)s^a$, with $g_a$ of order at least $-a/3$ in $z$. By \cref{def:fsp-wick} the left side of \eqref{eq:gam-deltasrc} is $\sum_{a,r\geq 0}\frac{c^r}{r!}g_a(z)\langle s^am^r\rangle_Cz^{(a+r)/2}$. As $C_{ss}=0$, in the matching sum of \cref{lem:fsp-ring}~(iii) every copy of $s$ is paired with a copy of $m$; so $\langle s^am^r\rangle_C=\frac{r!}{(r-a)!}(r-a-1)!!\,C_{sm}^{\,a}C_{mm}^{\,(r-a)/2}$ if $r-a\geq0$ is even ($(-1)!!:=1$), and $0$ otherwise. With $r=a+2j$ and $(2j-1)!!/(2j)!=1/(2^jj!)$,
\[
\sum_{r\geq 0}\frac{c^r}{r!}\langle s^am^r\rangle_C\,z^{(a+r)/2}=(cC_{sm}z)^a\exp\Bigl(\frac{c^2C_{mm}z}{2}\Bigr)=\delta^a\exp\Bigl(\frac{P''(0)\delta^2}{2z}\Bigr).
\]
Hence the left side is $\exp(P''(0)\delta^2/(2z))\sum_ag_a(z)\delta^a=\exp(P''(0)\delta^2/(2z))\exp(F_3(\delta)/z)A(\delta,z)$, which is the right side. All sums are finite at each order of $z$, since $g_a(z)\delta^a$ has order at least $2a/3$; and the substitution $s\mapsto\delta$ is compatible with sums, products and exponentials on series $\sum_ag_a(z)s^a$ with $g_a$ of order at least $-a/3$, a class containing $F_3/z$ and $A$ and closed under products.
\end{proof}

\begin{defn}[LMI data]
\label{def:gam-lmidata}
LMI data consist of a point $s_*\in\bbC$, a number $a_*\in\bbC^\times$, an integer $M_*$, an open disc $\bbD$ around $s_*$, a holomorphic function $\varphi$ on $\bbD$, and a series $\cF(\sigma,z)=\sum_{m\geq 0}\cF_m(\sigma)z^m$ of holomorphic functions on $\bbD$, such that $\cF_0$ has no zero on $\bbD$ (in particular $\cF(s_*,0)\neq 0$) and $a_*+(\sigma-s_*)\varphi''(\sigma)\neq 0$ for every $\sigma\in\bbD$. Put $\varphi_k:=\varphi^{(k)}(s_*)/k!$, put
\[
a(\sigma):=a_*(\sigma-s_*),
\]
which is exactly linear in $\sigma$, and $U_*:=\bbD\setminus\{s_*\}$. For $L\in U_*$ choose a branch of $\log a$ near $L$ and put
\[
\lambda(L):=-\varphi'(L)-a_*\log a(L),\qquad S(\sigma):=\lambda(L)\,\sigma+\varphi(\sigma)+a(\sigma)\log a(\sigma)-a(\sigma),
\]
\[
A(\sigma,z):=a(\sigma)^{M_*-1/2}\,\mathrm{Poch}_{M_*}\bigl(z/a(\sigma)\bigr)\exp\Bigl(\sum_{k\geq 1}\bk_k\bigl(z/a(\sigma)\bigr)^{2k-1}\Bigr)\,\cF(\sigma,z),
\]
that is, the expansion of $z^{M_*+a/z}\,\Gst(M_*+a/z)$ given by \cref{lem:fsp-stirling}~(iii) and \cref{lem:fsp-stirfe}, without its factor $(2\pi z)^{1/2}$ and without its phase factor $\re^{(a\log a-a)/z}$, times $\cF$; the power $a^{M_*-1/2}$ is taken with any branch near $L$. Let
\[
c^{\cF}(z;L)=\sum_{n\geq 0}c^{\cF}_n(L)z^n:=\SPh_L\bigl[S;\,A\bigr].
\]
Finally put $\delta:=-M_*z/a_*$, $\tilde\varphi(x):=\varphi(s_*+x)-\varphi(s_*)-\varphi'(s_*)x$, and define the chart parameter
\[
\mathsf Q(L):=\re^{-\lambda(L)/a_*}=a(L)\,\re^{\varphi'(L)/a_*}.
\]
\end{defn}

The point $L$ is critical for $S$: since the derivative of $a\log a-a$ is $a_*\log a$, $S'(\sigma)=\lambda(L)+\varphi'(\sigma)+a_*\log a(\sigma)$, which vanishes at $\sigma=L$. Moreover $S''(L)=\varphi''(L)+a_*^2/a(L)$, that is $S''(L)=(a_*+(L-s_*)\varphi''(L))/(L-s_*)$, which is non-zero on $U_*$; and $A(L,0)=a(L)^{M_*-1/2}\cF_0(L)\neq 0$; so $c^{\cF}(z;L)$ is defined for every $L\in U_*$.

% Section 9, part 2 of 4: Theorem LMI and its proof. Continued in part 3.

\begin{thm}[Theorem LMI: local Mellin inversion]\genusstatus{AI-proposed, unchecked}
\label{thm:gam-lmi}
Let LMI data be given (\cref{def:gam-lmidata}).
\begin{enumerate}
\item For $L\in U_*$ every coefficient $c^{\cF}_n(L)$ depends only on $L$, not on the branch of $\log a$ or on $\lambda(L)$, and $c^{\cF}_n$ is holomorphic on $U_*$.
\item Every $c^{\cF}_n$ extends holomorphically to $L=s_*$.
\item With these extensions,
\begin{equation}
\label{eq:gam-lmi}
c^{\cF}(z;s_*)=\exp\bigl(\tilde\varphi(\delta)/z\bigr)\,\frac{\cF(s_*+\delta,z)}{\cF(s_*,0)},\qquad \delta=-\frac{M_*z}{a_*},
\end{equation}
where both factors are re-expanded in $z$: $\tilde\varphi(\delta)/z=\sum_{k\geq 2}\varphi_k(-M_*/a_*)^kz^{k-1}$ and $\cF(s_*+\delta,z)=\sum_{n\geq 0}z^n\sum_{m+j=n}\cF_m^{(j)}(s_*)\,(-M_*/a_*)^j/j!$.
\end{enumerate}
\end{thm}

That $a$ is exactly linear in $\sigma$ and that $\varphi$ and $\cF$ do not involve $\lambda$ makes the chart data of Step~2 of the proof holomorphic at $L=s_*$; the condition $\cF(s_*,0)\neq 0$ is used for the normalisation at $s_*$.

\begin{proof}
\emph{Step 0: (i).} For $k\geq 2$ the Taylor coefficient of order $k$ at $L$ of $a\log a-a$ is $(-1)^ka_*/(k(k-1)(L-s_*)^{k-1})$; it involves no branch, and $\lambda(L)$ enters only the term of order one of $S$, which vanishes at the critical point. The normalised amplitude $A(L+t,z)/A(L,0)$ is the product of the binomial series $(1+t/(L-s_*))^{M_*-1/2}$, of the series $\mathrm{Poch}_{M_*}(z/a(L+t))\exp(\sum_k\bk_k(z/a(L+t))^{2k-1})$ in $t$, where $a(L+t)=a(L)(1+t/(L-s_*))$, and of $\cF(L+t,z)/\cF_0(L)$; it involves neither $\lambda$ nor a branch. By \cref{lem:fsp-univ}, $c^{\cF}_n(L)$ is a polynomial with rational coefficients in finitely many of these Taylor coefficients and in $-1/S''(L)=-(L-s_*)/(a_*+(L-s_*)\varphi''(L))$. Each of them is a holomorphic function of $L$ on $U_*$; the Taylor coefficients of $\varphi$, and those of the $\cF_m$ divided by $\cF_0(L)$, are holomorphic on all of $\bbD$, where $\cF_0$ has no zero. This proves (i).

\emph{Step 1: two variables.} Fix $L\in U_*$ and the branch of $\log a$ near $L$ used in $\lambda=\lambda(L)$. We use a local letter $\chi$ for the logarithm of the variable of the Euler integral of $\Gamma(M_*+a/z)$, and consider, in the two variables $(\sigma,\chi)$, the formal integrand with phase and amplitude
\[
\Phi_2(\sigma,\chi):=\lambda\sigma+\varphi(\sigma)+a(\sigma)\chi-\re^{\chi},\qquad \re^{M_*\chi}\,\cF(\sigma,z),
\]
at the point $(L,\chi_L)$, $\chi_L:=\log a(L)$ (the chosen branch). It is a formal integrand in the sense of \cref{sec:fsp}, motivated by Euler's integral $\int_0^\infty x^{M_*+a/z-1}\re^{-x/z}\,\dd x=z^{M_*+a/z}\,\Gamma(M_*+a/z)$, $a=a(\sigma)$, in the variable $x=\re^\chi$. The gradient $(\lambda+\varphi'(\sigma)+a_*\chi,\ a(\sigma)-\re^{\chi})$ vanishes at $(L,\chi_L)$ by the choice of $\lambda$, and, since $a''=0$, the Hessian there is
\[
\begin{pmatrix}\varphi''(L)&a_*\\ a_*&-a(L)\end{pmatrix},\qquad\text{with determinant }-a(L)\varphi''(L)-a_*^2=-a(L)S''(L)\neq 0.
\]
We apply \cref{lem:fsp-fubini} with the inner variable $\chi$ and the outer variable $\sigma$. The inner Hessian $-a(L)$ is invertible. Over the coefficient ring $\bbC\fps{\sigma-L}$ the inner critical point is $\chi(\sigma):=\chi_L+\log(1+(\sigma-L)/(L-s_*))$, the branch of $\log a(\sigma)$ that continues $\chi_L$, with $\re^{\chi(\sigma)}=a(\sigma)$; the inner critical value $\lambda\sigma+\varphi(\sigma)+a(\sigma)\chi(\sigma)-a(\sigma)$ is $S(\sigma)$; the inner Hessian at $\chi(\sigma)$ is $-a(\sigma)$, so that the factor of \cref{lem:fsp-fubini} is $(a(\sigma)/a(L))^{-1/2}$. The inner series is computed by \cref{lem:fsp-cov} with the map $\chi\mapsto x:=\re^{\chi}$, which sends $\chi(\sigma)$ to $a(\sigma)$ and has the Jacobian $a(\sigma)$ there; it turns the phase $a(\sigma)\log x-x$ into $a(\sigma)\chi-\re^{\chi}$ and the amplitude $x^{M_*-1}$ into $a(\sigma)^{-1}\re^{M_*\chi}$. With \cref{lem:fsp-stirling} for the unit $a=a(\sigma)$ of $\bbC\fps{\sigma-L}$, the integer $M=M_*$ and the measure $\dd x$, this gives
\begin{align*}
\SPn_{\chi(\sigma)}\bigl[a(\sigma)\chi-\re^{\chi};\,\re^{M_*\chi}\bigr]&=a(\sigma)\,\SPn_{a(\sigma)}\bigl[a(\sigma)\log x-x;\,x^{M_*-1}\bigr]\\
&=a(\sigma)^{M_*}\,\mathrm{Poch}_{M_*}\bigl(z/a(\sigma)\bigr)\exp\Bigl(\sum_{k\geq 1}\bk_k\bigl(z/a(\sigma)\bigr)^{2k-1}\Bigr).
\end{align*}
The factor $\cF(\sigma,z)\in\bbC\fps{\sigma-L}\fps{z}$ does not involve $\chi$ and is carried along by the linearity of the inner $\SPn$ over $\bbC\fps{\sigma-L}\fps{z}$ (\cref{def:fsp-wick}). So the outer amplitude is
\[
\Bigl(\frac{a(\sigma)}{a(L)}\Bigr)^{-1/2}a(\sigma)^{M_*}\,\mathrm{Poch}_{M_*}\bigl(z/a(\sigma)\bigr)\exp\Bigl(\sum_{k\geq1}\bk_k\bigl(z/a(\sigma)\bigr)^{2k-1}\Bigr)\cF(\sigma,z)=a(L)^{1/2}A(\sigma,z),
\]
and \cref{lem:fsp-fubini} gives
\begin{equation}
\label{eq:gam-step1}
\SPh_{(L,\chi_L)}\bigl[\Phi_2;\,\re^{M_*\chi}\cF\bigr]=\SPh_L\bigl[S;\,A\bigr]=c^{\cF}(z;L).
\end{equation}

\emph{Step 2: the chart.} Let $\mu$ be a second local letter and translate $\chi=\mu-\lambda/a_*$ (\cref{lem:fsp-cov}, Jacobian $1$). Put $\mu_c(L):=-\varphi'(L)/a_*$ and $p(L):=(L,\mu_c(L))$. The translation maps $p(L)$ to $(L,\chi_L)$, since $\mu_c(L)-\lambda/a_*=-\varphi'(L)/a_*+\varphi'(L)/a_*+\log a(L)=\chi_L$. Because $a$ is exactly $a_*(\sigma-s_*)$, we have $\lambda\sigma+a(\sigma)(\mu-\lambda/a_*)=\lambda s_*+a(\sigma)\mu$, and the translation turns the phase into
\[
\Phi_2(\sigma,\mu-\lambda/a_*)=\lambda s_*+\Psi_{\mathsf Q}(\sigma,\mu),\qquad \Psi_{\mathsf Q}(\sigma,\mu):=\varphi(\sigma)+a(\sigma)\mu-\mathsf Q\,\re^{\mu},
\]
where $\Psi_{\mathsf Q}$ is a phase local to this proof (as is $\Psi_0$ below), and the amplitude into $\re^{M_*(\mu-\lambda/a_*)}\cF=\mathsf Q^{M_*}\re^{M_*\mu}\cF$, with $\mathsf Q=\mathsf Q(L)=\re^{-\lambda/a_*}$ and $\mathsf Q^{M_*}:=\re^{-M_*\lambda/a_*}$. The constant $\lambda s_*$ of the phase changes only the critical value (\cref{lem:fsp-const}), and the constant factor $\mathsf Q^{M_*}$ of the amplitude cancels in $\SPh$. With \eqref{eq:gam-step1} we obtain, for $L\in U_*$,
\begin{equation}
\label{eq:gam-chart}
c^{\cF}(z;L)=\SPh_{p(L)}\bigl[\Psi_{\mathsf Q(L)};\,\re^{M_*\mu}\cF\bigr].
\end{equation}
Here $\lambda$ has disappeared, except through $\mathsf Q(L)=a(L)\re^{\varphi'(L)/a_*}$. This function is single-valued: a change of the branch of $\log a(L)$ changes $\lambda$ by an integer multiple of $2\pi\ri a_*$ and leaves $\re^{-\lambda/a_*}$ unchanged. It is holomorphic on $\bbD$, and it has a simple zero at $s_*$.

\emph{The family of critical points.} For every $L\in\bbD$, including $L=s_*$, the point $p(L)$ is a critical point of $\Psi_{\mathsf Q(L)}$: the derivative $\partial_\sigma\Psi=\varphi'(\sigma)+a_*\mu$ vanishes at $p(L)$, and $\partial_\mu\Psi=a(\sigma)-\mathsf Q(L)\re^{\mu}$ equals $a(L)-a(L)\re^{\varphi'(L)/a_*}\re^{-\varphi'(L)/a_*}=0$ there. Since $\mathsf Q(L)\re^{\mu_c(L)}=a(L)$, the Hessian at $p(L)$ is
\[
H(L)=\begin{pmatrix}\varphi''(L)&a_*\\ a_*&-a(L)\end{pmatrix},\qquad \det H(L)=-a_*\bigl(a_*+(L-s_*)\varphi''(L)\bigr)\neq 0\quad(L\in\bbD).
\]
At $L=s_*$ we have $\mathsf Q(s_*)=0$, $\Psi_0(\sigma,\mu)=\varphi(\sigma)+a(\sigma)\mu$ and $p(s_*)=(s_*,\mu_*)$ with $\mu_*:=-\varphi'(s_*)/a_*$.

\emph{Holomorphy: (ii).} In the variables $s=\sigma-L$ and $m=\mu-\mu_c(L)$ the phase $\Psi_{\mathsf Q(L)}$ is
\[
\Psi_{\mathsf Q(L)}(p(L))+\bigl[\varphi(L+s)-\varphi(L)-\varphi'(L)s\bigr]+a_*sm-a(L)\bigl(\re^{m}-1-m\bigr),
\]
because the terms of order one cancel; so its Taylor coefficients are the $\varphi^{(k)}(L)/k!$, the number $a_*$ and the numbers $-a(L)/k!$. The normalised amplitude is $\re^{M_*m}\cF(L+s,z)/\cF_0(L)$, with the Taylor coefficients $M_*^j/j!$ and $\cF_l^{(j)}(L)/(j!\,\cF_0(L))$; and the entries of $-H(L)^{-1}$ are holomorphic on $\bbD$, because $\det H(L)$ has no zero there. By \cref{lem:fsp-univ} the $n$-th coefficient of the right side of \eqref{eq:gam-chart} is a polynomial with rational coefficients in finitely many of these functions (compare \cref{lem:fsp-param}~(i)). It is therefore holomorphic on all of $\bbD$, and by \eqref{eq:gam-chart} it equals $c^{\cF}_n$ on $U_*$. This proves (ii). At $L=s_*$ the data are the Taylor data of $\Psi_0$ and of $\re^{M_*\mu}\cF$ at $(s_*,\mu_*)$ and the entries of $-H(s_*)^{-1}$; since the coefficients are given by the same universal polynomials,
\begin{equation}
\label{eq:gam-atQ0}
c^{\cF}(z;s_*)=\SPh_{(s_*,\mu_*)}\bigl[\Psi_0;\,\re^{M_*\mu}\cF\bigr].
\end{equation}

\emph{Step 3: the value.} In the variables $s:=\sigma-s_*$ and $m:=\mu-\mu_*$,
\[
\Psi_0=\varphi(s_*)+\tilde\varphi(s)+a_*sm
\]
exactly: the terms $\varphi'(s_*)s+a_*\mu_*s$ cancel, and there is no term in $m$ alone and no term $s^km$ with $k\geq 2$. The amplitude is $\re^{M_*\mu_*}\re^{M_*m}\cF(s_*+s,z)$. The constant $\varphi(s_*)$ of the phase does not matter (\cref{lem:fsp-const}), and \cref{lem:gam-deltasrc} with $P=\tilde\varphi\in s^2\bbC\fps{s}$, $c=M_*$ and $A(s,z)=\cF(s_*+s,z)$ gives $\SPn_{(s_*,\mu_*)}[\Psi_0;\,\re^{M_*\mu}\cF]=\re^{M_*\mu_*}\exp(\tilde\varphi(\delta)/z)\cF(s_*+\delta,z)$, whose constant term is $\re^{M_*\mu_*}\cF(s_*,0)$. Dividing by it, \eqref{eq:gam-atQ0} gives \eqref{eq:gam-lmi}.
\end{proof}

% Section 9, part 3 of 4: Theorem Gamma, the hatted series, the value of D. Continued in part 4.

\subsection{Theorem Gamma and the value of \texorpdfstring{$\Dn$}{D}}
\label{sec:gam-thm}

Recall $h=h_{\bw_0}$ from \eqref{eq:set-hw}, with $\bw_0=\{1-\zeta^j\}_{j=1,\dots,4}\cup\{-5\}$, and $h_n=\coef{z^n}h$. Put
\[
\mathsf T_k:=5^{1-2k}+\sum_{j=1}^4(\zeta^j-1)^{1-2k}\qquad(k\geq 1).
\]
Since $1-2k$ is odd, $\mathsf T_k=-\sum_{a\in\bw_0}a^{1-2k}$, and therefore $h(z)=\exp\bigl(\sum_{k\geq 1}\bk_k\mathsf T_kz^{2k-1}\bigr)$.

\begin{thm}[Theorem Gamma]\genusstatus{AI-proposed, unchecked}
\label{thm:gam-main}
Let $c(z;L)$ be the saddle series of \cref{def:gam-c}.
\begin{enumerate}
\item Every $c_n$ lies in $\AL$.
\item No $c_n$ has a pole at $L=1$, that is at $q=0$, and
\begin{equation}
\tag{$\Gamma$}
\label{eq:gam-Gamma}
c(z;1)=h(z),\qquad\text{that is, }c_n(1)=h_n\ \text{for every }n\geq 0.
\end{equation}
\item $c_n(\zeta L)=\zeta^{-n}c_n(L)$ for every $n$. Hence no $c_n$ has a pole at a fifth root of unity, $c_n$ lies in $\bbQ[L,L^{-1}]$, every exponent of $L$ that occurs in $c_n$ is congruent to $4n$ modulo $5$, and $c(z;\zeta^i)=h_{\bw_i}(z)=h(\zeta^{-i}z)$.
\item No $c_n$ has a pole at $L=0$, that is at $q=\infty$, and $c(z;0)=\Ginf(z)$, where
\[
\Ginf(z):=\exp\Bigl(\sum_{l\geq 1}(-1)^{l+1}\frac{B_{5l+1}(1/5)}{l(5l+1)}\,z^{5l}\Bigr),
\]
with $B_k(\theta)$ the Bernoulli polynomials; this is the series displayed on the right of (37) of \cite[Lemma~7.7]{GJR18} at $m=5$, $k=1$ and equivariant parameter $1$ in the notation there, read as a series in $z$ (the coefficients on the left of (37) are introduced in \cite[Corollary~7.6]{GJR18} through powers of $-z$); it is obtained there for their Landau--Ginzburg integral; the value $c(z;0)=\Ginf(z)$ is proved here independently (\cref{thm:gam-ginf}), and the comparison with \cite{GJR18} enters no proof. Consequently $c_n\in\bbQ[L]$ (\cref{cor:gam-poly}).
\end{enumerate}
\end{thm}

Only (ii) is used in the proof of the main theorem; (iii) and (iv) are by-products, and (iv) identifies the free coefficients (\cref{cor:gam-free}).

\begin{proof}
(i) is \cref{prop:gam-rational}. (iv) is proved in \cref{thm:gam-ginf} and \cref{cor:gam-poly} below; for the statements of (iv) these proofs use only (i), (iii) (and through it (ii)), \cref{thm:gam-lmi} and \cref{lem:fsp-dlmf5118} (the further statements of \cref{cor:gam-poly} also use \cref{thm:con-Pa}~(iv)).

(ii) We apply \cref{thm:gam-lmi} with $s_*=1$, $a_*=-1$, so that $a(\sigma)=1-\sigma$, and $M_*=0$. Let $\bbD$ be a disc around $1$ that contains none of $0,\zeta,\dots,\zeta^4$, fix branches of $\log(5\sigma)$ and of $\log(\zeta^j-\sigma)$, $j=1,\dots,4$, on $\bbD$, and put
\[
\varphi(\sigma):=5\sigma\log(5\sigma)-5\sigma+\sum_{j=1}^4\bigl[(\zeta^j-\sigma)\log(\zeta^j-\sigma)-(\zeta^j-\sigma)\bigr],
\]
\[
\cF(\sigma,z):=(5\sigma)^{1/2}\prod_{j=1}^4(\zeta^j-\sigma)^{-1/2}\exp\Bigl(\sum_{k\geq 1}\bk_kz^{2k-1}\Bigl[(5\sigma)^{1-2k}+\sum_{j=1}^4(\zeta^j-\sigma)^{1-2k}\Bigr]\Bigr),
\]
the phase and the amplitude of the five factors of the summand that do not degenerate at $\sigma=1$. Thus $\cF$ is $A_{\mathrm{St}}$ without its factor $(1-\sigma)^{-1/2}\exp(\sum_k\bk_kz^{2k-1}(1-\sigma)^{1-2k})$ of index $j=0$. The functions $\varphi$ and $\cF_m$ are holomorphic on $\bbD$, $\cF_0(1)=5^{1/2}\prod_{j\geq 1}(\zeta^j-1)^{-1/2}\neq 0$, and after shrinking $\bbD$ the function $\cF_0$ has no zero on $\bbD$ and $-1+(\sigma-1)\varphi''(\sigma)\neq 0$ there. The term $a\log a-a=(1-\sigma)\log(1-\sigma)-(1-\sigma)$ of \cref{def:gam-lmidata} is the term $j=0$ of $\mathsf S_0$, and, as $\mathrm{Poch}_0=1$, the Stirling form of $\Gamma(0+a/z)$ contributes $a^{-1/2}\exp(\sum_k\bk_k(z/a)^{2k-1})$, the factor $j=0$ of $A_{\mathrm{St}}$. Hence $S=\mathsf S_0$ and $A=A_{\mathrm{St}}$, and $\lambda(L)=-\varphi'(L)+\log(1-L)=-5\log(5L)+\sum_{j=0}^4\log(\zeta^j-L)$ is the logarithm of \cref{def:gam-c}, for the branches chosen. Therefore $c^{\cF}(z;L)=c(z;L)$ for $L\in U_*$. By \cref{thm:gam-lmi}~(ii) every $c_n$ extends holomorphically to $L=1$; since $c_n$ is rational by (i), it has no pole at $L=1$, and its value there is that of the extension. By \eqref{eq:gam-lmi} with $\delta=0$, and since the factor $(5\sigma)^{1/2}\prod_{j\geq 1}(\zeta^j-\sigma)^{-1/2}$ of $\cF$ does not involve $z$,
\[
c(z;1)=\frac{\cF(1,z)}{\cF(1,0)}=\exp\Bigl(\sum_{k\geq 1}\bk_k\mathsf T_kz^{2k-1}\Bigr)=h(z).
\]

(iii) Consider the Taylor data of $\mathsf S_0$, $\mathsf S_1$, $\mathsf S_{2k}$ at $L$ displayed after \cref{def:gam-c}. Since $p_k(\zeta L)=\zeta^{-k}p_k(L)$, the Taylor coefficients at $\zeta L$ are $\zeta^{-k}$, $\zeta^{-k}$ and $\zeta^{1-2k-r}$ times those at $L$. Substitute $t=\zeta t'$ (\cref{lem:fsp-cov}, a linear map with constant Jacobian). In the variable $t'$ the phase $\mathsf S_0(\zeta L+\zeta t')-\mathsf S_0(\zeta L)$ is $\zeta$ times $\mathsf S_0(L+t')-\mathsf S_0(L)$, the normalised factor $\exp(\mathsf S_1(\zeta L+\zeta t')-\mathsf S_1(\zeta L))$ equals $\exp(\mathsf S_1(L+t')-\mathsf S_1(L))$, and $z^{2k-1}\mathsf S_{2k}(\zeta L+\zeta t')=(\zeta^{-1}z)^{2k-1}\mathsf S_{2k}(L+t')$. So the normalised integrand at $\zeta L$ is that at $L$, with the phase multiplied by $\zeta$ and with $z$ replaced by $\zeta^{-1}z$ in the amplitude. \Cref{lem:fsp-homog} with the unit $\zeta$ gives $c(z;\zeta L)=c(\zeta^{-1}z;L)$, that is $c_n(\zeta L)=\zeta^{-n}c_n(L)$. Iterating, $c_n(L)=\zeta^{-in}c_n(\zeta^{-i}L)$; by (ii) the right side has no pole at $L=\zeta^i$, and its value there is $\zeta^{-in}h_n$. So $c(z;\zeta^i)=h(\zeta^{-i}z)$, which is $h_{\bw_i}(z)$ by \cref{lem:set-hrot}. By (i) the only possible poles of $c_n$ are $0$ and the fifth roots of unity; since there are none at the latter, $c_n\in\bbQ[L,L^{-1}]$. Writing $c_n=\sum_ea_eL^e$, the relation $c_n(\zeta L)=\zeta^{-n}c_n(L)$ gives $\zeta^e=\zeta^{-n}$ whenever $a_e\neq 0$, that is $e\equiv -n\equiv 4n$ modulo $5$.
\end{proof}

\begin{prop}[Flat and hatted saddle series]\genusstatus{AI-proposed, unchecked}
\label{prop:gam-hatflat}
For every $n$, $c_n(L)=L^{4n}\,\mathrm{ph}^{\mathrm{SP}}_n(L^{-5})$ as rational functions of $L$. Hence $\mathrm{ph}^{\mathrm{SP}}_n(1)=c_n(1)$ for every $n$, and
\[
\PhiSP(\zh,1)=h(\zh).
\]
\end{prop}

\begin{proof}
The identity of rational functions is \cref{thm:con-Pa}~(iv), which states that every $\mathrm{ph}^{\mathrm{SP}}_n$ is a polynomial in $u$. At $L=1$ it gives $\mathrm{ph}^{\mathrm{SP}}_n(1)=c_n(1)$, and $c_n(1)=h_n$ by \cref{thm:gam-main}~(ii); so $\PhiSP(\zh,1)=\sum_n\mathrm{ph}^{\mathrm{SP}}_n(1)\zh^n=h(\zh)$. The proof of \cref{thm:con-Pa} uses neither this proposition nor \cref{thm:gam-main}.
\end{proof}

\begin{cor}\genusstatus{AI-proposed, unchecked}
\label{cor:gam-equiv}
$\PhiSP(z,1)=h(z)$ holds if and only if the identity \eqref{eq:gam-Gamma} holds. With \cref{cor:con-A2}: $\Dn=1$ if and only if \eqref{eq:gam-Gamma} holds.
\end{cor}

\begin{proof}
By \cref{thm:con-Pa}~(iv), which uses neither \cref{thm:gam-main} nor \eqref{eq:gam-Gamma}, every $c_n$ is a Laurent polynomial in $L$ with $c_n(1)=\mathrm{ph}^{\mathrm{SP}}_n(1)$; so $c(z;1)=\PhiSP(z,1)$, and the first equivalence follows. The second is \cref{cor:con-A2}.
\end{proof}

\begin{thm}[Theorem C: the normalisation is canonical]\genusstatus{AI-proposed, unchecked}
\label{thm:gam-D}
In every conifold frame (\cref{def:set-frames}), $\Dn_0=1$ in $\bbQ\fps{z}$, and $\Dn=1$. Consequently $\Fg{g}$ holds for every $g\geq 2$.
\end{thm}

\begin{proof}
Fix a conifold frame. By \cref{thm:con-glue}, together with \cref{thm:con-Pb}, $\Dn_0(z)=\PhiSP(z,1)/h(z)$ exactly in $\bbQ\fps{z}$; the right side does not depend on the frame (\cref{def:con-PhiSP}). By \cref{prop:gam-hatflat}, which rests on \cref{thm:gam-main}~(ii), $\PhiSP(z,1)=h(z)$. Hence $\Dn_0=1$. By \cref{cor:gau-colours}, $\Dn_i(z)=\Dn_0(\zeta^{-i}z)=1$ for every colour $i$, so $\Dn=\diag(\Dn_i)=1$. Finally $\Fg{g}$, that is $\Dn_0=1$ modulo $z^{3g-2}$ (\cref{def:gau-Fg}), holds for every $g\geq 2$.
\end{proof}

The input (QRR) enters the proof of \cref{thm:gam-D} through \cref{thm:con-glue}, at colour $0$, and through \cref{cor:gau-colours}, whose proof uses the colour covariance of the rows (\cref{cor:gau-rowsNN}~(iii), resting on \cref{lem:gau-colrows}), that is (QRR) for the colours $i\neq0$; \cref{thm:con-glue} also uses the rows and their WKB form, which rest on (M), (G0) and the flatness part of (QRR) (\cref{def:set-rows}~(b)). The proof uses the $R$-matrix of the formal quintic, identified by (T) (\cref{thm:set-GT}), only through \cref{def:set-rows}~(b), $R|_{z=0}=1$ and (QRR).

\input{sections/s09_gamma_4}



\section{The main theorem}
\label{sec:mai}

This section assembles Theorem~A and derives Corollary~B; the consequences of $\Dn=1$ for the natural form and universality are in \cref{sec:cmp}. A conifold frame of \cref{def:set-frames}, given by $(\alpha,\beta)\in\bbQ^2$, is fixed, and $M=\Mo(\tau(u)\zh)$ is taken exactly in $u$. As in \cref{sec:flh}, $\Lambda_s\coloneqq\diag(s^{1-\nu_a})$ and $\Lambda_\varsigma\coloneqq\diag(\varsigma^{1-\nu_a})$ act on lowered components (local letters, not the leg map of \cref{def:mod-weights}), and $1_{\TFT}=\sum_i\Delta_i^{-1/2}\tilde e_i$.

\subsection{Theorem A and Corollary B}
\label{sec:mai-thm}

\begin{thm}[Theorem A]\genusstatus{AI-proposed, unchecked}
\label{thm:mai-main}
Let $g\geq2$, let $(\alpha,\beta)\in\bbQ^2$ give a conifold frame (\cref{def:set-frames}), let $\varpi_1=u+O(u^2)$ be the vanishing period and $\tc=\varpi_1/\varpi_0$. Let $M\coloneqq\Mo(\tau(u)\zh)$, exactly in $u$, and $R'\coloneqq\Qr M^{-1}$, where $\Qr$ is the matrix of the lowered rows of \cref{def:set-rows} expanded in the frame (\cref{def:set-confexp}). For a leg-free stable graph $\Gamma$ of genus $g$ let $T_\Gamma\coloneqq\mathrm{Cont}_\Gamma[E_{R'},q_{R'};M.\TFT]$ be its outer term, $\TFT$ the TFT of \cref{def:set-frob} expanded in the frame, $\TFT_{g,n}(\tilde e_i,\dots,\tilde e_i)=(\zeta^{4i}s\varpi_0)^{2g-2+n}$ (\cref{def:set-outer}). Then:
\begin{enumerate}
\item $T_\Gamma$ lies in $\bbQ\fps{u}$ for every $\Gamma$ with at least one edge;
\item for the vertex graph, $T_\Gamma\in c_g\,\tc^{2-2g}+\bbQ\fps{u}$, with
\[
c_g=\frac{(-1)^{g-1}B_{2g}\,5^{1-g}}{2g(2g-2)}
\]
as in \eqref{eq:set-cg}.
\end{enumerate}
The hatted terms are $G_\Gamma=u^{2g-2}T_\Gamma$ (\cref{lem:flh-hat}).
\end{thm}

A priori $T_\Gamma\in\bbQ(\zeta)\lser{s}$; the vertex graph is that of \cref{def:set-outer}.

\begin{proof}
Fix $g\geq2$ and the frame. By \cref{thm:gam-D}, $\Dn_0=1$ in $\bbQ\fps{z}$; in particular $\Fg{g}$ holds (\cref{def:gau-Fg}). By the comparison theorem, \cref{thm:flh-main}, $T_\Gamma=F_\Gamma$ for every leg-free stable graph $\Gamma$ of genus $g$, where $F_\Gamma$ is the term of $\Gamma$ in the flat sum of $\Pc.\Omega^{\mathrm{mod}}_{v(u)}$ (\cref{def:red-flatsum}). By \cref{thm:mod-MLK} the Model Lemma $\ML{h}$ holds for every $h\geq0$. So \cref{cor:red-can} applies: since $\ML{h}$ holds for $h\leq g-1$, $F_\Gamma$ lies in $\bbQ\fps{u}$ for every $\Gamma$ with an edge; since $\ML{g}$ holds as well, the vertex graph gives $F_\Gamma\in c_g\tc^{2-2g}+\bbQ\fps{u}$, the leaf point of $\Pc$ having the vanishing coordinate $T_{\Pc}=\tc$ (\cref{prop:tra-can}). This proves (i) and (ii). The last sentence is \cref{lem:flh-hat}.

The statements cited rest on these: \cref{thm:gam-D} on \cref{thm:con-glue}, Theorem~P (\cref{thm:con-Pa,thm:con-Pb}), \cref{thm:gam-main}~(ii), \cref{prop:gam-hatflat}, \cref{cor:gau-colours}; \cref{thm:flh-main} on \cref{thm:gau-R0}~(iv), \cref{cor:gau-colours}, \cref{prop:gau-symp}, \cref{lem:set-M0sympl}; \cref{prop:tra-can} on \cref{thm:gau-U,thm:gau-G,prop:gau-symp} and \cref{sec:tra}; all on \cref{prop:gau-quintic}, from \cref{prop:set-cyclicLR,thm:gau-cyc}.
\end{proof}

\emph{Where the inputs enter.} (QRR) enters through the rows: its flatness part (\cref{def:set-rows}~(b)) in \cref{prop:set-cyclicLR}~(iii) and \cref{cor:set-cyclicPF5}~(b), its value (\cref{thm:set-qrr}) in \cref{lem:gau-colrows} and the glue \cref{thm:con-glue}; (M) and (G0) through \cref{thm:set-mirror,prop:set-cyclicLR,cor:set-cyclicPF5}; (T) through \cref{thm:set-GT}, which identifies the $R$-matrix whose rows define $T_\Gamma$; \cite{PPZ15} through \cref{thm:set-G0} (as part of (G0)), \cref{def:set-graphsum}, \cref{lem:set-pushins} (in the reading of the topological-recursion input, \cref{prop:mod-Rtr}~(ii)), \cref{def:set-modcohft}, \cref{lem:set-modhom}, \cref{lem:set-dilaton}, \cref{def:set-outer}, \cref{sec:tra-dress}, \cref{thm:mod-tr} and \cref{sec:flh}; \cite{Eynard:2007kz}, \cite{DuninBarkowski:2012bw} (as recalled in \cite{Dunin-Barkowski:2015caa}) and \cite{ABDKS25} through \cref{thm:mod-MLK}, on the model curve $x=\Wv(p)$, $y=p$ only; and \cite[(1.7)]{AC96} through \cref{lem:set-dilaton}.

\emph{Conditionality.} Theorem~A holds modulo the quoted theorems (Q) of \cite{Eynard:2007kz}, \cite{DuninBarkowski:2012bw} (as recalled in \cite{Dunin-Barkowski:2015caa}), \cite{ABDKS25} and \cite{PPZ15} with the statement numbers listed in \cref{sec:int-inputs}, the flatness of the $R$-matrix and its value at the large-radius point (QRR), the genus-$0$ mirror theorem (M) and the genus-$0$ axioms (G0) of the twisted theory, Teleman's reconstruction with the axioms of the twisted theory in all genera (T), and the formula \cite[(1.7)]{AC96}; every other step is proved in the note. Topological recursion is used only on the model curve, never on a mirror curve of the quintic or of the formal quintic, and no finite computation enters the proof.

To expand $F_g$ (\cref{def:set-fq}) at the conifold point we use its graph-sum representative in $\Nring$ (\cref{def:set-confexp}). Write $\Psi_{\mathrm Q}\coloneqq\Qr|_{z=0}$, now for the rows as elements of $\Nring\otimes\bbQ(\zeta)$ (\cref{cor:gau-rowsNN}); its conifold expansion $\Psi_{\mathrm{cf}}\coloneqq\Cexp_{\alpha,\beta}(\Psi_{\mathrm Q})$ is the matrix $\Psi_{\mathrm Q}$ of \cref{lem:flh-inputs}, and $\Psi(u)\coloneqq\Lambda_s\Psi_{\mathrm{cf}}$ is its hatted form; $\mathfrak m_a$, $\mathfrak w_a$ are local letters.

\begin{lem}[the matrices $\Psi_{\mathrm Q}$ and $\Psi(u)$]\genusstatus{AI-proposed, unchecked}
\label{lem:mai-Psi}
\begin{enumerate}
\item In $\Nring\otimes\bbQ(\zeta)$, $(\Psi_{\mathrm Q})_{aj}=\zeta^{j(a+1)}\mathfrak m_a$, where
\[
(\mathfrak m_0,\dots,\mathfrak m_4)=\bigl(L/C_0,\;L^2/(C_0C_1),\;C_0C_1/L^2,\;C_0/L,\;1\bigr),
\]
and $\Psi_{\mathrm Q}^T\eta^{-1}\Psi_{\mathrm Q}=1$.
\item $\Psi(u)_{aj}=\zeta^{j(a+1)}\mathfrak w_a(u)$, where
\[
(\mathfrak w_0,\dots,\mathfrak w_4)=\bigl(1/\varpi_0,\;1/(\varpi_0\Dq\tc),\;\varpi_0\Dq\tc,\;\varpi_0,\;1\bigr).
\]
Hence $\Psi(u)\in GL_5(\bbQ(\zeta)\fps{u})$, $\Psi(0)=\Psi_0$, and $\Psi_{\mathrm{cf}}^T\eta^{-1}\Psi_{\mathrm{cf}}=\Psi(u)^T\eta^{-1}\Psi(u)=1$, equivalently $\Psi_{\mathrm{cf}}\Psi_{\mathrm{cf}}^T=\Psi(u)\Psi(u)^T=\eta$.
\end{enumerate}
\end{lem}

\begin{proof}
(i) At $z=0$, $\Qr^{(j)}_0(0)=\eta(1,\tilde e_j)=\Delta_j^{-1/2}=\zeta^jL/C_0$ (\cref{def:set-frob}); equivalently $\mathrm{ph}_0=1$ (\cref{def:set-rows}). By \cref{prop:set-cyclicLR}~(iii), $X_a=\zeta^{-j}\Qr^{(j)}_a$ satisfies $X_{a+1}=(C_1\kappa^{\mathrm{LR}}_a)^{-1}(z\Dq+\zeta^jL)X_a$ for $a\leq3$; at $z^0$ the term $z\Dq$ drops out, so $X_{a+1}(0)=\zeta^jL\,X_a(0)/(C_1\kappa^{\mathrm{LR}}_a)$, with $C_1\kappa^{\mathrm{LR}}=(C_1,L^5/(C_0^2C_1^2),C_1,C_0,C_0)$ and $X_0(0)=L/C_0$. This gives $\Qr^{(j)}_a(0)=\zeta^{j(a+1)}\mathfrak m_a$; only the relations of \cref{prop:set-cyclicLR} are used, so the identity holds in $\Nring\otimes\bbQ(\zeta)$ (\cref{cor:gau-rowsNN}). Since $(\eta^{-1})_{ab}=\frac15\delta_{a+b\equiv3}$ and $\mathfrak m_a\mathfrak m_{3-a}=1$ for every $a$,
\[
(\Psi_{\mathrm Q}^T\eta^{-1}\Psi_{\mathrm Q})_{ij}=\frac15\sum_a\zeta^{i(a+1)+j(4-a)}=\frac15\,\zeta^{i-j}\sum_a\zeta^{(i-j)a}=\delta_{ij}.
\]
(ii) The conifold expansion ($C_0\mapsto\varpi_0$, $C_1\mapsto\Dq\tc$, $L\mapsto s^{-1}$, \cref{lem:gau-confexp}) followed by the hatting factor $s^{1-\nu_a}$, $\nu=(0,-1,3,2,1)$, turns $\mathfrak m_a$ into $\mathfrak w_a$. At $u=0$, $\varpi_0(0)=1$ and $\Dq\tc(0)=-1$ (\cref{def:set-frames}) give $(\mathfrak w_a(0))_a=(1,-1,-1,1,1)=(\varepsilon_a)_a$, so $\Psi(0)=\Psi_0$ (\cref{def:set-M0}). The $\mathfrak w_a$ are units of $\bbQ\fps{u}$ with $\mathfrak w_a\mathfrak w_{3-a}=1$, which gives $\Psi(u)^T\eta^{-1}\Psi(u)=1$ as in (i); and $\Psi_{\mathrm{cf}}^T\eta^{-1}\Psi_{\mathrm{cf}}=\Psi(u)^T\Lambda_s^{-1}\eta^{-1}\Lambda_s^{-1}\Psi(u)=\Psi(u)^T\eta^{-1}\Psi(u)$, as $\Lambda_s\eta^{-1}\Lambda_s=\eta^{-1}$ ($\nu_a+\nu_{3-a}=2$). The equivalent forms follow from $A^{-1}=A^T\eta^{-1}$ for $A=\Psi_{\mathrm{cf}}$ and $A=\Psi(u)$.
\end{proof}

\begin{defn}[the genus-$g$ sum in the ring and its conifold expansion]
\label{def:mai-Fcf}
By \cref{cor:gau-rowsNN} the row coefficients $\Qr^{(i)}_{a,n}$ lie in $\Nring\otimes\bbQ(\zeta)$ and are free of $\logq$. By \cref{lem:mai-Psi}~(i), $\Psi_{\mathrm Q}^T\eta^{-1}\Psi_{\mathrm Q}=1$ in $\Nring\otimes\bbQ(\zeta)$. Put
\[
\tilde R\coloneqq\Psi_{\mathrm Q}^{-1}\Qr=\Psi_{\mathrm Q}^T\eta^{-1}\Qr ,
\]
the $R$-matrix in Teleman's colour-colour form; $\tilde R(0)=1$. For a matrix series $X(z)$ with $X(0)=1$ let $N_X(\psi,\psi')\coloneqq1-X^{-1}(\psi)X^{-1}(\psi')^T$ and
\[
E_X(\psi,\psi')\coloneqq\frac{N_X(\psi,\psi')-\frac12N_X(-\psi',\psi')-\frac12N_X(\psi,-\psi)}{\psi+\psi'} ,
\]
a power series for every $X$, since its numerator vanishes at $\psi=-\psi'$; it is symmetric, $E_X(\psi,\psi')^T=E_X(\psi',\psi)$, as $N_X(\psi,\psi')^T=N_X(\psi',\psi)$, and it is the edge of \cite[Def.~2.2]{PPZ15} when $X(-z)^TX(z)=1$, because then $N_X(-\psi',\psi')=N_X(\psi,-\psi)=0$ \cite[Rem.~2.1, (6)]{PPZ15}. (The symplecticity of $\tilde R$ is not known in $\Nring\otimes\bbQ(\zeta)$.) With the leaf point $q_{\tilde R}=\tilde R^{-1}(z)(-z1_{\TFT})$, $1_{\TFT}=\sum_i\Delta_i^{-1/2}\tilde e_i$, define
\[
F_g^{\Nring}\coloneqq\sum_\Gamma\mathrm{Cont}_\Gamma\bigl[E_{\tilde R},q_{\tilde R};\TFT\bigr]\in\Nring\otimes\bbQ(\zeta),
\]
the sum over the leg-free stable graphs of genus $g$ ($|\Aut\Gamma|^{-1}$ is part of $\mathrm{Cont}_\Gamma$). The \emph{conifold expansion} of $F_g$ in the frame $(\alpha,\beta)$ is
\[
F_g^{\mathrm{cf}}\coloneqq\Cexp_{\alpha,\beta}\bigl(F_g^{\Nring}\bigr)\in\bbQ\lser{s}.
\]
\end{defn}

The definition is justified as follows. First, $F_g^{\Nring}$ is a polynomial over $\bbQ$ in finitely many row coefficients, in $\Delta_i^{\pm1/2}=(\zeta^{4i}C_0/L)^{\pm1}$ and in $\zeta$: the contraction is finite because $q_{\tilde R}+z1_{\TFT}=z(1_{\TFT}-\tilde R^{-1}(z)1_{\TFT})$ is $O(z^2)$ (\cref{def:flh-outer}~(ii)), the coefficients of $\tilde R^{-1}$ are polynomials in those of $\tilde R$, and the correlators of $\TFT$ are rational multiples of $\Delta_i^{(2h-2+n)/2}$, $2h-2+n>0$. Second, $F_g^{\Nring}$ lies in $\Nring$. For $k\in(\bbZ/5)^\times$ let $\sigma_k$ be the automorphism $\zeta\mapsto\zeta^k$ of $\Nring\otimes\bbQ(\zeta)$. By \cref{cor:gau-rowsNN}~(iii), which rests on \cref{lem:gau-colrows}, $\Qr^{(i)}_{a,n}=\zeta^{i(a+1-n)}\Qr^{(0)}_{a,n}$; and the colour-$0$ representatives lie in $\Nring$: by \cref{thm:con-glue}, $\Dn_0\in1+z\bbQ\fps z$, so $A^{\mathrm{act}}_n=\sum_{m\leq n}A^{\mathrm{can}}_m\coef{z^{n-m}}\Dn_0^{-1}$ lies in $\bbQ[L,L^{-1}]$ (\cref{thm:gau-R0}~(ii), (iii)), and the representatives of \cref{cor:gau-rowsNN}~(i) are obtained from $A^{\mathrm{act}}_n(L)\,c_0^{-1}$ by a recursion with coefficients in $\Nring$. Moreover $\Delta_i^{1/2}=\zeta^{4i}C_0/L$; hence $\sigma_k(\Qr^{(i)}_{a,n})=\Qr^{(ki)}_{a,n}$, $\sigma_k(\Delta_i^{1/2})=\Delta_{ki}^{1/2}$, and, by the explicit form of $\Psi_{\mathrm Q}$ in \cref{lem:mai-Psi}, $\sigma_k(\tilde R_{ij})=\tilde R_{ki,kj}$. So $\sigma_k$ permutes the colours; as $\TFT$ is diagonal in the colour basis, $F_g^{\Nring}$ is a sum over all colourings of the vertices, which $\sigma_k$ permutes, so $\sigma_k(F_g^{\Nring})=F_g^{\Nring}$ for every $k$ and $F_g^{\Nring}\in(\Nring\otimes\bbQ(\zeta))^{\mathrm{Gal}}=\Nring$. Third, $\ev(F_g^{\Nring})=F_g$: under $\ev$, $\tilde R$ becomes the $R$-matrix of \cref{def:set-rows} in the normalised canonical basis, which is symplectic (\cref{def:set-rows}~(a)), so the image of $E_{\tilde R}$ is the edge of \cite[Def.~2.2]{PPZ15}; the graph sum is then the integral over $\overline{\mathcal M}_g$ of the genus-$g$ class of the unit-preserving action on $\TFT$ (\cref{lem:set-pushins}~(iii), unit $1_{\TFT}$), which is the potential by (T), with the normalisation (QRR) (\cref{thm:set-GT,thm:set-qrr}). A change of the primitive of $1$ in \cref{def:set-confexp} fixes the rows and $\Delta_i^{1/2}$, hence $F_g^{\Nring}$. Whether $\ev$ is injective on $\Nring$ is neither needed nor proved.

\begin{cor}[Corollary B: the conifold gap of the formal quintic]\genusstatus{AI-proposed, unchecked}
\label{cor:mai-gap}
Let $g\geq2$. In every conifold frame of \cref{def:set-frames},
\[
F_g^{\mathrm{cf}}\in c_g\,\tc^{2-2g}+\bbQ\fps{u}.
\]
Since $\tc=u+O(u^2)$ lies in $\bbQ\fps{u}$, $\bbQ\fps{u}=\bbQ\fps{\tc}$: the conifold expansion of $F_g$ contains no fractional power of $u$ and no power $\tc^k$ with $3-2g\leq k\leq-1$, and its leading term is $c_g\tc^{2-2g}$.
\end{cor}

\begin{proof}
Write $\Cexp$ for $\Cexp_{\alpha,\beta}$, extended $\bbQ(\zeta)$-linearly.

\emph{Step 1: expanding the graph sum.} By \cref{lem:gau-confexp}, $\Cexp$ is a ring homomorphism; $F_g^{\Nring}$ is a polynomial expression in finitely many coefficients of the edges of \cref{def:mai-Fcf}, of the leaf points and of $\TFT$, and the edge $E_X$, polynomial in the coefficients of $X^{-1}$ followed by an exact division by $\psi+\psi'$, commutes with ring homomorphisms. Hence
\[
F_g^{\mathrm{cf}}=\sum_\Gamma\mathrm{Cont}_\Gamma\bigl[E_{\Cexp(\tilde R)},q_{\Cexp(\tilde R)};\Cexp(\TFT)\bigr].
\]
Since $\Dn=1$ (\cref{thm:gam-D}), \cref{lem:flh-inputs}~(iii) gives $R'(-z)^T\eta^{-1}R'(z)=1$ for $R'=\Cexp(\Qr)M^{-1}$, and $M(-z)^TM(z)=1$ (\cref{lem:set-M0sympl}); so $\Cexp(\Qr)(-z)^T\eta^{-1}\Cexp(\Qr)(z)=1$. With $\Psi_{\mathrm{cf}}=\Cexp(\Psi_{\mathrm Q})$, so that $\Psi_{\mathrm{cf}}^T\eta^{-1}\Psi_{\mathrm{cf}}=1$ and $\Psi_{\mathrm{cf}}\Psi_{\mathrm{cf}}^T=\eta$ (\cref{lem:mai-Psi}), and $\Cexp(\tilde R)=\Psi_{\mathrm{cf}}^{-1}\Cexp(\Qr)$, this gives
\[
\Cexp(\tilde R)(-z)^T\,\Cexp(\tilde R)(z)=\Cexp(\Qr)(-z)^T(\Psi_{\mathrm{cf}}\Psi_{\mathrm{cf}}^T)^{-1}\Cexp(\Qr)(z)=1 :
\]
$\Cexp(\tilde R)$ is symplectic, and $E_{\Cexp(\tilde R)}$ is the edge of \cite[Def.~2.2]{PPZ15}. The classes of $\Cexp(\TFT)$ are those of $\TFT$ with $\Delta_i^{\pm1/2}$ replaced by $\Cexp(\Delta_i^{\pm1/2})$, and the axioms of a topological field theory with unit are polynomial identities in the $\Delta_i^{\pm1/2}$, so $\Cexp(\TFT)$ is one over $\bbQ(\zeta)\lser{s}$, with the unit $\Cexp(1_{\TFT})$, and $q_{\Cexp(\tilde R)}=\Cexp(\tilde R)^{-1}(z)\bigl(-z\,\Cexp(1_{\TFT})\bigr)$. By \cref{lem:set-pushins}~(iii), the right-hand side above is $\int_{\overline{\mathcal M}_g}\bigl(\Cexp(\tilde R).\Cexp(\TFT)\bigr)_{g,0}$, for the unit-preserving action of \cite[Def.~2.13]{PPZ15}.

\emph{Step 2: regrouping.} $M$ is symplectic with $M(0)=1$ (\cref{lem:set-M0sympl}), so $X\coloneqq\Cexp(\tilde R)M^{-1}$ is symplectic with $X(0)=1$, and $\Cexp(\Qr)M^{-1}=\Psi_{\mathrm{cf}}X$ is the outer matrix $R'$ of \cref{def:set-outer}. By \cite[Prop.~2.14]{PPZ15}, the unit-preserving action of \cite[Def.~2.13]{PPZ15} is a left group action, so
\[
\Cexp(\tilde R).\Cexp(\TFT)=X.\bigl(M.\Cexp(\TFT)\bigr).
\]
(The results of \cite{PPZ15} are stated over $\bbQ$; \cite[Props.~2.12 and~2.14]{PPZ15} hold over $\bbQ(\zeta)\lser s$ for the topological field theory $\Cexp(\TFT)$, whose unit has the components $y_i=\Cexp(\Delta_i^{-1/2})$, units of $\bbQ(\zeta)\lser s$, in the orthonormal basis $\tilde e_i$, and for the symplectic series $M$ and $X$, by the last paragraph of \cref{def:set-graphsum}.) Take the genus-$g$ parts without legs and integrate over $\overline{\mathcal M}_g$. The left-hand side gives $F_g^{\mathrm{cf}}$ by Step 1; the right-hand side gives $\sum_\Gamma\mathrm{Cont}_\Gamma[E_X,q_X;M.\Cexp(\TFT)]$ by \cref{lem:set-pushins}~(iii), applied to $M.\Cexp(\TFT)$, whose unit is that of $\Cexp(\TFT)$ by \cite[Prop.~2.12]{PPZ15}, and to the symplectic series $X$, $X(0)=1$; this is the contraction of \cref{def:flh-outer} with $\varkappa=1$. Since $\Psi_{\mathrm{cf}}$ is an isometry from the colour space onto $(\bbH,\eta)$ independent of $z$ which maps $1_{\TFT}$ to $1$ (\cref{lem:mai-Psi}~(ii), \cref{lem:flh-inputs}~(iv)), $R'(-\psi)^T\eta^{-1}R'(-\psi')=X(-\psi)^TX(-\psi')$ and $R'^{-1}(z)(-z1)=X^{-1}(z)(-z1_{\TFT})$; so $E_{R'}=E_X$ and $q_{R'}=q_X$: $R'^{-1}=X^{-1}\Psi_{\mathrm{cf}}^{-1}$, and $\Psi_{\mathrm{cf}}^{-1}\otimes\Psi_{\mathrm{cf}}^{-1}$ maps the inverse bivector of $(\bbH,\eta)$ to that of the colour space, $\Psi_{\mathrm{cf}}$ being an isometry. Each term is therefore $\mathrm{Cont}_\Gamma[E_{R'},q_{R'};M.\TFT]=T_\Gamma$, with $\TFT$ expanded in the frame as in \cref{def:set-outer}, and
\[
F_g^{\mathrm{cf}}=\sum_\Gamma T_\Gamma ,
\]
with no further factor $|\Aut\Gamma|^{-1}$.

\emph{Step 3.} By Theorem~A (\cref{thm:mai-main}), $T_\Gamma\in\bbQ\fps{u}$ for every $\Gamma$ with an edge, and the vertex graph gives $c_g\tc^{2-2g}+\bbQ\fps{u}$; the sum is finite, so $F_g^{\mathrm{cf}}\in c_g\tc^{2-2g}+\bbQ\fps{u}$.
\end{proof}

Corollary~B uses Theorem~A, the reconstruction (T) through the identification $\ev(F_g^{\Nring})=F_g$, \cite[\S0.2, (1); Def.~0.3, in particular (iii); (5); Rem.~2.1 and (6); Defs.~2.2 and~2.5, with the conventions stated after them; (8); Prop.~2.12; Def.~2.13; Prop.~2.14]{PPZ15}, and proved statements only; its conditionality is that of Theorem~A.

\begin{rmk}[the constant]
\label{rmk:mai-constant}
The constant $c_g$ is that of \cite[Conj.~1.5]{CGYZ26}, proved there for the quintic threefold \cite[Thm.~2]{CGYZ26}, in the same normalisation of the vanishing period; the two theories have different potentials, so nothing about the quintic follows. For the values see \cref{rmk:set-cg}; for the dictionary, the rescaling of the two periods of a frame, and the statement that $F_g$ lies in the algebra of \cite[Def.~1.1]{CGYZ26}, so that for an admissible pair $F_g^{\mathrm{cf}}$ is its conifold expansion in the sense of \cite[Def.~1.2]{CGYZ26}, see \cref{sec:sgn-lit}.
\draftnote{Numbers from arXiv:2609.13921v2, read in the HTML rendering; to be re-read in the PDF.}
\end{rmk}

\begin{rmk}[graph-locality]
\label{rmk:mai-graphlocal}
In the normalisation of \cref{rmk:mai-stored}, a computation that accompanied the work, not recorded in \cref{sec:evi}, found six edge graphs of genus $4$ singular while the singular part of the total was unchanged. If so, the singular part of the total does not determine those of the graph terms, and no argument that uses the total only through its singular part can prove Theorem~A; this is not used. The proof is graph-local, and it uses the odd part of $\log\cn$, which is fixed by (G) (\cref{thm:gau-G}) and not by symplecticity.
\end{rmk}

\begin{rmk}[the hypotheses cannot be dropped]
\label{rmk:mai-necessity}
Computations that accompanied the work, not recorded in \cref{sec:evi}, found singular edge graphs in the flat sums of gauges that violate (A), among them $P^\sharp$ of \cref{rmk:mai-stored}, and of gauges that are not regular in $u$ in the flat frame (and so are not dressings in the sense of \cref{def:tra-dressing}). This indicates, but does not prove here, that neither hypothesis of \cref{thm:red-main} can be dropped in general. For (Sp) no such example exists, because without it the edge of \cref{def:tra-edge} is not a power series. The gauge equation is not a hypothesis (\cref{rmk:mai-annih}).
\end{rmk}


% The final section: complements to the main theorem (label prefix cmp; subsection labels mai). Continued in part 2.
\section{Complements}
\label{sec:cmp}

This section records the consequences of $\Dn=1$ (\cref{thm:gam-D}) for the natural form and proves universality in flat form; none of it is used in the proofs of Theorem~A and Corollary~B. The conventions of \cref{sec:mai} are kept.

\subsection{Consequences of \texorpdfstring{$\Dn=1$}{D=1} for the natural form}
\label{sec:mai-nat}

Let $\hat\Qr(u,\zh)=\Lambda_s\Qr(u,s^4\zh)$ be the hatted rows of \cref{def:set-hat}. Its constant term is $\hat\Qr(u,0)=\Lambda_s\Psi_{\mathrm{cf}}=\Psi(u)$ (\cref{lem:mai-Psi}~(ii)), and $\tilde R(\zh;u)\coloneqq\Psi(u)^{-1}\hat\Qr(u,\zh)=\Psi(u)^T\eta^{-1}\hat\Qr(u,\zh)=\Psi_{\mathrm{cf}}^{-1}\Cexp(\Qr)(s^4\zh)$ is the conifold expansion of the matrix $\tilde R$ of \cref{def:mai-Fcf}, in the variable $\zh$.

\begin{lem}[the hatted gauge]\genusstatus{AI-proposed, unchecked}
\label{lem:mai-Ghat}
Write the graded canonical gauge as $\Pc=\sum_kz^k\sum_bp_{k,b}(u)E_{b+k,b}$ with $p_{k,b}\in\bbQ\fps{u}$ (\cref{thm:gau-U}, \cref{prop:gau-graded}, \cref{def:gau-can}), and put $\Ghat(u,\zh)\coloneqq\Lambda_s\Pc(u,s^4\zh)\Lambda_\varsigma^{-1}$. Then
\[
\Ghat_{ab}=\tau^{(1-\nu_b)/4}\sum_k\zh^k\,u^{\mu(b,k)}\,p_{k,b}(u)\,\delta_{a,b+k},\qquad\mu(b,k)\coloneqq\frac{4k+\nu_b-\nu_{b+k}}{5},
\]
where $\mu(b,k)$ is an integer with $\mu(b,k)\geq\max(0,(4k-4)/5)$, and $\mu(b,k)=0$ only for $k=0$ and for $(b,k)=(1,1)$, the entry $(2,1)$. Hence $\Ghat$ is a power series in $u$ whose coefficient of $u^n$ is a polynomial in $\zh$ of degree at most $\lfloor(5n+4)/4\rfloor$.
\end{lem}

\begin{proof}
Since $\varsigma=s\tau^{-1/4}$ (\cref{def:set-tau}), the entry $(b+k,b)$ of $\Lambda_s\Pc(u,s^4\zh)\Lambda_\varsigma^{-1}$ is $p_{k,b}\zh^k\tau^{(1-\nu_b)/4}$ times $s$ to the power $(1-\nu_{b+k})+4k-(1-\nu_b)=4k+\nu_b-\nu_{b+k}$. As $\nu_a\equiv-a\bmod5$, this exponent is $\equiv5k\equiv0\bmod5$, so it is $5\mu(b,k)$ and $s^{5\mu(b,k)}=u^{\mu(b,k)}$. The $\nu_a$ take values in $\{-1,\dots,3\}$, so $\nu_b-\nu_{b+k}\geq-4$ and $\mu(b,k)\geq(4k-4)/5$; for $k=0$ the exponent is $0$, and for $k\geq1$ it is $\geq0$. For $k\geq2$ it is positive; for $k=1$ it vanishes iff $\nu_{b+1}-\nu_b=4$, which happens only for $b=1$. At a given $u^n$ only the $k$ with $(4k-4)/5\leq n$ occur.
\end{proof}

\begin{cor}[the natural form after $\Dn=1$]\genusstatus{AI-proposed, unchecked}
\label{cor:mai-natfact}
In every conifold frame of \cref{def:set-frames}:
\begin{enumerate}
\item[(a)] $\hat\Qr(u,\zh)=\Ghat(u,\zh)\,\Rhat^{\mathrm m}(\tau(u)\zh)$. In particular $\hat\Qr$ is a power series in $u$ whose coefficients are power series in $\zh$, and $\Ghat$ is a power series in $u$ whose coefficients are polynomials in $\zh$ (\cref{lem:mai-Ghat}).
\item[(b)] $\hat\Qr(0,\zh)=\Rhat^{\mathrm m}(\zh)$: the conifold symbol of the formal quintic is that of the model. Hence $\tilde R(\zh;0)=\Mo(\zh)$ and $\Psi(0)=\Psi_0$.
\item[(c)] $\tilde R(\zh;u)=G(u,\zh)\,\Mo(\tau(u)\zh)$, with $\Mo$ on the right, where
\[
G\coloneqq\Psi(u)^{-1}\,\Ghat\,\Psi_0=\Psi_0^{-1}\,\Ghat(u,0)^{-1}\Ghat(u,\zh)\,\Psi_0 .
\]
Moreover $G(u,0)=G(0,\zh)=1$, the coefficient of every $u^n$ in $G$ is a polynomial in $\zh$, and $G(-\zh)^TG(\zh)=1$. The factor $\Mo(\tau\zh)$ is the same in every conifold frame; $G$ is built from the canonical gauge of the frame and may depend on it.
\end{enumerate}
\end{cor}

\begin{proof}
(a) By \cref{thm:gau-R0}~(iv), \cref{cor:gau-colours} and $\Dn=1$, $\Qr=\Pc\Qr^{\mathrm m}$, and $\Qr^{\mathrm m}=\Lambda_\varsigma^{-1}\Rhat^{\mathrm m}(z/\varsigma^4)$ with $z/\varsigma^4=\tau\zh$ (\cref{prop:set-modrows}, \cref{def:set-tau}, \cref{lem:set-canon}; \cref{lem:flh-inputs}~(i)). Multiplying by $\Lambda_s$ and inserting $\Lambda_\varsigma^{-1}\Lambda_\varsigma$ gives $\hat\Qr=\Ghat\,\Rhat^{\mathrm m}(\tau\zh)$. The coefficients of $\Rhat^{\mathrm m}(\tau\zh)$ lie in $\bbQ(\zeta)\fps{u}$, so the polynomiality of $\Ghat$ (\cref{lem:mai-Ghat}) makes $\hat\Qr$ a power series in $u$ with coefficients in $\bbQ(\zeta)\fps{\zh}$; in particular $\hat\Qr$ can be evaluated at $u=0$.

(b) At $\zh=0$, (a) gives $\Psi(u)=\Ghat(u,0)\Psi_0$, since $\Rhat^{\mathrm m}(0)=\Psi_0$ ($\Fs_b(0)=\varepsilon_b$, \cref{prop:set-modrows}); with \cref{lem:mai-Psi}, $\Ghat(0,0)=\Psi(0)\Psi_0^{-1}=1$. By \cref{lem:mai-Ghat} the only entry of $\Ghat$ at $u^0$ off the diagonal is the entry $(2,1)$, with coefficient $\zh\,p_{1,1}(0)$, and $p_{1,1}(0)=0$ by \cref{lem:gau-R1}: column $1$ of $\Pc(0,z)$ is $(\Pc)_{11}(0,z)e_1$. So $\Ghat(0,\zh)=1$, and $\tau(0)=1$ gives $\hat\Qr(0,\zh)=\Rhat^{\mathrm m}(\zh)$. Then $\tilde R(\zh;0)=\Psi_0^{-1}\Rhat^{\mathrm m}(\zh)=\Mo(\zh)$ by \cref{def:set-M0}.

(c) $\tilde R=\Psi(u)^{-1}\Ghat\,\Rhat^{\mathrm m}(\tau\zh)=[\Psi(u)^{-1}\Ghat\,\Psi_0][\Psi_0^{-1}\Rhat^{\mathrm m}(\tau\zh)]=G\,\Mo(\tau\zh)$. The second expression for $G$ follows from $\Psi(u)=\Ghat(u,0)\Psi_0$; $\Ghat(u,0)=\diag(\tau^{(1-\nu_b)/4}p_{0,b})$ is invertible over $\bbQ\fps{u}$, the $p_{0,b}$ being units (\cref{thm:gau-U}), so every $u^n$-coefficient of $G$ is a polynomial in $\zh$. $G(u,0)=1$ is clear, and $G(0,\zh)=1$ is (b). For the symplecticity, \cref{lem:flh-inputs}~(iii) with $\Dn=1$ and \cref{lem:set-M0sympl} give $\Qr(-z)^T\eta^{-1}\Qr(z)=1$, hence $\hat\Qr(-\zh)^T\eta^{-1}\hat\Qr(\zh)=1$ because $\Lambda_s\eta^{-1}\Lambda_s=\eta^{-1}$, and $\tilde R(-\zh)^T\tilde R(\zh)=1$ because $\Psi(u)^{-T}\Psi(u)^{-1}=\eta^{-1}$ (\cref{lem:mai-Psi}). Therefore $G(-\zh)^TG(\zh)=\Mo(-\tau\zh)^{-T}\Mo(\tau\zh)^{-1}=1$ by \cref{lem:set-M0sympl}. Neither $\Mo$ nor $\tau$ involves the frame.
\end{proof}

Only the symplecticity in (c) sees $\Dn$, through the even part of $\log\Dn_0$: for a scalar $\Dn=\cn(z)\in1+z^5\bbQ\fps{z^5}$ the proof gives the rest with $\Pc\cn(z)^{-1}$ in place of $\Pc$, as $z^5=u^4\zh^5$ (\cref{rmk:mai-stored}).

\begin{cor}[the actual gauge is canonical]\genusstatus{AI-proposed, unchecked}
\label{cor:mai-actual}
Let $P^{\mathrm{act}}\coloneqq\Qr\,(\Qr^{\mathrm m})^{-1}$, the matrix series that carries the model columns at $v(u)$ to the rows of the formal quintic, both with their exponentials stripped (\cref{thm:gau-R0}). Then $P^{\mathrm{act}}=\Pc$. In particular $P^{\mathrm{act}}$ is a graded regular gauge with coefficients in $\bbQ\fps{u}$, $P^{\mathrm{act}}_{11}(0,z)=1$, and $P^{\mathrm{act}}\YC(v(u))=\YCQ$.
\end{cor}

\begin{proof}
$\Qr^{\mathrm m}$ is invertible because $\Qr^{\mathrm m}(0)=\Psi^{\mathrm m}$ is (\cref{lem:flh-inputs}~(i)). By \cref{thm:gau-R0}~(iv) with \cref{cor:gau-colours}, $\Qr\Dn=\Pc\Qr^{\mathrm m}$, and $\Dn=1$ by \cref{thm:gam-D}; so $P^{\mathrm{act}}=\Pc$. The remaining assertions are \cref{def:gau-can} and \cref{thm:gau-G}.
\end{proof}

With \cref{thm:gam-D} this is Theorem~C (\cref{thm:int-D}); it is not circular, as the rows are normalised at $u=1$ and the canonical gauge at $u=0$.

\subsection{Universality}
\label{sec:mai-univ}

The flat half of the proof of Theorem~A, \cref{cor:red-can} with \cref{thm:mod-MLK}, uses the formal quintic only through its cyclic system and the canonical gauge, and applies to every self-dual cyclic system with one simple pole whose product is that of a pulled-back model; the comparison (\cref{thm:flh-main}) and the normalisation (\cref{thm:gam-D}) are specific to the formal quintic. The systems are over $\bbQ$ (\cref{def:gau-onepole}), and for a non-zero series $y\in u\bbQ\fps{u}$ (a local letter), $\tilde A_y$ denotes the model pulled back by $y$, with coefficients $(y',-y'/y,y',-y',-y')$ in $\bbQ\lser{u}$ and product $-y'^5/y$ (\cref{prop:gau-equalprod}); it has the one-pole shape if and only if $y'(0)\neq0$.

\begin{lem}[the pole order]\genusstatus{AI-proposed, unchecked}
\label{lem:mai-pole}
Let $A=\sum_ab_aE_{a,a+1}$ with $b_a\in\bbQ\lser{u}$ non-zero, $\Bprod\coloneqq\prod_ab_a$, and let $y\in u\bbQ\fps{u}$ be non-zero of order $n\geq1$.
\begin{enumerate}
\item If the gauge equation $z\,\partial_uP=AP-P\tilde A_y$ has a solution whose $z^0$-coefficient $P_0$ is non-zero, then $\Bprod=-y'^5/y$.
\item In that case $\ord_0\Bprod=4n-5\in\{-1,3,7,\dots\}$. If $\Bprod$ has a pole at $u=0$, then $n=1$, $\ord_0\Bprod=-1$ and $y'(0)^4=-\Res_{u=0}\Bprod$.
\end{enumerate}
\end{lem}

\begin{proof}
(i) At $z^0$ the gauge equation reads $AP_0=P_0\tilde A_y$, hence $A^5P_0=P_0\tilde A_y^5$, with $A^5=\Bprod\cdot1$ and $\tilde A_y^5=-(y'^5/y)\cdot1$; so $(\Bprod+y'^5/y)P_0=0$ over the field $\bbQ\lser{u}$, and $P_0\neq0$ gives the claim (the necessity clause of \cref{prop:gau-equalprod}). (ii) $y'^5/y$ has order $5(n-1)-n=4n-5$; a pole forces $n=1$, and then $y=y'(0)u+O(u^2)$ and $y'^5/y=y'(0)^4u^{-1}+O(1)$, so $\Res_{u=0}\Bprod=-y'(0)^4$.
\end{proof}

\begin{cor}[universality, flat form]\genusstatus{AI-proposed, unchecked}
\label{cor:mai-univ}
Let $v_A\in u\bbQ\fps{u}$ with $v_A'(0)\neq0$, and let $A$ be a self-dual one-pole system over $\bbQ\fps u$ (\cref{def:gau-onepole}) with product $-v_A'^5/v_A$ and residue $r$. Then $r=-1/\lambda^2$ with $\lambda\coloneqq-b_0(0)b_3(0)/v_A'(0)^2\in\bbQ^\times$ (a letter local to this corollary and its proof), because $b_1=\Bprod/(b_0^2b_3^2)$ by self-duality and $(u\Bprod)(0)=-v_A'(0)^4$ (\cref{lem:mai-pole}~(ii) with $y=v_A$). Let $P^A_{\mathrm{can}}$ be the canonical gauge from $\tilde A_{v_A}$ to $A$, that is, the regular gauge with $(P^A_{\mathrm{can}})_{11}(0,z)=1$, which exists, is unique, is graded and has rational coefficients, by the argument of \cref{def:gau-can}: a rational graded regular gauge $P$ exists (\cref{prop:gau-graded}~(i)), $P_{11}(0,z)$ is a unit of $\bbQ\fps{z^5}$ (\cref{prop:gau-graded}~(iii)), so $P\,P_{11}(0,z)^{-1}$ is a graded regular gauge (\cref{prop:gau-graded}~(ii)), and it is the only regular gauge with this normalisation and has rational coefficients (\cref{thm:gau-U}~(b), (d)); and let $P_A\coloneqq\lambda^{-1}P^A_{\mathrm{can}}$. Then for every $g\geq2$:
\begin{enumerate}
\item the term of every leg-free stable graph of genus $g$ with an edge in the flat sum of $P_A.\Omega^{\mathrm{mod}}_{v_A(u)}$ lies in $\bbQ\fps{u}$;
\item the term of the vertex graph lies in $c_g(T_A/\lambda)^{2-2g}+\bbQ\fps{u}$, where $T_A\coloneqq\int_0^ub_0\,\dd u$.
\end{enumerate}
\end{cor}

For the formal quintic in a conifold frame, $A$ is the system of \cref{prop:gau-quintic}, $v_A=v(u)$, $r=-1$, $\lambda=1$ ($b_0(0)=1$, $b_3(0)=-1$, $v'(0)=1$), $P_A=\Pc$ and $T_A=\tc$ (\cref{lem:gau-vanq}~(d)), and the corollary is \cref{cor:red-can} with \cref{thm:mod-MLK}. The sign of $\lambda$ is a convention: the proof uses only $\lambda^2=-1/r$, and the flat sums of $P$ and of $-P$ agree term by term, since $R$, $E_P$ and $\tilde q_T$ are unchanged, $d$ changes sign, and the vertex tensors of \cref{def:red-flatsum} change by $(-1)^{n_v}$ with $\sum_vn_v$ even.

\begin{proof}
For the pair $(A,\tilde A_{v_A})$ the residue of $\tilde A_{v_A}$ is $\tilde r=-1$. By \cref{lem:gau-R1}, column $1$ of $P(0,z)$ is $P_{11}(0,z)e_1$ and $rP_{22}(0,z)=\tilde rP_{11}(0,z)$ for every regular gauge $P$. The proof of \cref{prop:gau-symp} uses the residues only in its last step, the evaluation of the entry $(2,1)$ at $u=0$ by \cref{lem:gau-R1}; with the general form of that lemma it gives $P(-z)^T\eta P(z)=(\tilde r/r)\,P_{11}(0,z)P_{11}(0,-z)\,\eta$. For $P_A$ one has $P_{11}(0,z)=\lambda^{-1}$ and $\tilde r/r=\lambda^2$, so $P_A$ is symplectic. By \cref{thm:gau-G}~(c), $P^A_{\mathrm{can}}\YC(v_A(u))=Y^{\mathrm C}_A$, since $\YC(v_A(u))$ is the vanishing solution of $\tilde A_{v_A}$ (\cref{lem:gau-vanmod}, the chain rule, \cref{lem:gau-vanshape}, whose last claim holds for every $v\in u\bbQ\fps u$, and \cref{lem:gau-uniq}, as in the proof of \cref{thm:gau-G}); so $P_A\YC(v_A(u))=\lambda^{-1}Y^{\mathrm C}_A$ has no positive power of $z$ (\cref{lem:gau-vanq}), which is (A), and $T_{P_A}=T_A/\lambda$ by \cref{lem:gau-vanq}~(d). $P_A$ is graded with coefficients in $\bbQ\fps{u}$ and invertible constant term (\cref{thm:gau-U}, \cref{prop:gau-graded}), so it is a dressing (\cref{def:tra-dressing}) with (Sp) and (A). \Cref{thm:red-main} with \cref{thm:mod-MLK} gives (i) and (ii).
\end{proof}

\begin{rmk}[non-gauge dressings]
\label{rmk:mai-annih}
\Cref{thm:red-main} does not use the gauge equation: with \cref{thm:mod-MLK} it gives, for every dressing $P$ with (Sp) and (A) (\cref{def:tra-dressing}) and every $g\geq2$, regular edge graphs in the flat sum of $P.\Omega^{\mathrm{mod}}_{v(u)}$ and the vertex graph in $c_gT_P^{2-2g}+\bbQ\fps{u}$. Only the flat form is claimed: a hatted statement would need an identity like \cref{lem:flh-inputs}~(ii), available only for the formal quintic.
\end{rmk}

\begin{rmk}[a non-canonical normalisation]
\label{rmk:mai-stored}
Let $P^\sharp\coloneqq\Pc\,\cn(z)$, $\cn\in1+z^5\bbQ\fps{z^5}$, be the regular gauge of the other normalisation of \cref{sec:sgn-norm}, and write $\cn(z)=1+m_5z^5+m_{10}z^{10}+\cdots$ ($P^\sharp$ is a local letter). By \cref{thm:gau-U}, $P^\sharp$ is a graded regular gauge. By \cref{prop:gau-symp}, $P^\sharp(-z)^T\eta P^\sharp(z)=\cn(z)\cn(-z)\eta$ with $\cn(z)\cn(-z)=1-(m_5^2-2m_{10})z^{10}+\cdots$: $P^\sharp$ is symplectic modulo $z^{10}$, and symplecticity constrains only the even part of $\log\cn$. By \cref{thm:gau-G}, $P^\sharp\YC=\cn(z)\YCQ$, whose part of positive degree in $z$ begins with $m_5(z^5\YCQ)_{+}=-m_5(z^5e_1+\tc z^4e_0+\cdots)$; by \cref{thm:gau-G}~(b), $P^\sharp$ satisfies (A) only if $\cn=1$. So the gauge equation does not imply (A), by the freedom $\cn(z^5)$ of \cref{thm:gau-U}; (G) fixes $\cn$ altogether, in particular the odd part of $\log\cn$, which symplecticity leaves free. The value of $m_5$ is known only by computation.
\end{rmk}



\appendix
% Appendix: the proof of parts (i) and (ii) of prop:set-cyclicLR, the large-radius cyclic structure.
\section{The large-radius cyclic structure}
\label{sec:set-lrproof}

This appendix proves parts (i) and (ii) of \cref{prop:set-cyclicLR} from (M), (G0), \cref{lem:set-pf5} and the pairing of \cref{def:set-fq}, in the notation of \cref{sec:set-fq}. Every step is an identity over $\bbQ\fps q$ or $\bbQ\fps q[\logq]$ (in \cref{lem:set-periods} with the further variables $\mathfrak w$, $z^{-1}$, in \cref{lem:set-order3} over $\bbQ\lser q$); no truncation in $q$, analytic continuation or conifold frame is used. The letters $\mathfrak w$, $\cP$, $\mathfrak p_k$, $\tilde{\mathfrak p}_m$, $\mathsf p^{(a)}_k$ and $\chi$ are local. Before \cref{prop:set-cyclicLR}~(i) is proved, $\kappa^{\mathrm{LR}}_a$ denotes the structure constant of \cref{lem:set-grading}.

\begin{lem}[The $\bbZ/5$ grading]\genusstatus{AI-proposed, unchecked}
\label{lem:set-grading}
At $\lambda_i=\zeta^i$ and $t=TH$, $\langle\!\langle H^a,H^b,H^c\rangle\!\rangle=0$ unless $a+b+c\equiv3\bmod5$. Hence $H^a*H^b$ is a multiple of $H^{a+b}$, and $H*H^a=\kappa^{\mathrm{LR}}_aH^{a+1}$ with $\kappa^{\mathrm{LR}}_a\in1+q\bbQ\fps q$.
\end{lem}

\begin{proof}
By the divisor equation (G0)(b), with $\Qn=1$,
\[
\langle\!\langle H^a,H^b,H^c\rangle\!\rangle(TH)=\eta(H^aH^b,H^c)+\sum_{d\geq1}\re^{dT}\langle H^a,H^b,H^c\rangle_{0,3,d};
\]
in degree $0$ the insertions of $TH$ contribute nothing, as $\overline{\mathcal M}_{0,3+n}(\bbP^4,0)=\overline{\mathcal M}_{0,3+n}\times\bbP^4$ with the integrand pulled back from $\bbP^4$. By (G0)(a), before the specialisation the term of degree $d\geq1$ is a polynomial in the elementary symmetric functions $e_1,\dots,e_5$ of the weights, of weighted degree $m=a+b+c-3$. At $\lambda_j=\zeta^j$, $\prod_j(x-\zeta^j)=x^5-1$ gives $e_1=\dots=e_4=0$ and $e_5=1$, so the value is $0$ unless $5\mid m$. The classical term is $\eta(H^{a+b},H^c)=5\,\delta_{a+b+c\equiv3}$. Since $\eta^{cc'}$ is supported on $c+c'\equiv3$, $H^a*H^b$ is supported on $H^{a+b}$. The structure constant $\kappa^{\mathrm{LR}}_a$ is a power series in $\re^T$ with rational coefficients, hence lies in $\bbQ\fps q$, and $\kappa^{\mathrm{LR}}_a(0)=1$ since at $q=0$ the product is the classical one.
\end{proof}

\begin{lem}[Unit and Frobenius property]\genusstatus{AI-proposed, unchecked}
\label{lem:set-unitfrob}
$\kappa^{\mathrm{LR}}_0=\kappa^{\mathrm{LR}}_2=1$ and $\kappa^{\mathrm{LR}}_3=\kappa^{\mathrm{LR}}_4$.
\end{lem}

\begin{proof}
The class $1$ is the unit, so $\kappa^{\mathrm{LR}}_0=1$. By the Frobenius property, $5\kappa^{\mathrm{LR}}_a=\eta(H*H^a,H^{2-a})=\eta(H^a,H*H^{2-a})=5\kappa^{\mathrm{LR}}_{2-a}$ for all $a\in\bbZ/5$, which gives $\kappa^{\mathrm{LR}}_0=\kappa^{\mathrm{LR}}_2$ and $\kappa^{\mathrm{LR}}_3=\kappa^{\mathrm{LR}}_4$.
\end{proof}

\begin{lem}[The periods]\genusstatus{AI-proposed, unchecked}
\label{lem:set-periods}
Let $\mathfrak w$ be an indeterminate and
\[
\cP(\mathfrak w,q)\coloneqq\sum_{d\geq0}q^d\,\frac{\prod_{k=1}^{5d}(5\mathfrak w+k)}{\prod_{k=1}^d\bigl((\mathfrak w+k)^5-\mathfrak w^5\bigr)}\in\bbQ\fps q\fps{\mathfrak w},
\]
the denominators being polynomials in $\mathfrak w$ with the constant term $(d!)^5$. The map $\bbQ\fps q[\logq]\fps{\mathfrak w}\to\bbH\otimes\bbQ\fps q[\logq]\fps{z^{-1}}$, $\sum_ka_k\mathfrak w^k\mapsto\sum_ka_kH^kz^{-k}$ (with $H^k\coloneqq H^{k\bmod5}$), is an injective ring homomorphism, and $I/z$ is the image of $\re^{\logq\mathfrak w}\cP(\mathfrak w,q)$. Hence
\[
\frac Iz=\sum_{k\geq0}\mathfrak p_kH^kz^{-k},\qquad\mathfrak p_k\coloneqq\sum_{j=0}^k\frac{\logq^j}{j!}\,\tilde{\mathfrak p}_{k-j},\qquad\tilde{\mathfrak p}_m\coloneqq\coef{\mathfrak w^m}\cP\in\bbQ\fps q,\quad\tilde{\mathfrak p}_m(0)=\delta_{m0},
\]
with $\mathfrak p_0=C_0$ and $\mathfrak p_1=I_1$. Moreover $\PFfour\mathfrak p_k=0$ for $k=0,\dots,3$, and $\PFfour\mathfrak p_4=1$.
\end{lem}

\begin{proof}
The map is multiplicative because $H^5=1$, and injective because distinct powers of $\mathfrak w$ go to distinct powers of $z$; followed by the restriction to $p_i$ it is $G(\mathfrak w)\mapsto G(\zeta^i/z)$. As $\prod_j(\zeta^i-\zeta^j+kz)=(\zeta^i+kz)^5-1$, the powers $z^{-5d}$ cancel in each term $q^d$, and with $q^{\zeta^i/z}=\re^{\logq\zeta^i/z}$ the restriction of the image of $\re^{\logq\mathfrak w}\cP$ to $p_i$ is $I_i/z$; the five restrictions determine the element. Expanding gives $\mathfrak p_k$, and the term $q^0$ of $\cP$ is $1$. The constant term of $\cP$ in $\mathfrak w$ is $C_0$, and its coefficient of $\mathfrak w$ is, by logarithmic differentiation, $\sum_dq^d\frac{(5d)!}{(d!)^5}\bigl(\sum_{k=1}^{5d}\frac5k-\sum_{k=1}^d\frac5k\bigr)=\tilde I_1$; so $\mathfrak p_0=C_0$ and $\mathfrak p_1=I_1$.

\emph{The operator $\PFfive$.} On the term $q^d$ of $I_i$, which is $q^{d+\zeta^i/z}$ times a function of $z$, $z\Dq$ acts by $\zeta^i+dz$, so $(z\Dq)^5-1=\prod_j(z\Dq-\zeta^j)$ multiplies it by $\prod_j(\zeta^i-\zeta^j+dz)$; this cancels the factors $k=d$ of the denominator (and kills the term $d=0$), and leaves $q\prod_{k=1}^5(5\zeta^i+(5d-5+k)z)$ times the term $q^{d-1}$, which is what $q\prod_{k=1}^5(5z\Dq+kz)$ gives on the term $q^{d-1}$. Hence $\PFfive I_i=0$. By \cref{lem:set-pf5} and $I_i/z=\sum_k\zeta^{ik}\mathfrak p_kz^{-k}$,
\[
0=\PFfive(I_i/z)=\sum_{k\geq0}\zeta^{ik}\bigl(z^{5-k}\,\Dq\PFfour\mathfrak p_k-z^{-k}\mathfrak p_k\bigr),
\]
and the coefficients of $z^5,\dots,z^1$ give $\Dq\PFfour\mathfrak p_k=0$ for $k=0,\dots,4$.

\emph{Constants.} The constants of $\Dq$ on $\bbQ\fps q[\logq]$ are $\bbQ$: on $\bbQ\fps q$ this is clear from $\Dq q^n=nq^n$; and if $f=\sum_{j\leq N}a_j\logq^j$ with $a_N\neq0$, $N\geq1$ and $\Dq f=0$, then $a_N\in\bbQ$ and $\Dq a_{N-1}=-Na_N\neq0$, impossible as $\Dq$ maps $\bbQ\fps q$ into $q\bbQ\fps q$. Hence $\PFfour\mathfrak p_k\in\bbQ$ for $k\leq4$, equal to its part of order $q^0$. Taking this part commutes with $\Dq$, which acts on it as $\dd/\dd\logq$, and kills the multiples of $q$; by \eqref{eq:set-PF4} and $\tilde{\mathfrak p}_m(0)=\delta_{m0}$ it maps $\PFfour\mathfrak p_k$ to $(\dd/\dd\logq)^4(\logq^k/k!)$, which is $0$ for $k\leq3$ and $1$ for $k=4$.
\end{proof}

\begin{lem}[The Birkhoff recursion]\genusstatus{AI-proposed, unchecked}
\label{lem:set-birkhoff}
Put $\mathsf p^{(0)}_k\coloneqq\mathfrak p_k/C_0$ for $k\geq0$ and, recursively for $a=0,\dots,3$, $\mathsf p^{(a+1)}_k\coloneqq\partial_T\mathsf p^{(a)}_k/\kappa^{\mathrm{LR}}_a$ for $k\geq a+1$. Then, along the line $t=TH$,
\[
J_a=\sum_{k\geq a}\mathsf p^{(a)}_k\,H^kz^{a-k},\qquad\mathsf p^{(a)}_a=1\qquad(a=0,\dots,4),
\]
and $\kappa^{\mathrm{LR}}_a=\partial_T\mathsf p^{(a)}_{a+1}$ for $a=0,\dots,3$. In particular $\mathsf p^{(0)}_1=T$, $\kappa^{\mathrm{LR}}_1=\partial_T^2\mathsf p^{(0)}_2$, $\kappa^{\mathrm{LR}}_2=\partial_T\bigl(\partial_T^2\mathsf p^{(0)}_3/\kappa^{\mathrm{LR}}_1\bigr)$ and, since $\kappa^{\mathrm{LR}}_2=1$, $\kappa^{\mathrm{LR}}_3=\partial_T^2\bigl(\partial_T^2\mathsf p^{(0)}_4/\kappa^{\mathrm{LR}}_1\bigr)$.
\end{lem}

\begin{proof}
Induction on $a$. By (G0)(d), (M) and \cref{lem:set-periods}, $J_0=J/z=I/(zC_0)=\sum_k(\mathfrak p_k/C_0)H^kz^{-k}$ along the line, with $\mathsf p^{(0)}_0=1$; distinct $k$ carry distinct powers of $z$. By (G0)(d) with $b=1$ and \cref{lem:set-grading}, $z\partial_TJ_a=\kappa^{\mathrm{LR}}_aJ_{a+1}$ for $a\leq3$; only derivatives along the line enter. Differentiating the expansion of $J_a$ gives $z\partial_TJ_a=\sum_{k\geq a+1}(\partial_T\mathsf p^{(a)}_k)H^kz^{a+1-k}$, and since $J_{a+1}=H^{a+1}+O(z^{-1})$ the coefficient of $z^0$ gives $\kappa^{\mathrm{LR}}_a=\partial_T\mathsf p^{(a)}_{a+1}$; as $\kappa^{\mathrm{LR}}_a$ is a unit (\cref{lem:set-grading}), $J_{a+1}=z\partial_TJ_a/\kappa^{\mathrm{LR}}_a$ gives the expansion of $J_{a+1}$, with $\mathsf p^{(a+1)}_{a+1}=1$. Finally $\mathsf p^{(0)}_1=I_1/C_0=T$, $\kappa^{\mathrm{LR}}_0=\partial_TT=1$, and the closed forms follow from $\kappa^{\mathrm{LR}}_a=\partial_T\mathsf p^{(a)}_{a+1}$ and \cref{lem:set-unitfrob}.
\end{proof}

\begin{lem}[The Wronskian]\genusstatus{AI-proposed, unchecked}
\label{lem:set-wronskian}
$\Wr(\mathfrak p_0,\dots,\mathfrak p_3)\coloneqq\det\bigl(\Dq^j\mathfrak p_k\bigr)_{j,k=0,\dots,3}=L^{10}$.
\end{lem}

\begin{proof}
By \eqref{eq:set-PF4} the leading coefficient of $\PFfour$ is $1-5^5q=L^{-5}$ and its coefficient of $\Dq^3$ is $-6250q$. Differentiating the determinant row by row (the rows $j<3$ give two equal rows) and replacing $\Dq^4\mathfrak p_k$ in the last row by means of $\PFfour\mathfrak p_k=0$ ($k\leq3$, \cref{lem:set-periods}), whose terms of order $\leq2$ are combinations of the other rows, gives $\Dq\Wr=6250qL^5\,\Wr$; also $\Dq L^{10}=10L^9\Dq L=6250qL^5\,L^{10}$. So $\Wr\,L^{-10}$ is a constant of $\Dq$, hence (proof of \cref{lem:set-periods}) equal to its part of order $q^0$. Taking this part is a ring homomorphism to $\bbQ[\logq]$, on which $\Dq$ acts as $\dd/\dd\logq$; it maps $\Dq^j\mathfrak p_k$ to $\logq^{k-j}/(k-j)!$ for $k\geq j$ and to $0$ for $k<j$, so the determinant to $1$.
\end{proof}

\begin{lem}[Operators of order at most three]\genusstatus{AI-proposed, unchecked}
\label{lem:set-order3}
If $\Theta=\sum_{j=0}^3r_j\Dq^j$ with $r_j\in\bbQ\lser q$ satisfies $\Theta\mathfrak p_k=0$ for $k=0,\dots,3$, then $\Theta=0$.
\end{lem}

\begin{proof}
As the $r_j$ do not involve $\logq$ and $\Dq\re^{\logq\mathfrak w}=\mathfrak w\re^{\logq\mathfrak w}$, \cref{lem:set-periods} gives, coefficientwise in $\mathfrak w$,
\[
\sum_k(\Theta\mathfrak p_k)\mathfrak w^k=\Theta\bigl(\re^{\logq\mathfrak w}\cP\bigr)=\re^{\logq\mathfrak w}\,\cP_\Theta,\qquad\cP_\Theta\coloneqq\sum_{j=0}^3r_j(\Dq+\mathfrak w)^j\cP\in\bbQ\lser q\fps{\mathfrak w}.
\]
The hypothesis says that the left side is $O(\mathfrak w^4)$, hence so is $\cP_\Theta$. Now $\coef{\mathfrak w^n}\cP_\Theta=\sum_{j=0}^3m_{nj}r_j$ with $m_{nj}\coloneqq\sum_{i\leq\min(j,n)}\binom ji\Dq^{j-i}\tilde{\mathfrak p}_{n-i}\in\bbQ\fps q$, and $m_{nj}\equiv\delta_{nj}\bmod q$ because $\tilde{\mathfrak p}_m(0)=\delta_{m0}$ and $\Dq$ kills $q^0$. So $\det(m_{nj})$ is a unit, and $r_j=0$.
\end{proof}

\begin{proof}[Proof of \cref{prop:set-cyclicLR}~(i) and~(ii)]
(i) By \cref{lem:set-unitfrob}, $\kappa^{\mathrm{LR}}_0=\kappa^{\mathrm{LR}}_2=1$ and $\kappa^{\mathrm{LR}}_3=\kappa^{\mathrm{LR}}_4$.

\emph{The coefficient $\kappa^{\mathrm{LR}}_1$.} By \cref{lem:set-periods,lem:set-birkhoff}, $\mathfrak p_k=C_0\mathsf p^{(0)}_k$ with $\mathsf p^{(0)}_0=1$, $\mathsf p^{(0)}_1=T$. For any four functions $f_k$, $\det(\Dq^j(C_0f_k))=C_0^4\det(\Dq^jf_k)=C_0^4C_1^6\det(\partial_T^jf_k)$, since $\Dq^j(C_0f)-C_0\Dq^jf$ and $\Dq^j-C_1^j\partial_T^j$ are of lower order with coefficients independent of $f$. With $f=(1,T,\mathsf p^{(0)}_2,\mathsf p^{(0)}_3)$ and \cref{lem:set-birkhoff},
\[
\det(\partial_T^jf_k)=\partial_T^2\mathsf p^{(0)}_2\,\partial_T^3\mathsf p^{(0)}_3-\partial_T^3\mathsf p^{(0)}_2\,\partial_T^2\mathsf p^{(0)}_3=\bigl(\partial_T^2\mathsf p^{(0)}_2\bigr)^2\,\partial_T\Bigl(\frac{\partial_T^2\mathsf p^{(0)}_3}{\partial_T^2\mathsf p^{(0)}_2}\Bigr)=(\kappa^{\mathrm{LR}}_1)^2\kappa^{\mathrm{LR}}_2 .
\]
So \cref{lem:set-wronskian} and $\kappa^{\mathrm{LR}}_2=1$ give $(\kappa^{\mathrm{LR}}_1)^2=(L^5/C_0^2C_1^3)^2$ in the domain $\bbQ\fps q$; as $\kappa^{\mathrm{LR}}_1(0)=1=L(0)=C_0(0)=C_1(0)$ (\cref{lem:set-grading}), $\kappa^{\mathrm{LR}}_1=L^5/(C_0^2C_1^3)$.

\emph{The coefficient $\kappa^{\mathrm{LR}}_3$.} The operator $\Mcyc^{\mathrm{LR}}=\partial_T^2\circ(\kappa^{\mathrm{LR}}_1)^{-1}\circ\partial_T^2\circ C_0^{-1}$ has coefficients in $\bbQ\fps q$. By \cref{lem:set-birkhoff}, $\Mcyc^{\mathrm{LR}}\mathfrak p_0=\Mcyc^{\mathrm{LR}}\mathfrak p_1=0$ (as $\partial_T^21=\partial_T^2T=0$), $\Mcyc^{\mathrm{LR}}\mathfrak p_2=\partial_T^2\bigl((\kappa^{\mathrm{LR}}_1)^{-1}\kappa^{\mathrm{LR}}_1\bigr)=0$, $\Mcyc^{\mathrm{LR}}\mathfrak p_3=\partial_T\kappa^{\mathrm{LR}}_2=0$ and $\Mcyc^{\mathrm{LR}}\mathfrak p_4=\kappa^{\mathrm{LR}}_3$. In powers of $\Dq$, $\Mcyc^{\mathrm{LR}}=(\kappa^{\mathrm{LR}}_1C_0C_1^4)^{-1}\Dq^4+(\text{order}\leq3)$, and $\PFfour=L^{-5}\Dq^4+(\text{order}\leq3)$. So $\Mcyc^{\mathrm{LR}}-\chi\,\PFfour$, $\chi\coloneqq L^5/(\kappa^{\mathrm{LR}}_1C_0C_1^4)=C_0/C_1$, has order at most $3$ and kills $\mathfrak p_0,\dots,\mathfrak p_3$ (\cref{lem:set-periods}); by \cref{lem:set-order3} it vanishes. Hence $\kappa^{\mathrm{LR}}_3=\Mcyc^{\mathrm{LR}}\mathfrak p_4=\chi\,\PFfour\mathfrak p_4=\chi=C_0/C_1$, $\kappa^{\mathrm{LR}}_4=\kappa^{\mathrm{LR}}_3$, and
\begin{equation}
\label{eq:set-McycLR}
\Mcyc^{\mathrm{LR}}=\kappa^{\mathrm{LR}}_3\,\PFfour .
\end{equation}

(ii) For a solution $v$, the Frobenius property and the commutativity of $*$ give $z\partial_T\eta(H^a,v)=\eta(H^a,H*v)=\eta(H*H^a,v)=\kappa^{\mathrm{LR}}_a\eta(H^{a+1},v)$, as $\eta(H^a,\cdot)$ has constant coefficients; and $z\partial_T=C_1^{-1}z\Dq$.
\end{proof}

The coefficient $\kappa^{\mathrm{LR}}_3$ sees the relation $H^5=1$, through $\PFfour\mathfrak p_4=1$. Relations equivalent to (i) are stated in \cite[Props.~6.1--6.2]{GR19} and, for $a\leq3$, in \cite[Section~2.4.3, (17), and Section~2.4.4, (19)--(20)]{LP20}; the proof above does not use them.

% Appendix: the cyclic form of the fourth-order operator (proofs of thm:gau-cyc, cor:gau-cycsys and lem:gau-confexp); label prefix: gau.
\section{The cyclic form of the fourth-order operator}
\label{sec:gau-cycproof}

This appendix proves \cref{thm:gau-cyc,cor:gau-cycsys,lem:gau-confexp}, with the notation of \cref{def:gau-conc} and the coefficients $A_k$ of \cref{sec:gau-cyc}.

\begin{lem}[Constants of $\bbK{[\ell]}$]\genusstatus{AI-proposed, unchecked}
\label{lem:gau-const}
$\{f\in\bbK[\ell]:\Dq f=0\}=\bbQ$.
\end{lem}

\begin{proof}
For $f\in\bbK$, $\Dq f=-(1-u)f'=0$ gives $f'=0$, so $f\in\bbQ$. Let $f=\sum_{k\leq n}a_k\ell^k$ with $a_k\in\bbK$, $a_n\neq0$, $n\geq1$, and $\Dq f=0$. As $\Dq f=\sum_k\bigl(\Dq a_k+(k+1)a_{k+1}\Dq\ell\bigr)\ell^k$, the coefficient of $\ell^n$ gives $a_n\in\bbQ^\times$, and that of $\ell^{n-1}$ gives $\Dq a_{n-1}=na_n(1-u)/u$, i.e.\ $a_{n-1}'=-na_n/u$; impossible, since a derivative of a Laurent series has no term in $u^{-1}$.
\end{proof}

\begin{lem}[$\PFfour$ is self-adjoint]\genusstatus{AI-proposed, unchecked}
\label{lem:gau-adj}
$\PFfour^*=\PFfour$, as operators with coefficients in $\bbQ[u]$; hence as operators on every differential ring containing $\bbQ[u]$.
\end{lem}

\begin{proof}
The formal adjoint is the anti-automorphism of $\bbQ[u][\Dq]$ fixing the coefficients and sending $\Dq$ to $-\Dq$; it exists because it maps the relation $\Dq\circ a-a\circ\Dq=\Dq(a)$ to $a\circ(-\Dq)-(-\Dq)\circ a=\Dq(a)$. With $q=(1-u)/5^5$, $\Dq q=q$, write $\PFfour=\Dq^4-5^5q\circ P_4(\Dq)$, $P_4(\Dq):=\prod_{k=1}^4(\Dq+k/5)$. Then $(\Dq^4)^*=\Dq^4$ and $(q\circ P_4(\Dq))^*=P_4(-\Dq)\circ q=q\circ P_4(-\Dq-1)$, by $\Dq\circ q=q\circ(\Dq+1)$, and $P_4(-\Dq-1)=\prod_{k=1}^4\bigl(-(\Dq+\frac{5-k}5)\bigr)=P_4(\Dq)$ (four sign changes, $k\mapsto5-k$ permutes the factors).
\end{proof}

\begin{lem}[The concomitant]\genusstatus{AI-proposed, unchecked}
\label{lem:gau-conc}
\begin{enumerate}
\item For every $\cP=\sum_{k\leq4}a_k\Dq^k$ with coefficients in a differential ring $\cR_0$ (for instance $\bbK$ or $\bbQ[u]$) and all $f,g$ in a differential ring extending $\cR_0$: $g\,\cP f-f\,\cP^*g=\Dq\,\Conc_\cP(f,g)$.
\item For $f,g\in\mathcal S$, $\Conc(f,g)$ lies in $\bbQ$, and $\Conc$ is a $\bbQ$-bilinear alternating form on $\mathcal S$.
\item The jet matrix is $\Jet_{ab}=\sum_{k=a+b+1}^4(-1)^{k-1-a}\binom{k-1-a}{b}\Dq^{k-1-a-b}A_k$, explicitly
\[
\Jet=\begin{pmatrix}0&\frac25(1-u)&1-u&-u\\-\frac25(1-u)&0&u&0\\-(1-u)&-u&0&0\\u&0&0&0\end{pmatrix};
\]
so $\Jet_{ab}=0$ for $a+b\geq4$, $\Jet_{3-j,j}=(-1)^ju$, and $\det\Jet=u^4$.
\item For every basis $\mathbf x$ of $\mathcal S$, $\Conc_{\mathbf x}=\Wr_{\mathbf x}^T\Jet\,\Wr_{\mathbf x}$. In the standard basis
\[
\Conc_{\mathbf f}=\begin{pmatrix}0&-\frac75&2&\ast\\\frac75&0&0&-1\\-2&0&0&0\\-\ast&1&0&0\end{pmatrix},\qquad \operatorname{Pf}(\Conc_{\mathbf f})=2,\qquad\det\Conc_{\mathbf f}=4,
\]
where $\ast=\Conc(f_0,y_3)$ is not needed. Hence $\Conc$ is non-degenerate on $\mathcal S$, $\det\Wr_{\mathbf x}\neq0$ for every basis $\mathbf x$, and $u^2\det\Wr_{\mathbf f}\in\{2,-2\}$.
\end{enumerate}
\end{lem}

\begin{proof}
(i) Fix $k$ and put $w:=a_kg$, $\mathsf t_j:=(-1)^j\Dq^jw\,\Dq^{k-j}f$ (local letters). Then $\Dq\bigl[(-1)^j\Dq^jw\,\Dq^{k-1-j}f\bigr]=\mathsf t_j-\mathsf t_{j+1}$, and the sum over $j<k$ telescopes to $\mathsf t_0-\mathsf t_k=g\,a_k\Dq^kf-f\,(-\Dq)^k(a_kg)$; sum over $k$.

(ii) For $f,g\in\mathcal S$, (i) and \cref{lem:gau-adj} give $\Dq\Conc(f,g)=g\,\PFfour f-f\,\PFfour g=0$, so $\Conc(f,g)\in\bbQ$ by \cref{lem:gau-const}. For alternation, $Q:=\Conc(f,g)+\Conc(g,f)=\sum_{a,b}q_{ab}f_ag_b$ is a differential polynomial in indeterminates $f_a=\Dq^af$, $g_b=\Dq^bg$ with $\Dq Q=0$ identically, by (i) and \cref{lem:gau-adj}. If $Q\neq0$, let $s$ be the largest $a+b$ with $q_{ab}\neq0$; in $\Dq Q$ the coefficient of $f_0g_{s+1}$ is $q_{0s}$ and that of $f_cg_d$ ($c\geq1$, $c+d=s+1$) is $q_{c-1,d}+q_{c,d-1}$, so successively $q_{0,s}=q_{1,s-1}=\dots=q_{s,0}=0$, a contradiction.

(iii) Expand $\Dq^j(A_kg)=\sum_b\binom jb\Dq^{j-b}A_k\,\Dq^bg$ and put $a=k-1-j$; the entries follow with $\Dq(1-u)=1-u$, $\Dq u=-(1-u)$, e.g.\ $\Jet_{01}=-A_2+2\Dq A_3-3\Dq^2A_4=\frac25(1-u)$, $\Jet_{02}=A_3-3\Dq A_4=1-u$, $\Jet_{00}=A_1-\Dq A_2+\Dq^2A_3-\Dq^3A_4=0$. For $a+b=3$ only $k=4$ contributes, giving $(-1)^bu$, and nothing contributes for $a+b\geq4$; the antidiagonal $u,-u,u,-u$ and the even reversal of four letters give $\det\Jet=u^4$.

(iv) The identity is (iii) on the basis. By (ii) the entries of $\Conc_{\mathbf f}$ are rational and may be read at $u=0$. \emph{Power series.} At $u=0$, $\Jet=\frac25(E_{01}-E_{10})+E_{02}-E_{20}$, so $\Conc(f,g)=\bigl[\frac25(f\Dq g-g\Dq f)+(f\Dq^2g-g\Dq^2f)\bigr](0)$, and $(f,\Dq f,\Dq^2f)(0)$ is $(1,0,0)$, $(0,-1,-1)$, $(0,0,2)$ for $f_0,f_1,f_2$; hence $\Conc(f_0,f_1)=-\frac75$, $\Conc(f_0,f_2)=2$, $\Conc(f_1,f_2)=0$. \emph{Entries with $y_3$.} Write $\Dq^by_3=(\Dq^b\varpi_1)\ell+Y_b$, $Y_b\in\bbK$; then $\Conc(f,y_3)=\Conc(f,\varpi_1)\ell+\sum_{a,b}\Dq^af\,\Jet_{ab}Y_b$. The left side lies in $\bbQ$ by (ii), and the right side has degree at most $1$ in the indeterminate $\ell$ over $\bbK$; hence $\Conc(f,\varpi_1)=0$ for $f\in\mathcal S\cap\bbK$, and the second term is a constant of $\bbK$, equal to its coefficient of $u^0$. From $\Dq\ell=1-u^{-1}$, $\Dq^2\ell=u^{-1}-u^{-2}$, $\Dq^3\ell=-2u^{-3}+3u^{-2}-u^{-1}$ and $Y_1=\varpi_1\Dq\ell+\Dq\rho_3$, $Y_2=2\Dq\varpi_1\Dq\ell+\varpi_1\Dq^2\ell+\Dq^2\rho_3$, $Y_3=3\Dq^2\varpi_1\Dq\ell+3\Dq\varpi_1\Dq^2\ell+\varpi_1\Dq^3\ell+\Dq^3\rho_3$, with $\varpi_1=u+O(u^2)$, one gets $Y_0=\rho_3\in u^3\bbQ\fps u$, $Y_1=-1+O(u)$, $Y_2=u^{-1}+O(1)$, $Y_3=u^{-2}+O(u^{-1})$, and
\[
\begin{aligned}\sum_{a,b}\Dq^af\,\Jet_{ab}Y_b={}&f\bigl[\tfrac25(1-u)Y_1+(1-u)Y_2-uY_3\bigr]\\
&+\Dq f\bigl[-\tfrac25(1-u)Y_0+uY_2\bigr]+\Dq^2f\bigl[-(1-u)Y_0-uY_1\bigr]+\Dq^3f\,uY_0 ,\end{aligned}
\]
where the first bracket is a power series, the second is $1+O(u)$ and the others are $O(u)$. So $\Conc(f,y_3)=(\Dq f)(0)$ when $f(0)=0$: $\Conc(f_1,y_3)=-1$, $\Conc(f_2,y_3)=0$. Then $\operatorname{Pf}=c_{01}c_{23}-c_{02}c_{13}+c_{03}c_{12}=2$ and $\det=4$; $\det\Conc_{\mathbf x}=(\det\Wr_{\mathbf x})^2u^4$ gives $\det\Wr_{\mathbf x}\neq0$. Finally Abel's identity $\Dq\det\Wr_{\mathbf f}=-(A_3/A_4)\det\Wr_{\mathbf f}=2u^{-1}(1-u)\det\Wr_{\mathbf f}$ and $\Dq(u^2)=-2u(1-u)$ give $\Dq(u^2\det\Wr_{\mathbf f})=0$, so $u^2\det\Wr_{\mathbf f}\in\bbQ$ (\cref{lem:gau-const}), with square $\det\Conc_{\mathbf f}=4$.
\end{proof}

\begin{lem}[Transversality]\genusstatus{AI-proposed, unchecked}
\label{lem:gau-transv}
For every basis $\mathbf x$ of $\mathcal S$ and $a,b\in\{0,\dots,3\}$, in the fraction field of $\bbK[\ell]$,
\[
\sum_{i,j}(\Conc_{\mathbf x}^{-1})_{ij}\,\Dq^ax_i\,\Dq^bx_j=(\Jet^{-1})_{ab},
\]
where $(\Jet^{-1})_{ab}=0$ for $a+b\leq2$ and $(\Jet^{-1})_{03}=u^{-1}$. In particular the transversality sums of \cref{def:gau-conc} do not depend on the basis, $\Upsilon_0=\Upsilon_1=\Upsilon_2=0$ and $\Upsilon_3=u^{-1}$.
\end{lem}

\begin{proof}
As $\det\Wr_{\mathbf x}\neq0$, $\Conc_{\mathbf x}=\Wr_{\mathbf x}^T\Jet\Wr_{\mathbf x}$ gives $\Jet^{-1}=\Wr_{\mathbf x}\Conc_{\mathbf x}^{-1}\Wr_{\mathbf x}^T$, whose entry $(a,b)$ is the left side. With the exchange matrix $J_{ac}=\delta_{a+c,3}$, $\Jet J$ is upper triangular with diagonal $(-1)^{3-c}u$ (\cref{lem:gau-conc}~(iii)), hence so is $J\Jet^{-1}$, with diagonal entries the inverses: $(\Jet^{-1})_{ba}=0$ for $a+b<3$ and $(\Jet^{-1})_{03}=1/\Jet_{30}=u^{-1}$. The $\Upsilon_k$ are the entries $(0,k)$, and $\widetilde\Upsilon_k$ is built from the bivector $\sum_{i,j}(\Conc_{\mathbf x}^{-1})_{ij}x_i\otimes x_j$, which is basis independent since the basis $\mathbf xM$ has Gram matrix $M^T\Conc_{\mathbf x}M$.
\end{proof}

\begin{lem}[Isotropy, formal Picard--Lefschetz]\genusstatus{AI-proposed, unchecked}
\label{lem:gau-iso}
\begin{enumerate}
\item $\sh$ commutes with $\Dq$, fixes $f_0,f_1,f_2$ and maps $y_3$ to $y_3+\varpi_1$; so $\sh(\mathcal S)=\mathcal S$.
\item $\Conc(f,\varpi_1)=0$ for every $f\in\mathcal S\cap\bbK$, and $\Conc(y_3,\varpi_1)=1$.
\item On $\mathcal S\cap\bbK$ the form $\Conc$ has rank $2$ and radical $\bbQ\varpi_1$, and a plane in $\mathcal S\cap\bbK$ is isotropic if and only if it contains $\varpi_1$. In particular $\operatorname{span}(y_0,\varpi_1)$ is isotropic for every power-series solution $y_0$; and $\Conc(f_0,f_1+\gamma'f_2)=-\frac75+2\gamma'$ vanishes only for $\gamma'=\gamma=\frac7{10}$.
\end{enumerate}
\end{lem}

\begin{proof}
(i) $\Dq\circ\sh$ and $\sh\circ\Dq$ are additive and satisfy $X(fg)=X(f)\sh(g)+\sh(f)X(g)$; they agree on $\bbK$ and on $\ell$ (as $\Dq\ell\in\bbK$), hence on $\bbK[\ell]$. Clearly $\sh f_k=f_k$ and $\sh y_3=y_3+\varpi_1$. (ii) By (i), $\sh\Conc(f,g)=\Conc(\sh f,\sh g)$. For $f\in\mathcal S\cap\bbK$ the constant $\Conc(f,y_3)$ is fixed by $\sh$, so $\Conc(f,y_3)=\Conc(f,y_3+\varpi_1)$, i.e.\ $\Conc(f,\varpi_1)=0$; and $\Conc(y_3,\varpi_1)=-\Conc(f_1,y_3)-\gamma\Conc(f_2,y_3)=1$ (\cref{lem:gau-conc}~(iv)). (iii) The Gram matrix on $(f_0,f_1,f_2)$ is $\frac75(E_{10}-E_{01})+2(E_{02}-E_{20})$, of rank $2$ with kernel $\bbQ\varpi_1$; this gives the last assertion. A plane $\Lambda\not\ni\varpi_1$ satisfies $\mathcal S\cap\bbK=\Lambda\oplus\bbQ\varpi_1$, so $\Conc|_\Lambda$ has rank $2$; a plane containing $\varpi_1$ is isotropic by (ii).
\end{proof}

\begin{lem}[Special geometry]\genusstatus{AI-proposed, unchecked}
\label{lem:gau-sg}
Let $y_0\in\bbK^\times$ and $t\in\bbK$ with $\Dq t\neq0$.
\begin{enumerate}
\item[(a)] For every basis $\mathbf x$ of $\mathcal S$: $\widetilde\Upsilon_1=\widetilde\Upsilon_2=0$ and $\widetilde\Upsilon_3=1/(u\,y_0^2(\Dq t)^3)$.
\item[(b)] Suppose moreover that $y_0$ and $y_1:=ty_0$ lie in $\mathcal S\cap\bbK$ and span an isotropic plane; for instance $y_1=\varpi_1$ and $y_0$ any power-series solution not proportional to $\varpi_1$ (\cref{lem:gau-iso}). Then $\mathcal S$ has a basis $(x_0,\dots,x_3)$ with $x_0=y_0$, $x_1=y_1$, $\Conc(x_0,x_3)=\Conc(x_1,x_2)=1$ and $\Conc(x_0,x_2)=\Conc(x_1,x_3)=\Conc(x_2,x_3)=0$. With $\phi_i:=x_i/y_0$ (so $\phi_0=1$, $\phi_1=t$) and $\mathrm{Yuk}:=\de_t^2\phi_2$:
\begin{enumerate}
\item $\de_t^2\phi_3=-t\,\mathrm{Yuk}$; hence $\de_t^2(y/y_0)\in\mathrm{Yuk}\cdot\operatorname{span}_\bbQ(1,t)$ for every $y\in\mathcal S$;
\item $\mathrm{Yuk}=1/(u\,y_0^2(\Dq t)^3)$.
\end{enumerate}
\end{enumerate}
\end{lem}

\begin{proof}
(a) Put $\dot t:=\Dq t$ and $\Upsilon^{\Dq}_m:=\sum_{i,j}(\Conc_{\mathbf x}^{-1})_{ij}(x_i/y_0)\Dq^m(x_j/y_0)$ (local letters); $\Upsilon^{\Dq}_0=0$ by antisymmetry. The Leibniz rule gives $\Upsilon_k=\sum_m\binom km\,y_0\Dq^{k-m}y_0\,\Upsilon^{\Dq}_m$, a triangular system whose diagonal entries are the unit $y_0^2$; with $\Upsilon_1=\Upsilon_2=0$, $\Upsilon_3=u^{-1}$ (\cref{lem:gau-transv}) it gives $\Upsilon^{\Dq}_1=\Upsilon^{\Dq}_2=0$, $\Upsilon^{\Dq}_3=1/(uy_0^2)$. Likewise $\Dq=\dot t\de_t$, $\Dq^2=(\Dq\dot t)\de_t+\dot t^2\de_t^2$, $\Dq^3=(\Dq^2\dot t)\de_t+3\dot t(\Dq\dot t)\de_t^2+\dot t^3\de_t^3$ express $\Upsilon^{\Dq}_m$ ($m\leq3$) triangularly in $\widetilde\Upsilon_1,\dots,\widetilde\Upsilon_m$ with the unit $\dot t^m$ as coefficient of $\widetilde\Upsilon_m$; solving successively gives (a). Isotropy is not used in (a).

(b) $\Lambda:=\operatorname{span}(y_0,y_1)$ is an isotropic plane, since $t$ is not constant. As $\Conc$ is non-degenerate on $\mathcal S$ by \cref{lem:gau-conc}~(iv), $\Lambda^\perp=\Lambda$ and $x\mapsto\Conc(\cdot,x)|_\Lambda$ maps $\mathcal S$ onto $\Lambda^*$. Choose $x_3$ with $\Conc(y_0,x_3)=1$, $\Conc(y_1,x_3)=0$, and $x_2'$ with $\Conc(y_0,x_2')=0$, $\Conc(y_1,x_2')=1$, and put $x_2:=x_2'-\Conc(x_2',x_3)y_0$. The Gram matrix is $E_{03}+E_{12}-E_{21}-E_{30}$, so this is a basis with $\Conc_{\mathbf x}^{-1}=-\Conc_{\mathbf x}$, and $\widetilde\Upsilon_k=-\phi_0\de_t^k\phi_3+\phi_3\de_t^k\phi_0-\phi_1\de_t^k\phi_2+\phi_2\de_t^k\phi_1$. For $k\geq2$, $\de_t^k1=\de_t^kt=0$, so $\widetilde\Upsilon_k=-\de_t^k\phi_3-t\,\de_t^k\phi_2$: $k=2$ gives (i), and $k=3$ gives $\widetilde\Upsilon_3=\de_t(t\,\mathrm{Yuk})-t\,\de_t\mathrm{Yuk}=\mathrm{Yuk}$, which is (ii) by (a). For $y=\sum_ic_ix_i$, $\de_t^2(y/y_0)=(c_2-c_3t)\mathrm{Yuk}$.
\end{proof}

In a conifold frame, $\mathrm{Yuk}=1/(u\varpi_0^2(\Dq\tc)^3)$ is the conifold expansion of $\kappa^{\mathrm{LR}}_1$ of \cref{prop:set-cyclicLR}, and $\kappa_1=\mathrm{Yuk}\cdot\Dq\tc$ in \cref{prop:gau-quintic}.

\begin{proof}[Proof of \cref{thm:gau-cyc}]
Put $t:=y_1/y_0$, $\dot t:=\Dq t$. As $\Dq\circ\dot t^{-1}\circ\Dq=\dot t\,\de_t^2$,
\[
\Mcyc(y_0,y_1)=\frac{\dot t}{y_0}\circ\de_t^2\circ\bigl(u\,y_0^2\dot t^3\bigr)\circ\de_t^2\circ\frac1{y_0}.
\]
For $y\in\mathcal S$, \cref{lem:gau-sg}~(b) gives $u\,y_0^2\dot t^3\de_t^2(y/y_0)=c_2-c_3t$ with constants $c_2,c_3$, which $\de_t^2$ kills; so $\Mcyc(y_0,y_1)$ annihilates $\mathcal S$. Its coefficient of $\Dq^4$ is the product $y_0^{-1}\cdot\dot t^{-1}\cdot uy_0^2\dot t^2\cdot\dot t^{-1}\cdot y_0^{-1}=u=A_4$ of the multiplication factors. Hence $E:=\Mcyc(y_0,y_1)-\PFfour=\sum_{k\leq3}e_k\Dq^k$ annihilates $\mathcal S$, i.e.\ $(e_0,\dots,e_3)\Wr_{\mathbf f}=0$, and $E=0$ as $\det\Wr_{\mathbf f}\neq0$ (\cref{lem:gau-conc}~(iv)). Under $y_0\mapsto\mu y_0$, $y_1\mapsto\mu'y_1$ the multiplication factors of \eqref{eq:gau-Mcyc} acquire $\mu^{-1},\mu\mu'^{-1},\mu'^2,\mu\mu'^{-1},\mu^{-1}$, with product $1$; under $y_1\mapsto y_1+ay_0$, $\dot t$ is unchanged. The plane of the first statement contains $\varpi_1$ and is isotropic by \cref{lem:gau-iso}; in a conifold frame $\varpi_0$ is a power-series solution with $\varpi_0(0)=1$, not proportional to $\varpi_1$.
\end{proof}

\begin{proof}[Proof of \cref{cor:gau-cycsys}]
With $\dot t=\Dq t$ and $\kappa_1^{-1}=uy_0^2\dot t^2$, \eqref{eq:gau-Mcyc} gives the operator identity
\[
\kappa_4^{-1}\Dq\,\kappa_3^{-1}\Dq\,\kappa_2^{-1}\Dq\,\kappa_1^{-1}\Dq\,\kappa_0^{-1}\Dq\circ y_0^{-1}=y_0^{-1}\Dq\,y_0^{-1}\Dq\,\dot t^{-1}\Dq\,(uy_0^2\dot t^2)\Dq\,\dot t^{-1}\Dq\,y_0^{-1}=y_0^{-1}\Dq\circ\Mcyc(y_0,y_1).
\]
(i) Given $Y_{a+1}=z\Dq Y_a/\kappa_a$ for $a\leq3$, the closing relation $z\Dq Y_4=\kappa_4Y_0$ reads
\[
Y_0=z^5\kappa_4^{-1}\Dq\kappa_3^{-1}\Dq\cdots\kappa_0^{-1}\Dq\,Y_0\qquad(\text{$z$ commutes with $\Dq$});
\]
by the identity and \cref{thm:gau-cyc} this is $y_0Y_0=z^5\Dq\,\PFfour(y_0Y_0)$, i.e.\ $\PFfive(y_0Y_0)=0$ by \cref{lem:set-pf5}, an operator identity over $\bbQ[q][z]$ whose proof uses only $\Dq\circ q=q\circ(\Dq+1)$, valid in $\cR$. The converse is the same chain read backwards. (ii) For invertible $z$, $\Dq\mapsto\Dq+w/z$ fixing $\cR$ is a ring automorphism of $\cR[\Dq]$, since $(\Dq+w/z)\circ f-f\circ(\Dq+w/z)=\Dq(f)$; it maps $z\Dq$ to $z\Dq+w$, and applying it to the identities of (i) gives (ii).
\end{proof}

\begin{proof}[Proof of \cref{lem:gau-confexp}]
Write $\Cexp:=\Cexp_{\alpha,\beta}$ and $\Xi=-\log(1-u)$ (\cref{def:set-confexp}), so $\Dq\Xi=-1$ and $\Dq^k\Xi=0$ for $k\geq2$. By \cref{def:set-confexp}, $\Cexp(L)=s^{-1}$, $\Cexp(\logq)=-\Xi$, $\Cexp(c_j)=\Dq^j\varpi_0$ and $\Cexp(i_j)=\Dq^j\varpi_1-j\Dq^{j-1}\varpi_0+\Xi\Dq^j\varpi_0=\Dq^j(\varpi_1+\Xi\varpi_0)$ (Leibniz rule). $\Nring$ is the localisation of a polynomial ring at $L$, $c_0$ and $\mathsf W$, so the assignment extends uniquely once these go to units: $s^{-1}$ and $\varpi_0$ are units, and
\[
\Cexp(\mathsf W)=\varpi_0^2+\varpi_0\bigl(\Dq\varpi_1-\varpi_0+\Xi\Dq\varpi_0\bigr)-\Dq\varpi_0\bigl(\varpi_1+\Xi\varpi_0\bigr)=\varpi_0\Dq\varpi_1-\varpi_1\Dq\varpi_0=\varpi_0^2\Dq\tc
\]
is a unit since $\Dq\tc(0)=-\tc'(0)=-1$ (\cref{def:set-frames}). Also $\Cexp(I_1)=-\Xi\varpi_0+\varpi_1+\Xi\varpi_0=\varpi_1$.

$\Dq\circ\Cexp$ and $\Cexp\circ\Dq$ are $\Cexp$-derivations, so it suffices to compare them on the generators (on an inverted element, $\delta(x^{-1})=-\Cexp(x)^{-2}\delta(x)$). For $L$: $\Dq(s^{-1})=(s^{-6}-s^{-1})/5=\Cexp((L^6-L)/5)$; for $\logq$: $\Dq(-\Xi)=1$; for $c_0,c_1,c_2,i_0,i_1$ the images are shifted in $j$ by $\Dq$. $\Dq c_3$ is fixed in $\Nring$ by $\PFfour c_0=0$, and $\PFfour\varpi_0=0$ gives $\Cexp(\Dq c_3)=\Dq^4\varpi_0$. $\Dq i_2$ is fixed by $\Conc(c_0,I_1)=0$, with $\Dq^3I_1=c_3\logq+3c_2+\Dq i_2$; this relation is affine in $\Dq i_2$ with unit coefficient $\Jet_{03}c_0=-L^{-5}c_0$, and $\Cexp(\Dq^bI_1)=\Dq^b\varpi_1$ for $b\leq2$ by the cases treated. So its image differs from $\Conc(\varpi_0,\varpi_1)=0$ (\cref{lem:gau-iso}~(ii)) only in the last slot, $-u\varpi_0\bigl(\Cexp(\Dq^3I_1)-\Dq^3\varpi_1\bigr)=0$, whence $\Cexp(\Dq i_2)=\Dq^3\varpi_1+\Xi\Dq^3\varpi_0-3\Dq^2\varpi_0=\Dq\Cexp(i_2)$. The last assertion follows from $\Cexp(I_1/c_0)=\tc$ and $\Cexp(\mathsf W/c_0^2)=\Dq\tc$.
\end{proof}

% Appendix: u-adic convergence; the statement and proof of lem:tra-conv, used in sections 4, 5 and 7.

\section{$u$-adic convergence}
\label{sec:tra-convproof}

This appendix states and proves \cref{lem:tra-conv}, which justifies the infinite sums of \cref{sec:tra,sec:red,sec:flh}; the notation is that of \cref{sec:tra-tower}.

\begin{lem}[$u$-adic convergence]\genusstatus{AI-proposed, unchecked}
\label{lem:tra-conv}
Work in the $u$-adic setting. A flat model correlator is an ancestor correlator $\langle z^{k_1}H^{a_1}\cdots z^{k_n}H^{a_n}\rangle_h(v)$ of \cref{def:set-modcohft}, evaluated at $v = v(u)$; it vanishes unless $\sum_ik_i\leq\mathsf D_{h,n} := 3h-3+n$, because the class of the CohFT has no component of negative degree and $\dim\overline{\mathcal M}_{h,n} = \mathsf D_{h,n}$.
\begin{description}
\item[(a) Frames] For every family $(c_{(k,a)})$ in $\bbQ\fps{u}$ the family $(c_{(k,a)}\fr(k,a))$ is summable, the $z^jH$-coefficient of $\fr(k,a)$ having $u$-order at least $(4(k-j)+\nu_a+1)/5$; the decompositions of \cref{lem:tra-frames}~(c) and \cref{lem:tra-twoslot} hold with these sums.
\item[(b) Frame expansions inside correlators] Let $h\geq 0$, $n\geq 1$ and $2h-2+n>0$.
(b1) The correlator extends to a symmetric $\bbQ\fps{u}$-multilinear map $\langle\ \rangle_h\colon(\cH^u_+)^n\to\bbQ\lser{u}$,
\[
\langle X_1\cdots X_n\rangle_h := \sum_{\sum k_i\leq\mathsf D_{h,n}}\\prod_i(X_i)_{(k_i,a_i)}\,\langle z^{k_1}H^{a_1}\cdots z^{k_n}H^{a_n}\rangle_h(v(u)),
\]
a finite sum, which is of finite type in each slot.
(b2) Let $B := \{(k,a) : a\neq 1\}\cup\{H\}$, $t_{(k,a)} := \fr(k,a)$ and $t_H := H$. If $X_i = \sum_{\beta\in B}c_{i,\beta}t_\beta$ with $c_{i,\beta}\in\bbQ\fps{u}$, as is every element of $\sK$ (\cref{lem:tra-frames,cor:tra-slice}), then the family $(\prod_ic_{i,\beta_i}\langle t_{\beta_1}\cdots t_{\beta_n}\rangle_h)_{\beta\in B^n}$ is summable in $\bbQ\lser{u}$ and its sum is $\langle X_1\cdots X_n\rangle_h$. Every term is a correlator of homogeneous elements of $\sK$ over $\bbQ[v]$ evaluated at $v(u)$; no flat component $z^jH$, $j\geq 1$, of a frame occurs as an insertion of its own.
(b3) For $y\in z^2\cH^u_+$ the vertex factor $\sum_{m\geq 0}\frac1{m!}\langle X_1\cdots X_ny^m\rangle_h$ is a finite sum ($m\leq\mathsf D_{h,n}$ by (d)) and is of finite type in each $X_i$. If moreover $y\in\sK$, write $y = y^\circ+\delta H$ with $\delta := \Pi_0(y)$ and $y^\circ\in\sKo$ (\cref{cor:tra-slice}); then $y^\circ_{(k,a)} = y_{(k,a)} = 0$ for $k\leq 1$ and $a\neq 1$, and for each $m$ the expansion of $y^m = (y^\circ+\delta H)^m$ is finite, each of its terms being covered by (b2) with $n+m$ slots.
(b4) Let $\Gamma$ be a leg-free stable graph with at least one edge, so that every vertex carries $n_w\geq1$ half-edges (\cref{lem:red-genus}). Put on each edge a bivector $E_e\in\bbH\otimes\bbH\otimes\bbQ\fps{\psi,\psi'}\fps{u}$ with property (i) of \cref{lem:tra-twoslot}, and at each vertex $w$ a vertex factor as in (b3), of genus $g_w$, with one slot per half-edge and a leaf vector $y_w\in z^2\cH^u_+\cap\sK$. Then the contraction of these data along $\Gamma$ is a finite sum, and it is the sum of the summable family obtained by expanding every $E_e$ by \cref{lem:tra-twoslot} and every $y_w$ as in (b3). Every member of the family is a product of elements of $\bbQ\fps{u}$ and of correlators of homogeneous elements of $\sK$ in which every half-edge slot carries a frame.
\item[(c) Vanishing series; pairings] The coefficient $\coef{z^{-n}}\YC_b(v(u),z)$, given in closed form by \cref{lem:gau-vanshape}, has $u$-order at least $n-\lfloor n/5\rfloor\geq 4n/5$. Hence $\YC(v(u))\in\bbQ[z^{-1}]\fps{u}^5$, $J^{\mathrm C}=\eta^{-1}\YC$ lies in $\cH^u$ and $z^{-k}J^{\mathrm C}\in\cH^u_-$ for $k\geq 1$, and every pairing and product used in \cref{sec:tra} ($\Pi_i$ on $\cH^u_+$; $P\YC$, $P^{\mathrm{vec}}J^{\mathrm C}$ and $Rf$ for matrices $P$, $R$ over $\bbQ\fps{u,z}$; $\Omega_{\mathrm s}$ on $\cH^u$) is a ring operation in $\bbQ\lser{z}\fps{u}$, i.e.\ a sum that is finite modulo every power of $u$. (The contraction of \cref{lem:tra-E} is treated in the same way, in the ring $\bbQ\lser{\psi}\fps{u,\psi'}$, in the proof of that lemma.)
\item[(d) Leaves] A flat insertion $z^kH^b$ with $k\geq 2$ has $\wt(k,b)\geq 3/5$. At a vertex of genus $h$ with $n$ slots and $m$ further insertions of $\psi$-degree $\geq 2$, a non-zero correlator has slot degrees $\sum_ik_i\leq 3h-3+n-m$; in particular $m\leq 3h-3+n$.
\item[(e) Vertex resummation] Let $g\geq 2$ and $v'(0)\neq 0$. Let $\tilde q\in -z1+z^2\cH^u_+$ with $\Pi_i(\tilde q) = 0$ for $i\geq 1$ and $\delta := \Pi_0(\tilde q)-v\in u^2\bbQ\fps{u}$, and write $y := \tilde q+z1 = y^\circ+\delta H$ with $y^\circ\in\sKo$ (\cref{cor:tra-slice}). (e1) $\langle y^m\rangle_g = 0$ for $m>3g-3$. (e2) The family $(\frac{\delta^j}{l!\,j!}\langle(y^\circ)^lH^j\rangle_g)_{l,j\geq 0}$, every member with $(l,j)\neq(0,0)$ expanded in frames by (b2), is summable, and (here and below $H^j$ under a bracket denotes $j$ insertions of $H=H^1$, not the flat class $H^j$)
\[
\overline F_g(\tilde q) = \sum_m\frac1{m!}\langle y^m\rangle_g = \sum_{l,j\geq 0}\frac{\delta^j}{l!\,j!}\langle(y^\circ)^lH^j\rangle_g,
\]
where $\langle(y^\circ)^0H^0\rangle_g := \langle\ \rangle_g = \Fm_g(v)$. (e3) If $\ML{g}$ holds (\cref{def:red-ML}), then $\sum_j\frac{\delta^j}{j!}\langle H^j\rangle_g = c_g(v+\delta)^{2-2g}$ and $\sum_j\frac{\delta^j}{j!}\sum_{l\geq 1}\frac1{l!}\langle(y^\circ)^lH^j\rangle_g\in\bbQ\fps{u}$; hence $\overline F_g(\tilde q)\in c_g\,\Pi_0(\tilde q)^{2-2g}+\bbQ\fps{u}$.
\item[(f) Dilaton re-expansion] Let $2h-2+n>0$, $X_1,\dots,X_n\in\cH^u_+$, $\varkappa\in\bbQ\fps{u}^\times$, $y\in z^2\cH^u_+$ and $c := 1-1/\varkappa$. (f1) If $\varkappa(0) = 1$, the family $(\frac{c^{m_1}}{m_1!\,m_2!}\langle X_1\cdots X_n(z1)^{m_1}y^{m_2}\rangle_h)_{m_1,m_2\geq 0}$ is summable and
\[
\sum_{m\geq 0}\frac1{m!}\langle X_1\cdots X_n(cz1+y)^m\rangle_h = \varkappa^{2h-2+n}\sum_{m\geq 0}\frac1{m!}\langle X_1\cdots X_n(\varkappa y)^m\rangle_h ,
\]
the right side being a finite sum ($m\leq 3h-3+n$). (f2) If $\varkappa(0)\neq 1$ and $\langle X_1\cdots X_ny^{m}\rangle_h\neq 0$ for some $m$, the terms of the series on the left do not tend to $0$, so it has no $u$-adic limit.
\end{description}
Beyond \cref{conv:tra-forms}, \cref{lem:tra-sum,lem:tra-frames,cor:tra-slice,lem:tra-twoslot} (which precede \cref{sec:tra-dress} and do not use this lemma), \cref{lem:gau-vanshape} in (c) and the definitions of \cref{def:set-graphsum,def:set-modcohft}, the parts depend on the following statements only: (b) and (e) on the homogeneity of \cref{lem:set-modhom}; (f) on the dilaton equation of \cref{lem:set-dilaton}; (e3) on $\ML{g}$, as a hypothesis; (e) on $v'(0)\neq 0$. In the applications, $\delta\in u^2\bbQ\fps{u}$ in (e) is \cref{lem:tra-T}~(d), and $\varkappa(0) = 1$ in (f) holds for $\Pc$ by \cref{lem:gau-diag} (\cref{prop:tra-can}).
\end{lem}

\begin{proof}
(c) By \cref{lem:gau-vanshape} (for the residue solution of \cref{def:gau-van,lem:gau-vanmod}) the $z^{-n}$-coefficient is a multiple of $v^{n-\lfloor n/5\rfloor}$, and $\ord_uv(u)\geq 1$; so modulo $u^N$ only finitely many $n$ occur, $\YC(v(u),z)\in\bbQ[z^{-1}]\fps{u}^5$, and the statements on $J^{\mathrm C}$, $z^{-k}J^{\mathrm C}$ follow. The rest is \cref{conv:tra-forms}: $\bbQ\lser{z}\fps{u}$ is a ring containing $\bbQ\fps{u,z}$ and $\bbQ[z^{-1}]\fps{u}$, in which products, coefficient extraction, the splitting, $z\mapsto -z$ and $\Res_z$ are coefficientwise in $u$ and finite modulo every power of $u$.

(a) This is the proof of \cref{lem:tra-frames}~(c) and of \cref{lem:tra-twoslot}, the $u$-order being the exponent of \cref{lem:tra-frames}~(b) with $\ord_uv\geq 1$.

(b1) Only the finitely many flat slots with $\sum k_i\leq\mathsf D_{h,n}$ contribute. By \cref{lem:set-modhom} each flat correlator is $Cv^e$ with $C\in\bbQ$ and $e = 2-2h+\sum_i\wt(k_i,a_i)\geq 2-2h-n$, as $\wt(k,a)\geq\wt(0,1) = -1$; so at $v = v(u)$ all of them lie in $u^{e_0}\bbQ\fps{u}$ with $e_0 := \min(0,(2-2h-n)\ord_uv)$. Hence $\langle X_1\cdots X_n\rangle_h\in u^{e_0}\bbQ\fps{u}$ for all $X_i\in\cH^u_+$; it is symmetric and $\bbQ\fps{u}$-multilinear, and with $X_2,\dots,X_n$ fixed it is of finite type in $X_1$: modulo $u^N$ it depends only on the $(X_1)_{(k,a)}$ with $k\leq\mathsf D_{h,n}$, taken modulo $u^{N-e_0}$.

(b2) For $\beta_i\in B$ put $\kappa(\beta_i) := k$ if $\beta_i = (k,a)$ and $\kappa(H) := 0$. By (b1) the term of $\beta = (\beta_1,\dots,\beta_n)$ is the finite sum over flat slots of $\prod_ic_{i,\beta_i}(t_{\beta_i})_{(k_i,a_i)}\langle\text{slot}\rangle_h(v(u))$. A frame $t_{\beta_i}$ with $\kappa(\beta_i)>\mathsf D_{h,n}$ enters a slot with $k_i\leq\mathsf D_{h,n}$ only through its $H$-coefficient at $z^{k_i}$, of $u$-order at least $(4(\kappa(\beta_i)-\mathsf D_{h,n})+1)/5$ by (a). So the term has $u$-order at least $e_0+\sum_i\max(0,(4(\kappa(\beta_i)-\mathsf D_{h,n})+1)/5)$, and for each $N$ only finitely many $\beta$ give a term outside $u^N\bbQ\fps{u}$: the family is summable. By \cref{lem:tra-sum} in the first slot (finite type by (b1); the family $(c_{1,\beta}t_\beta)_\beta$ is summable by (a)), $\langle X_1\cdots X_n\rangle_h = \sum_{\beta_1}c_{1,\beta_1}\langle t_{\beta_1}X_2\cdots X_n\rangle_h$; then in the second slot for each $\beta_1$, and so on. This expresses $\langle X_1\cdots X_n\rangle_h$ as an iterated sum of the family, which equals its unordered sum by regrouping. Each $t_\beta$ is a frame or $H$, a homogeneous element of $\sK$ over $\bbQ[v]$ evaluated at $v(u)$ (\cref{lem:tra-frames}~(b), (d)); the expansion is indexed by $B$, so a flat $z^jH$ with $j\geq 1$ never appears as an insertion.

(b3) By (d), $\langle X_1\cdots X_ny^m\rangle_h = 0$ for $m>\mathsf D_{h,n}$, since all components of $y$ have $\psi$-degree $\geq 2$; for each $m$ the term is multilinear, hence of finite type in each $X_i$ by (b1). If $y\in\sK$, then $y^\circ := y-\delta H\in\sKo$ by \cref{cor:tra-slice}, its frame coordinates are $y^\circ_{(k,a)} = y_{(k,a)}$ for $a\neq 1$, which vanish for $k\leq 1$ because $y\in z^2\cH^u_+$, and the multinomial expansion of $(y^\circ+\delta H)^m$ is finite.

(b4) For an edge between the vertices $w$ and $w'$ only its coefficients of $\psi^k\psi'^l$ with $k\leq\mathsf D_w := 3g_w-3+n_w$ and $l\leq\mathsf D_{w'}$ matter, by (d); so the contraction is a finite sum of products. As a function of one edge bivector, the others fixed, it is linear and, by (b1), modulo $u^N$ it depends on finitely many coefficients of that bivector, taken modulo a fixed power of $u$: it is of finite type. Expanding every edge by \cref{lem:tra-twoslot} and every leaf as in (b3) gives a family indexed by the frame labels of all half-edges and leaves; by the bound of (b2) at each vertex it is summable, and by \cref{lem:tra-sum}, applied edge by edge, its sum is the contraction. The edge coefficients $E^{\alpha\beta}_e$ lie in $\bbQ\fps{u}$, and each half-edge slot carries a frame.

(d) $\wt(k,b) = (4k+\nu_b-4)/5\geq(8-1-4)/5 = 3/5$ for $k\geq 2$, since $\nu_b\geq -1$. A correlator with $n$ slots of degrees $k_i$ and $m$ insertions of degrees $l_j\geq 2$ at genus $h$ vanishes unless $\sum k_i+\sum l_j\leq 3h-3+n+m$; so $\sum k_i\leq 3h-3+n+m-2m = 3h-3+n-m$, and $m\leq 3h-3+n$ because $\sum k_i\geq 0$.

(e) (e1) is (d) with no slot. (e2) For fixed $m$, multilinearity and symmetry give
\[
\langle y^m\rangle_g = \sum_{l+j=m}\frac{m!}{l!\,j!}\,\delta^j\langle(y^\circ)^lH^j\rangle_g .
\]
Expand $y^\circ$ in frames; only frames with $k\geq 2$ occur (b3), and they have $\wt(k,a)\geq(8+0-4)/5 = 4/5$, as $a\neq 1$ gives $\nu_a\geq 0$. By \cref{lem:set-modhom} a correlator of $l$ such frames and $j$ insertions $H$ ($\wt(0,1) = -1$) is $Cv^e$ with $e\geq 2-2g+4l/5-j$. Since $\ord_uv = 1$ (here $v'(0)\neq 0$ is used: for negative $e$, $\ord_uv^e = e\ord_uv$) and $\ord_u\delta^j\geq 2j$, every term of the expanded $(l,j)$-family has $u$-order at least $2-2g+j+4l/5$. So for each $N$ only finitely many $(l,j)$ contribute outside $u^N\bbQ\fps{u}$, and for each $(l,j)$ the frame family is summable by (b2); the family over $(l,j)$ and the frame labels is summable. Grouping it by $m = l+j$ gives $\sum_m\frac1{m!}\langle y^m\rangle_g$, which is $\overline F_g(\tilde q)$ by the definition of $\overline F_g$ (\cref{def:set-graphsum,def:set-modcohft}); grouping it by $(l,j)$ gives the right side. (e3) For $j\geq 1$, $\ML{g}$~(i) with $X_1 = \dots = X_j = H$ ($\Pi_0(H) = 1$) says that $\langle H^j\rangle_g-(\Fm_g)^{(j)}(v)$ lies in $\bbQ[v]$; both terms are homogeneous of exponent $2-2g-j<0$ (\cref{lem:set-modhom}; $\Fm_g = c_gv^{2-2g}$ by $\ML{g}$~(ii)), so the difference is $0$; for $j = 0$, $\ML{g}$~(ii) gives $\langle\ \rangle_g = c_gv^{2-2g}$. Hence
\[
\sum_j\frac{\delta^j}{j!}\langle H^j\rangle_g = c_gv^{2-2g}\sum_j\binom{2-2g}{j}\Bigl(\frac{\delta}{v}\Bigr)^j = c_g(v+\delta)^{2-2g}:
\]
$\delta/v$ lies in $u\bbQ\fps{u}$, as $v$ is $u$ times a unit and $\delta\in u^2\bbQ\fps{u}$, and the binomial identity $(1+x)^{2-2g} = \sum_j\binom{2-2g}{j}x^j$ of $\bbQ\fps{x}$ is preserved by the continuous substitution $x\mapsto\delta/v$. For $l\geq 1$, each term of the frame expansion is a correlator of homogeneous elements of $\sK$ containing at least one frame, on which $\Pi_0$ vanishes; by $\ML{g}$~(i) it lies in $\bbQ[v]$, so the term lies in $\bbQ\fps{u}$, and so does the sum of the summable family of these terms.

(f) The dilaton equation (\cref{lem:set-dilaton}~(i), which applies to the model by \cref{lem:set-dilaton}~(ii), in the form $\langle z1\,Y_1\cdots Y_r\rangle_h = (2h-2+r)\langle Y_1\cdots Y_r\rangle_h$ for $2h-2+r>0$), applied $m_1$ times, gives $\langle X(z1)^{m_1}y^{m_2}\rangle_h = (N)_{m_1}\langle Xy^{m_2}\rangle_h$, where $X := X_1\cdots X_n$, $N := 2h-2+n+m_2\geq 1$ and $(N)_{m_1} := N(N+1)\cdots(N+m_1-1)$; every intermediate step has $2h-2+r>0$. By multilinearity and symmetry the $m$-th term on the left is $\sum_{m_1+m_2 = m}\frac{c^{m_1}}{m_1!\,m_2!}\langle X(z1)^{m_1}y^{m_2}\rangle_h$, and for fixed $m_2$ the series in $m_1$ is $\frac{1}{m_2!}\langle Xy^{m_2}\rangle_h$ times the binomial series $\sum_{m_1}\binom{N+m_1-1}{m_1}c^{m_1}$. (f1) If $\varkappa(0) = 1$, then $\ord_uc\geq 1$ and the binomial series converges to $(1-c)^{-N} = \varkappa^N$. The family is summable: $\langle Xy^{m_2}\rangle_h = 0$ for $m_2>3h-3+n$ by (d), $\ord_uc^{m_1}\geq m_1$, and the finitely many correlators $\langle Xy^{m_2}\rangle_h$ lie in one $u^{N_0}\bbQ\fps{u}$ by (b1). Regrouping by $m_2$ gives $\sum_{m_2}\varkappa^{2h-2+n+m_2}\langle Xy^{m_2}\rangle_h/m_2!$, which is the right side. (f2) Let $c_0 := c(0)\neq 0$, $A_{m_2} := \langle Xy^{m_2}\rangle_h/m_2!$, let $o$ be the least $u$-order of the non-zero $A_{m_2}$ (finitely many, by (d)), $S$ the set of $m_2$ with $\ord_uA_{m_2} = o$, $a_{m_2} := \coef{u^o}A_{m_2}$ and $N_{m_2} := 2h-2+n+m_2$. Since $\ord_uc\geq 0$ and $\coef{u^0}c^k = c_0^k$, for $m\geq\max S$ the coefficient of $u^o$ in the $m$-th term $T_m$ is
\[
\coef{u^o}T_m = c_0^m\sum_{m_2\in S}c_0^{-m_2}a_{m_2}\binom{N_{m_2}+m-m_2-1}{N_{m_2}-1}.
\]
Each binomial coefficient is a polynomial in $m$ of degree $N_{m_2}-1$, and these degrees are distinct; so the sum is a non-zero polynomial in $m$, and $\coef{u^o}T_m\neq 0$ for all large $m$. Thus $T_m$ does not tend to $0$, and the series has no $u$-adic limit.
\end{proof}

In the applications of (f) (\cref{lem:tra-F}), $\varkappa = d$ and $y = T_{\geq 2}$, so that $cz1+y = q_T+z1$ and $\varkappa y = \tilde q_T+z1$ by \eqref{eq:tra-leaf}.

% Appendix (formal stationary phase), part 1 of 5; label prefix: fsp. Each part inputs the next one at its end.
% No new macros; local letters are declared where they occur.
\section{Formal stationary phase}
\label{sec:fsp}

This appendix proves the properties of formal stationary phase used in \cref{sec:con,sec:gam}. Coefficients lie in a commutative $\bbQ$-algebra $A$, the stationary phase at a critical point is the Wick contraction of the Taylor data there, and nothing is evaluated at $z\neq0$. Only the branch part of \cref{lem:fsp-refl}, used in \cref{lem:con-form}~(d), is not about formal series. Here $t,x,y,v,w$ denote integration or outer variables, $\xi$ the variable of a generating function, $q$ and $\varepsilon$ formal parameters, and $\mathsf H$ a Hessian matrix.

%-------------------------------------------------------------------------------
\subsection{The weighted ring and the stationary-phase functional}
\label{sec:fsp-def}

\begin{defn}[Weighted ring and Gaussian functional]
\label{def:fsp-ring}
Let $t=(t_1,\dots,t_n)$ be variables and $z$ a formal variable, and put $\wt(t^\alpha z^j):=|\alpha|/2+j$ for $\alpha\in\bbZ_{\geq0}^n$, $j\in\bbZ$. Let $\cR_n(A)$ be the set of formal sums $\sum c_{\alpha,j}t^\alpha z^j$, $c_{\alpha,j}\in A$, such that for every $W$ only finitely many monomials of weight $\leq W$ have $c_{\alpha,j}\neq0$. A family $(P_\lambda)$ in $\cR_n(A)$ is \emph{summable} if for every $W$ only finitely many $P_\lambda$ have a non-zero term of weight $\leq W$; its sum is taken termwise. For a symmetric $C\in A^{n\times n}$ put $\langle t^\alpha\rangle_C:=\de_\xi^\alpha\exp(\xi^{\mathrm T}C\xi/2)|_{\xi=0}$, so that $\sum_\alpha\langle t^\alpha\rangle_C\,\xi^\alpha/\alpha!=\exp(\xi^{\mathrm T}C\xi/2)$, extended $A$-linearly to $A[t]$. The \emph{Gaussian functional} with covariance $zC$ is the $A$-linear map $E_{zC}\colon\cR_n(A)\to A\lser{z}$, $E_{zC}[t^\alpha z^j]:=\langle t^\alpha\rangle_C\,z^{j+|\alpha|/2}$, applied termwise; $E_{zC}[1]=1$.
\emph{Outer variables.} If variables $y=(y_1,\dots,y_m)$ of weight $1/2$ are added and $C$ has entries in $A\fps{y}$, the partial functional $E^t_{zC}\colon\cR_{n+m}(A)\to\cR_m(A)$, $t^\alpha y^\beta z^j\mapsto\langle t^\alpha\rangle_C\,y^\beta z^{j+|\alpha|/2}$, acts on $t$ only.
\emph{Parameters.} Over $A=A_0\fps{q}$ we work in $\cR_n(A_0)\fps{q}\supset\cR_n(A)$, coefficientwise in $q$, in the following sense. For a symmetric matrix $C$ over $A$, which may depend on $q$, and $P=\sum_dq^dP_d$ with $P_d\in\cR_n(A_0)$, put $E_{zC}[P]:=\sum_dq^dE_{zC}[P_d]\in A_0\lser z\fps q$, where $E_{zC}[P_d]\in A\lser z$; on $\cR_n(A)$ this is the functional defined above. Modulo $q^N$ it is the Gaussian functional over $B_N:=A_0[q]/(q^N)$ on $\cR_n(B_N)=\cR_n(A_0)\fps q/(q^N)$. Identities in $\cR_n(A_0)\fps q$, in particular between exponentials, are meant modulo every $q^N$, in $\cR_n(B_N)$, where an element divisible by $q$ is nilpotent and $\re^{Y+X}=\re^Y\re^X$ for $Y$ of weight $\geq1/2$ and $X$ nilpotent. Likewise, for an amplitude $a\in A_0\fps{x-c_0}\lser z\fps q$ and a critical point $c\in c_0+(qA)^n$, the series $\SPn_c[F;a]\in A_0\lser z\fps q$ of \cref{def:fsp-wick} is defined modulo every $q^N$, where $a(c+t)$ is an amplitude over $B_N$. Alternatively $q$ gets the weight $1$, and $\cR_{n;q}(A_0)$ is the set of sums $\sum c\,t^\alpha q^dz^j$ with the same finiteness condition for the weight $|\alpha|/2+d+j$.
\end{defn}

\begin{lem}[The weighted ring]\genusstatus{AI-proposed, unchecked}
\label{lem:fsp-ring}
(i) $\cR_n(A)$ is a ring containing $A\fps{t}\lser{z}$; multiplication by a fixed element, $\de/\de t_i$ and multiplication by $t_i$ preserve summable families and commute with their sums.

(ii) If every term of $P\in\cR_n(A)$ has weight $\geq1/2$, the family $(P^k/k!)_k$ is summable, $\exp(P):=\sum_kP^k/k!$ lies in $\cR_n(A)$, and the usual rules for exponentials hold.

(iii) $\langle t^\alpha\rangle_C=0$ for $|\alpha|$ odd; for $|\alpha|$ even it is the sum over the perfect matchings of $|\alpha|$ labelled points, $\alpha_i$ of them carrying the index $i$, of the products of the $C_{ij}$ over the matched pairs; and $\langle t_it^\beta\rangle_C=\sum_jC_{ij}\beta_j\langle t^{\beta-e_j}\rangle_C$ (Wick recursion).

(iv) $E_{zC}$ and $E^t_{zC}$ are well defined, do not lower the weight, and commute with sums of summable families; $E_{zC}$, and $E^t_{zC}$ for constant $C$, preserve the weight, $E_{zC}[fP]=fE_{zC}[P]$ for $f\in A\lser{z}$, and $E^t_{zC}[gP]=g\,E^t_{zC}[P]$ for $g\in\cR_m(A)$, a series in the outer variables and $z$, and $P\in\cR_{n+m}(A)$.

(v) For $J\in A^{n\times n}$ the substitution $t=Jy$ maps $\cR_n(A)$ to itself preserving the weight, and $E_{zC'}[P(Jy)]=E_{zJC'J^{\mathrm T}}[P]$ for all $P\in\cR_n(A)$ and symmetric $C'$.
\end{lem}

\begin{proof}
(i) The weights of the non-zero terms of an element are bounded below, so in a product the coefficient of a monomial of weight $W$ is a finite sum over pairs of terms; the same count gives the finiteness condition and the assertions on summable families. An element $\sum_{j\geq j_0}a_jz^j$ of $A\fps{t}\lser{z}$ has in weight $\leq W$ only terms with $j_0\leq j\leq W$, $|\alpha|\leq2(W-j_0)$. The operators $\de/\de t_i$ and $t_i\cdot$ shift weights by $\mp1/2$. (ii) $P^k$ has weight $\geq k/2$. (iii) $\exp(\xi^{\mathrm T}C\xi/2)$ is even in $\xi$, and $\de_{\xi_i}\exp(\xi^{\mathrm T}C\xi/2)=(C\xi)_i\exp(\xi^{\mathrm T}C\xi/2)$; applying $\de_\xi^\beta$ at $\xi=0$ with the Leibniz rule gives the recursion (one derivative $\de_{\xi_j}$ must hit $\xi_j$, in $\beta_j$ ways). The recursion determines all moments from $\langle1\rangle_C=1$, and the matching sum satisfies it too (match $i$ with one of the remaining entries). (iv) A monomial of weight $w$ goes to a multiple of $z^w$, odd moments vanishing, and the coefficient of $z^k$ receives contributions only from the finitely many terms of weight $k$; this gives the assertions on $E_{zC}$. For $E^t_{zC}$, $\langle t^\alpha\rangle_C\in A\fps{y}$ is a polynomial in the entries of $C$, so the image of $t^\alpha y^\beta z^j$ is a sum of monomials $y^{\beta+\gamma}z^{j+|\alpha|/2}$ of weight $\geq\wt(t^\alpha y^\beta z^j)$ ($\gamma=0$ if $C$ is constant), and the coefficient of $y^\delta z^k$ receives contributions only from the finitely many terms with $\beta\leq\delta$, $j+|\alpha|/2=k$, which have weight $\leq k+|\delta|/2$. For $g\in\cR_m(A)$ let $g_\gamma$ be its monomials; the family $(g_\gamma P)_\gamma$ is summable with sum $gP$, and $E^t_{zC}[g_\gamma P]=g_\gamma E^t_{zC}[P]$ by the definition on monomials. (v) By the generating function, $\langle\re^{\xi\cdot Jy}\rangle_{C'}=\langle\re^{(J^{\mathrm T}\xi)\cdot y}\rangle_{C'}=\exp(\xi^{\mathrm T}JC'J^{\mathrm T}\xi/2)$, i.e.\ $\langle(Jy)^\alpha\rangle_{C'}=\langle t^\alpha\rangle_{JC'J^{\mathrm T}}$, which is the claim for $P=t^\alpha z^j$; the substitution maps a monomial of weight $w$ to a finite sum of monomials of weight $w$, and (iv) extends the identity.
\end{proof}

\begin{defn}[The normalised stationary phase]
\label{def:fsp-wick}
Let $c\in A^n$. For a \emph{phase} $F\in A\fps{x-c}$, $c$ is a \emph{non-degenerate critical point over $A$} if $\nabla F(c)=0$ and $\mathsf H:=\Hess F(c)\in GL_n(A)$. With $t:=x-c$ write $F(c+t)=F(c)+\frac12t^{\mathrm T}\mathsf Ht+F_{\geq3}(t)$, $F_{\geq3}\in(t)^3$, and $C:=-\mathsf H^{-1}$. An \emph{amplitude} is a series $a\in A\fps{x-c}\lser{z}$, or any $a$ with $\exp(F_{\geq3}/z)\,a(c+t)\in\cR_n(A)$; every term of $F_{\geq3}/z$ has weight $\geq1/2$. The \emph{normalised stationary phase} is
\[
\SPn_c[F;a]:=E_{zC}\bigl[\exp\bigl(F_{\geq3}(t)/z\bigr)\,a(c+t)\bigr]\in A\lser{z}.
\]
Equivalently, substitute $t=z^{1/2}\tilde t$, expand in powers of $z^{1/2}$ and apply $\langle\cdot\rangle_C$ in $\tilde t$ termwise. The map $a\mapsto\SPn_c[F;a]$ is $A\lser{z}$-linear, in particular $A\fps{z}$-linear (\cref{lem:fsp-ring}~(iv)). If $a=\sum_{m\geq0}a_mz^m$, the constant term of $\SPn_c[F;a]$ is $a_0(c)$; if it is a unit, $\SPh_c[F;a]:=\SPn_c[F;a]/a_0(c)\in1+zA\fps{z}$. The \emph{full symbol} is
\[
\mathrm{SP}_c\bigl[\re^{F/z}a\,\dd^nx\bigr]:=\bigl[\re^{F(c)/z}\bigr]\,(2\pi z)^{n/2}\,\bigl[\det(-\mathsf H)^{-1/2}\bigr]\,\SPn_c[F;a],
\]
with formal symbols in brackets ($[\re^{F_1/z}][\re^{F_2/z}]=[\re^{(F_1+F_2)/z}]$ for $F_1,F_2\in A$; the square root, in an extension of $A$, is part of the datum). An identity of full symbols is an identity of normalised series together with an explicit identity of symbols under a stated rule for the square roots. Only normalised series are canonical; the note normalises series by their constant terms, so only the $z$-dependence of the symbols matters, and it is always explicit. The convention $\re^{F/z}$ becomes the convention $\re^{-\cW/z}$, resp.\ $\re^{-\hat\cW/\zh}$, of \cref{sec:con} through $F=-\cW$, resp.\ $F=-\hat\cW$ with $z$ renamed $\zh$.
\end{defn}

% Formal stationary phase, part 2 of 5: universality, constants, homogeneity (end of sec:fsp-def); exact forms, Gaussian
% perturbation, products, formal implicit functions (sec:fsp-basic). Continued in part 3.

\begin{lem}[Universality of the coefficients]\genusstatus{AI-proposed, unchecked}
\label{lem:fsp-univ}
Let $c$ be a non-degenerate critical point of $F$ over $A$ and $a(c+t)=\sum_{m\geq m_0}a_m(t)z^m$, $a_m\in A\fps{t}$. Then
\begin{equation}
\label{eq:fsp-univ}
\coef{z^k}\,\SPn_c[F;a]=\sum_{m=m_0}^{k}\ \sum_{r=0}^{2(k-m)}\frac1{r!}\,\bigl\langle\,[F_{\geq3}^{\,r}\,a_m]_{2(k-m+r)}\bigr\rangle_C ,
\end{equation}
$[P]_d$ denoting the homogeneous part of degree $d$ in $t$. So $\coef{z^k}\SPn_c[F;a]$ is a polynomial over $\bbQ$ in the entries of $C$, the Taylor coefficients of $F$ at $c$ of orders $3,\dots,2(k-m_0)+2$ and those of the $a_m$ of orders $\leq2(k-m)$, linear in the latter. Consequently: (i) a ring homomorphism $\rho\colon A\to A'$, applied to the coefficients, maps $\SPn_c[F;a]$ to $\SPn_{\rho(c)}[\rho F;\rho a]$; (ii) if $A$ is a ring of holomorphic functions or of formal power series in some parameters, every coefficient is such a function, resp.\ series.
\end{lem}

\begin{proof}
The terms $F_{\geq3}^{\,r}a_mz^{m-r}/r!$ form a summable family with sum $\exp(F_{\geq3}/z)\,a(c+t)$, so $E_{zC}$ applies termwise. A monomial $t^\alpha$ of $F_{\geq3}^{\,r}a_m$ contributes to $z^k$ if and only if $|\alpha|=2(k-m+r)$; as $F_{\geq3}^{\,r}\in(t)^{3r}$, this forces $r\leq2(k-m)$ and $m\leq k$. A contributing term with factors of $F_{\geq3}$ of degrees $l_1,\dots,l_r\geq3$ and a monomial of $a_m$ of degree $j$ has $l_1+\dots+l_r+j=2(k-m+r)$, so $l_i\leq2(k-m)+2$ and $j\leq2(k-m)$. Moments are polynomials in the entries of $C$ (\cref{lem:fsp-ring}~(iii)). For (i), $\rho$ maps the Taylor data at $c$ to those at $\rho(c)$, and $C$ to $-\rho(\mathsf H)^{-1}$; (ii) follows.
\end{proof}

\begin{lem}[Constants]\genusstatus{AI-proposed, unchecked}
\label{lem:fsp-const}
For $\kappa\in A$, $\SPn_c[F+\kappa;a]=\SPn_c[F;a]$, and the full symbol of $\re^{(F+\kappa)/z}a\,\dd^nx$ is $[\re^{\kappa/z}]$ times that of $\re^{F/z}a\,\dd^nx$; so $\SPn_c[F;a]$ depends only on the Taylor coefficients of order $\geq2$ of $F$ at $c$ and on the Taylor series of $a$ at $c$.
\end{lem}

\begin{proof}
$\mathsf H$ and $F_{\geq3}$ do not change; the last assertion is \eqref{eq:fsp-univ}.
\end{proof}

\begin{lem}[Homogeneity]\genusstatus{AI-proposed, unchecked}
\label{lem:fsp-homog}
For $\mu\in A^\times$, $\SPn_c[\mu F;a(\cdot\,;z)](z)=\SPn_c[F;a(\cdot\,;\mu z')](z')$ at $z'=z/\mu$, i.e.\ after the isomorphism $A\lser{z}\to A\lser{z'}$, $z\mapsto\mu z'$; and, with the branch rule $\det(-\mu\mathsf H)^{1/2}:=\mu^{n/2}\det(-\mathsf H)^{1/2}$,
\[
[\re^{\mu F(c)/z}]\,(2\pi z)^{n/2}\,[\det(-\mu\mathsf H)^{-1/2}]=[\re^{F(c)/z'}]\,(2\pi z')^{n/2}\,[\det(-\mathsf H)^{-1/2}].
\]
\end{lem}

\begin{proof}
For $\mu F$ the Hessian is $\mu\mathsf H$, the covariance $C/\mu$ and the cubic part $\mu F_{\geq3}$. As $\langle t^\alpha\rangle_C$ is homogeneous of degree $|\alpha|/2$ in $C$, $E_{zC/\mu}[t^\alpha z^j]=\mu^{-|\alpha|/2}\langle t^\alpha\rangle_Cz^{j+|\alpha|/2}$, which at $z=\mu z'$ equals $\langle t^\alpha\rangle_C\mu^jz'^{\,j+|\alpha|/2}=E_{z'C}[t^\alpha(\mu z')^j]$. Both sides are linear, weight-preserving and compatible with summable sums, so $E_{zC/\mu}[P(t,z)]|_{z=\mu z'}=E_{z'C}[P(t,\mu z')]$ on $\cR_n(A)$; take $P=\exp(\mu F_{\geq3}/z)a(c+t;z)$, for which $P(t,\mu z')=\exp(F_{\geq3}/z')a(c+t;\mu z')$. The symbols agree term by term.
\end{proof}

%-------------------------------------------------------------------------------
\subsection{Exact forms, Gaussian perturbation, implicit functions and parameters}
\label{sec:fsp-basic}

\begin{lem}[Exact forms]\genusstatus{AI-proposed, unchecked}
\label{lem:fsp-ibp}
Let $c$ be a non-degenerate critical point of $F$ over $A$ and $h_1,\dots,h_n$ amplitudes with $\exp(F_{\geq3}/z)h_i\in\cR_n(A)$ (e.g.\ $h_i\in A\fps{x-c}\lser{z}$). Then $\SPn_c\bigl[F;\sum_i(\de_ih_i+z^{-1}(\de_iF)h_i)\bigr]=0$, i.e.\ $\mathrm{SP}_c[\dd(\re^{F/z}\iota)]=0$ for $\iota=\sum_i(-1)^{i-1}h_i\,\dd x_1\cdots\widehat{\dd x_i}\cdots\dd x_n$. The same holds for every sum of such expressions that is summable in $\cR_n(A)$ after multiplication by $\exp(F_{\geq3}/z)$.
\end{lem}

\begin{proof}
First, for every $g\in\cR_n(A)$ and every $i$,
\begin{equation}
\label{eq:fsp-gibp}
E_{zC}\bigl[\de_ig+z^{-1}(\mathsf Ht)_i\,g\bigr]=0\qquad\text{(Gaussian integration by parts).}
\end{equation}
Multiplying by $C$ and using $C\mathsf H=-1$, \eqref{eq:fsp-gibp} for all $i$ is equivalent to $E_{zC}[t_kg]=z\sum_jC_{kj}E_{zC}[\de_jg]$ for all $k$. For $g=t^\beta z^l$ both sides equal $z^{l+(|\beta|+1)/2}\sum_jC_{kj}\beta_j\langle t^{\beta-e_j}\rangle_C$ by the Wick recursion, and both sides are linear and compatible with summable sums. Now put $g_i:=\exp(F_{\geq3}/z)h_i$. Then $\de_ig_i=\exp(F_{\geq3}/z)(\de_ih_i+z^{-1}(\de_iF_{\geq3})h_i)$ and $\de_iF(c+t)=(\mathsf Ht)_i+\de_iF_{\geq3}(t)$, so $\exp(F_{\geq3}/z)\sum_i(\de_ih_i+z^{-1}(\de_iF)h_i)=\sum_i(\de_ig_i+z^{-1}(\mathsf Ht)_ig_i)$, whose image under $E_{zC}$ vanishes; summable sums are treated termwise.
\end{proof}

\begin{lem}[Gaussian perturbation]\genusstatus{AI-proposed, unchecked}
\label{lem:fsp-pg}
Let $R$ be a commutative $\bbQ$-algebra, complete and separated for an ideal $I$, $\mathsf H\in GL_n(R)$ symmetric, $C:=-\mathsf H^{-1}$, $\mathsf K$ symmetric with entries in $I$, and $C_{\mathsf K}:=-(\mathsf H+\mathsf K)^{-1}$.

(i) For $P\in R[t]$, $\langle\exp(\frac12t^{\mathrm T}\mathsf Kt)P\rangle_C=\det(1+\mathsf H^{-1}\mathsf K)^{-1/2}\langle P\rangle_{C_{\mathsf K}}$, both sides converging $I$-adically, the power being the binomial series.

(ii) (Infinitesimal form.) For any symmetric $\mathsf K'$ over $R$, a variable $s$, $C_s:=-(\mathsf H+s\mathsf K')^{-1}$ and $d_s:=\det(1+s\mathsf H^{-1}\mathsf K')^{-1/2}$ in $R\fps{s}$: $\frac{\dd}{\dd s}(d_s\langle P\rangle_{C_s})=d_s\langle\frac12t^{\mathrm T}\mathsf K't\,P\rangle_{C_s}$.

(iii) (With $z$ and outer variables.) Let $\mathsf H\in GL_n(A)$ be symmetric, $C=-\mathsf H^{-1}$, $y=(y_1,\dots,y_m)$ outer variables, $\mathsf K$ symmetric with entries in the ideal $(y)$ of $A\fps{y}$, and $\mathsf H':=\mathsf H+\mathsf K$. Then every term of $t^{\mathrm T}\mathsf Kt/(2z)$ has weight $\geq1/2$, and for $P\in\cR_{n+m}(A)$
\[
E^t_{zC}\bigl[\exp\bigl(t^{\mathrm T}\mathsf Kt/(2z)\bigr)P\bigr]=\bigl(\det\mathsf H'/\det\mathsf H\bigr)^{-1/2}E^t_{-z\mathsf H'^{-1}}[P]\qquad\text{in }\cR_m(A).
\]
The same holds over $A_0\fps{q}$ for $\mathsf K$ with entries in $(q)$ and $P\in\cR_n(A_0)\fps{q}$, coefficientwise in $q$ (in the sense of \cref{def:fsp-ring}).
\end{lem}

\begin{proof}
(i), (ii). By the generating function $\langle P\rangle_C=[\exp(\frac12\sum_{ij}C_{ij}\de_i\de_j)P](0)$ (on $t^\alpha$ the operator removes the variables in pairs). Let $Q:=\frac12t^{\mathrm T}\mathsf K't$ and $\rho_s(P):=d_s\langle P\rangle_{C_s}$. As $\dd C_s/\dd s=C_s\mathsf K'C_s$, $\frac{\dd}{\dd s}\langle P\rangle_{C_s}=\frac12\sum_{ij}(C_s\mathsf K'C_s)_{ij}\langle\de_i\de_jP\rangle_{C_s}$; applying the Wick recursion $\langle t_iP\rangle_C=\sum_lC_{il}\langle\de_lP\rangle_C$ twice,
\[
\langle QP\rangle_{C_s}=\tfrac12\tr(\mathsf K'C_s)\,\langle P\rangle_{C_s}+\tfrac12\textstyle\sum_{lm}(C_s\mathsf K'C_s)_{lm}\langle\de_l\de_mP\rangle_{C_s};
\]
and $\frac{\dd}{\dd s}\log d_s=-\frac12\tr((\mathsf H+s\mathsf K')^{-1}\mathsf K')=\frac12\tr(C_s\mathsf K')$. Together $\frac{\dd}{\dd s}\rho_s(P)=\rho_s(QP)$, which is (ii). For (i) take $\mathsf K'=\mathsf K$ and $\Lambda_s(P):=\langle\exp(sQ)P\rangle_C=\sum_ks^k\langle Q^kP\rangle_C/k!$; then $\frac{\dd}{\dd s}\Lambda_s(P)=\Lambda_s(QP)$ and $\Lambda_0=\rho_0$, so all $s$-derivatives at $0$ agree ($\langle Q^kP\rangle_C$) and $\Lambda_s(P)=\rho_s(P)$ in $R\fps{s}$. The $s^k$-coefficients of both sides, of $d_s$ and of $C_s=-\sum_r(-s)^r(\mathsf H^{-1}\mathsf K)^r\mathsf H^{-1}$ lie in $I^k$, and $s=1$ is a ring homomorphism from such series to the complete ring $R$; it gives (i).
(iii) The terms of $t^{\mathrm T}\mathsf Kt/(2z)$ are multiples of $t_it_jy^\gamma z^{-1}$, $|\gamma|\geq1$. For $P=t^\alpha y^\beta z^j$, the term $(t^{\mathrm T}\mathsf Kt/(2z))^kt^\alpha/k!$ has degree $2k+|\alpha|$ in $t$, so $E^t_{zC}[\exp(t^{\mathrm T}\mathsf Kt/(2z))t^\alpha y^\beta z^j]=\langle\exp(\frac12t^{\mathrm T}\mathsf Kt)t^\alpha\rangle_Cy^\beta z^{j+|\alpha|/2}$, which by (i) over $R=A\fps{y}$, $I=(y)$ (with $\det(1+\mathsf H^{-1}\mathsf K)=\det\mathsf H'/\det\mathsf H$, $C_{\mathsf K}=-\mathsf H'^{-1}$) is the right side for this $P$. Both sides are linear and compatible with summable sums (the exponential has terms of weight $\geq0$ only), so the identity extends to $\cR_{n+m}(A)$. Over $A_0\fps q$ argue modulo $q^N$: over $B_N=A_0[q]/(q^N)$ the ideal $(q)$ is nilpotent, (i) applies with $R=B_N$ and $I=(q)$, and the monomial argument above gives the identity on $\cR_n(B_N)$.
\end{proof}

\begin{lem}[Products]\genusstatus{AI-proposed, unchecked}
\label{lem:fsp-prod}
Split the variables as $(t,y)$ and let $C=\diag(C_t,C_y)$ with constant blocks. Then, for $P\in\cR_{n+m}(A)$, $E_{zC}[P]=E^y_{zC_y}\bigl[E^t_{zC_t}[P]\bigr]$. Hence, for phases and amplitudes in separate variables,
\[
\SPn_{(c_1,c_2)}\bigl[F_1(t)+F_2(y);\,a_1(t)\,a_2(y)\bigr]=\SPn_{c_1}[F_1;a_1]\,\SPn_{c_2}[F_2;a_2],
\]
and the full symbols multiply, with $\det(-\mathsf H)^{1/2}:=\det(-\mathsf H_1)^{1/2}\det(-\mathsf H_2)^{1/2}$.
\end{lem}

\begin{proof}
For block-diagonal $C$ the generating function factorises, so $\langle t^\alpha y^\beta\rangle_C=\langle t^\alpha\rangle_{C_t}\langle y^\beta\rangle_{C_y}$; both sides are linear and compatible with summable sums. For separated data the integrand is the product of $\exp(F_{1,\geq3}/z)a_1$, a function of $t$, and $\exp(F_{2,\geq3}/z)a_2$, a function of $y$; and $E^t[P_1(t)P_2(y)]=E^t[P_1]\,P_2$ with $E^t[P_1]\in A\lser{z}$.
\end{proof}

\begin{lem}[Formal implicit function theorem]\genusstatus{AI-proposed, unchecked}
\label{lem:fsp-ift}
Let $B$ be a commutative $\bbQ$-algebra, complete and separated for an ideal $I$ (e.g.\ $B=A\fps{y}$, $I=(y)$, or $B=A_0\fps{q}$, $I=(q)$), and $F\in B\fps{t}$ with $\nabla F(0)\in I^n$ and $\mathsf H_0:=\Hess F(0)\in GL_n(B)$. Then there is a unique $c\in I^n$ with $\nabla F(c)=0$, and $\Hess F(c)$ is invertible.
\end{lem}

\begin{proof}
Substitution of elements of $I$ into power series converges, $B$ being complete. For $c\in I^n$ put $\Psi(c):=c-\mathsf H_0^{-1}\nabla F(c)$; it lies in $I^n$, because $\nabla F(c)=\nabla F(0)+\mathsf H_0c+O(c^2)$ and $\nabla F(0)\in I^n$. Writing $\nabla F(c)-\nabla F(c')=\mathsf N(c,c')(c-c')$ with difference quotients, $\mathsf N(c,c')\equiv\mathsf H_0$ modulo $I$, one gets $\Psi(c)-\Psi(c')=-\mathsf H_0^{-1}(\mathsf N(c,c')-\mathsf H_0)(c-c')$, a matrix with entries in $I$ times $c-c'$; so $c-c'\in I^r$ implies $\Psi(c)-\Psi(c')\in I^{r+1}$. The iterates $c_{k+1}:=\Psi(c_k)$, $c_0=0$, satisfy $c_{k+1}-c_k\in I^{k+1}$, so they converge to some $c\in I^n$ with $\Psi(c)=c$, i.e.\ $\nabla F(c)=0$. Two such zeros $c,c'$ satisfy $c-c'=\Psi(c)-\Psi(c')$, hence $c-c'\in\bigcap_rI^r=0$. Finally $\Hess F(c)\equiv\mathsf H_0$ modulo $I$, and a matrix over $B$ invertible modulo $I$ is invertible.
\end{proof}

% Formal stationary phase, part 3 of 5 (sec:fsp-basic, continued): translations, moving the base point, parameters,
% derivatives in a parameter, the stationary phase at a generic base point. Continued in part 4.

\begin{lem}[Translations]\genusstatus{AI-proposed, unchecked}
\label{lem:fsp-transl}
Let $c$ be a non-degenerate critical point of $F$ over $A$, $a$ an amplitude and $p\in A^n$ small in one of two senses: (a) $A=A_0\fps{q}$, $F$ and $c$ are defined over $A_0$, $a\in A_0\fps{x-c}\lser{z}\fps{q}$ and $p\in(qA)^n$; or (b) $p\in(zA\fps{z})^n$ and $a\in A\fps{x-c}\lser{z}$. Then the translate $b_p(t):=\exp\bigl((F(c+t+p)-F(c+t))/z\bigr)a(c+t+p)$ (the formal analogue of $x\mapsto x+p$ applied to $\re^{F/z}a$) is an amplitude, in $A_0\fps{t}\lser{z}\fps{q}$, resp.\ $A\fps{t}\lser{z}$, and $\SPn_c[F;b_p]=\SPn_c[F;a]$.
\end{lem}

\begin{proof}
In case (a), $F(c+t+p)-F(c+t)$ has coefficients in $qA$, so the exponential converges $q$-adically, each $q$-coefficient lying in $A_0\fps{t}[z^{-1}]$. In case (b), with $p=zv$, $(F(c+t+zv)-F(c+t))/z=\nabla F(c+t)\cdot v+O(z)$ lies in the ideal $(t,z)$ of $A\fps{t,z}$ because $\nabla F(c+t)=O(t)$, so the exponential converges in $A\fps{t,z}$. Let $s$ be a variable and $B(s):=b_{sp}$. At fixed degree in $t$ and powers of $z$ (and $q$), the coefficient of $B(s)$ is a polynomial in $s$: in case (a) each $p_i$ carries a factor $q$; in case (b) each factor $\nabla F(c+t)\cdot sv$ of the exponential has degree $\geq1$ in $t$, while every other factor, and every term of the Taylor expansion of $a(c+t+szv)$ that involves $s$, carries a positive power of $z$. By \cref{lem:fsp-univ} the coefficients of $\SPn_c[F;B(s)]$ are polynomials in $s$, with derivative $\SPn_c[F;\dd B/\dd s]$. Differentiating $B(s)$ and comparing with $\de_iB(s)$, $\de_i=\de/\de t_i$, gives $\dd B/\dd s=\sum_ip_i\bigl(\de_iB+z^{-1}(\de_iF)(c+t)B\bigr)$; by \cref{lem:fsp-ibp} with $h_i:=p_iB(s)$ (coefficientwise in $q$ in case (a)) the derivative vanishes. A polynomial over a $\bbQ$-algebra with zero derivative is constant: $\SPn_c[F;B(1)]=\SPn_c[F;B(0)]$.
\end{proof}

\begin{cor}[Moving the base point]\genusstatus{AI-proposed, unchecked}
\label{cor:fsp-movebase}
Let $A:=A_0\fps{q}$, let $F_0\in A_0\fps{x-c_0}$ have the non-degenerate critical point $c_0$ over $A_0$, with Hessian $\mathsf H_0$, let $F:=F_0+F_1$ with $F_1\in qA_0\fps{x-c_0}\fps{q}$, and $a\in A_0\fps{x-c_0}\lser{z}\fps{q}$. Then $F$ has a unique critical point $c\in c_0+(qA)^n$, non-degenerate with Hessian $\mathsf H_c$. For amplitudes $b$ put
\[
\mathcal S_q[F;b]:=\exp\bigl((F(c)-F_0(c_0))/z\bigr)\,(\det\mathsf H_c/\det\mathsf H_0)^{-1/2}\,\SPn_c[F;b]
\]
(the power is a binomial series, as $\det\mathsf H_c/\det\mathsf H_0\in1+qA$, and the exponential converges, as $(F(c)-F_0(c_0))/z$ lies in $qA_0\fps{q}\,z^{-1}$). Then $\SPn_{c_0}[F_0;\re^{F_1/z}a]=\mathcal S_q[F;a]$ in $A_0\lser{z}\fps{q}$, the left side being computed coefficientwise in $q$. With the branch rule $\det(-\mathsf H_c)^{1/2}:=\det(-\mathsf H_0)^{1/2}(\det\mathsf H_c/\det\mathsf H_0)^{1/2}$, $\mathcal S_q[F;a]$ is $\mathrm{SP}_c[\re^{F/z}a\,\dd^nx]$ divided by the $q$-independent symbols $[\re^{F_0(c_0)/z}](2\pi z)^{n/2}[\det(-\mathsf H_0)^{-1/2}]$.
\end{cor}

\begin{proof}
\Cref{lem:fsp-ift} over $B=A$, $I=(q)$, gives $c$: $\nabla F(c_0)=\nabla F_1(c_0)\in(q)^n$, and $\Hess F(c_0)\in\mathsf H_0+q\,\mathfrak{gl}_n(A)$ is invertible. Put $p:=c-c_0$. By \cref{lem:fsp-transl}~(a) for $F_0$ and the amplitude $\re^{F_1/z}a$, the left side is $\SPn_{c_0}[F_0;b_p]$, $b_p(t)=\exp((F(c+t)-F_0(c_0+t))/z)\,a(c+t)$. With $F(c+t)=F(c)+\frac12t^{\mathrm T}\mathsf H_ct+F_{c,\geq3}(t)$ ($c$ is critical) and $F_0(c_0+t)=F_0(c_0)+\frac12t^{\mathrm T}\mathsf H_0t+F_{0,\geq3}(t)$, in $\cR_n(A_0)\fps{q}$, that is, modulo every $q^N$ (\cref{def:fsp-ring}; both sides are the exponential of $(F(c+t)-F_0(c_0)-\frac12t^{\mathrm T}\mathsf H_0t)/z$),
\[
\re^{F_{0,\geq3}/z}\,\re^{(F(c+t)-F_0(c_0+t))/z}=\re^{(F(c)-F_0(c_0))/z}\,\re^{t^{\mathrm T}(\mathsf H_c-\mathsf H_0)t/(2z)}\,\re^{F_{c,\geq3}/z}.
\]
The scalar factor comes out of $E_{-z\mathsf H_0^{-1}}$. As $\mathsf H_c-\mathsf H_0$ has entries in $qA$, \cref{lem:fsp-pg}~(iii) gives, for $P\in\cR_n(A_0)\fps{q}$,
\[
E_{-z\mathsf H_0^{-1}}\bigl[\re^{t^{\mathrm T}(\mathsf H_c-\mathsf H_0)t/(2z)}P\bigr]=(\det\mathsf H_c/\det\mathsf H_0)^{-1/2}\,E_{-z\mathsf H_c^{-1}}[P],
\]
and $E_{-z\mathsf H_c^{-1}}[\re^{F_{c,\geq3}/z}a(c+t)]=\SPn_c[F;a]$.
\end{proof}

\begin{lem}[Parameters]\genusstatus{AI-proposed, unchecked}
\label{lem:fsp-param}
(i) (Dependence on parameters.) If the Taylor coefficients at $c$ of $F$ (orders $\geq2$) and of $a$ lie in a $\bbQ$-subalgebra $A'\subset A$ and $\det\mathsf H\in A'^\times$, all coefficients of $\SPn_c[F;a]$ lie in $A'$; rational data give rational coefficients, with poles only where the data have poles or $\det\mathsf H=0$. Over $A=A_0\fps{q}$, $q=(q_1,\dots,q_p)$, if $F|_{q=0}$ has the non-degenerate critical point $c_0$, then $F$ has a unique critical point $c(q)\in c_0+(q)A^n$, non-degenerate, and the coefficients of $\SPn_{c(q)}[F;a]$, $F(c(q))$ and $\det\Hess F(c(q))$ lie in $A_0\fps{q}$. If $F$, $a$ are holomorphic near $(c_0,q_0)$ and $c(q)$ is a holomorphic family of non-degenerate critical points with $c(q_0)=c_0$, these are holomorphic near $q_0$.

(ii) (Taylor expansion and evaluation.) In the holomorphic case, the Taylor series at $q_0$ of each coefficient is the corresponding coefficient over $\bbC\fps{q-q_0}$ for the Taylor-expanded data, at the critical point of the formal case, and its value at $q_0$ is that of $\SPn_{c_0}[F(\cdot\,;q_0);a(\cdot\,;q_0)]$. In the formal case, $q\mapsto0$ maps $\SPn_{c(q)}[F;a]$ to $\SPn_{c_0}[F|_{q=0};a|_{q=0}]$.

(iii) (Perturbation in a parameter.) \Cref{cor:fsp-movebase}: the $q$-expansion of the stationary phase at the moving critical point, with its two prefactors, is that at $c_0$ with $\re^{F_1/z}$ in the amplitude; for $F_1=qg$ and $a$ free of $q$, $\coef{q^d}\SPn_{c_0}[F_0;\re^{qg/z}a]=\SPn_{c_0}[F_0;(g/z)^da]/d!$.

(iv) (A $z$-dependent part of the phase.) Let $c$ be a non-degenerate critical point of $F$ over $A$, $a\in A\fps{x-c}\lser z$ an amplitude (free of $\varepsilon$), $g\in A\fps{x-c}$, $g(c)=0$, $\varepsilon$ a variable and $F_\varepsilon:=F+\varepsilon g$, with its unique critical point $c(\varepsilon)\in c+(\varepsilon A\fps{\varepsilon})^n$. Then $\mathcal S_\varepsilon[F_\varepsilon;a]=\SPn_c[F;\re^{\varepsilon g/z}a]$; giving $\varepsilon$ the weight $1$, this series has finitely many non-zero terms $\varepsilon^dz^k$ of each bounded weight $d+k$, the substitution $\varepsilon\mapsto z$ is well defined on such series, and $[\mathcal S_\varepsilon[F_\varepsilon;a]]_{\varepsilon=z}=\SPn_c[F;\re^ga]$. So the stationary phase with the $z$-dependent phase $F+zg$ is defined and equals the one with $zg$ moved into the amplitude as $\re^g$.
\end{lem}

\begin{proof}
(i) The first assertion is \eqref{eq:fsp-univ}, as $C=-\operatorname{adj}(\mathsf H)/\det\mathsf H$. In the formal case \cref{lem:fsp-ift} over $A$, $I=(q)$, gives $c(q)$, and the data at $c(q)$ lie in $A$. In the holomorphic case $\de_x^\alpha F(c(q);q)$, $\de_x^\alpha a(c(q);q)$ and the covariance are holomorphic in $q$. (ii) The Taylor map and the evaluations are ring homomorphisms (\cref{lem:fsp-univ}~(i)); the Taylor series of $c(q)$ is a critical point of the expanded phase, hence the formal one by uniqueness (\cref{lem:fsp-ift}). (iii) $\re^{qg/z}a=\sum_dq^d(g/z)^da/d!$. (iv) \Cref{lem:fsp-ift} over $A\fps{\varepsilon}$ gives $c(\varepsilon)$, as $\nabla F_\varepsilon(c)=\varepsilon\nabla g(c)$; the first identity is \cref{cor:fsp-movebase} with $q=\varepsilon$, $F_0=F$, $F_1=\varepsilon g$. In $\cR_{n;\varepsilon}(A)$ every term of $\varepsilon g/z$ and of $F_{\geq3}/z$ has weight $\geq1/2$ ($g=O(t)$), so $\exp(F_{\geq3}/z)\exp(\varepsilon g/z)a(c+t)$ lies there; it is the same formal sum as in $\cR_n(A)\fps{\varepsilon}$, $E_{zC}$ preserves the weight, and each coefficient of $\varepsilon^dz^k$ is the same finite sum in both readings, which gives the finiteness. The substitution $t^\alpha\varepsilon^dz^j\mapsto t^\alpha z^{d+j}$ maps $\cR_{n;\varepsilon}(A)$ to $\cR_n(A)$ (each output monomial comes from finitely many inputs of its weight), is a ring homomorphism compatible with summable sums, maps $\exp(\varepsilon g/z)$ to $\re^g$, and commutes with $E_{zC}$ (both composites send $t^\alpha\varepsilon^dz^j$ to $\langle t^\alpha\rangle_Cz^{d+j+|\alpha|/2}$).
\end{proof}

\begin{lem}[Derivatives in a parameter]\genusstatus{AI-proposed, unchecked}
\label{lem:fsp-dq}
In the situation of \cref{cor:fsp-movebase}, for every amplitude $a\in A_0\fps{x-c_0}\lser{z}\fps{q}$,
\[
\frac{\dd}{\dd q}\,\mathcal S_q[F;a]=\mathcal S_q\bigl[F;\,z^{-1}(\de_qF)\,a+\de_qa\bigr]\qquad\text{in }A_0\lser{z}\fps{q},
\]
i.e.\ $\frac{\dd}{\dd q}\mathrm{SP}[\re^{F/z}a\,\dd^nx]=\mathrm{SP}[\de_q(\re^{F/z}a)\,\dd^nx]$; the same holds for $\mathrm D:=(q_0+q)\,\dd/\dd q$, $q_0\in A_0$, in place of $\de_q$. If $\mathrm DF=\varphi$ does not depend on $q$ and $\mathrm Da=0$, then $P(z\mathrm D)\,\mathcal S_q[F;a]=\mathcal S_q[F;P(\varphi)a]$ for every polynomial $P$ over $A_0[z]$.
\end{lem}

\begin{proof}
By \cref{cor:fsp-movebase}, $\mathcal S_q[F;a]=\SPn_{c_0}[F_0;\re^{F_1/z}a]$ with a $q$-independent functional, applied coefficientwise in $q$; so $\dd/\dd q$ acts on the amplitude, $\frac{\dd}{\dd q}(\re^{F_1/z}a)=\re^{F_1/z}(z^{-1}(\de_qF)a+\de_qa)$, and \cref{cor:fsp-movebase} for this amplitude gives the claim. For $\mathrm D$ multiply by $q_0+q$. If $\mathrm DF=\varphi$ is $q$-independent and $\mathrm Da=0$, then $\mathrm D(\varphi^ka)=0$ and $z\mathrm D\,\mathcal S_q[F;\varphi^ka]=\mathcal S_q[F;\varphi^{k+1}a]$; induction and linearity over $A_0[z]$ give the clause for $P$ (without $\mathrm Da=0$, terms $z\mathrm D(\varphi^ka)$ appear).
\end{proof}

\begin{lem}[The stationary phase at a generic base point]\genusstatus{AI-proposed, unchecked}
\label{lem:fsp-generic}
Let $\delta(\sigma):=(\sigma^6-\sigma)/5$, so that $\AL=\bbQ[L,\delta(L)^{-1}]$ and $\Dq L=\delta(L)$, and let $\bbA_\sigma:=\bbQ[\sigma,\delta(\sigma)^{-1}]$ be the same ring in a variable $\sigma$; for $f\in\bbA_\sigma$ let $f(L+t)\in\AL\fps{t}$ be its Taylor expansion at $\sigma=L$. Let $S$ be a phase at $\sigma=L$ over $\AL$ with $S'(L)=0$ and $S''(L+t)=-1/\delta(L+t)$, $\varepsilon$ a variable and $F(\sigma;\varepsilon):=S(\sigma)+\varepsilon\sigma$. Let $\mathsf r(\sigma):=5\sigma/(\sigma^5-1)\in\bbA_\sigma^\times$, and for $\sigma_\ast\in\{L,\sigma_\varepsilon\}$ let $(\mathsf r(\sigma)/\mathsf r(\sigma_\ast))^{1/2}$ be the square root with constant term $1$ of the expansion of $\mathsf r(\sigma)/\mathsf r(\sigma_\ast)$ at $\sigma_\ast$.

(i) $F(\cdot\,;\varepsilon)$ has a unique critical point $\sigma_\varepsilon\in L+\varepsilon\AL\fps{\varepsilon}$, non-degenerate, and $\sigma_\varepsilon=\exp(\varepsilon\Dq)(L)$, where $\exp(\varepsilon\Dq):=\sum_m\varepsilon^m\Dq^m/m!$; the ring homomorphism $\AL\to\AL\fps{\varepsilon}$, $L\mapsto\sigma_\varepsilon$, is $\exp(\varepsilon\Dq)$.

(ii) For $b\in\bbA_\sigma\fps{z}$, with $\exp(\varepsilon\Dq)$ acting coefficientwise in $z$,
\[
\SPn_{\sigma_\varepsilon}\bigl[F(\cdot\,;\varepsilon);(\mathsf r(\sigma)/\mathsf r(\sigma_\varepsilon))^{1/2}b\bigr]=\exp(\varepsilon\Dq)\,\SPn_L\bigl[S;(\mathsf r(\sigma)/\mathsf r(L))^{1/2}b\bigr].
\]

(iii) For $a:=(\mathsf r(\sigma)/\mathsf r(L))^{1/2}b$, $b\in\bbA_\sigma\fps{z}$, and $\Dq$ acting coefficientwise in $z$,
\begin{equation}
\label{eq:fsp-generic}
(z\Dq+L)\bigl(L\cdot\SPn_L[S;a]\bigr)=L\cdot\SPn_L[S;\sigma a]\qquad\text{in }\AL\lser{z},
\end{equation}
where $\sigma a$ is the amplitude $(L+t)\,a(L+t,z)$.
\end{lem}

\begin{proof}
(i) $F(L+t;\varepsilon)=S(L)+\varepsilon L+\varepsilon t+\sum_{k\geq2}S^{(k)}(L)t^k/k!$ has gradient $\varepsilon$ and Hessian $-1/\delta(L)\in\AL^\times$ at $t=0$, so \cref{lem:fsp-ift} over $\AL\fps{\varepsilon}$ gives $\sigma_\varepsilon=L+t_\varepsilon$, $t_\varepsilon\in\varepsilon\AL\fps{\varepsilon}$. For $f\in\bbA_\sigma$, the substitution $t=t_\varepsilon$ in $f(L+t)$ is the ring homomorphism $\bbA_\sigma\to\AL\fps{\varepsilon}$, $\sigma\mapsto\sigma_\varepsilon$ (well defined, $\delta(\sigma_\varepsilon)\equiv\delta(L)$ modulo $\varepsilon$ being a unit): both are ring homomorphisms agreeing on $\sigma$. So $S''(\sigma_\varepsilon)=-1/\delta(\sigma_\varepsilon)$, and differentiating $S'(\sigma_\varepsilon)+\varepsilon=0$ gives $\sigma_\varepsilon'=\delta(\sigma_\varepsilon)$, $\sigma_0=L$. On the other hand $\exp(\varepsilon\Dq)$ is a ring homomorphism (Leibniz rule for $\Dq^m$) with $\frac{\dd}{\dd\varepsilon}\exp(\varepsilon\Dq)(f)=\exp(\varepsilon\Dq)(\Dq f)$, so $\tilde\sigma_\varepsilon:=\exp(\varepsilon\Dq)(L)$ satisfies $\tilde\sigma_\varepsilon'=\delta(\tilde\sigma_\varepsilon)$, $\tilde\sigma_0=L$. A solution $y\in\AL\fps{\varepsilon}$ of $y'=\delta(y)$, $y(0)=L$, is unique, since $(m+1)y_{m+1}=\coef{\varepsilon^m}\delta(y)$ involves only $y_0,\dots,y_m$. So $\sigma_\varepsilon=\tilde\sigma_\varepsilon$, and the two homomorphisms agree on $L$, hence on $\AL$.
(ii) For $f\in\bbA_\sigma$, the expansion $f(\sigma_\varepsilon+t)$ is $\exp(\varepsilon\Dq)$ applied coefficientwise to $f(L+t)$: both are ring homomorphisms $\bbA_\sigma\to\AL\fps{\varepsilon}\fps{t}$ sending $\sigma$ to $\sigma_\varepsilon+t$, by (i). This applies to the derivatives of $-1/\delta$ (the Taylor coefficients of order $\geq2$ of $F(\cdot\,;\varepsilon)$ at $\sigma_\varepsilon$), to the covariance $\delta(\sigma_\varepsilon)$, to the coefficients of $b$ and to $\mathsf r$; the binomial series has rational coefficients. \Cref{lem:fsp-univ}~(i) and \cref{lem:fsp-const} give (ii).
(iii) Put $N:=\SPn_L[S;a]$ and $N^\sigma:=\SPn_L[S;\sigma a]$. The amplitude $a$ is free of $\varepsilon$, and at $\sigma_\varepsilon$ it equals $(\mathsf r(\sigma_\varepsilon)/\mathsf r(L))^{1/2}(\mathsf r(\sigma)/\mathsf r(\sigma_\varepsilon))^{1/2}b$, the first factor being the constant $(\exp(\varepsilon\Dq)(\mathsf r(L))/\mathsf r(L))^{1/2}$. By linearity and (ii), and since the Hessian at $\sigma_\varepsilon$ is $\exp(\varepsilon\Dq)(-1/\delta(L))$, the series of \cref{cor:fsp-movebase} (base ring $\AL$, parameter $\varepsilon$, $F_0=S$, $F_1=\varepsilon\sigma$) is
\[
\mathcal S_\varepsilon[F;a]=\re^{(F(\sigma_\varepsilon;\varepsilon)-S(L))/z}\Bigl(\frac{\exp(\varepsilon\Dq)(\delta(L)\mathsf r(L))}{\delta(L)\mathsf r(L)}\Bigr)^{1/2}\exp(\varepsilon\Dq)(N)=\re^{(F(\sigma_\varepsilon;\varepsilon)-S(L))/z}\,\frac{\sigma_\varepsilon}L\,\exp(\varepsilon\Dq)(N),
\]
because $\delta(L)\mathsf r(L)=L^2$ and two elements of $1+\varepsilon\AL\fps{\varepsilon}$ with equal squares are equal. By \cref{lem:fsp-dq} with $\de_\varepsilon F=\sigma$, $\de_\varepsilon a=0$, the derivative at $\varepsilon=0$ of the left side is $z^{-1}N^\sigma$. On the right, $\frac{\dd}{\dd\varepsilon}F(\sigma_\varepsilon;\varepsilon)=\sigma_\varepsilon$ by criticality, so the three factors contribute $L/z$, $\Dq L/L$ and $\Dq N$. Hence $(L/z)N+(\Dq L/L)N+\Dq N=z^{-1}N^\sigma$; multiply by $zL$.
\end{proof}

\input{sections/a_formalsp_4}



% Appendix: signs and normalisations (label prefix sgn). Conventions and dictionaries only.
% Continued in part 2.
\section{Signs and normalisations}
\label{sec:sgn}

This appendix collects the sign conventions and normalisations, and the dictionaries to the sources quoted. Every entry is a convention, or a dictionary derived at the place cited; the body refers to it only for orientation. The comparison with \cite{CGYZ26} at the end is proved there, for the definitions of \cite{CGYZ26} as recalled there, and enters no proof of the body.

%-------------------------------------------------------------------------------
\subsection{The \texorpdfstring{$R$}{R}-matrix and its action}
\label{sec:sgn-R}

\emph{Givental's convention.} $R(z)=1+O(z)$ is taken in Givental's convention, $S\sim R\,\re^{u/z}$, with legs $R^{-1}(\psi)$. That the columns $R(z)\tilde e_i\,\re^{u_i/z}$ solve $z\de_Tv=H*v$ along $t=T(q)H$ is not part of the convention but the quoted property (QRR) (\cref{def:set-rows}~(b), \cref{thm:set-qrr}); and $R(-z)^T\eta R(z)=\eta$ in the flat basis (\cref{def:set-rows}~(a)); equivalently $R(z)^{-1}=R(-z)^*$, the adjoint with respect to $\eta$.

\emph{The action.} The action is that of \cite[Def.~2.2]{PPZ15}, also the convention of Givental and Teleman by the remark after it: the legs carry $R^{-1}(\psi)$, and the edge is $E_R(\psi,\psi')=\bigl(\eta^{-1}-R^{-1}(\psi)\,\eta^{-1}R^{-1}(\psi')^T\bigr)/(\psi+\psi')$, a power series because $R$ is symplectic \cite[Remark~2.1, (6)]{PPZ15}. The action of \cite[\S2.3]{Dunin-Barkowski:2015caa} is the same up to the factor $1/|\Aut\Gamma|$, omitted there and present in \cite[(2-26)]{DuninBarkowski:2012bw} and \cite[Def.~2.2]{PPZ15}, and up to the sign of the dilaton leaves (legs $R^{-1}(\psi)$ acting on column vectors, the row index being the outgoing colour); by \cite[Remark~2.1]{Dunin-Barkowski:2015caa} Givental's matrix $R$ is the transpose of theirs. The unit-preserving action is $R.\Omega=R\,\ttr_R\,\Omega$ with the translation $\ttr_R(z)=z\,1-z\,R^{-1}(z)\,1$ of \cite[Def.~2.13]{PPZ15}, the translation $T_b$ of \cite[Prop.~2.12]{PPZ15}; so the dilaton leaves carry $+\psi\bigl(1-R^{-1}(\psi)1\bigr)$. This is also the sign of the dilaton leaves of topological recursion (\cref{prop:mod-Rtr}) and of \cite[\S2.3]{DuninBarkowski:2012bw}; \cite[\S2.3]{Dunin-Barkowski:2015caa} prints the opposite sign. $\mathrm{Cont}_\Gamma$ includes the factor $|\Aut\Gamma|^{-1}$ (\cref{def:set-graphsum}).

\emph{The colour basis.} In the basis of idempotents $e_i=\Delta_i^{-1/2}\tilde e_i$, $\eta(e_i,e_j)=\delta_{ij}/\Delta_i$, a vertex of genus $g_v$ and colour $i$ carries $\Delta_i^{g_v-1}$; the edge between colours $i$ and $j$ is
\[
\frac{\delta_{ij}\Delta_i-\Delta_i^{1/2}\Delta_j^{1/2}\sum_{a,b}\eta^{ab}\Qr^{(i)}_a(-\psi)\Qr^{(j)}_b(-\psi')}{\psi+\psi'},\qquad\eta^{ab}=\tfrac15\delta_{a+b\equiv3\,(5)};
\]
and the dilaton leaf of colour $i$ in degree $k\geq2$ is $(-1)^k\Delta_i^{1/2}\coef{z^{k-1}}\Qr^{(i)}_0(z)$. These are \cite[Defs.~2.2, 2.13]{PPZ15} in the basis $e_i$; the proofs use the form of \cite{PPZ15}.

\emph{Lowered and vector forms.} Solutions $\phi$ and gauges are written in lowered components $Y_a=\eta(H^a,\phi)$ (\cref{conv:tra-forms}). As a matrix $\eta$ is $5$ times the permutation matrix of $a\leftrightarrow3-a$; a gauge $P$ acting on lowered components acts on vectors by $P^{\mathrm{vec}}=\eta^{-1}P\eta$, and the dilaton factor is $d:=(P_0)_{33}=(P_0^{\mathrm{vec}})_{00}$ (\cref{def:tra-dressing}). The condition (A) is stated in lowered form, \cref{lem:tra-E} and the leaf point of \cref{lem:tra-T} in vector form; parts (b), (d) of \cref{lem:tra-T} are computed in lowered form.

%-------------------------------------------------------------------------------
\subsection{Variables, colours and quantum Riemann--Roch}
\label{sec:sgn-var}

\emph{Operators and points.} $\Dq=q\,\dd/\dd q=-(1-u)\,\dd/\dd u$, so $z\Dq=-(1-u)z\de_u$; along the mirror coordinate $\de_T=C_1^{-1}\Dq$ (\cref{prop:set-cyclicLR}); in \cref{sec:gau-cyc}, $\de_t=(\Dq t)^{-1}\Dq$ (\cref{def:gau-conc}); $\theta_u=u\,\dd/\dd u$, $\theta_{\zh}=\zh\,\dd/\dd\zh$ (\cref{lem:gau-HPF}). The points are those of \cref{def:set-points}.

\emph{Colours.} Colour $i\in\bbZ/5$ is the fixed point $p_i$ of the torus, of weight $\lambda_i=\zeta^i$; the eigenvalue $\zeta^iL$ of $C_1H*$, so that $\Dq u_i=\zeta^iL$ (\cref{cor:set-cyclicPF5}); and the critical point $\sigma=\zeta^iL$ of the Mellin--Barnes summand (\cref{def:gam-c}). The colour covariances are $\Phi_i(\zh,u)=\Phi_0(\zeta^{-i}\zh,u)$ (\cref{lem:gau-colrows}), $\Dn_i(z)=\Dn_0(\zeta^{-i}z)$ (\cref{cor:gau-colours}) and $h_{\bw_i}(z)=h(\zeta^{-i}z)$ (\cref{lem:set-hrot}). The square roots are fixed by $\zeta^{1/2}:=\zeta^3$: $\Delta_i^{1/2}=\zeta^{4i}C_0/L$, and $(W^0_i)^{-1/2}=\zeta^i$ for $W^0_i=\zeta^{3i}$ (\cref{def:set-frob}).

\emph{Flat and hatted.} The flat theory uses $(u,z)$, the classes $H^a$ and the rows $\Qr$; the hatted theory (\cref{def:set-hat}) uses $(u,\zh)$, $z=s^4\zh$, the classes $\hat H^a=s^{1-\nu_a}H^a$, the rows $\hat\Qr_a=s^{1-\nu_a}\Qr_a$, the hatted TFT $\hat\TFT_{g,n}(\tilde e_i,\dots,\tilde e_i)=(\zeta^{4i}\hat x)^{2g-2+n}$ with $\hat x=\varpi_0$, the unit $\hat1=s\,1$ and the rescaling $\Omega_{g,n}\mapsto s^{-(2g-2+n)}\Omega_{g,n}$; $\eta$ is preserved, since $\nu_a+\nu_b=2$ on the $\eta$-pairs $(0,3)$, $(1,2)$, $(4,4)$; and $G_\Gamma=u^{2g-2}T_\Gamma$ (\cref{lem:flh-hat}).

\emph{Quantum Riemann--Roch.} The tangent directions at $p_i$ contribute their weights $\lambda_i-\lambda_j$, $j\neq i$, and the Euler-class twist minus its fibre weight $5\lambda_i$ \cite[(10)]{LP20}; so $\bw_i=\{\zeta^i-\zeta^j\}_{j\neq i}\cup\{-5\zeta^i\}$ and $h_1=-\frac3{20}$ (\cref{thm:set-qrr,rmk:set-qrrsign}), used in the form $\Phi_i(z,1)=h_{\bw_i}(z)$ (\cref{eq:set-qrr}). The Bernoulli numbers and polynomials are those of \cref{def:fsp-stirling}, $\xi\,\re^{\theta\xi}/(\re^\xi-1)=\sum_nB_n(\theta)\xi^n/n!$ \cite[24.2.3]{DLMF}, $B_n=B_n(0)$; in $h_{\bw}$, $\Gst$ and Stirling's series only the numbers $B_{2k}$, $k\geq1$, enter, so there the conventions $\xi/(\re^\xi-1)$ and $\xi/(1-\re^{-\xi})$ give the same formulas ($B_1=-\frac12$ occurs only inside the proof of \cref{lem:fsp-stirfe}). The polynomials $B_n(\theta)$ of \cref{lem:fsp-dlmf5118} and the values $B_{5l+1}(\frac15)$ in $\Ginf$ (\cref{thm:gam-ginf}) depend on the convention, which is that of \cite[24.2.3]{DLMF}; the other one would replace $B_n(\theta)$ by $B_n(\theta+1)=B_n(\theta)+n\theta^{n-1}$.

%-------------------------------------------------------------------------------
\subsection{Normalisations of the model}
\label{sec:sgn-model}

The note uses one normalisation of the model, $\Wv(p)=p^4/4+v/p$ with the integrands $p^m\re^{-\Wv(p)/z}\dd p$ and the curve $x=\Wv(p)$, $y=p$ (\cref{def:set-model,def:mod-curve}); its branch points are $\varsigma_b=\zeta^{-b}\varsigma$, $\varsigma^5=v$, with $x''(\varsigma_b)=5\varsigma_b^2$ and the pairwise distinct critical values $x(\varsigma_b)=\frac54\zeta^b\varsigma^4$. \Cref{tab:sgn-model} gives the other normalisations that occur in the sources quoted.

\begin{table}[ht]
\caption{Normalisations of the model.}
\label{tab:sgn-model}
\centering
\footnotesize
\renewcommand{\arraystretch}{1.2}
\begin{tabular}{@{}>{\raggedright\arraybackslash}p{0.24\textwidth}>{\raggedright\arraybackslash}p{0.4\textwidth}>{\raggedright\arraybackslash}p{0.3\textwidth}@{}}
\hline
form & relation to $\Wv$ & use \\
\hline
\cite[(8-1)]{Dunin-Barkowski:2015caa}, $n=4$: $f=p^4+a_2p^2+a_3p+a_4+\frac{a_5}{p-a_1}$ & at $(a_1,\dots,a_5)=(0,0,0,0,4v)$, $f=4\Wv$; critical points $p^5=v$, critical values $5p_c^4$ & the curve to which the quoted theorem is applied in \cref{thm:mod-tr}, transferred to $x=\Wv$ by \cref{lem:mod-scale}; \cite[\S8]{Dunin-Barkowski:2015caa} itself is not used \\
the curve $x=p^4+V/p$, $y=p$ (cf.\ \cite[(7.2)]{ABDKS25}) & $V=4v$: $x=4\Wv$; $\omega_{h,n}[\Wv]=4^{2h-2+n}\,\omega_{h,n}[p^4+V/p]|_{V=4v}$, and every pairing scales by $4^{2h-2+n}$ & not used: \cite[Theorem~7.1]{ABDKS25} is applied to $x=\Wv$ itself (\cref{lem:mod-abdks}) \\
$\lambda=x^4+u'/x$, integrand $\re^{\lambda/z'}x^k\dd x$ & $x=p$, $u'=4v$, $z'=-4z$: $\re^{\lambda/z'}=\re^{-\Wv/z}$; at $v=1$ the row $b$ of colour $i$ of \cref{def:set-M0} is $\varepsilon_b\zeta^{i(b+1)}\SPn_{x=1}[\lambda;x^{\nu_b}](-4\zeta^{-i}t)$ (\cref{prop:set-modrows,lem:con-S2}) & dictionary only \\
TFT of topological recursion & $\Delta_j=-5\varsigma_j^2$ in the basis $e^{\mathrm{TR}}_j$ of \cite[(2-30)]{Dunin-Barkowski:2015caa}, with $x=x(\varsigma_j)-w_j^2/2$ and the uniform orientation $\frac{\dd w_j}{\dd p}(\varsigma_j)=\sqrt{-5}\,\varsigma_j$ (one root for all $j$); frozen basis $\bar e_j=e^{\mathrm{TR}}_j/\sqrt{-5}$ & \cref{prop:mod-Rtr,thm:mod-tr} \\
model TFT & defined by $\hat\Delta_i^{1/2}:=\zeta^{4i}\varsigma$ in the normalised canonical basis, unit $\sum_i\hat\Delta_i^{-1/2}\tilde e_i$ & \cref{def:set-modcohft} \\
topological recursion versus the model & $\omega_{h,n}=(-5)^{h-1}\sum\langle\prod_i\psi^{d_i}H^{a_i}\rangle_h\prod_i\dd\xi_{a_i,d_i}$; $\Fm_h=(-5)^{1-h}\FTR_h$ & \cref{thm:mod-tr,cor:mod-trF} \\
Hurwitz space & \cite[(8-1)]{Dunin-Barkowski:2015caa} with $n=4$ on the line $(a_1,\dots,a_5)=(0,0,0,0,4v)$; called $H_{0;1,4}$ in \cite{CvdLPS23} & the identification is not used \\
$H$ and $\se$ & $H=\dd/\dd v$, $\se=z\,1+vH=\fr(1,0)$; $\langle H\rangle_1=-1/(12v)$, $\langle\se\rangle_1=\frac18$; replacing $H$ and $\se$ by $-H$ and $-\se$ changes both signs & \cref{lem:tra-frames}; sign anchors, used in no proof \\
\hline
\end{tabular}
\end{table}

% Signs and normalisations, part 2: topological recursion, legs, the tower, normalisations, stationary phase, the literature.

%-------------------------------------------------------------------------------
\subsection{Topological recursion}
\label{sec:sgn-tr}

Topological recursion is used in \cref{sec:mod} only, on the curve $x=\Wv(p)$, $y=p$, with the kernel of \cref{conv:mod-kernel}; $\iota_b$ is the local involution of $x$ at $\varsigma_b$. \Cref{tab:sgn-tr} compares it with the printed conventions of the sources quoted.

\begin{table}[ht]
\caption{Conventions of topological recursion.}
\label{tab:sgn-tr}
\centering
\footnotesize
\renewcommand{\arraystretch}{1.2}
\begin{tabular}{@{}>{\raggedright\arraybackslash}p{0.15\textwidth}>{\raggedright\arraybackslash}p{0.4\textwidth}>{\raggedright\arraybackslash}p{0.39\textwidth}@{}}
\hline
object & this note (\cref{conv:mod-kernel}) & sources \\
\hline
kernel & $K_b(p_0,q)=\frac12\int_{\iota_bq}^{q}\frac{\omega_{0,2}(p_0,\cdot)}{(y(q)-y(\iota_bq))\,\dd x(q)}$ near $\varsigma_b$, with $\omega_{0,1}=y\,\dd x$ and $\omega_{0,2}=\frac{\dd p_1\dd p_2}{(p_1-p_2)^2}$ & equal to the kernel of \cite[(2-26)--(2-29)]{Dunin-Barkowski:2015caa} and of \cite[(2.1)]{Hock25}; the printed kernel of \cite[(4-3), (4-9)]{Eynard:2007kz} equals that of \cite[(1.1)]{ABDKS25} with $\omega^{(0)}_{1,0}=-y\,\dd x$, and both are minus this kernel \\
correlators & $\omega_{h,n}$ & $\omega_{h,n}=(-1)^n\omega^{\mathrm{ABDKS}}_{h,n}$; $\omega_{h,n}=(-1)^nW^{\mathrm{EO}}_{h,n}$ for $2h-2+n\geq0$ \\
$\omega_{0,3}$ & $-\sum_b\Res_{q=\varsigma_b}\frac{\omega_{0,2}(q,p_1)\,\omega_{0,2}(q,p_2)\,\omega_{0,2}(q,p_3)}{\dd x(q)\,\dd y(q)}$ & minus \cite[Theorem~4.1, (4-20)]{Eynard:2007kz} \\
free energies & $(2-2h)\FTR_h=\sum_b\Res_{\varsigma_b}\Phi_{y\mathrm dx}\,\omega_{h,1}$, $\dd\Phi_{y\mathrm dx}=y\,\dd x$; on the model curve $\FTR_h=\frac{B_{2h}}{2h(2h-2)}v^{2-2h}$ for $h\geq2$ (\cref{thm:mod-gap}) & $\FTR_h=-F^{\mathrm{print}}_h$, with $F^{\mathrm{print}}_h$ of \cite[(4-10)]{Eynard:2007kz} \\
dilaton equation & $\sum_b\Res_{q=\varsigma_b}\Phi_{y\mathrm dx}(q)\,\omega_{h,n+1}(q,p_1,\dots,p_n)=(2-2h-n)\,\omega_{h,n}(p_1,\dots,p_n)$ for $2h-2+n>0$ & \cite[Theorem~4.7, (4-30)]{Eynard:2007kz}, with the factor $2h+n-2$ for $W^{\mathrm{EO}}$; the residues at the $p_i$ and the term $\delta_{g,0}\delta_{k,1}y\,\dd x$ printed there drop out in this range \\
variation in $v$ & \cref{thm:mod-var}, at fixed $x$, along a ray $\mathfrak r$ & \cite[Theorem~5.1, (5-14)]{Eynard:2007kz} with \S5.3, (5-27)--(5-30); for $\partial_v$ at fixed $x$ the function $\Lambda$ of \cite[(5-13)]{Eynard:2007kz} (a letter of that paper) is $-1$ on the ray, with the printed signs of \cite[(5-1), (5-13)]{Eynard:2007kz} \\
Laplace transforms & \cite[(2-31), (2-32)]{Dunin-Barkowski:2015caa} read as formal Laplace series, the finite part at $i=j$ giving the $\delta_{ij}$ of $R^{-1}=1+O(\vartheta^{-1})$ & \cref{lem:mod-laplace} \\
sign anchors & the curve $x=-p^2/2$, $y=p$ gives the intersection numbers of $\psi$ classes; $x=p+1/p$, $y=p$ gives $\FTR_2=-\frac1{240}$ (\cref{conv:mod-kernel}) & classical; used in no proof \\
\hline
\end{tabular}
\end{table}

%-------------------------------------------------------------------------------
\subsection{Legs}
\label{sec:sgn-legs}

The dictionary of record is the leg dictionary of \cref{prop:mod-LD}, with the legs $\phi^b_k$ of \cref{def:mod-legs} and the residue pairing of \cref{def:mod-pair}: the leg of $v^m\psi^kH^b$ is $-\varrho_b^{-1}v^m\phi^b_k$, $\varrho=(-1,1,-1,-3,-2)$. In particular the unit has the leg $y=p$, acting in the string equation (used in no proof) by $+\dd/\dd x$; the dilaton shift $q=-z\,1$ has the leg $\int y\,\dd x=-\phi^0_1$; $H$ has the logarithmic leg $-\log p$, paired along a ray $\mathfrak r$; and $\se=z\,1+vH$ has the leg $-p^5/5$. For legs with a pole at $p=0$ the residue pairing and the pairing on a circle $|p|=R>|\varsigma|$ differ by the residue at $p=0$, in general non-zero.

%-------------------------------------------------------------------------------
\subsection{The vanishing tower and regularity}
\label{sec:sgn-tower}

The note uses one tower of covectors, defined from the vanishing solution $\YC$ of the model (\cref{def:tra-tower,def:gau-van}):
\[
\Pi_i(q):=-\coef{z^i}\bigl\langle\YC(-z),q(z)\bigr\rangle,\qquad i\geq0;
\]
$\Pi_0$ is the vanishing covector, and $\Pi_i(z^kH^b)=\Pi_0(z^{k-i}H^b)$, which is zero for $k<i$. On insertions, $\Pi_0(z^kH^b)=(-1)^{k+1}\coef{z^{-k}}\YC_b$, which vanishes unless $b+k\equiv1\bmod5$; for instance $\Pi_0(H)=1$ and $\Pi_0(z\,1)=-v$.

\emph{Weights and regularity.} The $v$-exponent of a model correlator is $e=2-2h+\sum_i\wt(k_i,a_i)$, $\wt(k,a)=(4k+\nu_a-4)/5$ (\cref{lem:set-modhom}); in the graded setting of \cref{conv:tra-forms}, $v$ has weight $5$ and $z^kH^a$ the weight $4k+\nu_a+1=5+5\wt(k,a)$. A model correlator is regular if $e\geq0$, i.e.\ if it lies in $\bbQ[v]$; a graph term is regular if it lies in $\bbQ\fps{u}$, so for $G_\Gamma=u^{2g-2}T_\Gamma$ regularity of $T_\Gamma$ means $G_\Gamma\in u^{2g-2}\bbQ\fps{u}$. ``Regular'' never means ``zero''.

%-------------------------------------------------------------------------------
\subsection{Normalisations of gauges and series}
\label{sec:sgn-norm}

The normalisations used throughout are: $(\Pc)_{11}(0,z)=1$ (\cref{def:gau-can}); $\varpi_1=u+O(u^2)$, $\varpi_0(0)=1$, $\tc'(0)=1$ (\cref{lem:set-frobbasis,def:set-frames}); $v(u)=u+O(u^2)$ (\cref{def:set-tau}); $\Fs_0(0)=1$ (\cref{prop:set-modrows}); $\mathrm{ph}_0=1$ (\cref{def:set-rows}); $c_0(L)=1$ (\cref{def:gam-c}); $\SPh=1+O(z)$ (\cref{def:fsp-wick}); and $d(0)=1$ for the dilaton factor of $\Pc$ (\cref{lem:gau-diag}).

The other normalisation of the gauge, fixed by $p_{5n,0}(0)=0$ for $n\geq1$, where $P_k=\sum_bp_{k,b}E_{b+k,b}$ (\cref{def:gau-onepole}), equals $\Pc\,\cn(z)$ with $\cn(z)=1+m_5z^5+m_{10}z^{10}+\dots$ (\cref{rmk:mai-stored}); it enters no proof.

%-------------------------------------------------------------------------------
\subsection{Stationary phase}
\label{sec:sgn-sp}

\emph{Phases.} Stationary-phase integrands are written $\re^{F/z}$ in \cref{sec:fsp}; the saddle series of \cref{def:gam-c} has the phase $F=\mathsf S_0$; the flat integrand of \cref{sec:con} is written $\re^{-\cW/z}$ and the hatted one $\re^{-\hat\cW/\zh}$ (\cref{lem:con-form}, \cref{def:con-PhiSP}). They correspond under $F=-\cW$, resp.\ $F=-\hat\cW$ with $z$ renamed $\zh$ (\cref{def:fsp-wick}).

\emph{Measures.} In the Mellin--Barnes summands (\cref{sec:con-form}), $\Gamma(1+5\sigma/z)=\int y^{5\sigma/z}\re^{-y}\dd y$ carries $\dd y$ and $\Gamma(w)=\int x^w\re^{-x}\dd x/x$ carries $\dd x/x$; an insertion $x^m$ with $\dd x/x$ is the case $M=m$ of \cref{lem:fsp-stirling}, with $\dd x$ the case $M=m+1$.

\emph{Normalised and full stationary phase.} The series $\SPn$ of \cref{def:fsp-wick} is normalised by the Gaussian, so its constant term is the value $a_0(c)$ of the $z^0$-part of the amplitude at the critical point; $\SPh=\SPn/a_0(c)$; the full symbol is $(2\pi z)^{n/2}\det(-\mathsf H)^{-1/2}\re^{F(c)/z}\SPn$, $\mathsf H$ the Hessian at $c$. Only the series $\SPh$ are canonical: a prefactor such as $(2\pi z/5)^{1/2}$ cancels in a comparison (\cref{def:con-PhiSP}).

\emph{Branches.} Roots and logarithms occur only in the symbols: $\det(-J^{\mathrm T}\mathsf HJ)^{1/2}:=\det J\cdot\det(-\mathsf H)^{1/2}$ under a change of variables with Jacobian $J$ (\cref{lem:fsp-cov}), $\det(-\mu\mathsf H)^{1/2}:=\mu^{n/2}\det(-\mathsf H)^{1/2}$ under a rescaling by $\mu$ (\cref{lem:fsp-homog}), $\log(a/z):=\log a-\log z$ and $\Gst(1+w):=w\,\Gst(w)$ (\cref{def:fsp-stirling}). The saddle series $c(z;L)$ does not depend on the branches (\cref{prop:gam-rational}).

\emph{$q$ and $-q$.} In every form of the Mellin--Barnes summand the critical equation is $-q(5\sigma)^5=\prod_j(\zeta^j-\sigma)$, i.e.\ $u\sigma^5=1$, so $\sigma=L$ is the critical point of colour $0$; the logarithm $\lambda$ of $-q$ is forced by criticality (\cref{def:gam-c}), and $c(z;L)$ depends on neither $\lambda$ nor the sign convention for $q$. The two summands $(-q)^{\sigma/z}\Gamma(1+5\sigma/z)\prod_j\Gamma((\zeta^j-\sigma)/z)$ and $q^{\sigma/z}\Gamma(1+5\sigma/z)/\prod_j\Gamma(1+(\sigma-\zeta^j)/z)$ give the same normalised series (\cref{lem:con-form}). In Theorem~LMI the chart parameter $\mathsf Q=\re^{-\lambda/a_*}$ is $-q$ at $s_*=1$ and $(-q)^{-1/5}$ at $s_*=0$ (\cref{def:gam-lmidata,thm:gam-main,thm:gam-ginf}).

\emph{Exponentials.} In the comparison of the actual rows with the canonical gauge the same exponential $\re^{u_i/z}$ is removed on both sides (\cref{thm:gau-R0}); the hatted phase gives $\re^{(u_0-u_0^{\mathrm c})/z}$, off by $\re^{-u_0^{\mathrm c}/z}$, so \cref{sec:con} uses normalised series only.

%-------------------------------------------------------------------------------
\subsection{Conventions of the literature}
\label{sec:sgn-lit}

\Cref{tab:sgn-lit} records the conventions of the sources in their own letters; they are used only through the translations given at the places cited.

\begin{table}[ht]
\caption{Conventions of the literature and of the note.}
\label{tab:sgn-lit}
\centering
\footnotesize
\renewcommand{\arraystretch}{1.2}
\begin{tabular}{@{}>{\raggedright\arraybackslash}p{0.14\textwidth}>{\raggedright\arraybackslash}p{0.47\textwidth}>{\raggedright\arraybackslash}p{0.33\textwidth}@{}}
\hline
topic & literature & this note \\
\hline
$R$-matrix action & \cite[Def.~2.2]{PPZ15}; \cite[\S2.3]{Dunin-Barkowski:2015caa} (transposed matrix, opposite dilaton sign) & \cite{PPZ15} throughout (\cref{sec:sgn-R}) \\
characters & \cite[\S2.2, Theorem~1]{Giv01}: tangent characters, minus sign; \cite[9.1~Theorem]{MR1901075}: cotangent characters, plus sign & equal, since only odd powers occur \\
twist weight & \cite[(17)]{GJR18}: $-5h_\alpha$; \cite[(20)]{GJR17}: $t-5H$; \cite[Corollary~1]{MR2276766}: the symbol $\exp\bigl(+\sum_m\bk_m(z/(\lambda+\rho))^{2m-1}\bigr)$ of an Euler-class twist; \cite[(26)]{LP20}: $(5\lambda_i)/e(\mathbb E^*\tensor(5\lambda_i))$ & all agree with $-5\zeta^i\in\bw_i$ (\cref{thm:set-qrr,rmk:set-qrrsign}) \\
asymptotic series & \cite[Theorem~4 (i)]{ZZ08}, \cite[\S2.4.4, (18)]{LP20}: $\Phi_s(0)=\delta_{0,s}$ & (QRR): the value $h_{\bw_i}$ at $q=0$ \\
Landau--Ginzburg charts & \cite[\S7, Remark~7.2]{GJR18}: critical point $(x_i)_\alpha=L_\alpha-\lambda_i$, $y_\alpha=mL_\alpha$, measure $\dd x/x$ with $x^{k/m}$; \cite[\S10]{MR1901075}: chart $x_0=Q/(x_1\cdots x_n)$, measure $\dd x$; \cite[Lemma~4.20]{MR2553377}: $\re^{-W/z}$ at $z=-1$, measure $\dd w/w$; \cite[\S7, (11)]{MR2276766}: $\dd x$ with $(2\pi z(\lambda+\rho))^{-1/2}$ & the summand of \cref{def:gam-c}: critical point $x_j=\zeta^j-L$ (opposite sign) and $y=5L$, measures $\dd y$ and $\dd x_j/x_j$ (\cref{sec:sgn-sp}) \\
gap constants & \cite[(1.2)]{HKQ09}: $\frac{(-1)^{g-1}B_{2g}}{2g(2g-2)}t_D^{2-2g}$; \cite[Conjecture~1.5]{CGYZ26} (quintic): $\frac{(-1)^{g-1}B_{2g}}{5^{g-1}2g(2g-2)}\tau_c^{2-2g}$; \cite[Definition~1.8]{CGYZ26} has $(-1)^g$ with a constant $\kappa$ and is not used; \cite[(1.5)]{Br25} (local $\bbP^2$): $\frac{3^{g-1}B_{2g}}{2g(2g-2)}t_{\mathrm{CF}}^{2-2g}$ & $c_g=\frac{(-1)^{g-1}B_{2g}}{2g(2g-2)5^{g-1}}$ for $\varpi_1=u+O(u^2)$; rescaling $\varpi_1$ by $\sqrt5$ gives the form of \cite[(1.2)]{HKQ09} (see below); $\Fm_h=(-5)^{1-h}\FTR_h=c_hv^{2-2h}$ \\
mirror theorem & \cite[Corollary~11.7]{MR1408320} prints $l_\alpha$ for $\lambda_\alpha$ & \cite[Theorem~11.6]{MR1408320} is used (\cref{thm:set-mirror}) \\
\hline
\end{tabular}
\end{table}

\emph{The comparison with \cite{CGYZ26}} (\cref{rmk:int-constant,rmk:mai-constant}). In our notation, and as read in version~2 of the preprint, \cite[Def.~1.1]{CGYZ26} works over $\bbQ(q)$ with $\Dq$, the periods $C_0$, $I_1$ of $\PFfour$ (the genus-$0$ periods of the formal quintic are those of the quintic, \cref{def:set-genus0}), $T=I_1/C_0$ and the differential $\bbQ(q)$-algebra generated by $C_0^{\pm1}$, $C_1^{\pm1}$. A pair of periods at the conifold point, one not vanishing and one vanishing simply, is admissible \cite[Def.~1.2]{CGYZ26} if substituting them for $C_0$, and the derivative of their quotient for $C_1$, respects the relations of that algebra; this defines their conifold expansion, in $\Delta=1-5^5q$. Their pair \cite[Ex.~1.3]{CGYZ26} is $(f_0,\varpi_1)$, the frame $(\alpha,\beta)=(0,0)$ with $\Delta=u$ and their flat coordinate $\tc$; their family \cite[Rem.~1.4]{CGYZ26} with rational coefficients is \cref{def:set-frames}; and $\varpi_1=u+O(u^2)$ is the normalisation of the pair of \cite[Ex.~1.3]{CGYZ26} and of \cite[Def.~1.8\,(2); \S4.2, Prop.-Def.~4.5\,(1)]{CGYZ26}; \cite[Conj.~1.5]{CGYZ26} is read with this normalisation, on which its constant depends (see the rescaling below). Both expansions are differential homomorphisms for $\Dq$ sending $C_0\mapsto\varpi_0$ and $C_1\mapsto\Dq\tc$ and expanding at $u=0$, so they agree on the differential subring of $\Nring$ generated by $c_0^{\pm1}$ and $(\mathsf W/c_0^2)^{\pm1}$ over $\bbQ[L^5,L^{-5}]$, wherever theirs is defined. The constant of \cite[Conj.~1.5]{CGYZ26} is $(-1)^{g-1}B_{2g}/(5^{g-1}\,2g(2g-2))=c_g$; the sign convention of \cite[Def.~1.8]{CGYZ26} is not used.

\emph{Rescaling the periods.} Let $\varkappa_0,\varkappa_1\in\bbC^\times$ (local letters), let $\Cexp'$ be the conifold expansion of \cref{def:set-confexp} with $(\varpi_0,\varpi_1)$ replaced by $(\varkappa_0\varpi_0,\varkappa_1\varpi_1)$, and $\tc'\coloneqq(\varkappa_1/\varkappa_0)\tc$. Then
\[
\Cexp'(F_g^{\Nring})-c_g\,\varkappa_1^{2g-2}\,\tc'^{\,2-2g}\in\bbC\fps{u},
\]
so the constant does not depend on $\varkappa_0$ and is $c_g$ if and only if $\varkappa_1^{2g-2}=1$. Indeed, the $\bbC$-linear endomorphism $\mathfrak k$ (local) of $\Nring\otimes\bbC$ with $c_j\mapsto\varkappa_0c_j$, $i_j\mapsto\varkappa_1i_j+(\varkappa_1-\varkappa_0)(\logq\,c_j+jc_{j-1})$, $\logq\mapsto\logq$, $L\mapsto L$ commutes with $\Dq$ ($\PFfour$ is linear and the concomitant relation bilinear), satisfies $\Cexp'=\Cexp\circ\mathfrak k$, and sends $C_0$, $I_1$, $C_1$, $\mathsf W$ to $\varkappa_0C_0$, $\varkappa_1I_1$, $(\varkappa_1/\varkappa_0)C_1$, $\varkappa_0\varkappa_1\mathsf W$. It fixes the flat amplitudes $A_n\in\bbQ(\zeta)[L,L^{-1}]$ (\cref{thm:gau-R0}~(iii)), so by the forward relations of \cref{prop:set-cyclicLR}, with coefficients $C_1\kappa^{\mathrm{LR}}_a=(C_1,L^5/(C_0^2C_1^2),C_1,C_0,C_0)$, it multiplies the rows and $\Psi_{\mathrm Q}$ by $\diag(\varkappa_0^{-1},\varkappa_1^{-1},\varkappa_1,\varkappa_0,1)$. Hence $\mathfrak k(\tilde R)=\tilde R$, while $\Delta_i^{1/2}\mapsto\varkappa_0\Delta_i^{1/2}$ and $1_{\TFT}\mapsto\varkappa_0^{-1}1_{\TFT}$: $\mathfrak k$ fixes the edges, multiplies every leaf insertion by $\varkappa_0^{-1}$ and every vertex with $n_v$ half-edges and $l_v$ leaf insertions by $\varkappa_0^{2g_v-2+n_v+l_v}$, so every term by $\varkappa_0^{2g-2}$ (as in the proof of \cref{lem:flh-comm}); so $\Cexp'(F_g^{\Nring})=\varkappa_0^{2g-2}F_g^{\mathrm{cf}}$, and Corollary~B gives the claim. Consequently, for every pair of the family of \cite[Rem.~1.4]{CGYZ26} normalised as in \cite[Prop.-Def.~4.5\,(1)]{CGYZ26}, the formal quintic has the conifold gap in the sense of Corollary~B, with the constant of \cite[Conj.~1.5]{CGYZ26}; by the membership statement below this is also the gap for the conifold expansion of \cite[Def.~1.2]{CGYZ26}, when the pair is admissible.

\emph{Membership.} The comparison needs that $F_g^{\Nring}$ lies in the differential subring $\mathscr S$ of $\Nring$ generated by $c_0^{\pm1}$ and $(\mathsf W/c_0^2)^{\pm1}$ over $\bbQ[L^5,L^{-5}]$, on which the two expansions agree; then $F_g=\ev(F_g^{\Nring})$ lies in the algebra of \cite[Def.~1.1]{CGYZ26} as recalled above, since $\ev$ maps $L^5$, $c_0$ and $\mathsf W/c_0^2$ to $(1-5^5q)^{-1}$, $C_0$ and $C_1$; and for every frame whose pair $(\varpi_0,\varpi_1)$ is admissible, which is asserted in \cite[Ex.~1.3, Rem.~1.4]{CGYZ26} and is not proved here, $F_g^{\mathrm{cf}}$ is the conifold expansion of $F_g$ in the sense of \cite[Def.~1.2]{CGYZ26}: that expansion composed with $\ev$, and $\Cexp_{\alpha,\beta}$, are homomorphisms of differential rings on $\mathscr S$ which agree on its generators. The injectivity of $\ev$ is not used. The membership holds, under the inputs of \cref{thm:gam-D} (the letters $\mathscr S$, $\mathsf B$, $\mathsf B_5$, $x$, $\pi_j$, $\varrho$ are local). (1) $F_g^{\Nring}$ is free of $\logq$ (\cref{cor:gau-rowsNN}~(ii)) and fixed by the automorphisms $i_j\mapsto i_j-kc_j$, $\logq\mapsto\logq+k$, $k\in\bbQ$ (\cref{def:set-confexp}). With $x:=i_0/c_0$ and $\pi_j:=i_j-xc_j$ ($j=1,2$), the $\logq$-free subring of $\Nring$ is the polynomial ring $\mathsf B[x]$ over $\mathsf B:=\bbQ[L^{\pm1}][c_0^{\pm1},c_1,c_2,c_3,\pi_1,\pi_2][\mathsf W^{-1}]$, because $\mathsf W=c_0^2+c_0\pi_1$; the automorphisms fix $\mathsf B$ and send $x$ to $x-k$, so $F_g^{\Nring}\in\mathsf B$. (2) Let $\varrho$ be the automorphism of $\Nring\otimes\bbQ(\zeta)$ with $L\mapsto\zeta L$ which fixes $\zeta$ and the other generators; it commutes with $\Dq$, since $\Dq L=L(L^5-1)/5$ and $\Dq c_3$, $\Dq i_2$ involve $L$ only through $L^5$. By \cref{thm:gam-D}, $A^{\mathrm{act}}_n=A^{\mathrm{can}}_n=L^{4n+1}\mathrm{ph}^{\mathrm{can}}_n(L^{-5})$ (\cref{thm:gau-R0}~(ii), (iii)), so $\varrho$ multiplies $\Qr^{(0)}_{0,n}=A^{\mathrm{act}}_nc_0^{-1}$ by $\zeta^{1-n}$ and, by the recursion of \cref{cor:gau-rowsNN}~(i), whose coefficients it fixes, $\Qr^{(0)}_{a,n}$ by $\zeta^{a+1-n}$. Hence $\varrho(\Qr^{(i)}_{a,n})=\Qr^{(i+1)}_{a,n}$, $\varrho(\Delta_i^{1/2})=\Delta_{i+1}^{1/2}$ and $\varrho\bigl((\Psi_{\mathrm Q})_{aj}\bigr)=(\Psi_{\mathrm Q})_{a,j+1}$ (\cref{lem:mai-Psi}~(i)): $\varrho$ permutes the colours and fixes $F_g^{\Nring}$. As $\mathsf W$ is free of $L$, $\mathsf B=\bigoplus_{k=0}^4L^k\mathsf B_5$, where $\mathsf B_5$ (local) is the ring displayed in (3), and $\varrho$ acts on the $k$-th summand by $\zeta^k$; so $F_g^{\Nring}\in\mathsf B_5$. (3) $\mathsf B_5:=\bbQ[L^{\pm5}][c_0^{\pm1},c_1,c_2,c_3,\pi_1,\pi_2][\mathsf W^{-1}]\subset\mathscr S$, as $c_j=\Dq^jc_0$, $\pi_1=(\mathsf W-c_0^2)/c_0$ and $\pi_2=(\Dq\mathsf W-2c_0c_1)/c_0$.

\emph{Differences.} The rows involve $L\notin\bbQ(q)$, yet $F_g^{\mathrm{cf}}$ is a Laurent series in $u$ by Corollary~B. The expansion of \cite{CGYZ26} is defined on functions under admissibility, $\Cexp$ on the presented ring $\Nring$ (\cref{lem:gau-confexp}).


% Evidence appendix, part 1 of 3 (continued in parts 2 and 3); label prefix: evi. No new macros.
\section{Evidence: finite computations}
\label{sec:evi}
\noindent\genusstatus{numerically supported} The finite computations recorded in this appendix are exact or modular computations; they support the statements at the genera, orders and primes given, and are not proofs. The code and outputs are in \texttt{code/} of this project.


Nothing in this appendix is a step of a proof: the proofs of the note are complete without it, and no proof refers to it. It records in tables some finite computations made while the statements were developed, with exact arithmetic over $\bbQ$ or $\bbQ(\zeta)$ or arithmetic modulo the primes named (a zero test modulo a prime is probabilistic); conifold frames are labelled by $(\alpha,\beta)$ as in \cref{def:set-frames}, $\varpi_0=f_0+\alpha\varpi_1+\beta f_2$.

\subsection{Graph sums of genus 2 to 5}
\label{sec:evi-graphs}

The outer terms $T_\Gamma=u^{2-2g}G_\Gamma(u)$ (\cref{def:set-outer,lem:flh-hat}), with $M=\Mo(\tau(u)\zh)$, were computed for every leg-free stable graph of genus $2$ to $5$, modulo $2^{89}-1$ and $2^{61}-1$ with rational reconstruction, in the frames $(0,0)$ and $(1/3,-2/7)$; an edge graph is regular if and only if $G_\Gamma=O(u^{2g-2})$. At genus $4$ and $5$ the outer sum was obtained by regrouping the plain graph sum along the edges of $M$, a regrouping checked against genuine vertex tensors only at genus $2$ and $3$. With the bare symbol $\Mo(\zh)$ all $149$ edge graphs of genus $4$ with at most five edges are singular, and with $\tau$ replaced by $\tau+u^5$ exactly the three one-edge graphs of genus $4$ fail. The value $c_g$ and the gap of the total do not depend on $M$; only the regularity of the edge graphs tests the factorisation. The results are in \cref{tab:evi-g2,tab:evi-g3,tab:evi-g4,tab:evi-g5}.

\begin{table}[ht]
\centering\small\setlength{\tabcolsep}{4pt}
\begin{tabular}{l|r|r|r}
graph of genus $2$ & plain, $\coef{\tc^{-2}}$ & plain, $\coef{\tc^{-1}}$ & natural, at $u=0$\\ \hline
one vertex of genus $2$ & $1807/3072000$ & $-791/15360000$ & $\tfrac{1}{1200}\tc^{-2}+O(1)$\\
genus-$1$ vertex with a loop & $-811/256000$ & $-1771/3840000$ & $479/756000$\\
genus-$0$ vertex with two loops & $-57/64000$ & $41/12800$ & $1/2250$\\
two genus-$1$ vertices, one edge & $843/1024000$ & $743/15360000$ & $-127/2016000$\\
genus $0$ with a loop, edge to genus $1$ & $3/51200$ & $-41/153600$ & $0$\\
two genus-$0$ vertices, three edges & $87/32000$ & $37/160000$ & $0$\\
two genus-$0$ vertices with loops, an edge & $9/12800$ & $-173/64000$ & $0$\\ \hline
sum & $1/1200$ & $0$ & \\
\end{tabular}
\caption{Genus $2$, per stable graph, summed over colourings, in the frame $(0,0)$, modulo both primes: the polar coefficients of the plain graph sum of \cref{def:set-graphsum} (vertices the TFT $\TFT$), and the terms $T_\Gamma$ of the natural form at $u=0$ (in the frame $(1/3,-2/7)$, modulo $2^{89}-1$, the edge graphs are regular as well). The $115$ coloured graphs give identical tables at both primes.}
\label{tab:evi-g2}
\end{table}

\begin{table}[ht]
\centering\small
\begin{tabular}{l|p{0.68\textwidth}}
stable graphs & $42$; with $0,1,\dots,6$ edges: $1,2,5,9,12,8,5$\\
edge graphs & all $41$ regular: frame $(0,0)$ modulo both primes, frame $(1/3,-2/7)$ modulo $2^{89}-1$\\
vertex graph, $(0,0)$ & $u^{-4}\bigl(\tfrac{1}{25200}-\tfrac{u}{9000}+\tfrac{407u^2}{3780000}-\tfrac{151u^3}{3780000}+\dots\bigr)=\tfrac{1}{25200}\tc^{-4}+O(1)$\\
vertex graph, $(1/3,-2/7)$ & $u^{-4}\bigl(\tfrac{1}{25200}-\tfrac{11u}{189000}+\dots\bigr)$, with the same polar part\\
\end{tabular}
\caption{Genus $3$, the natural form.}
\label{tab:evi-g3}
\end{table}

\begin{table}[ht]
\centering\small
\begin{tabular}{l|p{0.72\textwidth}}
stable graphs & $379$; with $0,1,\dots,9$ edges: $1,3,7,21,43,75,89,81,42,17$\\
edge graphs & all $378$ regular, in both frames, modulo both primes\\
total, $(0,0)$ & $u^{-6}\bigl(\tfrac{1}{180000}-\tfrac{7u}{300000}+\tfrac{701u^2}{18000000}-\tfrac{1459u^3}{45000000}+\tfrac{923639u^4}{67500000000}-\tfrac{171583u^5}{67500000000}\bigr)$ $=\tfrac{1}{180000}\tc^{-6}+O(1)$\\
total, $(1/3,-2/7)$ & $u^{-6}\bigl(\tfrac{1}{180000}-\tfrac{11u}{900000}-\tfrac{79u^2}{378000000}+\dots\bigr)$, with the same polar part\\
\end{tabular}
\caption{Genus $4$, the natural form; $1/180000=c_4$.}
\label{tab:evi-g4}
\end{table}

\begin{table}[ht]
\centering\small
\begin{tabular}{l|p{0.72\textwidth}}
stable graphs & $4555$; with $0,1,\dots,12$ edges: $1,3,11,34,100,239,492,784,1002,926,632,260,71$\\
edge graphs & all $4554$ regular: frame $(0,0)$ modulo both primes, frame $(1/3,-2/7)$ modulo $2^{89}-1$\\
total, $(0,0)$ & $u^{-8}\bigl(\tfrac{1}{660000}-\tfrac{7u}{825000}+\tfrac{199u^2}{9900000}-\tfrac{717u^3}{27500000}+\tfrac{449357u^4}{22500000000}-\tfrac{5593217u^5}{618750000000}$ $+\tfrac{14686042363u^6}{6496875000000000}-\tfrac{8370858511u^7}{32484375000000000}\bigr)=\tfrac{1}{660000}\tc^{-8}+O(1)$, by the Chinese remainder theorem from both primes\\
total, $(1/3,-2/7)$ & $u^{-8}\bigl(\tfrac{1}{660000}-\tfrac{u}{225000}+\tfrac{323u^2}{207900000}+\tfrac{66841u^3}{15592500000}-\dots\bigr)$, modulo both primes: different from the first at every order $u^1,\dots,u^7$, with the same polar part\\
\end{tabular}
\caption{Genus $5$, the natural form; $1/660000=c_5$.}
\label{tab:evi-g5}
\end{table}

% Evidence appendix, part 2 of 3: the normalisation. Continued in part 3.

\subsection{The normalisation and the identity \texorpdfstring{$(\Gamma)$}{(Gamma)}}
\label{sec:evi-normal}

\Cref{tab:evi-windows} lists finite windows in which $\Dn=1$, the identity \eqref{eq:gam-Gamma} and some identities entering their proofs were tested directly, independently of the proofs of \cref{thm:gam-D} and \cref{thm:gam-main}. Since $\Phican=\PhiSP$ in every conifold frame (\cref{thm:con-Pb}) and $\PhiSP$ involves no frame, further conifold frames give no independent evidence for $\Dn=1$.

\begin{table}[ht]
\centering\small
\begin{tabular}{>{\raggedright\arraybackslash}p{0.4\textwidth}|>{\raggedright\arraybackslash}p{0.3\textwidth}|>{\raggedright\arraybackslash}p{0.22\textwidth}}
statement & range & arithmetic\\ \hline
$\Dn_0=1$, order by order at the conifold point & $z^{204}$ ($g\leq 69$), two implementations; $z^{125}$, a third & exact, frame $(0,0)$\\
$\Dn_0=1$ & $z^{1004}$ ($g\leq 335$) & modulo $2^{89}-1$, $2^{61}-1$\\
the actual gauge is $\Pc$ (\cref{cor:mai-actual}) & $z^{35}$; $z^{150}$ & two primes, two frames\\
\eqref{eq:gam-Gamma}: $c_n(1)=h_n$ & $n\leq 60$ (Laurent expansion at $L=1$); $n\leq 40$ (tracked precision); $n\leq 30$ (fit at $32$ points) & exact over $\bbQ$\\
\eqref{eq:gam-Gamma} & $z^{500}$ & modulo $2^{127}-1$, $2^{107}-1$\\
the chart of \cref{thm:gam-lmi} at $L=1$ equals $h$ & $z^{60}$ & exact over $\bbQ(\zeta)$\\
$\PhiSP=\Phican$ (\cref{thm:con-Pb}) & $\zh^{80}$, $u^{75}$ & modulo $2^{61}-1$ (two frames), $2^{89}-1$\\
$\cV=\varpi_0(\YCQ)_0$ (\cref{prop:con-res} with (G)) & $u^{40}$; $u^{30}$ & exact, two frames; two primes\\
$c(z;0)=\Ginf$ (\cref{thm:gam-ginf}) & $z^{60}$ (chart at $L=0$); $z^{80}$ (the value) & exact\\
$\fc^{\mathrm{can}}_m=\coef{z^{5m}}\Ginf$ (\cref{cor:gam-free}) & $m\leq 25$; $m\leq 100$ & exact; two primes\\
\end{tabular}
\caption{Finite windows of $\Dn=1$ and of the identity \eqref{eq:gam-Gamma}.}
\label{tab:evi-windows}
\end{table}

% Evidence appendix, part 3 of 3: the model. Last part of the appendix.

\subsection{The model}
\label{sec:evi-model}

\Cref{tab:evi-model} lists the main checks of \cref{sec:mod}. Further, the formula of \cref{rmk:mod-hock} was checked for $h\leq 12$ exactly, and the singular lists of the Model Lemma at genus $0$ to $3$ all vanish, while deliberately wrong frames give non-zero values. Off $\sK$ (\cref{rmk:mod-offK}), exactly at $v=1$, $\langle\psi^5H\rangle_3=53/1209600$, whereas the one-point law of \cref{cor:mod-onepoint} would give $-1/604800$.

\begin{table}[ht]
\centering\small
\begin{tabular}{>{\raggedright\arraybackslash}p{0.4\textwidth}|>{\raggedright\arraybackslash}p{0.36\textwidth}|>{\raggedright\arraybackslash}p{0.16\textwidth}}
statement & range & arithmetic\\ \hline
topological recursion on $x=\Wv$ is $(-5)^{h-1}$ times the model CohFT, i.e.\ the model correlators are $(-5)^{1-h}$ times those of topological recursion (\cref{thm:mod-tr}) & $72120$ colour-basis correlators with $\psi$-deficit $\leq 6$ at $(h,n)=(0,3)$, $(0,4)$, $(0,5)$, $(1,1)$, $(1,2)$, $(1,3)$, $(2,1)$, $(2,2)$, $(3,1)$ ($4665$ non-zero); on $x=\kappa\Wv$, $\kappa=1,4,-1$: $6495$ each & two primes\\
the R-matrix of the curve $x=\kappa\Wv$ is $\Mo(z/(\kappa\varsigma^4))$ (\cref{prop:mod-Rtr}) & all $25$ entries through $z^{10}$, $\kappa=1,4,-1$, $\varsigma=1,2$ & exact, $\bbQ(\zeta)$\\
the dilaton-leaf sign of \cite[Def.~2.13]{PPZ15} & $(h,n)=(1,1)$, four points, $\kappa=1,4$; the opposite sign fails at all four & exact\\
$\FTR_h(v=1)=B_{2h}/(2h(2h-2))=-\tfrac{1}{240},\tfrac{1}{1008},-\tfrac{1}{1440},\tfrac{1}{1056}$ & $h=2,\dots,5$ & $\bbQ(\zeta)$ ($h\leq 3$); two primes\\
the period identity of \cref{prop:mod-period}; $\Sigma_h=-4^hB_{2h}/(2h)$ & $h\leq 7$; $h\leq 20$ & exact\\
the lattice theorem (\cref{thm:mod-lattice}) & weight $\leq 400$; $\leq 1000$; $16200$ raw-leg kernel vectors & exact; two primes\\
the leg dictionary (\cref{prop:mod-LD}) & $35000$ pairings per prime; $52$ frame one-point functions, $h=2,3$ & two primes; $\bbQ(\zeta)$\\
the Model Lemma on $\sK$ (\cref{thm:mod-MLK}) & $5289$ correlators with $h\leq 4$, including $1302$ at $(h,n)=(2,4)$; $1395$ singular, all zero & exact; two primes\\
\end{tabular}
\caption{Checks on the model.}
\label{tab:evi-model}
\end{table}




\bibliographystyle{amsalpha}
\bibliography{refs}

\end{document}
